\documentclass[UTF8, 12pt, fontset=fandol, oneside]{ctexart}
\usepackage{researchnotes}
\providecommand{\xlongequal}[1]{\stackrel{#1}{=}}

\graphicspath{{math/数分习题/}}

\title{第五章\quad 级数}
\author{}
\date{}

\begin{document}

\maketitle

\noindent\textbf{导读}\quad 级数是一个工具，有完善的理论，是数学分析课程三大主干内容之一，也是考研重点内容之一. 本章内容适合各类读者，过于理论性的问题（如一致收敛等），非数学院系学生一般不作过高要求.

\section*{§5.1\quad 数项级数}

\subsection*{一、求和问题}

级数求和的问题，一般来说，是一个困难问题，没有什么好办法. 因为部分和
\[
  S_n=\sum_{k=1}^n a_k
\]
随 $n$ 增大时，项数越来越多，除非能化为已知级数，人们只能设法把 $S_n$ 写成紧缩形式，才便于求极限. 本段主要讨论把 $S_n$ 转化为紧缩形式的几种常用方法，以及用于列求极限的方法. 至于用 Abel 第二定理化为幂级数求和问题，我们将在 §5.3 专门讨论.

\noindent\textbf{a.\quad 利用已知级数}

\noindent\textbf{例 5.1.1}\quad 计算
\[
  \frac12+\frac3{2^2}+\frac5{2^3}+\cdots+\frac{2n-1}{2^n}+\cdots.
\]

\noindent\textbf{解}\quad
\[
\begin{aligned}
S_n=2S_n-S_n
={}&1+\frac32+\frac5{2^2}+\cdots+\frac{2n-1}{2^{n-1}}
      -\frac12-\frac3{2^2}-\cdots-\frac{2n-3}{2^{n-1}}-\frac{2n-1}{2^n}\\
={}&1+1+\frac12+\cdots+\frac1{2^{n-2}}-\frac{2n-1}{2^n}\\
={}&1+\frac{1-\dfrac1{2^{n-1}}}{1-\dfrac12}-\frac{2n-1}{2^n}.
\end{aligned}
\]
故原级数的和
\[
  S=\lim_{n\to\infty}S_n=3.
\]

\noindent\textbf{练习}\quad 计算 $\displaystyle\sum_{n=1}^\infty ne^{-nx}$.

\noindent\textbf{提示}\quad 计算 $(1-e^{-x})S_n$.

\noindent\textbf{再提示}\quad
\[
\begin{aligned}
(1-e^{-x})\sum_{k=1}^n ke^{-kx}
={}&(e^{-x}+2e^{-2x}+3e^{-3x}+\cdots+ne^{-nx})\\
&-\left[e^{-2x}+2e^{-3x}+\cdots+(n-1)e^{-nx}+ne^{-(n+1)x}\right]\\
={}&e^{-x}+e^{-2x}+\cdots+e^{-nx}-\frac{n}{e^{(n+1)x}}
\to \frac{e^{-x}}{1-e^{-x}}\quad(n\to\infty).
\end{aligned}
\]
故
\[
  \lim_{n\to\infty}\sum_{k=1}^n ke^{-kx}
  =\frac{e^{-x}}{(1-e^{-x})^2}.
\]

\noindent\textbf{例 5.1.2}\quad 设级数 $\displaystyle\sum_{n=1}^\infty a_n$ 收敛于 $A$（有限数），试证：
\[
  \lim_{n\to\infty}\frac1n\left[a_n+2a_{n-1}+\cdots+(n-1)a_2+na_1\right]=A.
\qquad(1)
\]
（华东师范大学）

\noindent\textbf{证}\quad 记 $S_n=\displaystyle\sum_{k=1}^n a_k$，则
\[
  \lim_{n\to\infty}S_n=\lim_{n\to\infty}\sum_{k=1}^n a_k\stackrel{\text{已知}}{=}A.
\qquad(2)
\]
且容易看出
\[
  a_n+2a_{n-1}+\cdots+(n-1)a_2+na_1=S_1+S_2+\cdots+S_n.
\qquad(3)
\]
因为：在 $S_1+S_2+\cdots+S_n$ 中，按 $a_k$ 同类项合并，即得
\[
  a_n+2a_{n-1}+\cdots+(n-1)a_2+na_1.
\]
或（写成三角矩阵）将 $S_k(k=1,2,\cdots,n)$ 展开，各摆一行，排成（如下）三角矩阵：
\[
\begin{pmatrix}
a_1&0&0&\cdots&0\\
a_1&a_2&0&\cdots&0\\
a_1&a_2&a_3&\cdots&0\\
\vdots&\vdots&\vdots& &\vdots\\
a_1&a_2&a_3&\cdots&a_n
\end{pmatrix},
\]
那么式 (3)：左端 =“列和”相加，右端 =“行和”相加. 故式 (3) 成立.
\[
\begin{aligned}
\lim_{n\to\infty}\frac1n\left[a_n+2a_{n-1}+\cdots+(n-1)a_2+na_1\right]
&\stackrel{\text{用 Stolz 公式}}{=}
\lim_{n\to\infty}\frac{S_1+S_2+\cdots+S_n}{n}\\
&=\lim_{n\to\infty}S_n\stackrel{\text{式 }(2)}{=}A.
\end{aligned}
\]
式 (1) 获证.

\noindent\textbf{b.\quad 连锁消去法}

\noindent\textbf{例 5.1.3}\quad 设 $0<x<1$，求如下级数之和：
\[
  \sum_{n=0}^\infty \frac{x^{2^n}}{1-x^{2^{n+1}}}.
\]
（国外赛题）

\noindent\textbf{解}\quad
\[
\begin{aligned}
\sum_{k=0}^n \frac{x^{2^k}}{1-x^{2^{k+1}}}
&=\sum_{k=0}^n\left(\frac1{1-x^{2^k}}-\frac1{1-x^{2^{k+1}}}\right)\\
&=\left(\frac1{1-x}-\frac1{1-x^2}\right)
 +\left(\frac1{1-x^2}-\frac1{1-x^{2^2}}\right)+\cdots\\
&\quad+\left(\frac1{1-x^{2^{n-1}}}-\frac1{1-x^{2^n}}\right)
 +\left(\frac1{1-x^{2^n}}-\frac1{1-x^{2^{n+1}}}\right)\\
&=\frac1{1-x}-\frac1{1-x^{2^{n+1}}}.
\end{aligned}
\]
因此
\[
  \sum_{n=0}^\infty \frac{x^{2^n}}{1-x^{2^{n+1}}}
  =\lim_{n\to\infty}\left(\frac1{1-x}-\frac1{1-x^{2^{n+1}}}\right)
  =\frac{x}{1-x}.
\]

\noindent\textbf{例 5.1.4}\quad 求如下级数之和：
\[
  1)\ \sum_{k=1}^\infty \arctan\frac1{2k^2},\qquad
  2)\ \sum_{k=2}^\infty \arctan\frac2{4k^2-4k+1}.
\]

\noindent\textbf{提示}\quad 利用公式
\[
  \arctan x-\arctan y=\arctan\frac{x-y}{1+xy},\qquad
  \arctan\frac1{2k^2}=\arctan\frac1{2k-1}-\arctan\frac1{2k+1}.
\]

这种连锁消去法，还可以是多项相消，如

\noindent\textbf{例 5.1.5}\quad 计算
\[
  \sum_{n=1}^\infty\left(\sqrt n-2\sqrt{n+1}+\sqrt{n+2}\right).
\]

\noindent\textbf{解}\quad
\[
\begin{aligned}
S_n={}&(1-2\sqrt2+\sqrt3)+(\sqrt2-2\sqrt3+\sqrt4)\\
&+(\sqrt3-2\sqrt4+\sqrt5)+(\sqrt4-2\sqrt5+\sqrt6)+\cdots\\
&+\left(\sqrt{n-2}-2\sqrt{n-1}+\sqrt n\right)
 +\left(\sqrt{n-1}-2\sqrt n+\sqrt{n+1}\right)\\
&+\left(\sqrt n-2\sqrt{n+1}+\sqrt{n+2}\right)\\
={}&1-\sqrt2-\sqrt{n+1}+\sqrt{n+2}
=1-\sqrt2+\frac1{\sqrt{n+1}+\sqrt{n+2}}\\
&\to 1-\sqrt2\quad(\text{当 }n\to\infty\text{ 时}).
\end{aligned}
\]

\noindent\textbf{练习}\quad 计算
\[
  \sum_{n=1}^\infty \frac1{n(n+1)(n+2)}.\qquad\left(\frac14\right)
\]

\noindent\textbf{提示}\quad （参看例 4.5.6 里的拆分法）
\[
  S_n=\frac12\sum_{k=1}^n\left(\frac1k-\frac2{k+1}+\frac1{k+2}\right).
\]

\noindent\textbf{例 5.1.6}\quad 已知级数 $\displaystyle\sum_{n=1}^\infty a_n=A$，$\displaystyle\lim_{n\to\infty}na_n=0$，求级数之和
\[
  \sum_{n=1}^\infty n(a_n-a_{n+1}).
\]
（华东师范大学）

\noindent\textbf{解}\quad
\[
\begin{aligned}
\sum_{k=1}^n k(a_k-a_{k+1})
&=(a_1-a_2)+2(a_2-a_3)+3(a_3-a_4)+\cdots\\
&\quad+(n-1)(a_{n-1}-a_n)+n(a_n-a_{n+1})\\
&=\sum_{i=1}^n a_i-na_{n+1}\to A+0=A\quad(n\to\infty).
\end{aligned}
\]

\noindent\textbf{c.\quad 方程式法}

\noindent\textbf{要点}\quad 建立 $S_n$ 的方程式，从而求出 $S_n$.

\noindent\textbf{例 5.1.7}\quad 计算
\[
  q\cos\alpha+q^2\cos2\alpha+\cdots+q^n\cos n\alpha+\cdots\qquad(|q|<1).
\]

\noindent\textbf{解}\quad 记
\[
  S_n=q\cos\alpha+q^2\cos2\alpha+\cdots+q^n\cos n\alpha
  =\sum_{k=1}^n q^k\cos k\alpha.
\]
两边同乘 $2q\cos\alpha$，得
\[
\begin{aligned}
2q\cos\alpha\cdot S_n
&=\sum_{k=1}^n 2q^{k+1}\cos\alpha\cos k\alpha\\
&=\sum_{k=1}^n q^{k+1}[\cos(k+1)\alpha+\cos(k-1)\alpha],
\end{aligned}
\]
即
\[
  2q\cos\alpha\cdot S_n
  =\left(q^{n+1}\cos(n+1)\alpha+S_n-q\cos\alpha\right)
   +\left(q^2+q^2S_n-q^{n+2}\cos n\alpha\right),
\]
解此方程便得
\[
  S_n=\frac{q^{n+2}\cos n\alpha-q^{n+1}\cos(n+1)\alpha+q\cos\alpha-q^2}
  {1+q^2-2q\cos\alpha}
  \to \frac{q\cos\alpha-q^2}{1+q^2-2q\cos\alpha}\quad(n\to\infty).
\]

\noindent\textbf{注}\quad 本例亦可由如下复数和式取实部得到：
\[
  \sum_{k=0}^n z^k=\frac{1-z^{n+1}}{1-z}.
\]

\noindent\textbf{d.\quad 利用子列的极限}

\noindent\textbf{要点}\quad 我们知道，若 $\{S_{2n}\}$ 与 $\{S_{2n+1}\}$ 有相同极限 $S$，则 $\lim\limits_{n\to\infty}S_n=S$. 因此对于级数 $\displaystyle\sum_{n=1}^\infty a_n$，若通项 $a_n\to0$（当 $n\to\infty$ 时），则部分和的子列 $\{S_{2n}\}$ 收敛于 $S$，意味着 $\{S_{2n+1}\}$ 也收敛于 $S$，从而 $\displaystyle\sum_{n=1}^\infty a_n=S$. 我们把 $\{S_{2n}\}$ 与 $\{S_{2n+1}\}$ 称为互补子列. 这个原理可推广到一般情形：若 $\displaystyle\sum_{n=1}^\infty a_n$ 的通项 $a_n\to0$（当 $n\to\infty$ 时），$\{S_n\}$ 的子列 $\{S_{pn}\}_{n=1}^\infty\to S$（$p$ 是某个正整数），则 $\displaystyle\sum_{n=1}^\infty a_n=S$. 我们把这种方法称为子列方法.

\noindent\textbf{$\star$例 5.1.8}\quad 计算
\[
  1+\frac12+\left(\frac13-1\right)+\frac14+\frac15+\left(\frac16-\frac12\right)
  +\frac17+\frac18+\left(\frac19-\frac13\right)+\cdots.
\]

\noindent\textbf{解}\quad 此级数通项趋向零，因此只要求 $S_{3n}$ 的极限，注意公式
\[
  1+\frac12+\frac13+\cdots+\frac1n=C+\ln n+\varepsilon_n,
\]
其中 $C$ 为 Euler 常数（见例 1.2.11），$\varepsilon_n\to0$（当 $n\to\infty$ 时）. 因此，对原级数，
\[
\begin{aligned}
S_{3n}
&=1+\frac12+\frac13+\cdots+\frac1{3n}-1-\frac12-\cdots-\frac1n\\
&=\ln(3n)-\ln n+\varepsilon_{3n}-\varepsilon_n\to\ln3\quad(n\to\infty).
\end{aligned}
\]
故原级数和
\[
  S=\ln3.
\]

\noindent\textbf{练习}\quad 求
\[
  1-\frac12-\frac14+\frac13-\frac16-\frac18+\cdots+\frac1{2n-1}
  -\frac1{4n-2}-\frac1{4n}+\cdots
\]
之和.（中国科学技术大学）

\noindent\textbf{解}\quad 原式
\[
\begin{aligned}
&=\sum_{n=1}^\infty\left(\frac1{2n-1}-\frac1{4n-2}-\frac1{4n}\right)
=\sum_{n=1}^\infty\left(\frac1{4n-2}-\frac1{4n}\right)\\
&=\frac12\lim_{n\to\infty}\sum_{k=1}^n\left(\frac1{2k-1}-\frac1{2k}\right)
=\frac12\lim_{n\to\infty}\sum_{k=1}^n\left[\left(\frac1{2k-1}+\frac1{2k}\right)-2\cdot\frac1{2k}\right]\\
&=\frac12\lim_{n\to\infty}\left(\sum_{k=1}^{2n}\frac1k-\sum_{k=1}^n\frac1k\right)\quad(\text{再利用 Euler 常数})\\
&=\frac12\lim_{n\to\infty}\left[(\ln(2n)+C+\varepsilon_{2n})-(\ln n+C+\varepsilon_n)\right]
=\frac12\ln2.
\end{aligned}
\]

\noindent\textbf{*例 5.1.9}\quad 将级数
\[
  1-\frac12+\frac13-\frac14+\frac15-\cdots
\]
的各项重新安排，使先依次出现 $p$ 个正项，再出现 $q$ 个负项，然后如此交替，试证新级数的和为
\[
  \ln2+\frac12\ln\frac pq.
\]

\noindent\textbf{证}\quad 因为通项趋向零，根据上述子列求和法，对新级数我们只要求子列 $\{S_{(p+q)n}\}_{n=1}^\infty$ 的极限也就够了. 新级数前 $(p+q)n$ 项的和
\[
\begin{aligned}
S_{(p+q)n}
={}&1+\frac13+\cdots+\frac1{2p-1}-\frac12-\frac14-\cdots-\frac1{2q}\\
&+\frac1{2p+1}+\frac1{2p+3}+\cdots+\frac1{4p-1}
  -\frac1{2q+2}-\frac1{2q+4}-\cdots-\frac1{4q}\\
&+\cdots+\frac1{2np-(2p-1)}+\frac1{2np-(2p-3)}+\cdots+\frac1{2np-1}\\
&-\frac1{2nq-(2q-2)}-\frac1{2nq-(2q-4)}-\cdots-\frac1{2nq}\\
={}&1+\frac13+\frac15+\cdots+\frac1{2np-1}
  -\frac12-\frac14-\cdots-\frac1{2nq}\\
={}&1+\frac12+\frac13+\frac14+\frac15+\cdots+\frac1{2np}
  -\left(\frac12+\frac14+\cdots+\frac1{2np}\right)
  -\frac12-\frac14-\cdots-\frac1{2nq}\\
={}&1+\frac12+\cdots+\frac1{2np}
  -\frac12\left(1+\frac12+\cdots+\frac1{np}\right)
  -\frac12\left(1+\frac12+\cdots+\frac1{nq}\right).
\end{aligned}
\]
注意对于调和级数，有公式
\[
  1+\frac12+\cdots+\frac1n=C+\ln n+\varepsilon_n,
\]
其中 $C$ 为 Euler 常数，$\varepsilon_n\to0$（当 $n\to\infty$ 时），故
\[
\begin{aligned}
S_{(p+q)n}
&=C+\ln(2np)+\varepsilon_{2np}
 -\frac12[C+\ln(np)+\varepsilon_{np}]
 -\frac12[C+\ln(nq)+\varepsilon_{nq}]\\
&\to \ln2+\frac12\ln\frac pq\quad(n\to\infty).
\end{aligned}
\]

\noindent\textbf{e.\quad 先求 $S_n'(x)$ 的紧缩形式}

\noindent\textbf{$\star$例 5.1.10}\quad 设 $x\in[0,\pi]$，试求级数
\[
  \sum_{n=1}^\infty \frac{\sin nx}{n}
\]
的和函数.

\noindent\textbf{解}\quad 若 $x=0$ 或 $\pi$，显然级数和为 $0$.

现设 $0<x\le\pi$. 记
\[
  S_n(x)=\sum_{k=1}^n\frac{\sin kx}{k},
\]
则
\[
\begin{aligned}
S_n'(x)
&=\left(\sum_{k=1}^n\frac{\sin kx}{k}\right)'
=\sum_{k=1}^n\cos kx
=\frac1{2\sin\dfrac x2}\sum_{k=1}^n 2\sin\frac x2\cos kx\\
&=\frac1{2\sin\dfrac x2}\sum_{k=1}^n
\left[\sin\left(k+\frac12\right)x-\sin\left(k-\frac12\right)x\right]\\
&=\frac1{2\sin\dfrac x2}\left(\sin\frac{2n+1}2x-\sin\frac x2\right)
=\frac{\sin\left(n+\frac12\right)x}{2\sin\dfrac x2}-\frac12.
\end{aligned}
\]
于是
\[
\begin{aligned}
S_n(x)=S_n(x)-S_n(\pi)
&=-\int_x^\pi S_n'(t)\,dt\\
&=-\frac12\int_x^\pi\frac1{\sin\dfrac t2}\sin\left(n+\frac12\right)t\,dt
  +\frac12(\pi-x).
\end{aligned}
\]
利用 Riemann 引理，$n\to\infty$ 时上式第一项趋向零. 所以，级数和
\[
  S(x)=
  \begin{cases}
  0, & x=0,\ x=\pi,\\[2mm]
  \dfrac12(\pi-x), & 0<x<\pi.
  \end{cases}
\]

\subsection*{$\star$二、级数敛散性的判断}

\noindent\textbf{a.\quad Cauchy 准则及其应用}

\noindent\textbf{要点}\quad 1）Cauchy 准则. 级数 $\displaystyle\sum_{n=1}^\infty a_n$ 收敛的充要条件是：$\forall\varepsilon>0$，$\exists N>0$，当 $n>N$ 时，
\[
  \left|\sum_{k=n+1}^{n+p}a_k\right|<\varepsilon\qquad(\forall p\in\mathbb{N}).
\]
值得注意的是，此条件意味着
\[
  \sum_{k=n+1}^{n+p}a_k\to0\quad(\text{当 }n\to\infty\text{ 时关于 }p\in\mathbb{N}\text{ 一致})
\]
（$\mathbb{N}$ 是自然数的集合）；而不只是 $\forall p$，有 $\displaystyle\sum_{k=n+1}^{n+p}a_k\to0(n\to\infty)$（见例 5.1.13）.

2）Cauchy 准则的否定形式. 级数 $\displaystyle\sum_{n=1}^\infty a_n$ 发散的充要条件是：$\exists\varepsilon_0>0$，$\forall N>0$，$\exists n>N$ 及某自然数 $p$，使得
\[
  \left|\sum_{k=n+1}^{n+p}a_k\right|\ge\varepsilon_0.
\]

\noindent\textbf{例 5.1.11}\quad 证明级数 $\displaystyle\sum_{n=1}^\infty\frac1n$ 发散.

\noindent\textbf{证}\quad 取 $\varepsilon_0=\dfrac12>0$，则 $\forall n\in\mathbb{N}$，取 $p=n$ 时，恒有
\[
  \left|\sum_{k=n+1}^{2n}\frac1k\right|
  \ge \frac n{2n}=\frac12=\varepsilon_0>0,
\]
故 $\displaystyle\sum_{n=1}^\infty\frac1n$ 发散.

\noindent\textbf{$\star$例 5.1.12}\quad 设 $a_n>0$，$S_n=a_1+a_2+\cdots+a_n$，级数 $\displaystyle\sum_{n=1}^\infty a_n=\infty$. 试证 $\displaystyle\sum_{n=1}^\infty\frac{a_n}{S_n}$ 发散.（武汉大学）

\noindent\textbf{证}\quad 因 $a_n>0$，$S_n\uparrow$，所以
\[
  \sum_{k=n+1}^{n+p}\frac{a_k}{S_k}
  \ge \frac{\sum_{k=n+1}^{n+p}a_k}{S_{n+p}}
  =\frac{S_{n+p}-S_n}{S_{n+p}}
  =1-\frac{S_n}{S_{n+p}}.
\]
因为 $S_n\to+\infty$，故 $\forall n$，当 $p\in\mathbb{N}$ 充分大时，有 $\dfrac{S_n}{S_{n+p}}<\dfrac12$，从而
\[
  \sum_{k=n+1}^{n+p}\frac{a_k}{S_k}\ge 1-\frac12=\frac12.
\]
所以 $\displaystyle\sum_{n=1}^\infty\frac{a_n}{S_n}$ 发散.

\noindent\textbf{$\star$例 5.1.13}\quad 如果
\[
  \lim_{n\to\infty}a_{n+1}=0,\quad
  \lim_{n\to\infty}(a_{n+1}+a_{n+2})=0,\quad\cdots,\quad
  \lim_{n\to\infty}(a_{n+1}+a_{n+2}+\cdots+a_{n+p})=0,
\]
试问级数 $\displaystyle\sum_{n=1}^\infty a_n$ 是否一定收敛？（“是”或“不一定”，要说明理由.）（华中科技大学）

\noindent\textbf{解}\quad 不一定. 例如上面例 5.1.11：
\[
  \sum_{n=1}^\infty\frac1n,
\]
虽然 $\forall p\in\mathbb{N}$，
\[
  0<\frac1{n+1}+\cdots+\frac1{n+p}<\frac p{n+1}\to0\quad(n\to\infty),
\]
但 $\displaystyle\sum_{n=1}^\infty\frac1n$ 发散.

\noindent\textbf{$\star$例 5.1.14}\quad 证明：级数 $\displaystyle\sum_{n=1}^\infty a_n$ 收敛的充分必要条件是：对于任意的正整数序列 $p_1,p_2,\cdots,p_k,\cdots$ 及自然数的任意子列 $\{n_k\}$，皆有
\[
  \lim_{k\to\infty}(a_{n_k+1}+a_{n_k+2}+\cdots+a_{n_k+p_k})=0.
\]
（中山大学）

\noindent\textbf{证}\quad $1^\circ$ 必要性. 因为 $\displaystyle\sum_{n=1}^\infty a_n$ 收敛，所以 $\forall\varepsilon>0$，$\exists N>0$，当 $n>N$ 时，
\[
  \left|\sum_{k=n+1}^{n+p}a_k\right|<\varepsilon\qquad(\forall p\in\mathbb{N})
\]
成立. 由 $n_k\ge k$ 知，当 $k>N$ 时，有
\[
  |a_{n_k+1}+\cdots+a_{n_k+p_k}|<\varepsilon.
\]
故
\[
  \lim_{k\to\infty}(a_{n_k+1}+\cdots+a_{n_k+p_k})=0.
\]

$2^\circ$ 充分性.（反证法）若 $\displaystyle\sum_{n=1}^\infty a_n$ 发散，则 $\exists\varepsilon_0>0$，$\forall N>0$，$\exists n>N$ 及 $p\in\mathbb{N}$ 使得
\[
  |a_{n+1}+\cdots+a_{n+p}|\ge\varepsilon_0.
\]
特别地，对 $N_1=1$，$\exists n_1>1,p_1\in\mathbb{N}$ 使得
\[
  |a_{n_1+1}+\cdots+a_{n_1+p_1}|\ge\varepsilon_0;
\]
对 $N_2=\max\{n_1,2\}$，$\exists n_2>N_2,p_2\in\mathbb{N}$，使得
\[
  |a_{n_2+1}+\cdots+a_{n_2+p_2}|\ge\varepsilon_0.
\]
如此，我们得到自然数的子列 $\{n_k\}$ 以及 $\{p_k\}$，使得恒有
\[
  |a_{n_k+1}+\cdots+a_{n_k+p_k}|\ge\varepsilon_0\quad(k=1,2,\cdots),
\]
与已知条件矛盾. 证毕.

值得注意的是，Cauchy 准则不仅能用于级数敛散性的判别，还可导出收敛级数的其他性质，如例 5.1.38 和例 5.1.39.

\noindent\textbf{b.\quad 正项级数敛散性的判定}

\noindent\textbf{要点}\quad 判断级数 $\sum a_n$ 的敛散性，通常有如下方法：

1）若通项 $a_n\not\to0$（当 $n\to\infty$ 时），则 $\sum a_n$ 发散.

2）判阶法：如果 $a_n\to0$（当 $n\to+\infty$ 时），并且相对 $\dfrac1n$ 来讲，它是 $p$ 阶的无穷小量，那么当 $p>1$ 时，级数 $\sum a_n$ 收敛；当 $p\le1$ 时，$\sum a_n$ 发散.

3）D'Alembert 判别法（亦称比式判别法）：对正项级数 $\displaystyle\sum_{n=1}^\infty a_n$，

若 $\exists q>0$ 和 $N>0$，使得 $\forall n>N$，有 $\dfrac{a_{n+1}}{a_n}\le q<1$，则级数 $\displaystyle\sum_{n=1}^\infty a_n$ 收敛；

若 $\exists N>0$，$n>N$ 时，恒有 $\dfrac{a_{n+1}}{a_n}\ge1$，则级数 $\displaystyle\sum_{n=1}^\infty a_n$ 发散.

特别地，若 $\displaystyle\lim_{n\to\infty}\frac{a_{n+1}}{a_n}=\ell$，则当 $\ell<1$ 时，正项级数 $\displaystyle\sum_{n=1}^\infty a_n$ 收敛；当 $\ell>1$ 时，$\displaystyle\sum_{n=1}^\infty a_n$ 发散.

或：若 $\displaystyle\overline{\lim_{n\to\infty}}\frac{a_{n+1}}{a_n}=\ell$，则当 $\ell<1$ 时，正项级数 $\displaystyle\sum_{n=1}^\infty a_n$ 收敛；若 $\displaystyle\underline{\lim_{n\to\infty}}\frac{a_{n+1}}{a_n}=\ell>1$，则 $\displaystyle\sum_{n=1}^\infty a_n$ 发散.

4）根式判别法（也称 Cauchy 判别法）：若 $a_n\ge0$，$\exists q,N>0$，$\forall n>N$，有 $\sqrt[n]{a_n}\le q<1$，则级数 $\displaystyle\sum_{n=1}^\infty a_n$ 收敛；

若 $\exists N>0$，$\forall n>N$，有 $\sqrt[n]{a_n}\ge1$，则级数 $\displaystyle\sum_{n=1}^\infty a_n$ 发散.

特别地，若 $\displaystyle\lim_{n\to\infty}\sqrt[n]{a_n}=\ell$，则当 $\ell<1$ 时，$\displaystyle\sum_{n=1}^\infty a_n$ 收敛；当 $\ell>1$ 时，$\displaystyle\sum_{n=1}^\infty a_n$ 发散.

或：若 $\displaystyle\overline{\lim_{n\to\infty}}\sqrt[n]{a_n}=\ell$，则当 $\ell<1$ 时，正项级数 $\displaystyle\sum_{n=1}^\infty a_n$ 收敛；若 $\displaystyle\underline{\lim_{n\to\infty}}\sqrt[n]{a_n}=\ell>1$，则 $\displaystyle\sum_{n=1}^\infty a_n$ 发散.

\noindent\textbf{注}\quad 1）凡是用比式判别法能判别的，用根式判别法也一定能判别.

2）比式判别法和根式判别法都是基于跟等比级数比较. 凡是递减速度不低于某等比 $(0<q<1)$ 级数的级数，必然收敛. 虽然使用方便，但适用面较窄. 技巧性较大，适用面较宽的是比较判别法.

5）比较判别法：$\displaystyle\sum_{n=1}^\infty a_n$ 和 $\displaystyle\sum_{n=1}^\infty b_n$ 是两正项级数，若从某项开始恒有 $a_n\le b_n$，若 $\displaystyle\sum_{n=1}^\infty b_n$ 收敛，则 $\displaystyle\sum_{n=1}^\infty a_n$ 也收敛；反之，若 $\displaystyle\sum_{n=1}^\infty a_n$ 发散，则 $\displaystyle\sum_{n=1}^\infty b_n$ 也发散.

使用此法主要用缩放法：给定一个级数 $\displaystyle\sum_{n=1}^\infty a_n$，如果通项 $a_n$ 十分复杂，欲证级数收敛，应将 $a_n$ 化简放大为 $b_n$，使得 $\displaystyle\sum_{n=1}^\infty b_n$ 收敛；欲证级数 $\displaystyle\sum_{n=1}^\infty a_n$ 发散，应将 $a_n$ 化简缩小为 $c_n$，使得 $\displaystyle\sum_{n=1}^\infty c_n$ 发散.

\noindent\textbf{注}\quad 特别，若正项级数 $\displaystyle\sum_{n=1}^\infty a_n$ 与 $\displaystyle\sum_{n=1}^\infty b_n$ 满足
\[
  \lim_{n\to\infty}\frac{a_n}{b_n}=\ell\quad(0<\ell<+\infty),
\]
则 $\displaystyle\sum_{n=1}^\infty a_n$ 与 $\displaystyle\sum_{n=1}^\infty b_n$ 同时敛散.

6）Cauchy 积分判别法：若 $[1,+\infty)$ 上，$f(x)\downarrow$，且 $f(x)\ge0$，则级数 $\displaystyle\sum_{n=1}^\infty f(n)$ 与
\[
  \int_1^{+\infty} f(x)\,dx
\]
同时敛散.

7）考虑部分和 $\displaystyle\sum_{k=1}^n a_k$ 是否关于 $n$ 有界. 有界则收敛，无界则发散.

\noindent\textbf{利用判阶法及比较判别法}

\noindent\textbf{$\star$例 5.1.15}\quad 若
\[
  \lim_{n\to\infty}\left(n^{2\sin\frac1n}a_n\right)=1,
\]
判断级数 $\displaystyle\sum_{n=1}^\infty a_n$ 是否收敛？试证之.（上海交通大学）

\noindent\textbf{解}\quad 已知
\[
  \frac{a_n}{n^{-2\sin\frac1n}}=n^{2\sin\frac1n}a_n\to1\quad(n\to\infty),
\]
且
\[
  0\le n^{-2\sin\frac1n}=\left(\frac1{n^2}\right)^{\frac{\sin(1/n)}{1/n}}
  \le \left(\frac1{n^2}\right)^{3/4}\quad(\text{当 }n\text{ 充分大时}),
\]
所以 $n^{-2\sin\frac1n}$ 为无穷小量，$a_n$ 与 $n^{-2\sin\frac1n}$ 为等价无穷小量，故 $\sum a_n$ 与 $\sum n^{-2\sin\frac1n}$ 同时敛散. 另由 $\sum\dfrac1{n^{3/2}}$ 收敛，知 $\sum n^{-2\sin\frac1n}$ 收敛，从而级数 $\sum a_n$ 收敛.

\noindent\textbf{$\star$例 5.1.16}\quad 设
\[
  a_n=\left(1-\frac{p\ln n}{n}\right)^n,
\]
讨论 $\sum a_n$ 的敛散性.

\noindent\textbf{分析}\quad
\[
  a_n=e^{\ln\left(1-\frac{p\ln n}{n}\right)^n}
  =e^{n\ln\left(1-\frac{p\ln n}{n}\right)},
\]
而当 $n\to\infty$ 时，$\dfrac{\ln n}{n}\to0$，因此
\[
  \ln\left(1-\frac{p\ln n}{n}\right)\sim-\frac{p\ln n}{n},
\]
从而可以设想 $a_n\sim e^{n\left(-\frac{p\ln n}{n}\right)}=n^{-p}$.

\noindent\textbf{证 I}\quad
\[
\begin{aligned}
\lim_{n\to\infty}\ln(n^p a_n)
&=\lim_{n\to\infty}\ln\left[n^p\left(1-\frac{p\ln n}{n}\right)^n\right]\\
&=\lim_{n\to\infty}\left[p\ln n+n\ln\left(1-\frac{p\ln n}{n}\right)\right]\\
&=\lim_{n\to\infty}\frac1{1/n}\left[-\frac1n p\ln\frac1n+\ln\left(1+\frac pn\ln\frac1n\right)\right]\\
&=\lim_{x\to0^+}\frac1x[\ln(1+px\ln x)-px\ln x]\quad(\text{应用 L'Hospital 法则})\\
&=\lim_{x\to0^+}\frac{-p^2(x\ln x+x\ln x)}{1+px\ln x}=0.
\end{aligned}
\]
故 $\displaystyle\lim_{n\to\infty}n^p a_n=1$，$a_n\sim n^{-p}$，所以级数当 $p>1$ 时收敛，当 $p\le1$ 时发散.

\noindent\textbf{证 II}\quad 利用带 Peano 余项的 Taylor 公式：
\[
  \ln(1+x)=x-\frac12x^2+o(x^2)\quad(x\to0).
\]
\[
  a_n=e^{n\ln\left(1-\frac{p\ln n}{n}\right)}
  =e^{n\left[-\frac{p\ln n}{n}+o\left(\left(\frac{\ln n}{n}\right)^{3/2}\right)\right]}
  =n^{-p}\cdot e^{o\left(\frac{(\ln n)^{3/2}}{n^{1/2}}\right)}
  \sim n^{-p}\quad(n\to\infty),
\]
因此 $\sum a_n$ 当且仅当 $p>1$ 时收敛.

\noindent\textbf{例 5.1.17}\quad 证明：

1）级数 $\displaystyle\sum_{n=1}^\infty \frac1{(\ln n)^{\ln n}}$ 收敛；（武汉大学）

2）级数 $\displaystyle\sum_{n=1}^\infty \frac1{n^{1+1/n}}$ 发散.

\noindent\textbf{证}\quad 1）$n$ 充分大时，
\[
  \frac1{(\ln n)^{\ln n}}\le
  \frac1{(e^{1+\alpha})^{\ln n}}=\frac1{n^{1+\alpha}}\quad(\alpha>0).
\]
由比较判别法，结论自明.

2）（利用不等式：几何平均 $\le$ 算术平均，和“对数不等式”：
\[
  \frac{x}{x+1}\le\ln(1+x)\le x\quad(x>-1)
\]
（见例 1.1.8）.）
\[
  n^{1/n}=\sqrt[n]{n}=\sqrt[n]{n\cdot1\cdots1}
  \le \frac{n+(n-1)}n=2\cdot\frac{2n-1}{2n}
  \le2\cdot\frac{2n}{2n+1}\le2\ln(1+2n)\quad(n>1).
\]
于是
\[
  n^{1+\frac1n}\le2n\ln(1+2n),
\qquad
  \frac1{n^{1+\frac1n}}\ge\frac1{2n\ln(1+2n)}.
\qquad(1)
\]
由
\[
  \int_2^{+\infty}\frac1{(1+2x)\ln(1+2x)}\,dx
\]
发散，且
\[
  0<\frac1{(1+2x)\ln(1+2x)}
  <\frac{1+2x}{2x}\frac1{(1+2x)\ln(1+2x)},
\]
知
\[
  \int_2^{+\infty}\frac1{2x\ln(1+2x)}\,dx
\]
发散. 再由 Cauchy 积分判别法，$\displaystyle\sum_{n=2}^\infty\frac1{2n\ln(1+2n)}$ 发散.

于是，利用比较判别法，由不等式 (1) 知：
\[
  \sum_{n=1}^\infty\frac1{n^{1+\frac1n}}
\]
发散.

\noindent\textbf{练习 1}\quad 设 $\{a_n\},\{b_n\}$ 是各项均为正的数列，满足：
\[
  \lim_{n\to\infty}\frac{b_n}{n}=0
\qquad(1)
\]
及
\[
  \lim_{n\to\infty} b_n\left(\frac{a_n}{a_{n+1}}-1\right)=\lambda>0.
\qquad(2)
\]
求证：级数 $\displaystyle\sum_{n=1}^\infty a_n$ 收敛.（中国科学技术大学）

\noindent\textbf{提示}\quad 利用条件 (1) 和 (2)，
\[
  \lim_{n\to\infty}n\left(\frac{a_n}{a_{n+1}}-1\right)
  =\lim_{n\to\infty}
  \frac{b_n\left(\dfrac{a_n}{a_{n+1}}-1\right)}{\dfrac{b_n}{n}}
  =+\infty.
\]
因此有
\[
  \frac{a_n}{a_{n-1}}\le\left(\frac{n-1}{n}\right)^2,\qquad
  \frac{a_{n+1}}{a_n}\le\left(\frac n{n+1}\right)^2.
\]

\noindent\textbf{证}\quad 将 $N$ 固定，则 $\forall n>N$，有
\[
\begin{aligned}
a_n
&=\frac{a_n}{a_{n-1}}\cdot\frac{a_{n-1}}{a_{n-2}}\cdots
  \frac{a_{N+2}}{a_{N+1}}\cdot\frac{a_{N+1}}{a_N}\cdot a_N\\
&\le\left(\frac{n-1}{n}\cdot\frac{n-2}{n-1}\cdots
  \frac{N+1}{N+2}\cdot\frac N{N+1}\right)^2a_N
=\frac{N^2a_N}{n^2}.
\end{aligned}
\]
（利用比较判别法）由 $\dfrac{a_N}{N^2}\sum_{n=N+1}^\infty\dfrac1{n^2}$ 收敛知 $\displaystyle\sum_{n=N+1}^\infty a_n$ 收敛，从而 $\displaystyle\sum_{n=1}^\infty a_n$ 也收敛.

\noindent\textbf{练习 2}\quad 证明：级数
\[
  \sum_{n=1}^\infty\frac{\sqrt[n]{n}-1}{n^\alpha}
\]
当 $\alpha>0$ 时收敛，当 $\alpha\le0$ 时发散.（南开大学）

\noindent\textbf{提示}\quad （当 $n\to\infty$ 时）作为无穷小量：
\[
  \sqrt[n]{n}-1
  \quad\text{与}\quad
  \ln[1+(\sqrt[n]{n}-1)]=\frac{\ln n}{n}
\]
等价，从而 $\dfrac{\sqrt[n]{n}-1}{n^\alpha}$ 与 $\dfrac{\ln n}{n^{\alpha+1}}$ 等价.

\noindent\textbf{证}\quad $1^\circ$ 当 $\alpha>0$ 时，
\[
  \lim_{n\to\infty}\frac{\ln n}{n^{\alpha/2}}
  =\lim_{x\to\infty}\frac{\ln x}{x^{\alpha/2}}
  \stackrel{\text{L'Hospital 法则}}{=}
  \lim_{x\to\infty}\frac1x\cdot\frac2\alpha\cdot\frac1{x^{\alpha/2-1}}=0.
\qquad(1)
\]
因此
\[
  \lim_{n\to\infty}
  \frac{\dfrac{\sqrt[n]{n}-1}{n^\alpha}}{\dfrac1{n^{\alpha/2+1}}}
  =\lim_{n\to\infty}\frac{\ln n}{n^{\alpha/2}}
  \stackrel{\text{式 }(1)}{=}0.
\qquad(2)
\]
式 (2) 说明 $\dfrac{\sqrt[n]{n}-1}{n^\alpha}$ 是比 $\dfrac1{n^{\alpha/2+1}}$ 高阶的无穷小量. 而已知 $\displaystyle\sum_{n=1}^\infty\dfrac1{n^{\alpha/2+1}}$ 收敛，故 $\displaystyle\sum_{n=1}^\infty\dfrac{\sqrt[n]{n}-1}{n^\alpha}$ 收敛.

$2^\circ$ 当 $\alpha\le0$ 时，
\[
  \frac{\ln n}{n^{\alpha+1}}=\frac{\ln n}{n^\alpha}\cdot\frac1n\ge\frac1n,
\]
因 $\displaystyle\sum\frac1n$ 发散，故 $\displaystyle\sum_{n=1}^\infty\dfrac{\sqrt[n]{n}-1}{n^\alpha}$ 发散.

\noindent\textbf{练习 3}\quad 若正项级数 $\displaystyle\sum_{n=1}^\infty a_n$ 收敛，求证：

1）$\displaystyle\sum_{n=1}^\infty a_n^p$ 收敛 $(p>1)$；

2）$\displaystyle\sum_{n=1}^\infty\frac{\sqrt[k]{a_n}}n$ 收敛 $(2<k\in\mathbb{N})$.（南开大学）

\noindent\textbf{提示}\quad 1）$\displaystyle\sum_{n=1}^\infty a_n$ 收敛，可知 $a_n$ 为无穷小 $(n\to\infty)$，故对 $0<a_n<1$（$n$ 充分大），当 $p>1$ 时，有 $0<a_n^p<a_n$.

2）用 $k$ 个正数的均值不等式.

\noindent\textbf{再提示}\quad
\[
\begin{aligned}
\frac{\sqrt[k]{a_n}}n
&=\sqrt[k]{a_n\cdot
\frac1{n^{k/(k-1)}}\cdots\frac1{n^{k/(k-1)}}}\\
&\le\frac1k\left[a_n+\frac{k-1}{n^{k/(k-1)}}\right].
\end{aligned}
\]
而级数
\[
  \frac1k\sum_{n=1}^\infty\left[a_n+\frac{k-1}{n^{k/(k-1)}}\right]
\]
收敛.

\noindent\textbf{练习 4}\quad 判断级数
\[
  \sum_{n=1}^\infty\left[e-\left(1+\frac1{1!}+\frac1{2!}+\cdots+\frac1{n!}\right)\right]
\]
的敛散性.（南开大学）

\noindent\textbf{提示}\quad
\[
  0<e-\left(1+\frac1{1!}+\frac1{2!}+\cdots+\frac1{n!}\right)
  =R_n(x)\big|_{x=1}
\]
（$e^x$ 在 $x=1$ 处 Taylor 公式）
\[
  =\frac{e^\xi}{(n+1)!}\le\frac e{(n+1)!}\qquad(0<\xi<1),
\]
而 $\displaystyle\sum_{n=1}^\infty\dfrac e{(n+1)!}$ 收敛.

\noindent\textbf{例 5.1.18}\quad 设 $0<p_1<p_2<\cdots<p_n<\cdots$，求证：$\displaystyle\sum_{n=1}^\infty\frac1{p_n}$ 收敛的充要条件为如下级数收敛：
\[
  \sum_{n=1}^\infty\frac{n}{p_1+p_2+\cdots+p_n}.
\]

\noindent\textbf{提示}\quad 当 $n\ge2$ 时，
\[
  p_1+\cdots+p_n
  \ge p_{\left[\frac n2\right]}+p_{\left[\frac n2\right]+1}+\cdots+p_n
  \ge\left[\frac n2\right]p_{\left[\frac n2\right]}
  \ge\frac n4p_{\left[\frac n2\right]}>0,
\]
\[
  0\le\frac1{p_n}\le
  \frac{n}{p_1+p_2+\cdots+p_n}\le
  \frac4{p_{\left[\frac n2\right]}}.
\]
并注意
\[
  \sum_{n=2}^\infty\frac1{p_{\left[\frac n2\right]}}
  =\frac1{p_1}+\frac1{p_1}+\frac1{p_2}+\frac1{p_2}+\cdots.
\]

\noindent\textbf{例 5.1.19}\quad 设 $a_n>0\downarrow$，试证 $\displaystyle\sum_{n=1}^\infty a_n$ 与 $\displaystyle\sum_{n=0}^\infty 2^n a_{2^n}$ 同时敛散.

\noindent\textbf{证}\quad 因为对正项级数，任意加括号不改变敛散性，因此由
\[
\begin{aligned}
\sum_{n=1}^\infty a_n
&=a_1+(a_2+a_3)+(a_4+a_5+a_6+a_7)+(a_8+\cdots+a_{15})+\cdots\\
&\le a_1+2a_2+4a_4+8a_8+\cdots
=\sum_{n=0}^\infty 2^n a_{2^n}
\end{aligned}
\]
知，当级数 $\displaystyle\sum_{n=1}^\infty a_n$ 发散时，$\displaystyle\sum_{n=0}^\infty 2^n a_{2^n}$ 亦发散. 另外由
\[
\begin{aligned}
\sum_{n=1}^\infty a_n
&=a_1+a_2+(a_3+a_4)+(a_5+\cdots+a_8)+(a_9+\cdots+a_{16})+\cdots\\
&\ge a_1+a_2+2a_4+2^2a_{2^3}+2^3a_{2^4}+\cdots
=a_1+\frac12\sum_{n=1}^\infty2^n a_{2^n}
\end{aligned}
\]
知，当级数 $\displaystyle\sum_{n=1}^\infty a_n$ 收敛时，级数 $\displaystyle\sum_{n=0}^\infty2^n a_{2^n}$ 亦收敛. 总之两级数同时敛散.

\noindent\textbf{例 5.1.20}\quad 证明 Kummer 判别法：假设 $a_n>0,b_n>0(n=1,2,\cdots)$.

1）若 $\exists\alpha>0$，使得
\[
  \frac{b_n}{b_{n+1}}a_n-a_{n+1}\ge\alpha\qquad(n=1,2,\cdots),
\qquad(1)
\]
则级数 $\displaystyle\sum_{n=1}^\infty b_n$ 收敛；

2）若 $\displaystyle\sum_{n=1}^\infty\frac1{a_n}$ 发散，且
\[
  \frac{b_n}{b_{n+1}}a_n-a_{n+1}\le0\qquad(n=1,2,\cdots),
\qquad(2)
\]
则级数 $\displaystyle\sum_{n=1}^\infty b_n$ 发散.

\noindent\textbf{证}\quad 1）由式 (1) 知
\[
  b_na_n-b_{n+1}a_{n+1}\ge\alpha b_{n+1}>0,
\qquad(3)
\]
故 $b_na_n\downarrow$. 又因 $b_na_n>0$，所以 $\{b_na_n\}$ 收敛，从而级数
\[
  \sum_{n=1}^\infty(b_na_n-b_{n+1}a_{n+1})
\]
亦收敛. 再根据式 (3)，用比较判别法，知 $\displaystyle\sum_{n=1}^\infty b_n$ 亦收敛.

2）由式 (2) 知
\[
  \frac{1/a_{n+1}}{1/a_n}\le\frac{b_{n+1}}{b_n}\qquad(n=1,2,\cdots),
\]
故由 $\displaystyle\sum_{n=1}^\infty\frac1{a_n}$ 发散知 $\displaystyle\sum_{n=1}^\infty b_n$ 亦发散.

\noindent\textbf{$\star$例 5.1.21}\quad 若正项级数 $\displaystyle\sum_{n=1}^\infty a_n$ 收敛，且
\[
  e^{a_n}=a_n+e^{a_n+b_n}\qquad(n=1,2,\cdots),
\]
证明 $\displaystyle\sum_{n=1}^\infty b_n$ 收敛.（华东师范大学）

\noindent\textbf{提示}\quad （用比较判别法的极限形式.）
\[
  b_n=\ln(e^{a_n}-a_n)-a_n,\qquad
  \sum_{n=1}^\infty b_n=\sum_{n=1}^\infty\ln(e^{a_n}-a_n)-\sum_{n=1}^\infty a_n,
\]
\[
  \ln(e^{a_n}-a_n)=o(a_n).
\]

\noindent\textbf{例 5.1.22}\quad 研究级数 $\displaystyle\sum_{n=1}^\infty\dfrac1{x_n^2}$ 的敛散性，这里 $x_n$ 是方程 $x=\tan x$ 的正根，并且按递增的顺序编号.（国外赛题）

\noindent\textbf{提示}\quad
\[
  x_n\in\left(\frac\pi2+(n-1)\pi,\frac\pi2+n\pi\right),\qquad
  \frac1{x_n^2}\le\frac1{n^2}.
\]

\noindent\textbf{练习 1}\quad 设 $n$ 是一正整数，试证：

1）方程 $x^n+nx-1=0$ 在 $(-1,+\infty)$ 内有唯一正实根 $x_n$；

2）当 $\alpha>1$ 时，级数 $\displaystyle\sum_{n=1}^\infty x_n^\alpha$ 收敛.（中国科学院）

\noindent\textbf{提示}\quad 1）函数 $F(x)=x^n+nx-1$ 的导数 $F'(x)=n(x^{n-1}+1)>0$（当 $x>-1$ 时）. 而 $F(0)=-1$，$F\left(\dfrac1n\right)=\dfrac1{n^n}>0$.

2）如此 $0<x_n^\alpha\le\dfrac1{n^\alpha}$（$\forall n\in\mathbb{N}$），且 $\alpha>1$ 时，$\displaystyle\sum_{n=1}^\infty\dfrac1{n^\alpha}<+\infty$.

\noindent\textbf{练习 2}\quad 设 $\displaystyle\sum_{n=1}^\infty a_n$ 为收敛的正项级数. 求证：级数
\[
  \sum_{n=1}^\infty a_n^{1-\frac1n}
\]
收敛.（中国科学技术大学）

\noindent\textbf{提示}\quad 将 $a_n^{1-\frac1n}$ 变形，应用例 4.4.9 的不等式.

\noindent\textbf{证}\quad $\forall b>0$，$p,q>0$：$\dfrac1p+\dfrac1q=1$，
\[
  a_n^{1-\frac1n}=\frac1b(a_n^{\frac{n-1}n}\cdot b)
  \le\frac1b\left[\frac1p(a_n^{\frac{n-1}n})^p+\frac1q b^q\right]
\]
（利用例 4.4.9 的不等式）. 取 $p=\dfrac n{n-1}$，$b=\dfrac12$，则 $\dfrac1q=1-\dfrac1p=\dfrac1n<1$. 故
\[
  a_n^{1-\frac1n}<2a_n+\frac1{2^{n-1}}.
\]
因 $\displaystyle\sum_{n=1}^\infty\left(2a_n+\frac1{2^{n-1}}\right)$ 收敛，故 $\displaystyle\sum_{n=1}^\infty a_n^{1-\frac1n}$ 收敛.

\noindent\textbf{练习 3}\quad 已知 $\varphi(x)$ 是 $(-\infty,+\infty)$ 上周期为 $1$ 的连续函数，
\[
  \int_0^1\varphi(x)\,dx=0.
\qquad(1)
\]
设
\[
  a_n=\int_0^1 e^x\varphi(nx)\,dx\qquad(n=1,2,\cdots),
\qquad(2)
\]
求证：级数 $\displaystyle\sum_{n=1}^\infty a_n^2$ 收敛.（南开大学）

\noindent\textbf{提示}\quad 只需证明：$|a_n|\le\dfrac Mn$（$M$ 为常数）. 为此，引入
\[
  \Phi(x)=\int_0^x\varphi(t)\,dt.
\]

\noindent\textbf{证}\quad
\[
\begin{aligned}
a_n
&=\int_0^1 e^x\varphi(nx)\,dx
=\frac1n\int_0^1 e^x\,d\Phi(nx)\\
&=\frac1n e^x\Phi(nx)\bigg|_{x=0}^{x=1}
  -\frac1n\int_0^1 e^x\Phi(nx)\,dx.
\end{aligned}
\]
此项为 $0$. $\Phi(x)$ 连续、有界（记界为 $M_1>0$）：$|\Phi(x)|\le M_1$（$\forall x\in\mathbb{R}$），则
\[
  |a_n|\le\frac1n\int_0^1e^x|\Phi(nx)|\,dx
  \le\frac1nM_1(e-1),
\qquad
  a_n^2\le\frac1{n^2}M_1^2(e-1)^2.
\]
而 $\displaystyle\sum_{n=1}^\infty\frac1{n^2}M_1^2(e-1)^2$ 收敛，故 $\displaystyle\sum_{n=1}^\infty a_n^2$ 收敛.

\noindent\textbf{练习 4}\quad 设 $\{\alpha_n\}$ 及 $\{\delta_n\}$ 是两个（非负）无穷实数列，满足条件：
\[
  \text{(i)}\quad \alpha_{n+1}\le(1+\delta_n)\alpha_n+\delta_n\quad(n\ge1);
  \qquad
  \text{(ii)}\quad \sum_{n=1}^\infty\delta_n<+\infty.
\]
试证：

1）数列 $\left\{\displaystyle\prod_{i=1}^n(1+\delta_i)\right\}$ 收敛；

2）数列 $\{\alpha_n\}$ 有界；

3）数列 $\{\alpha_n\}$ 收敛.（中国科学院）

\noindent\textbf{分析}\quad 1）对 $\displaystyle\prod_{i=1}^n(1+\delta_i)$ 取对数，利用 $\ln(1+\delta_k)<\delta_k$，放大求上界；由 $\left\{\displaystyle\sum_{k=1}^n\ln(1+\delta_k)\right\}$ 递增有界得收敛.

2）将条件 (i) 的不等式改写成 $\alpha_{n+1}+1\le(1+\delta_n)(\alpha_n+1)$，进行迭代，推出 $\{\alpha_n\}$ 有界.

3）当 $n$ 增大时 $\sup_{k\ge n}|\alpha_k|$ 递减有下限，因此 $\displaystyle\lim_{n\to\infty}\sup_{k\ge n}|\alpha_k|$ 存在，记为 $\alpha$，可证：$\lim\alpha_n=\alpha$.

\noindent\textbf{证}\quad 1）因 $\delta_k\ge0$，$\displaystyle\sum_{k=1}^n\ln(1+\delta_k)$ 递增，且
\[
  \sum_{k=1}^n\ln(1+\delta_k)\le\sum_{k=1}^n\delta_k\le\sum_{n=1}^\infty\delta_n<+\infty,
\]
有上界. 故
\[
  \lim_{n\to\infty}\prod_{k=1}^n(1+\delta_k)
  =\lim_{n\to\infty}e^{\sum_{k=1}^n\ln(1+\delta_k)}
  =e^{\lim_{n\to\infty}\sum_{k=1}^n\ln(1+\delta_k)}
  \stackrel{\text{极限存在}}{=}\prod_{n=1}^\infty(1+\delta_n)<+\infty.
\qquad(1)
\]
2）由已知条件 (i) 可得
\[
  \alpha_{n+1}+1\le(1+\delta_n)(\alpha_n+1)
  \le(1+\delta_n)(1+\delta_{n-1})(\alpha_{n-1}+1)\le\cdots
\]
\[
  \le(\alpha_1+1)\prod_{k=1}^n(1+\delta_k)
  \stackrel{\text{式 }(1)}{\le}
  (\alpha_1+1)\prod_{n=1}^\infty(1+\delta_n)<+\infty\qquad(\forall n\in\mathbb{N}).
\]
因此 $\{\alpha_n\}$ 有上界：
\[
  \exists M>0:\ \forall n\in\mathbb{N},\ 0\le\alpha_n\le M.
\qquad(2)
\]
3）因 $\{\alpha_n\}$ 有上界，$\sup_{k\ge n}|\alpha_k|<+\infty$；又 $\sup_{k\ge n}|\alpha_k|$ 随 $n$ 增加而递减，而 $\alpha_k\ge0$，从而 $\sup_{k\ge n}|\alpha_k|\ge0$，递减且有下界，知
\[
  \lim_{n\to\infty}\sup_{k\ge n}|\alpha_k|
\]
存在，记为 $\alpha(\alpha\ge0)$.

记 $\beta_n=\sup_{k\ge n}|\alpha_k|$，那么 $\beta_n\downarrow\alpha$：$\forall\varepsilon>0$，$\exists N>0$，当 $n>N$ 时，有
\[
  0\le\alpha\le\beta_n<\alpha+\varepsilon.
\qquad(3)
\]
另一方面，因 $\displaystyle\sum_{n=1}^\infty\delta_n<+\infty$，$\forall\varepsilon>0$，$\exists N_1>0$，当 $n>N_1$ 时，有
\[
  \left(\sum_{k=n}^\infty\delta_k\right)(M+1)<\varepsilon.
\qquad(4)
\]
由条件 (i)：$\alpha_{n+1}\le(1+\delta_n)\alpha_n+\delta_n$ 可得
\[
  \alpha_{n+1}-\alpha_n\le\delta_n(\alpha_n+1)
  \stackrel{\text{式 }(2)}{\le}\delta_n(M+1),
\qquad(5)
\]
即
\[
  \alpha_{n+1}-\alpha_n\le\delta_n(M+1)\qquad(n=1,2,\cdots).
\]
当 $m>n>N_1$ 时，将不等式 (5) 从 $n$ 至 $(m-1)$ 相加，得
\[
  \alpha_m-\alpha_n\le\left(\sum_{k=n}^{m-1}\delta_k\right)(M+1)
  \le\left(\sum_{k=n}^{\infty}\delta_k\right)(M+1)
  \stackrel{\text{式 }(4)}{<}\varepsilon,
\]
亦即 $\alpha_m-\varepsilon\le\alpha_n$.

注意：当 $m$ 换为更大的整数时，此式保持成立，因此得 $\beta_m-\varepsilon\le\alpha_n$，再将 $m(>N_1)$ 固定，令 $m\to\infty$，取极限，得 $\alpha-\varepsilon\le\alpha_n$. 联系式 (1)，当 $n>\max\{N,N_1\}$ 时，有
\[
  \alpha-\varepsilon\le\alpha_n\le\sup_{k\ge n}|\alpha_k|=\beta_n<\alpha+\varepsilon.
\]
即 $|\alpha-\alpha_n|\le\varepsilon$，亦即 $\lim\alpha_n=\alpha$. 证毕.

\noindent\textbf{利用 D'Alembert 判别法}

\noindent\textbf{例 5.1.23}\quad 试证如下级数收敛：
\[
  \sqrt{2-\sqrt2}+\sqrt{2-\sqrt{2+\sqrt2}}
  +\sqrt{2-\sqrt{2+\sqrt{2+\sqrt2}}}+\cdots.
\]

\noindent\textbf{证}\quad 记
\[
  A_1=\sqrt2,\quad A_2=\sqrt{2+\sqrt2},\quad
  A_3=\sqrt{2+\sqrt{2+\sqrt2}},\cdots,
\]
则易知 $A_n\to2(n\to\infty)$，
\[
\begin{aligned}
\lim_{n\to\infty}\frac{a_{n+1}}{a_n}
&=\lim_{n\to\infty}
\frac{\sqrt{2-\sqrt{2+A_n}}}{\sqrt{2-A_n}}
=\lim_{x\to2}\frac{\sqrt{2-\sqrt{2+x}}}{\sqrt{2-x}}\\
&=\sqrt{\lim_{x\to2}\frac{2-\sqrt{2+x}}{2-x}}
\stackrel{\text{L'Hospital 法则}}{=}\frac12<1.
\end{aligned}
\]
利用 D'Alembert 判别法，知级数收敛.

\noindent\textbf{例 5.1.24}\quad 级数 $\displaystyle\sum_{n=1}^\infty u_n^{-1}$ 收敛吗？这里 $u_1=1,u_2=2,u_n=u_{n-2}+u_{n-1}(n\ge3)$.（国外赛题）

\noindent\textbf{提示}\quad （D'Alembert 判别法）由数学归纳法易得
\[
  \frac{u_{n+1}^{-1}}{u_n^{-1}}=\frac{u_n}{u_{n+1}}\le\frac23<1,
\]
故收敛.

\noindent\textbf{例 5.1.25}\quad 证明：若 $f(x)$ 为单调减少的正值函数，又设
\[
  \lim_{x\to+\infty}\frac{e^xf(e^x)}{f(x)}=\lambda,
\qquad(1)
\]
则 $\lambda<1$ 时级数 $\displaystyle\sum_{n=1}^\infty f(n)$ 收敛，$\lambda>1$ 时级数 $\displaystyle\sum_{n=1}^\infty f(n)$ 发散.

\noindent\textbf{分析}\quad 因 $f(x)>0,\downarrow$，根据 Cauchy 积分判别法，$\displaystyle\sum_{n=1}^\infty f(n)$ 与积分 $\displaystyle\int_1^{+\infty}f(x)\,dx$ 同时敛散. 要考察正值函数的反常积分 $\displaystyle\int_1^{+\infty}f(x)\,dx$ 的敛散性，只需要取一序列 $x_1<x_2<\cdots<x_n<\cdots$，使 $x_n\to+\infty$，看极限 $\displaystyle\lim_{n\to\infty}\int_1^{x_n}f(x)\,dx$ 是否存在.

\noindent\textbf{证}\quad 已知 $\displaystyle\lim_{x\to+\infty}\frac{e^xf(e^x)}{f(x)}=\lambda$.

$1^\circ$ 若 $\lambda>1$，则 $\exists A>1$，使得 $x>A$ 时，
\[
  \frac{e^xf(e^x)}{f(x)}>1,
\]
即 $e^xf(e^x)>f(x)$. 从而 $\forall x_{n-1}<x_n$，有
\[
  \int_{x_{n-1}}^{x_n} e^xf(e^x)\,dx>\int_{x_{n-1}}^{x_n}f(x)\,dx.
\qquad(2)
\]
（因 $f(x)$ 单调，积分有意义.）左端积分作变量替换，令 $e^x=t$，于是 (2) 成为
\[
  \int_{e^{x_{n-1}}}^{e^{x_n}}f(t)\,dt>\int_{x_{n-1}}^{x_n}f(x)\,dx.
\qquad(3)
\]
由此可见，若取序列 $\{x_n\}$ 如下：
\[
  x_1=1,\quad x_2=e,\quad x_3=e^{x_2},\cdots,\quad
  x_n=e^{x_{n-1}},\quad x_{n+1}=e^{x_n},\cdots,
\]
将积分变量 $t$ 仍写成 $x$，则式 (3) 可改写成
\[
  \int_{x_n}^{x_{n+1}}f(x)\,dx>\int_{x_{n-1}}^{x_n}f(x)\,dx
  \qquad(n=2,3,\cdots).
\qquad(4)
\]
于是
\[
  \int_1^{x_n}f(x)\,dx
  =\sum_{k=2}^n\int_{x_{k-1}}^{x_k}f(x)\,dx
  >(n-1)\int_1^e f(x)\,dx\to+\infty\quad(n\to\infty).
\]
所以 $\lambda>1$ 时，$\displaystyle\sum_{n=1}^\infty f(n)$ 发散.

$2^\circ$ 若 $\lambda<1$，取实数 $q:\lambda<q<1$. 由已知条件 (1)，对 $q$ 而言，$\exists A>1$，使得 $x>A$ 时，有
\[
  \frac{e^xf(e^x)}{f(x)}<q.
\]
即 $e^xf(e^x)<qf(x)$. 采用上面同样的方法进行推理可得
\[
\begin{aligned}
\int_1^{x_n}f(x)\,dx
&=\sum_{k=2}^n\int_{x_{k-1}}^{x_k}f(x)\,dx\\
&<\sum_{k=2}^n q^{k-1}\int_1^e f(x)\,dx
<\frac{\int_1^e f(x)\,dx}{1-q}<+\infty,
\end{aligned}
\]
故 $\displaystyle\int_1^{+\infty}f(x)\,dx$ 收敛，从而 $\displaystyle\sum_{n=1}^\infty f(n)$ 收敛.

\noindent\textbf{利用部分和有界}

\noindent\textbf{例 5.1.26}\quad 设 $\displaystyle\sum_{n=1}^\infty a_n$ 为正项级数，满足：
\[
  1)\quad \sum_{k=1}^n(a_k-a_n)\text{ 对 }n\text{ 有界};\qquad
  2)\quad a_n\downarrow0.
\]
试证级数 $\displaystyle\sum_{n=1}^\infty a_n$ 收敛.

\noindent\textbf{证}\quad 要证明正项级数 $\displaystyle\sum_{n=1}^\infty a_n$ 收敛，只要证明 $\exists M>0$，使得 $\forall n\in\mathbb{N}$ 有 $\displaystyle\sum_{k=1}^n a_k\le M$.

已知 $\displaystyle\sum_{k=1}^n(a_k-a_n)$ 有界，所以 $\exists M>0$，使得
\[
  \sum_{k=1}^n(a_k-a_n)\le M\qquad(\forall n\in\mathbb{N}).
\qquad(1)
\]
现任意固定一个 $n\in\mathbb{N}$，取 $m>n$，于是利用条件 2）及式 (1) 得
\[
  \sum_{k=1}^n a_k-na_m
  =\sum_{k=1}^n(a_k-a_m)
  \le\sum_{k=1}^m(a_k-a_m)\le M.
\qquad(2)
\]
此式对任意 $m>n$ 皆成立. 令 $m\to\infty$，因 $na_m\to0$，故式 (2) 成为 $\displaystyle\sum_{k=1}^n a_k\le M$. 由 $n$ 的任意性，知 $\displaystyle\sum_{k=1}^\infty a_k$ 收敛.

\noindent\textbf{$\star$例 5.1.27}\quad 若 $\displaystyle\sum_{n=1}^\infty a_n$ 收敛，且 $a_n>0$，则当 $p>\dfrac12$ 时，$\displaystyle\sum_{n=1}^\infty\frac{\sqrt{a_n}}{n^p}$ 收敛.（东北师范大学，郑州大学）

\noindent\textbf{提示}\quad 应用 Cauchy 不等式：
\[
  \sum_{k=1}^n\frac{\sqrt{a_k}}{k^p}
  \le\left(\sum_{k=1}^n a_k\right)^{1/2}
  \left(\sum_{k=1}^n\frac1{k^{2p}}\right)^{1/2}
\]
或不等式
\[
  \frac{\sqrt{a_n}}{n^p}\le\frac12\left(a_n+\frac1{n^{2p}}\right).
\]

\noindent\textbf{$\star$例 5.1.28}\quad 设 $\{a_n\}(n\ge1)$ 是正实数序列，证明：若级数 $\displaystyle\sum_{n=1}^\infty\dfrac1{a_n}$ 收敛，则级数
\[
  \sum_{n=1}^\infty
  \frac{n^2}{(a_1+a_2+\cdots+a_n)^2a_n}
\]
也收敛.（中南大学）

\noindent\textbf{证}\quad 我们希望证明部分和
\[
  S_n=\sum_{k=1}^n\frac{k^2}{(a_1+\cdots+a_k)^2a_k}
\]
有界. 记 $A_n=a_1+\cdots+a_n(n\ge1)$，$A_0=0$，于是
\[
\begin{aligned}
S_n
&=\sum_{k=1}^n\frac{k^2}{A_k^2}(A_k-A_{k-1})
\le\frac1{a_1}+\sum_{k=2}^n\frac{k^2}{A_kA_{k-1}}(A_k-A_{k-1})\\
&=\frac1{a_1}+\sum_{k=2}^n\frac{k^2}{A_{k-1}}-\sum_{k=2}^n\frac{k^2}{A_k}\\
 =\frac1{a_1}+\sum_{k=1}^{n-1}\frac{(k+1)^2}{A_k}-\sum_{k=2}^n\frac{k^2}{A_k}\\
&=\frac1{a_1}+2\sum_{k=2}^{n-1}\frac{k}{A_k}
  +\sum_{k=2}^{n-1}\frac1{A_k}
  +\frac4{A_1}-\frac{n^2}{A_n}\\
&\le\frac5{a_1}+2\sum_{k=1}^n\frac{k}{A_k}
  +\sum_{k=1}^n\frac1{a_k}.
\end{aligned}
\]
右端第二项用 Cauchy 不等式放大：
\[
\begin{aligned}
\sum_{k=1}^n\frac{k}{A_k}
&=\sum_{k=1}^n\frac{k}{A_k}\sqrt{a_k}\cdot\frac1{\sqrt{a_k}}\\
&\le
\left[\sum_{k=1}^n\left(\frac{k}{A_k}\sqrt{a_k}\right)^2\right]^{1/2}
\left[\sum_{k=1}^n\left(\frac1{\sqrt{a_k}}\right)^2\right]^{1/2}\\
&=\left(S_n\sum_{k=1}^n\frac1{a_k}\right)^{1/2}.
\end{aligned}
\]
得
\[
  S_n\le\frac5{a_1}+2\left(S_n\sum_{k=1}^n\frac1{a_k}\right)^{1/2}
  +\sum_{k=1}^n\frac1{a_k}.
\]
这是关于 $S_n$ 的不等式，解此不等式得
\[
  \sqrt{S_n}\le
  \sqrt{\sum_{k=1}^n\frac1{a_k}}+
  \sqrt{\frac5{a_1}+2\sum_{k=1}^n\frac1{a_k}}.
\]
因为 $\displaystyle\sum_{k=1}^\infty\dfrac1{a_k}$ 收敛，可知 $S_n$ 有界，原级数收敛.

\noindent\textbf{例 5.1.29}\quad 设 $a_n>0$，试证如下级数收敛：
\[
  \sum_{n=1}^\infty
  \frac{a_n}{(1+a_1)(1+a_2)\cdots(1+a_n)}.
\]

\noindent\textbf{提示}\quad 用数学归纳法或连锁消去法可证：
\[
  0\le S_n=
  1-\frac1{(1+a_1)(1+a_2)\cdots(1+a_n)}<1.
\]

\noindent\textbf{$\star$例 5.1.30}\quad 假设 $a_n>0(n=1,2,\cdots)$，$S_n=\displaystyle\sum_{k=1}^n a_k$，试问：级数
\[
  \sum_{n=1}^\infty\frac{a_n}{S_n^\lambda}
\]
与 $\displaystyle\sum_{k=1}^\infty a_k$ 的敛散性之间有何关系？给出证明.（南开大学）

\noindent\textbf{解}\quad $1^\circ$ 当 $\lambda=1$ 时，两者同时敛散. 因为：

1）若 $\displaystyle\sum_{k=1}^\infty a_k$ 发散，则由例 5.1.12 知 $\displaystyle\sum_{n=1}^\infty\dfrac{a_n}{S_n}$ 亦发散；

2）若 $\displaystyle\sum_{k=1}^\infty a_k$ 收敛，因为
\[
  0<\frac{a_n}{S_n}\le\frac1{S_1}a_n,
\]
可知 $\displaystyle\sum_{n=1}^\infty\dfrac{a_n}{S_n}$ 也收敛.

$2^\circ$ 当 $\lambda>1$ 时，不论 $\displaystyle\sum_{k=1}^\infty a_k$ 是否收敛，$\displaystyle\sum_{n=1}^\infty\dfrac{a_n}{S_n^\lambda}$ 恒收敛. 下面证明.

\noindent\textbf{证法 I}\quad
\[
  0<\frac{a_n}{S_n^\lambda}
  =\frac{S_n-S_{n-1}}{S_n^\lambda}
  =\int_{S_{n-1}}^{S_n}\frac1{S_n^\lambda}\,dx
  \le\int_{S_{n-1}}^{S_n}\frac1{x^\lambda}\,dx
  \qquad(n=2,3,\cdots).
\qquad(1)
\]
累加起来，$\forall N>2$，
\[
  0<\sum_{n=2}^N\frac{a_n}{S_n^\lambda}
  \le\int_{a_1}^{+\infty}\frac1{x^\lambda}\,dx<+\infty
  \qquad(\lambda>1).
\]
说明部分和有界，故正项级数 $\displaystyle\sum_{n=1}^\infty\dfrac{a_n}{S_n^\lambda}$ 收敛.

\noindent\textbf{证法 II}\quad （利用 Lagrange 定理.）当 $\lambda>1$ 时，
\[
  \frac1{1-\lambda}\left(\frac1{S_n^{\lambda-1}}-\frac1{S_{n-1}^{\lambda-1}}\right)
  \stackrel{\text{Lagrange 定理}}{=}
  \frac1{\xi_n^\lambda}(S_n-S_{n-1})
  \ge\frac{S_n-S_{n-1}}{S_n^\lambda}
  =\frac{a_n}{S_n^\lambda},
\]
其中 $\xi_n\in(S_{n-1},S_n)$. 因此
\[
\begin{aligned}
0<\sum_{n=1}^N\frac{a_n}{S_n^\lambda}
&\le\frac1{1-\lambda}\sum_{n=2}^N\left(\frac1{S_n^{\lambda-1}}-\frac1{S_{n-1}^{\lambda-1}}\right)\\
&=\frac1{1-\lambda}\left(\frac1{S_N^{\lambda-1}}-\frac1{S_1^{\lambda-1}}\right)\\
&\le\frac1{\lambda-1}\cdot\frac1{S_1^{\lambda-1}},
\end{aligned}
\]
此式说明部分和有界，故正项级数 $\displaystyle\sum_{n=1}^\infty\dfrac{a_n}{S_n^\lambda}$ 收敛.

$3^\circ$ 当 $0<\lambda<1$ 时，两者同时敛散. 这是因为：

1）若 $\displaystyle\sum_{k=1}^\infty a_k$ 收敛于 $S$，则
\[
\begin{aligned}
0<\sum_{n=2}^N\frac{a_n}{S_n^\lambda}
&=\sum_{n=2}^N\frac{S_n-S_{n-1}}{S_n^\lambda}
=\sum_{n=2}^N\frac1{S_n^\lambda}\int_{S_{n-1}}^{S_n}dx\\
&\le\sum_{n=2}^N\int_{S_{n-1}}^{S_n}\frac1{x^\lambda}\,dx
=\frac1{1-\lambda}\sum_{n=2}^N\int_{S_{n-1}}^{S_n}dx^{1-\lambda}\\
&=\frac1{1-\lambda}(S_N^{1-\lambda}-a_1^{1-\lambda})
\le\frac{S}{1-\lambda},
\end{aligned}
\]
说明部分和有界，故 $\displaystyle\sum_{n=1}^\infty\dfrac{a_n}{S_n^\lambda}$ 收敛.

2）若 $\displaystyle\sum_{k=1}^\infty a_k$ 发散，因为
\[
  \sum_{n=1}^N\frac{a_n}{S_n^\lambda}
  \ge\frac1{S_N^\lambda}\sum_{n=1}^N a_n
  =\frac1{S_N^\lambda}S_N
  =S_N^{1-\lambda}\stackrel{0<\lambda<1}{\longrightarrow}+\infty\quad(N\to\infty),
\]
故 $\displaystyle\sum_{n=1}^\infty\dfrac{a_n}{S_n^\lambda}$ 发散.

$4^\circ$ 若 $\lambda=0$，则
\[
  \sum_{n=1}^\infty\frac{a_n}{S_n^\lambda}
  =\sum_{k=1}^\infty a_k
\]
（二者归一，当然同时敛散）.

$5^\circ$ 若 $\lambda<0$，则 $S_n=\displaystyle\sum_{k=1}^n a_k$ 与
\[
  \sum_{n=1}^\infty\frac{a_n}{S_n^\lambda}
  =\sum_{n=1}^\infty a_nS_n^{-\lambda}
\]
同时敛散. 因为：

1）若 $\displaystyle\sum_{k=1}^\infty a_k$ 收敛，则 $S_n=\displaystyle\sum_{k=1}^n a_k$ 有界：$\exists M>0$，$\forall n>0$：$0<S_n<M$. 因此
\[
  0<\sum_{n=1}^N a_nS_n^{-\lambda}
  \le\sum_{n=1}^N a_nM^{-\lambda}
  =M^{-\lambda}\sum_{n=1}^N a_n
  =M^{-\lambda}S_N\le M^{-\lambda}S,
\]
即 $\displaystyle\sum_{n=1}^N a_nS_n^{-\lambda}$ 有界，可知 $\displaystyle\sum_{n=1}^\infty\dfrac{a_n}{S_n^\lambda}$ 收敛.

2）若 $\displaystyle\sum_{k=1}^n a_k$ 发散，则当 $n$ 充分大时，$S_n>1$，$S_n^{-\lambda}>1(-\lambda>0)$. 故
\[
  \sum_{n=1}^N\frac{a_n}{S_n^\lambda}
  =\sum_{n=1}^N a_nS_n^{-\lambda}
  \ge\sum_{n=1}^N a_n\to+\infty\quad(N\to\infty).
\]
因此 $\displaystyle\sum_{n=1}^\infty\dfrac{a_n}{S_n^\lambda}$ 发散.

\noindent\textbf{c.\quad 变号级数敛散性的判断}

\noindent\textbf{要点}\quad 设 $\sum a_n$ 为变号级数，判断 $\sum a_n$ 的敛散性，通常方法是

1）对 $\sum |a_n|$ 应用 D'Alembert 判别法或根式判别法，若 $\sum |a_n|$ 收敛，则 $\sum a_n$ 绝对收敛. 用这两种判别法时，若 $\sum |a_n|$ 发散，则意味着 $a_n\not\to0(n\to\infty)$，从而 $\sum a_n$ 亦发散.

2）应用 Leibniz 定理：若 $a_n\ge0,a_n\downarrow0$，则
\[
  \sum_{n=1}^\infty(-1)^{n-1}a_n
\]
收敛.

3）应用 Abel 判别法或 Dirichlet 判别法：

\noindent\textbf{Abel 判别法：}若 $\displaystyle\sum_{n=1}^\infty a_n$ 收敛，$\{b_n\}$ 单调有界，则 $\displaystyle\sum_{n=1}^\infty a_nb_n$ 收敛.

\noindent\textbf{Dirichlet 判别法：}若 $\left\{\displaystyle\sum_{k=1}^n a_k\right\}$ 有界，$b_n\nearrow0$（或 $b_n\searrow0$），则 $\displaystyle\sum_{n=1}^\infty a_nb_n$ 收敛.

4）应用 Cauchy 准则（或利用 Abel 变换等）. 另外注意，证明条件收敛时，必须同时证明两点：一是 $\sum a_n$ 收敛，二是 $\sum |a_n|$ 发散.

\noindent\textbf{$\star$例 5.1.31}\quad 证明级数
\[
  1-\frac13\left(1+\frac12\right)
  +\frac15\left(1+\frac12+\frac13\right)
  -\frac17\left(1+\frac12+\frac13+\frac14\right)+\cdots
\]
是收敛的.（上海师范大学）

\noindent\textbf{证}\quad 因为
\[
\begin{aligned}
|a_n|
&=\frac1{2n-1}\left(1+\frac12+\cdots+\frac1n\right)\\
&=\frac{(2n-1)+2}{(2n-1)(2n+1)}
  \left(1+\frac12+\cdots+\frac1n\right)\\
&=\frac1{2n+1}\left[\left(1+\frac2{2n-1}\right)
  \left(1+\frac12+\cdots+\frac1n\right)\right]\\
&>\frac1{2n+1}\left(1+\frac12+\cdots+\frac1n+\frac1{n+1}\right)
=|a_{n+1}|\qquad(n=1,2,\cdots).
\end{aligned}
\]
即 $|a_n|\downarrow$. 又
\[
  |a_n|=\frac1{2n-1}\left(1+\frac12+\cdots+\frac1n\right)
  =\frac1{2n-1}(C+\ln n+\varepsilon_n)\to0\quad(n\to\infty),
\]
其中 $C$ 为 Euler 常数，故原级数收敛（Leibniz 定理）.

\noindent\textbf{例 5.1.32}\quad 证明
\[
  \sum_{n=1}^{\infty}\frac{(-1)^{[\sqrt n]}}n
\]
收敛（方括号表示取整数部分）.

\noindent\textbf{提示}\quad 将级数中相邻并且符号相同的项合并为一项，组成一新的交错级数：
\[
  \sum_{n=1}^{\infty}(-1)^n
  \left[\frac1{n^2}+\frac1{n^2+1}+\cdots+\frac1{(n+1)^2-1}\right].
\]
注意到方括号内共有 $2n+1$ 项，其中前 $n$ 项之和与后 $n+1$ 项之和，分别夹在 $\dfrac1{n+1}$ 与 $\dfrac1n$ 之间，因此
\[
  \frac2{n+1}<\frac1{n^2}+\frac1{n^2+1}
  +\cdots+\frac1{(n+1)^2-1}<\frac2n.
\]
知新级数为 Leibniz 级数，故收敛. 原级数的任一部分和总是夹在新级数某相邻的两部分和之间，所以原级数也收敛.

\noindent\textbf{练习}\quad 有兴趣的读者不妨类似讨论级数
\[
  \sum_{n=1}^{\infty}\frac1n(-1)^{[\sqrt[3]n]}
\]
的敛散性.

\noindent\textbf{$\star$例 5.1.33}\quad 讨论级数
\[
  \sum_{n=1}^{\infty}a_n
  =\frac1{1^p}-\frac1{2^q}+\frac1{3^p}-\frac1{4^q}
  +\cdots+\frac1{(2n-1)^p}-\frac1{(2n)^q}+\cdots
\qquad(1)
\]
$(p>0,q>0)$ 的绝对收敛与条件收敛性.（复旦大学）

\noindent\textbf{解}\quad $1^\circ$ 若 $p,q>1$，则 $\displaystyle\sum_{n=1}^{\infty}a_n$ 绝对收敛（因为例如 $p>q$，则 $\displaystyle\sum_{n=1}^{\infty}|a_n|$ 以 $\displaystyle\sum_{n=1}^{\infty}\frac1{n^q}$ $(q>1)$ 为优级数）.

$2^\circ$ 若 $0<p=q\le1$，应用 Leibniz 定理知级数收敛，且条件收敛.

$3^\circ$ 当 $p,q>0$ 时（此时通项 $\to0$），原级数 (1) 跟级数
\[
  \sum_{n=1}^{\infty}\left[\frac1{(2n-1)^p}-\frac1{(2n)^q}\right]
\qquad(2)
\]
同时敛散.

若 $p>1,0<q\le1$ 或 $q>1,0<p\le1$，级数
\[
  \sum_{n=1}^{\infty}\frac1{(2n-1)^p},\qquad
  \sum_{n=1}^{\infty}\frac1{(2n)^q}
\]
一收敛一发散，故原级数发散.

若 $0<p<q<1$，则
\[
  \frac1{(2n-1)^p}-\frac1{(2n)^q}>0,
\]
且与 $\dfrac1{(2n-1)^p}$ 同阶 $(n\to\infty)$，故级数 (2) 发散，从而 (1) 发散.

同理可证，若 $0<q<p<1$，则原级数发散.

\noindent\textbf{例 5.1.34}\quad 研究级数
\[
  \sum_{n=1}^{\infty}\frac{(-1)^{n-1}}{n^{p+\frac1n}}
\]
的绝对收敛与条件收敛性.（辽宁大学）

\noindent\textbf{解}\quad $1^\circ$ 当 $p\le0$ 时，通项 $\nrightarrow0$ $(n\to\infty)$，原级数发散.

$2^\circ$ 当 $p>1$ 时，因 $\dfrac1{n^{p+\frac1n}}<\dfrac1{n^p}$，故原级数绝对收敛.

$3^\circ$ 当 $0<p\le1$ 时，
\[
  \sum_{n=1}^{\infty}\frac{(-1)^{n-1}}{n^p}
\]
收敛，$\dfrac1{n^{1/n}}$ 单调有界，应用 Abel 判别法知原级数收敛. 因为
\[
  \left|\frac{(-1)^{n-1}}{n^{p+\frac1n}}\right|\bigg/\frac1{n^p}\to1\quad(n\to\infty),
\]
故原级数条件收敛.

\noindent\textbf{new 练习 1}\quad 设 $f(x)$ 在区间 $[-1,1]$ 上有二阶连续导数，且 $\displaystyle\lim_{x\to0}\frac{f(x)}x=0$. 证明：数项级数
\[
  \sum_{n=1}^{\infty}f\left(\frac1n\right)
\]
绝对收敛.（华中科技大学）

\noindent\textbf{证}\quad （利用 Taylor 公式.）
\[
  \lim_{x\to0}\frac{f(x)}x=0
  \Rightarrow \lim_{x\to0}f(x)=0\Rightarrow f(0)=0
  \Rightarrow f'(0)=\lim_{x\to0}\frac{f(x)-0}{x-0}=0.
\]
因 $f''(x)$ 在 $[-1,1]$ 上连续，故 $\exists M>0$，使得
$|f''(x)|\le M,\ \forall x\in[-1,1]$. 因此
\[
  \left|f\left(\frac1n\right)\right|
  \stackrel{\text{Taylor 公式}}{\le}
  \left|f(0)+f'(0)\left(\frac1n\right)
  +\frac{f''(\xi)}{2!}\left(\frac1n\right)^2\right|
  \le\frac M2\cdot\frac1{n^2},
  \qquad 0<\xi<\frac1n.
\]
故 $\displaystyle\sum_{n=1}^{\infty}f\left(\dfrac1n\right)$ 绝对收敛.

\noindent\textbf{new 练习 2}\quad 设 $x_1=\dfrac12$，
$x_{n+1}=\dfrac12-\dfrac{x_n^2}2\ (n=1,2,\cdots)$.

1）试证极限 $\displaystyle\lim_{n\to\infty}x_n=A$ 存在，有限；

2）证明级数 $\displaystyle\sum_{n=1}^{\infty}(x_n-A)$ 绝对收敛.（华中师范大学）

\noindent\textbf{证}\quad 1）由 $0<x_n\le\dfrac12\ (\forall n\in\mathbb N)$ 可知 $|x_n+x_{n-1}|\le1$，故
\[
  |x_{n+1}-x_n|
  =\frac12|x_n^2-x_{n-1}^2|
  \le\frac12|x_n-x_{n-1}|.
\]
（由压缩映象原理知 $\{x_n\}$ 收敛.）从递推式取极限得
$A=\dfrac12-\dfrac{A^2}2$，解得 $A=\sqrt2-1$.

2）对 $\displaystyle\sum_{n=1}^{\infty}|x_n-A|$ 应用 D'Alembert 判别法：
\[
\begin{aligned}
  \frac{|x_{n+1}-A|}{|x_n-A|}
  &=\frac{\left|\dfrac12-\dfrac{x_n^2}2-\left(\dfrac12-\dfrac{A^2}2\right)\right|}
          {|x_n-A|}\\
  &=\frac12\cdot\frac{|x_n^2-A^2|}{|x_n-A|}
  =\frac12|x_n+A|\to |A|<1\quad(n\to\infty).
\end{aligned}
\]
故 $\displaystyle\sum_{n=1}^{\infty}(x_n-A)$ 绝对收敛获证.

\noindent\textbf{d.\quad Abel 变换}

\noindent\textbf{要点}\quad 所谓 Abel 变换，就是借助于
$A_n=\displaystyle\sum_{k=1}^n a_k$，将和式
$\displaystyle\sum_{k=n+1}^m a_kb_k$ 或
$\displaystyle\sum_{k=1}^n a_kb_k$ 进行改写. 下面将相关公式列为例 5.1.35，以便后面引用和查询.

\noindent\textbf{new 例 5.1.35}\quad 试证 Abel 变换公式：

1）
\[
  \sum_{k=n+1}^m a_kb_k
  =-A_nb_{n+1}+\sum_{k=n+1}^{m-1}A_k(b_k-b_{k+1})+A_mb_m.
\tag{A}
\]
若 $|A_n|\le M(\forall n\in\mathbb N),\ m>n\ge1,\ b_n\le0,\nearrow$
（或 $b_n\ge0,\searrow$），则
\[
  \left|\sum_{k=n+1}^m a_kb_k\right|\le2M|b_{n+1}|.
\tag{B}
\]
若只知道 $b_n$ 单调，则
\[
  \left|\sum_{k=n+1}^m a_kb_k\right|
  \le2M(|b_{n+1}|+|b_m|).
\tag{C}
\]

2）对 $\displaystyle\sum_{k=1}^n a_kb_k$，有
\[
  \sum_{k=1}^n a_kb_k=A_1b_1+\sum_{k=2}^n(A_k-A_{k-1})b_k.
\tag{$A_1$}
\]
若 $\forall n\in\mathbb N,\ |A_n|\le M$，且 $b_n\ge0,\searrow$（或 $b_n\le0,\nearrow$），则
\[
  \left|\sum_{k=1}^n a_kb_k\right|\le M|b_1|.
\tag{$B_1$}
\]
若只知道 $b_n$ 单调，则
\[
  \left|\sum_{k=1}^n a_kb_k\right|\le M(|b_1|+2|b_n|).
\tag{$C_1$}
\]

3）由此直接得出 Abel 判别法和 Dirichlet 判别法.

\noindent\textbf{证}\quad 1）$(m>n\ge1)$ 用
$A_n=\displaystyle\sum_{k=1}^n a_k$ 代入得
\[
\begin{aligned}
  \sum_{k=n+1}^m a_kb_k
  &=\sum_{k=n+1}^m(A_k-A_{k-1})b_k\\
  &=-A_nb_{n+1}+\sum_{k=n+1}^{m-1}A_k(b_k-b_{k+1})+A_mb_m,
\end{aligned}
\]
此即式 (A). 当 $|A_n|\le M(\forall n\in\mathbb N)$ 时，
\[
  \sum_{k=n+1}^m a_kb_k
  \le M|b_{n+1}|+M\sum_{k=n+1}^{m-1}|b_k-b_{k+1}|+M|b_m|.
\]
若 $\{b_n\}$ 单调下降且 $b_n\ge0$（或 $\{b_n\}$ 单调上升且 $b_n\le0$），则 $\{|b_n-b_{n+1}|\}$ 保持同号（因而 $\sum|b_k-b_{k+1}|=\left|\sum(b_k-b_{k+1})\right|$），故有
\[
  \left|\sum_{k=n+1}^m a_kb_k\right|
  \le M(|b_{n+1}|+|b_{n+1}-b_m|+|b_m|)
  =2M|b_{n+1}|.
\tag{B}
\]
若只知道 $\{b_n\}$ 单调，则
\[
  \left|\sum_{k=n+1}^m a_kb_k\right|
  \le M(|b_{n+1}|+|b_{n+1}|+|b_m|+|b_m|)
  \le2M(|b_{n+1}|+|b_m|).
\tag{C}
\]

2）$\displaystyle\sum_{k=1}^n a_kb_k$ 与
$\displaystyle\sum_{k=n+1}^m a_kb_k$ 不同之处在首项：前者为
$a_1=A_1$，后者为 $a_{n+1}=A_{n+1}-A_n$. 若令
$a_0=A_0=0$，则 $a_1=A_1-A_0$ 成为 $a_{n+1}=A_{n+1}-A_n$ 当 $n=0$ 的情形. 这时式 (B) 和 (C) 推导过程有效. 又 $n=0$ 时 $A_0=0$，因而公式 (A) 右端第一项 $-A_nb_{n+1}|_{n=0}=0$，故对应的公式 $(A_1)$ 比公式 (A) 少（此）一项. 公式 $(B_1)$ 和 $(C_1)$ 都是由公式 (A) 导出. 同样，因为 $A_0=0$ 使得对应的公式 $(B_1)$ 和 $(C_1)$ 分别比公式 (B) 和 (C) 各少了一个 $b_1$.

3）（Dirichlet 判别法）若 $A_n=\displaystyle\sum_{k=1}^n a_k$ 有界（即存在 $M$ 使得 $|A_n|\le M(\forall n\in\mathbb N)$），$|b_n|$ 单调趋于 $0$，则
$\displaystyle\sum_{k=1}^{\infty}a_kb_k$ 收敛.

（下证 Dirichlet 判别法.）因为 $|b_n|$ 单调趋于 $0$，$\forall\varepsilon>0$，$\exists N>0$，当 $\forall m>n>N$ 时，有 $|b_n|<\dfrac{\varepsilon}{2M}$，利用上面的公式 (B) 得
\[
  \left|\sum_{k=n+1}^m a_kb_k\right|
  \le2M|b_{n+1}|<2M\cdot\frac{\varepsilon}{2M}=\varepsilon.
\]
根据 Cauchy 准则，级数 $\displaystyle\sum_{k=1}^{\infty}a_kb_k$ 收敛.

（Abel 判别法）若 $\displaystyle\sum_{k=1}^{\infty}a_k$ 收敛，$|b_n|$ 单调有界，则
$\displaystyle\sum_{k=1}^{\infty}a_kb_k$ 收敛.

（下证 Abel 判别法.）因为 $|b_n|$ 单调有界，所以
$\displaystyle\lim_{n\to\infty}b_n=b$ 存在，于是 $|b_n-b|$ 单调 $\to0$ $(n\to\infty)$. 因此
\[
  \sum_{k=1}^n a_kb_k
  =\sum_{k=1}^n a_k(b_k-b)+b\sum_{k=1}^n a_k.
\]
因为 $\displaystyle\sum_{k=1}^{\infty}a_k$ 收敛，其部分和有界，$|b_k-b|$ 单调趋于零，根据 Dirichlet 判别法，
$\displaystyle\sum_{k=1}^{\infty}a_k(b_k-b)$ 收敛；右端第二项已知
$\displaystyle\sum_{k=1}^{\infty}a_k$ 收敛，故
$\displaystyle\sum_{k=1}^{\infty}a_kb_k$ 收敛.

\noindent\textbf{new 练习 1}\quad 讨论下面级数的敛散性：
\[
  \sum_{n=1}^{\infty}\left(1+\frac12+\frac13+\cdots+\frac1n\right)\frac{\sin nx}{n}.
\]
（浙江大学）

\noindent\textbf{解}\quad （用 Dirichlet 判别法.）当 $x=2n\pi$ 时，级数成为零级数，当然收敛.

当 $x\ne2n\pi$ 时，记 $a_k=\sin kx$，则
\[
\begin{aligned}
  \left|\sum_{k=1}^n a_k\right|
  &=\left|\sum_{k=1}^n\sin kx\right|
   =\left|\left(\sum_{k=1}^n2\sin kx\cdot\sin\frac x2\right)\bigg/\left(2\sin\frac x2\right)\right|\\
  &=\left|\left(\sum_{k=1}^n\cos\frac{2k-1}{2}x-\cos\frac{2k+1}{2}x\right)\bigg/\left(2\sin\frac x2\right)\right|\\
  &\le\frac1{\left|\sin\frac x2\right|}
  \qquad(\text{当 }x\text{ 固定时，对 }n\text{ 保持有界}).
\end{aligned}
\]
又
\[
  b_n=\frac1n\left(1+\frac12+\frac13+\cdots+\frac1n\right)
  =\frac{C+\ln n+\varepsilon_n}{n}\to0
\]
（$C$ 为 Euler 常数）. 而
\[
  \left(\frac{C+\ln x}{x}\right)'=\frac{1-C-\ln x}{x^2}<0
  \qquad(\text{当 }x>e\text{ 时}),
\]
所以 $b_n\searrow0$，故利用 Dirichlet 判别法知原级数收敛. 证毕.

\noindent\textbf{new 练习 2}\quad 设 $\{a_n\}$ 和 $\{b_n\}$ 是两个有界数列，满足
$\alpha a_{n+1}+\beta a_n=b_n,\ n\ge1$. 若 $\alpha>\beta>0$，且 $b_n$ 单调，试证 $\{a_n\}$ 收敛.（仿华中师范大学）

\noindent\textbf{提示}\quad 因为
\[
  a_{n+1}=-\frac{\beta}{\alpha}a_n+\frac1\alpha b_n,
\]
若记 $q=\dfrac{\beta}{\alpha}$，则 $0<q<1$. 利用数学归纳法或回推法，易证 $\{a_n\}$ 的通项：
\[
  a_{n+1}=\frac1\alpha\sum_{k=0}^{n-1}(-q)^k b_{n-k}+(-q)^n a_1\quad(n\ge1).
\]
再用 Abel 判别法即得.

\noindent\textbf{证}\quad
\[
\begin{aligned}
  a_{n+1}
  &=-qa_n+\frac1\alpha b_n
   =-q\left(-qa_{n-1}+\frac1\alpha b_{n-1}\right)+\frac1\alpha b_n\\
  &=(-q)^2a_{n-1}+\frac{-q}{\alpha}b_{n-1}+\frac1\alpha b_n\\
  &=(-q)^2\left(-qa_{n-2}+\frac1\alpha b_{n-2}\right)
    +\frac{-q}{\alpha}b_{n-1}+\frac1\alpha b_n\\
  &=(-q)^3a_{n-2}+\frac{(-q)^2}{\alpha}b_{n-2}
    +\frac{-q}{\alpha}b_{n-1}+\frac1\alpha b_n=\cdots\\
  &=\frac1\alpha\sum_{k=0}^{n-1}(-q)^k b_{n-k}+(-q)^n a_1.
\end{aligned}
\]
因 $\{b_n\}$ 单调有界，而 $\displaystyle\sum_{k=0}^{\infty}(-1)^kq^k$ 收敛，利用 Abel 判别法，知 $\{a_n\}$ 收敛.

\noindent\textbf{练习 3}\quad 设级数 $\displaystyle\sum_{n=1}^{\infty}a_n$ 收敛，$\displaystyle\sum_{n=1}^{\infty}(b_{n+1}-b_n)$ 绝对收敛，试证级数
$\displaystyle\sum_{n=1}^{\infty}a_nb_n$ 也收敛.

\noindent\textbf{证}\quad 因为 $\displaystyle\sum(b_{n+1}-b_n)$ 绝对收敛，所以
$\displaystyle\sum(b_{n+1}-b_n)$ 收敛，从而 $b_n-b_1\to A$，
$b_n\to A+b_1$，所以 $\{b_n\}$ 有界，$|b_n|\le M$. 因为
$\displaystyle\sum_{n=1}^{\infty}a_n$ 收敛，$\displaystyle\sum|b_{n+1}-b_n|$ 收敛，根据 Cauchy 准则，$\forall\varepsilon>0$，$\exists N>0$，当 $n>N$ 时，
\[
  \left|\sum_{k=n+1}^{n+p}a_k\right|<\frac{\varepsilon}{1+M},
  \qquad
  \sum_{k=n+1}^{n+p}|b_{k+1}-b_k|<1\quad(\forall p\in\mathbb N).
\]
记 $S_{n+i}=\displaystyle\sum_{k=n+1}^{n+i}a_k\ (i=1,2,\cdots,p)$，则
\[
\begin{aligned}
  \left|\sum_{k=n+1}^{n+p}a_kb_k\right|
  &=|a_{n+1}b_{n+1}+a_{n+2}b_{n+2}+\cdots+a_{n+p}b_{n+p}|\\
  &=|S_{n+1}b_{n+1}+(S_{n+2}-S_{n+1})b_{n+2}
     +\cdots+(S_{n+p}-S_{n+p-1})b_{n+p}|\\
  &=|S_{n+1}(b_{n+1}-b_{n+2})+\cdots
    +S_{n+p-1}(b_{n+p-1}-b_{n+p})+S_{n+p}b_{n+p}|\\
  &\le |S_{n+1}||b_{n+1}-b_{n+2}|+\cdots
    +|S_{n+p-1}||b_{n+p-1}-b_{n+p}|+|S_{n+p}||b_{n+p}|\\
  &\le \frac{\varepsilon}{1+M}
    \left(\sum_{k=n+1}^{n+p}|b_{k+1}-b_k|+|b_{n+p}|\right)
  \le\frac{\varepsilon}{1+M}(1+M)=\varepsilon.
\end{aligned}
\]
所以 $\displaystyle\sum_{n=1}^{\infty}a_nb_n$ 收敛.

\subsection*{三、级数敛散性的应用}

\noindent\textbf{a.\quad 收敛性的应用}

\noindent\textbf{$\star$例 5.1.36}\quad 设
$x_n=\dfrac{n^n}{n!3^n}$，求 $\displaystyle\lim_{n\to\infty}x_n$.（上海交通大学，华中师范大学）

\noindent\textbf{解}\quad 将 $x_n$ 看成级数 $\displaystyle\sum_{n=1}^{\infty}x_n$ 的通项，因
\[
  \frac{x_n}{x_{n-1}}=\frac13\left(1+\frac1{n-1}\right)^{n-1}
  \to\frac e3<1\quad(n\to\infty).
\]
因此，级数 $\sum x_n$ 收敛，$\displaystyle\lim_{n\to\infty}x_n=0$.

类似可证
\[
  \lim_{n\to\infty}\frac{a^n}{n!}=0,\qquad
  \lim_{n\to\infty}\frac{(2n)!}{a^{n!}}=0\quad(a>1).
\]
（此类考题很多.）

\noindent\textbf{new 练习}\quad 设 $a_n>0(n=1,2,\cdots)$，
$\displaystyle\sum_{n=1}^{\infty}a_n=1$，$A_n=\displaystyle\sum_{k=1}^n a_k$，求极限：
\[
  \lim_{n\to\infty}\frac{e^{A_n}-e^{A_{n-1}}}{A_n^e-A_{n-1}^e}.
\]
（华中科技大学）

\noindent\textbf{解}\quad 因级数 $\displaystyle\sum_{n=1}^{\infty}a_n$ 收敛，故通项 $a_n\to0\ (n\to\infty)$.
\[
\begin{aligned}
  \frac{e^{A_n}-e^{A_{n-1}}}{A_n^e-A_{n-1}^e}
  &=\frac{e^{A_{n-1}}(e^{A_n-A_{n-1}}-1)}
          {A_{n-1}^e\left[\left(A_n/A_{n-1}\right)^e-1\right]}\\
  &=\frac{e^{A_{n-1}}}{A_{n-1}^e}\cdot
    \frac{e^{a_n}-1}{\left(1+\dfrac{a_n}{A_{n-1}}\right)^e-1}.
\end{aligned}
\]
其中
\[
  \frac{e^{a_n}-1}{\left(1+\dfrac{a_n}{A_{n-1}}\right)^e-1}
  =
  \frac{(e^{a_n}-1)/a_n}
       {\left[\left(1+\dfrac{a_n}{A_{n-1}}\right)^e-1\right]/
        \left[\left(\dfrac{a_n}{A_{n-1}}\right)\cdot A_{n-1}\right]}
  \to\frac1e.
\]
而
\[
  \frac{e^{A_{n-1}}}{A_{n-1}^e}\to e,
\]
故所求极限为 $1$.

\noindent\textbf{$\star$例 5.1.37}\quad 设
\[
  x_n=1+\frac1{\sqrt2}+\cdots+\frac1{\sqrt n}-2\sqrt n,
\]
试证 $\displaystyle\lim_{n\to\infty}x_n$ 存在.

\noindent\textbf{证}\quad 因为
\[
  x_n=\sum_{k=1}^n(x_k-x_{k-1})\quad(\text{记 }x_0=0)
\]
是 $\displaystyle\sum_{k=1}^{\infty}(x_k-x_{k-1})$ 的部分和，而
\[
  x_k-x_{k-1}=\frac1{\sqrt k}-2(\sqrt k-\sqrt{k-1})
  =\frac1{\sqrt k}-\frac2{\sqrt k+\sqrt{k-1}}
  =-\frac1{\sqrt k(\sqrt k+\sqrt{k-1})^2}
  =O\left(\frac1{k^{3/2}}\right),
\]
所以 $\displaystyle\sum(x_k-x_{k-1})$ 收敛，$\displaystyle\lim_{n\to\infty}x_n$ 存在.

\noindent\textbf{例 5.1.38}\quad 求极限
\[
  \lim_{n\to\infty}\left(\frac1{p^{n+1}}+\frac1{p^{n+2}}+\cdots+\frac1{p^{2n}}\right)
  \qquad(p>1).
\]

\noindent\textbf{提示}\quad 考虑级数 $\displaystyle\sum_{n=1}^{\infty}\dfrac1{p^n}$，利用 Cauchy 准则.

\noindent\textbf{$\star$例 5.1.39}\quad 设级数 $\displaystyle\sum_{n=1}^{\infty}a_n$ 收敛，$a_n>0$，$a_n\downarrow$，试证
$\displaystyle\lim_{n\to\infty}na_n=0$.

\noindent\textbf{证}\quad （要证明 $\lim_{n\to\infty}na_n=0$，即 $\forall\varepsilon>0$，要证 $\exists N>0$，使得 $n>N$ 时，有 $0\le na_n<\varepsilon$.）
因 $\displaystyle\sum_{n=1}^{\infty}a_n(a_n>0)$ 收敛，根据 Cauchy 准则，$\forall\varepsilon>0$，$\exists N>0$，$n>N$ 时，
\[
  0<a_{N+1}+a_{N+2}+\cdots+a_n<\frac{\varepsilon}{2}.
\qquad(1)
\]
但 $a_n\downarrow$，故
\[
  (n-N)a_n\le a_{N+1}+\cdots+a_n<\frac{\varepsilon}{2}.
\]
特别，令 $n=2N$，得 $(2N-N)a_{2N}<\dfrac{\varepsilon}{2}$. 故当 $n>2N$ 时，
\[
  na_n=(n-N)a_n+(2N-N)a_n
  <(n-N)a_n+(2N-N)a_{2N}<\frac{\varepsilon}{2}+\frac{\varepsilon}{2}
  =\varepsilon.
\]
故 $\displaystyle\lim_{n\to\infty}na_n=0$.

\noindent\textbf{注}\quad 本例说明递减正项级数要收敛，其通项必须是比 $\dfrac1n$ 高阶的无穷小量，但注意此条件并不充分.

\noindent\textbf{new 练习}\quad 设数列 $\{a_n\}$ 满足条件：

1）$0<a_k\le100a_n\ (n=k+1,k+2,\cdots)$；

2）$\displaystyle\sum_{n=1}^{\infty}a_n$ 收敛.

求证：$\displaystyle\lim_{n\to\infty}na_n=0$.（华东师范大学）

\noindent\textbf{提示}\quad 利用 Cauchy 准则，由条件 2）知，$\forall\varepsilon>0$，$\exists N\in\mathbb N$，$\forall n>N$，有
\[
  0<\sum_{k=n+1}^{2n}a_k<\frac{\varepsilon}{100}.
\qquad(1)
\]
再利用条件 1），可证：$n>N$ 时，有 $|na_n|<\varepsilon$.

\noindent\textbf{再提示}\quad 因 $a_n<100a_{n+1},a_n<100a_{n+2},\cdots,a_n<100a_{2n}$，故
\[
  |na_n|\le\left|100\sum_{k=n+1}^{2n}a_k\right|
  \stackrel{\text{式 }(1)}{<}100\cdot\frac{\varepsilon}{100}=\varepsilon.
\]

\noindent\textbf{例 5.1.40}\quad 设正项级数 $\displaystyle\sum_{n=1}^{\infty}a_n$ 收敛，试证：
\[
  \lim_{n\to\infty}\frac{\sum_{k=1}^n ka_k}{n}=0.
\]

\noindent\textbf{证}\quad 记 $S=\displaystyle\sum_{n=1}^{\infty}a_n$，$S_n=\displaystyle\sum_{k=1}^n a_k$，则 $S_n\to S$（当 $n\to\infty$ 时）. 利用 Abel 变换，
\[
  \sum_{k=1}^n ka_k=nS_n-\sum_{k=1}^{n-1}S_k,
\]
从而
\[
  \lim_{n\to\infty}\frac{\sum_{k=1}^n ka_k}{n}
  \stackrel{\text{Abel 变换}}{=}
  \lim_{n\to\infty}\left(S_n-\frac{S_1+S_2+\cdots+S_{n-1}}n\right)
  =S-S=0.
\]

\noindent\textbf{new 练习}\quad 设正项级数 $\displaystyle\sum_{n=1}^{\infty}a_n$ 收敛，求证：
\[
  \lim_{n\to\infty}\frac{n^2}{\dfrac1{a_1}+\dfrac1{a_2}+\cdots+\dfrac1{a_n}}=0.
\]
（北京大学）

\noindent\textbf{证}\quad 利用已知不等式（见例 3.4.8），
\[
  0\le\frac{n}{\dfrac1{a_1}+\dfrac1{a_2}+\cdots+\dfrac1{a_n}}
  \le\sqrt[n]{a_1a_2\cdots a_n}\le\frac{a_1+a_2+\cdots+a_n}{n}.
\]
以 $n$ 同乘各项，得
\[
  0<\frac{n^2}{\dfrac1{a_1}+\dfrac1{a_2}+\cdots+\dfrac1{a_n}}
  \le n\sqrt[n]{a_1a_2\cdots a_n}
  \le n\frac{a_1+a_2+\cdots+a_n}{n}.
\qquad(1)
\]
式 (1) 右边的不等式，对每个 $k$，用 $ka_k$ 替代 $a_k(k=1,2,\cdots,n)$，再同乘 $\dfrac1{\sqrt[n]{n!}}$ 得
\[
  0<\frac{n^2}{\dfrac1{a_1}+\dfrac1{a_2}+\cdots+\dfrac1{a_n}}
  \le\frac n{\sqrt[n]{n!}}\sqrt[n]{(a_1)(2a_2)\cdots(na_n)}
  \le\frac n{\sqrt[n]{n!}}\cdot\frac{a_1+2a_2+\cdots+na_n}{n}.
\qquad(2)
\]
其中
\[
  \lim_{n\to\infty}\frac{a_1+2a_2+\cdots+na_n}{n}
  =\lim_{n\to\infty}\frac{\sum_{k=1}^nka_k}{n}
  \stackrel{\text{例 5.1.40}}{=}0.
\]
而 $\displaystyle\lim_{n\to\infty}\dfrac n{\sqrt[n]{n!}}=e$（见例 1.3.14 的 4）. 根据式 (2)，利用求极限的两边夹法则，待证的极限为零.

\noindent\textbf{例 5.1.41}\quad 试证：若 $\displaystyle\sum_{n=1}^{\infty}a_n$ 收敛，$a_n>0$，$|a_n-a_{n+1}|\downarrow$，则 $a_n\downarrow0$，且
\[
  \lim_{n\to\infty}\left(\frac1{a_{n+1}}-\frac1{a_n}\right)=+\infty.
\]

\noindent\textbf{证}\quad $1^\circ$ 因 $|a_n-a_{n+1}|\downarrow$，且 $a_n-a_{n+1}\to0$（因 $\sum a_n$ 收敛，知 $a_n\to0$），故 $a_n-a_{n+1}\ge0$，即 $a_n\ge a_{n+1}$，因此 $a_n\downarrow0$.

$2^\circ$ 要证 $n\to\infty$ 时，
\[
  \frac1{a_{n+1}}-\frac1{a_n}
  =\frac{a_n-a_{n+1}}{a_na_{n+1}}\to+\infty,
\]
即要证明 $\dfrac{a_na_{n+1}}{a_n-a_{n+1}}\to0$. 事实上，
\[
\begin{aligned}
0\le\frac{a_na_{n+1}}{a_n-a_{n+1}}
&\le\frac{a_n^2}{a_n-a_{n+1}}
=\frac1{a_n-a_{n+1}}\sum_{k=n}^{\infty}(a_k^2-a_{k+1}^2)\\
&\le\sum_{k=n}^{\infty}\frac{a_k^2-a_{k+1}^2}{a_k-a_{k+1}}
=\sum_{k=n}^{\infty}(a_k+a_{k+1})
=R_{n-1}+R_n\to0\quad(n\to\infty),
\end{aligned}
\]
其中 $R_{n-1}=\displaystyle\sum_{k=n}^{\infty}a_k$ 为收敛级数
$\displaystyle\sum_{n=1}^{\infty}a_n$ 的余和.

\noindent\textbf{条件收敛的应用}

\noindent\textbf{例 5.1.42}\quad 研究级数
$\displaystyle\sum_{n=1}^{\infty}\dfrac{\sin n}{n}$，把这个级数的前 $n$ 项和分成两项：
\[
  S_n=\sum_{k=1}^n\frac{\sin k}{k}=S_n^++S_n^-,
\qquad(1)
\]
其中 $S_n^+$ 和 $S_n^-$ 分别是正项之和与负项之和. 证明：
$\displaystyle\lim_{n\to\infty}\dfrac{S_n^+}{S_n^-}$ 存在并求其值.（国外赛题）

\noindent\textbf{提示}\quad $1^\circ$ 利用 Dirichlet 判别法知，
$\displaystyle\sum_{n=1}^{\infty}\dfrac{\sin n}{n}$ 收敛.

$2^\circ$ 由
\[
  \sum_{k=1}^n\frac{|\sin k|}{k}
  \ge\sum_{k=1}^n\frac{\sin^2 k}{k}
  =\frac12\sum_{k=1}^n\frac1k-\frac12\sum_{k=1}^n\frac{\cos2k}{k}
\]
可知，$\displaystyle\sum_{n=1}^{\infty}\dfrac{\sin n}{n}$ 非绝对收敛.

$3^\circ$ 由 (1) 式知
\[
  S_n^-=\frac{S_n^++S_n^--(S_n^+-S_n^-)}2
  =\frac{S_n}2-\frac{\sigma_n}2\to-\infty,
  \qquad \sigma_n=\sum_{k=1}^n\frac{|\sin k|}{k}\to+\infty.
\]

$4^\circ$
\[
  \frac{S_n^+}{S_n^-}
  =\frac{S_n^++S_n^-}{S_n^-}-\frac{S_n^-}{S_n^-}
  =\frac{S_n}{S_n^-}-1\to-1\quad(n\to\infty).
\]

\noindent\textbf{注}\quad 本题结论对任一条件收敛级数都成立.

\noindent\textbf{b.\quad 发散性的应用}

\noindent\textbf{$\star$例 5.1.43}\quad 假设 $\displaystyle\sum_{n=1}^{\infty}a_n$ 发散，且 $|a_n|$ 是正的不增数列，试证：
\[
  \lim_{n\to\infty}\frac{a_2+a_4+\cdots+a_{2n}}{a_1+a_3+\cdots+a_{2n-1}}=1.
\qquad(1)
\]
（东北师范大学）

\noindent\textbf{证}\quad 因 $a_1\ge a_2\ge a_3\ge a_4\ge\cdots\ge a_{2n-1}\ge a_{2n}\ge\cdots\ge0$，故
\[
  a_1+a_3+\cdots+a_{2n-1}
  \ge a_2+a_4+\cdots+a_{2n}
  \ge a_3+a_5+\cdots+a_{2n-1}.
\qquad(2)
\]
从而
\[
  1\ge\frac{a_2+a_4+\cdots+a_{2n}}{a_1+a_3+\cdots+a_{2n-1}}
  \ge1-\frac{a_1}{a_1+a_3+\cdots+a_{2n-1}}\to1\quad(n\to\infty).
\qquad(3)
\]
最后的极限是因 $\displaystyle\sum_{n=1}^{\infty}a_n$ 发散，由式 (2)，
\[
  a_1+a_3+\cdots+a_{2n-1}
  \ge\frac12(a_2+a_3+\cdots+a_{2n})
  =\frac12(S_{2n}-a_1)\to+\infty\quad(n\to\infty).
\]
式 (3) 表明式 (1) 成立.

\noindent\textbf{new $\star$练习}\quad 设 $|a_n|$ 单调递减收敛于 $0$，级数 $\displaystyle\sum_{n=1}^{\infty}|a_n\sin n|$ 发散，试证：

1）$\displaystyle\sum_{n=1}^{\infty}a_n\sin n$ 收敛；

2）$\displaystyle\lim_{n\to\infty}\frac{u_n}{v_n}=1$，其中
\[
  u_n=\sum_{k=1}^n(|a_k\sin k|+a_k\sin k),\qquad
  v_n=\sum_{k=1}^n(|a_k\sin k|-a_k\sin k).
\]
（南开大学）

\noindent\textbf{提示}\quad 1）（应用 Dirichlet 判别法.）注意
\[
  \left|\sum_{k=1}^n\sin k\right|\le\frac1{\left|\sin\frac12\right|}
\]
（参看例 5.1.35 练习 1 的解）.

2）记
$A_n=\displaystyle\sum_{k=1}^n|a_k\sin k|$，
$B_n=\displaystyle\sum_{k=1}^n a_k\sin k$，则
\[
  \frac{u_n}{v_n}
  =\frac{A_n+B_n}{A_n-B_n}
  =\frac{1+B_n/A_n}{1-B_n/A_n}\to1\quad(n\to\infty),
\]
（因已证 $\displaystyle\sum_{n=1}^{\infty}a_n\sin n$ 收敛，而
$\displaystyle\sum_{n=1}^{\infty}|a_n\sin n|$ 发散.）

\noindent\textbf{*例 5.1.44}\quad 设

1）$a_k>0(k=1,2,\cdots)$；
2）$\displaystyle\lim_{k\to\infty}a_k=0$；
3）$\displaystyle\sum_{k=1}^{\infty}a_k$ 发散.

证明：$\{S_n-[S_n]\}\ (n=1,2,\cdots)$ 在 $[0,1]$ 中稠密，其中
$S_n=\displaystyle\sum_{k=1}^n a_k$，$[S_n]$ 为 $S_n$ 的整数部分.（兰州大学）

\noindent\textbf{分析}\quad 问题等价于数列 $\{S_n\}\nearrow+\infty$，
$S_n-S_{n-1}=a_n\to0$（当 $n\to\infty$ 时），求证 $S_n$ 的小数部分
$\alpha_n=S_n-[S_n]$ 在 $[0,1]$ 中稠密. 即 $\forall(\alpha,\beta)\subset[0,1]$，$\exists\alpha_n$，使得
$\alpha_n\in(\alpha,\beta)$. 因为 $S_n$ 挨个地走过每个整数区间 $[k,k+1]$，且“步子”$S_n-S_{n-1}=a_n$ 无限变小. 这意味着它的小数部分 $\alpha_n$ 一遍又一遍地从左到右走过区间 $[0,1)$，向右的“步子”同样无限变小. 可见“步子”小到比指定区间 $(\alpha,\beta)$ 的长度还小时，就必有 $\alpha_n$ 落入 $(\alpha,\beta)$ 中.

\noindent\textbf{证}\quad 设 $(\alpha,\beta)\subset[0,1]$ 为任一小区间. 因为 $a_k\to0(k\to\infty)$，所以，$\exists N>0$，$k>N$ 时，
\[
  0<a_k<\beta-\alpha.
\qquad(1)
\]
取一 $n_0>N$，对 $S_{n_0}$，可取充分大的正整数 $m$，使得 $S_{n_0}<m+\alpha$. 因为 $S_n\to+\infty$，所以 $\exists n_1>n_0>N$，使得 $S_{n_1}>m+\beta(>m+\alpha>S_{n_0})$，于是必 $\exists n:n_0<n<n_1$ 使得 $S_n\in(m+\alpha,m+\beta)$，$S_n-[S_n]\in(\alpha,\beta)$. 否则，必存在某个 $k>N$，使得 $S_{k-1}<m+\alpha$，$S_k>m+\beta$，从而 $a_k=S_k-S_{k-1}>\beta-\alpha$，与式 (1) 矛盾. 证毕.

最后（不作重点）我们介绍一个有趣的应用.

\noindent\textbf{※例 5.1.45}\quad 试在 $[0,1]$ 上构造一个函数，使之在 $[0,1]$ 上单调，在有理点上间断，在无理点上连续.

\noindent\textbf{解}\quad $(0,1)$ 内全体有理点，可排成一个序列
\[
  \{x_n\}_{n=1}^{\infty}
  =\left\{\frac12,\frac13,\frac23,\frac14,\frac34,
  \frac15,\frac25,\frac35,\frac45,\cdots\right\}.
\qquad(1)
\]
设 $\displaystyle\sum_{n=1}^{\infty}c_n$ 是某一个正项收敛级数，$c_n>0$，$\forall x\in(0,1)$，定义
\[
  f(x)=\sum_{x_n<x}c_n.
\]
则 $f(x)$ 递增. 任意有理点 $\dfrac nm\in(0,1)$ 对应序列中某项 $x_{n_0}$，于是
\[
  f\left(\frac nm+0\right)-f\left(\frac nm-0\right)\ge c_{n_0}>0,
\]
故在有理点上间断. 若 $x_0$ 为无理点，则由 $\sum c_n$ 收敛，$\forall\varepsilon>0$，可取 $N$ 使 $\sum_{N+1}^{\infty}c_n<\varepsilon$. 再取 $\delta>0$ 使 $|x_i-x_0|\ge\delta\ (i=1,\cdots,N)$，则 $|x-x_0|<\delta$ 时
\[
  |f(x)-f(x_0)|
  \le\sum_{x_n\text{ 在 }x\text{ 与 }x_0\text{ 之间}}c_n
  \le\sum_{N+1}^{\infty}c_n<\varepsilon.
\]
故 $f(x)$ 在无理点上连续.

\section*{四、级数问题的若干反例}

\noindent\textbf{例 5.1.46}\quad 试写出一正项级数 $\displaystyle\sum_{n=1}^{\infty}a_n$，使得

1）$\displaystyle\sum_{n=1}^{\infty}a_n$ 收敛；\qquad
2）$a_n\ne o\left(\dfrac1n\right)$.（国外赛题）

\noindent\textbf{解}\quad 取
\[
  a_n=
  \begin{cases}
  \dfrac1n, & n\text{ 为整数平方数},\\[0.4em]
  \dfrac1{n^2}, & \text{否则}.
  \end{cases}
\]
显然 $a_n\ne o(1/n)$. 且
\[
  \sum a_n\le \sum_{n=1}^{\infty}\frac1{n^2}
  +\sum_{k=1}^{\infty}\frac1{k^2}<+\infty.
\]
故此级数收敛. 本例说明例 5.1.39 中单调性条件去掉之后，结论可能不成立.

\noindent\textbf{例 5.1.47}\quad 举出一个发散的交错级数，使其通项趋向零.（国外赛题）

\noindent\textbf{解}\quad 例如
\[
  1-\frac12+\frac1{3^2}-\frac14+\frac1{5^2}-\cdots
  +\frac1{(2n-1)^2}-\frac1{2n}+\cdots.
\]
其通项趋于零，但偶数部分为调和级数的一半，奇数平方部分收敛，故交错部分和发散. 本例说明 Leibniz 级数的三个条件中减少了单调性条件，定理就不再成立.

\noindent\textbf{例 5.1.48}\quad 举出一个收敛级数 $\displaystyle\sum a_n$，使得级数 $\displaystyle\sum a_n^3$ 发散.（国外赛题）

\noindent\textbf{解}\quad 令
\[
  \sum a_n=1-1+\frac1{\sqrt[3]2}-\frac1{2\sqrt[3]2}-\frac1{2\sqrt[3]2}
  +\cdots+\frac1{\sqrt[3]k}
  -\underbrace{\frac1{k\sqrt[3]k}-\cdots-\frac1{k\sqrt[3]k}}_{k\text{项}}+\cdots.
\]
每一组和为 $0$，故 $\sum a_n$ 收敛. 但
\[
  \sum a_n^3
  =1-1+\frac12-\frac1{2^3\cdot2}-\frac1{2^3\cdot2}
  +\cdots+\frac1k-\underbrace{\frac1{k^3\cdot k}-\cdots-\frac1{k^3\cdot k}}_{k\text{项}}+\cdots
\]
发散，因为部分和的子列
\[
  1+\frac12+\cdots+\frac1k-1-\frac1{2^3}-\cdots-\frac1{k^3}\to+\infty.
\]

\noindent\textbf{例 5.1.49}\quad 证明：若 $\forall |x_n|\to0(n\to\infty)$，有 $\displaystyle\sum a_nx_n$ 收敛，则 $\displaystyle\sum a_n$ 绝对收敛.

\noindent\textbf{证}\quad 反证. 若 $\displaystyle\sum |a_n|=+\infty$，则可取
$0=m_0<m_1<\cdots<m_n<\cdots$，使得
\[
  \sum_{i=m_{n-1}+1}^{m_n}|a_i|\ge n\qquad(n=1,2,\cdots).
\]
令
\[
  x_i=\frac{\operatorname{sgn}(a_i)}n
  \qquad(m_{n-1}<i\le m_n).
\]
则 $x_i\to0$，但每一片段
\[
  \sum_{i=m_{n-1}+1}^{m_n}a_ix_i
  =\sum_{i=m_{n-1}+1}^{m_n}\frac{|a_i|}{n}\ge1,
\]
故 $\sum a_nx_n$ 发散，矛盾.

\section*{五、数项级数与反常积分的关系}

\noindent\textbf{a.\quad 关于收敛性}

\noindent\textbf{要点}\quad 设 $a=A_0<A_1<\cdots<A_n<\cdots$，$A_n\to\infty$，$f(x)>0$，则
\[
  \int_a^{+\infty}f(x)\,dx
  \quad\text{与}\quad
  \sum_{n=1}^{\infty}\int_{A_{n-1}}^{A_n}f(x)\,dx
\]
同时敛散，并且收敛时两者大小相等.

\noindent\textbf{$\star$例 5.1.50}\quad 讨论
\[
  \int_0^{+\infty}\frac{dx}{1+x^\alpha|\sin x|^\beta}
\]
的收敛性，其中 $\alpha>\beta>1$.（复旦大学）

\noindent\textbf{证}\quad 取 $A_n=n\pi$，并用 $|\sin x|$ 的周期性，将积分分解为
\[
  \sum_{n=0}^{\infty}\int_0^{\pi/2}
  \frac{dt}{1+(n\pi+t)^\alpha\sin^\beta t}
  +\sum_{n=0}^{\infty}\int_{\pi/2}^{\pi}
  \frac{dt}{1+(n\pi+t)^\alpha\sin^\beta t}.
\]
当 $t\in[0,\pi/2]$ 时，$\sin t\ge\dfrac2\pi t$，故第一列积分
\[
  I_n\le\frac c{n^{\alpha/\beta}}.
\]
因 $\alpha/\beta>1$，$\sum I_n$ 收敛. 第二列积分作变换 $v=\pi-t$ 同理收敛，故原积分收敛.

\noindent\textbf{new 例 5.1.51}\quad 设 $f(x)$ 在 $[a,+\infty)$ 的内闭区间上 Riemann 可积 $(a>0)$，试证：反常积分 $\displaystyle\int_a^{+\infty}f(x)\,dx$ 绝对可积的充分必要条件是：对于任意满足 $x_0=a,\ x_n\to+\infty$ 的单调递增数列 $\{x_n\}$，级数
\[
  \sum_{n=0}^{+\infty}\int_{x_n}^{x_{n+1}}f(x)\,dx
\]
绝对收敛.（北京大学）

\noindent\textbf{证}\quad 必要性显然. 充分性可利用 Riemann 可积函数在小区间上取样，使每小段积分的符号与 $|f|$ 的积分一致；再在 $[a,A]$ 上取有限覆盖并逐段延伸，得到
\[
  \sum\left|\int_{x_k}^{x_{k+1}}f(x)\,dx\right|
  =\sum\int_{x_k}^{x_{k+1}}|f(x)|\,dx,
\]
从而推出 $\displaystyle\int_a^{+\infty}|f(x)|\,dx$ 收敛.

\noindent\textbf{练习}\quad 证明
\[
  \int_0^1\left|x\sin\frac1{x^2}-\frac1{x^3}\cos\frac1{x^2}\right|\,dx
\]
发散.

\noindent\textbf{证}\quad
\[
  \left|x\sin\frac1{x^2}-\frac1{x^3}\cos\frac1{x^2}\right|
  \ge\frac1{x^3}\left|\cos\frac1{x^2}\right|-x.
\]
右端第一项在适当积分小区间上可估计为
\[
  \frac14\sum_{k=1}^{\infty}\ln\left(1+\frac2{6k-1}\right)=+\infty,
\]
而 $\int_0^1x\,dx=\dfrac12$，故原积分发散.

\noindent\textbf{b.\quad “和”值的计算与估计}

\noindent\textbf{*例 5.1.52}\quad 求
\[
  \int_E e^{-x/2}\frac{|\sin x-\cos x|}{\sqrt{\sin x}}\,dx,
\]
其中 $E$ 为区间 $(0,+\infty)$ 中使被积表达式有意义的一切 $x$ 值所成之集合.

\noindent\textbf{解}\quad
\[
  E=\{x>0\mid 2k\pi<x<(2k+1)\pi,\ k\in\mathbb N\}.
\]
令 $x=t+2k\pi$，则
\[
  \int_E e^{-x/2}\frac{|\sin x-\cos x|}{\sqrt{\sin x}}\,dx
  =\sum_{k=0}^{\infty}e^{-k\pi}\int_0^\pi
  e^{-t/2}\frac{|\sin t-\cos t|}{\sqrt{\sin t}}\,dt.
\]
又
\[
  \int e^{-t/2}\frac{\sin t-\cos t}{\sqrt{\sin t}}\,dt
  =-2e^{-t/2}\sqrt{\sin t}+C,
\]
按 $[0,\pi/4]$ 与 $[\pi/4,\pi]$ 分段去绝对值，得
\[
  \int_E e^{-x/2}\frac{|\sin x-\cos x|}{\sqrt{\sin x}}\,dx
  =\frac{2\sqrt[4]8\,e^{-\pi/8}}{1-e^{-\pi}}.
\]

\noindent\textbf{*例 5.1.53}\quad 若 $f(x)>0,\ f(x)\downarrow$，$\displaystyle\sum f(n)$ 收敛，试证对余和 $R_n=\displaystyle\sum_{k=n+1}^{\infty}f(k)$ 有
\[
  \int_{n+1}^{+\infty}f(x)\,dx\le R_n\le f(n+1)+\int_{n+1}^{+\infty}f(x)\,dx.
\qquad(1)
\]

\noindent\textbf{证}\quad 因为 $f(x)\downarrow$，
\[
  \sum_{k=1}^{\infty}f(n+k)
  \ge\sum_{k=1}^{\infty}\int_{n+k}^{n+k+1}f(x)\,dx
  \ge\sum_{k=1}^{\infty}f(n+k+1).
\]
即 $R_n\ge\int_{n+1}^{+\infty}f(x)\,dx\ge R_n-f(n+1)$，移项即得.

\noindent\textbf{*例 5.1.54}\quad 试证：
\[
  \frac1{n\ln n}-\sum_{k=n}^{\infty}\frac1{k^2\ln k}
  \sim\frac1{n(\ln n)^2}\qquad(n\to\infty).
\]

\noindent\textbf{证}\quad 取 $f(x)=\dfrac1{x^2\ln x}$ 并应用上题，再对
$\displaystyle\int_n^{+\infty}\dfrac{dx}{x^2\ln x}$ 反复分部积分：
\[
  \int_n^{+\infty}\frac{dx}{x^2\ln x}
  =\frac1{n\ln n}-\frac1{n(\ln n)^2}
  +\frac{2\theta_n}{n(\ln n)^3}\qquad(0<\theta_n<1).
\]
代入夹逼估计即得结论.

\noindent\textbf{c.\quad 反常积分作为级数的极限}

\noindent\textbf{$\star$例 5.1.55}\quad 设单调函数 $f(x)$ 在 $x\ge0$ 有定义，并且反常积分 $\displaystyle\int_0^{+\infty}f(x)\,dx$ 存在，试证明
\[
  \lim_{h\to0+}h[f(h)+f(2h)+\cdots]
  =\int_0^{+\infty}f(x)\,dx.
\]

\noindent\textbf{证}\quad 例如 $f(x)\nearrow$，$\forall h>0$，有
\[
  \int_0^{+\infty}f(x)\,dx
  \le\sum_{k=1}^{\infty}hf(kh)
  \le\int_h^{+\infty}f(x)\,dx.
\]
令 $h\to0+$ 即得结论；$f(x)\downarrow$ 类似.

\noindent\textbf{*例 5.1.56}\quad 计算
\[
  \lim_{t\to1^-}(1-t)\left(\frac t{1+t}+\frac{t^2}{1+t^2}+\cdots\right).
\]

\noindent\textbf{解}\quad 令 $h=-\ln t$，则
\[
\begin{aligned}
\text{原式}
&=\lim_{h\to0+}\frac{1-e^{-h}}h\cdot
h\sum_{n=1}^{\infty}\frac{e^{-nh}}{1+e^{-nh}}\\
&=1\cdot\int_0^{+\infty}\frac{e^{-x}}{1+e^{-x}}\,dx=\ln2.
\end{aligned}
\]

\noindent\textbf{new 例 5.1.57}\quad 假设 $a_k>0,\ b_n>0$，
$\lim a_n=a$，$\displaystyle\sum b_n=1$. 证明：
\[
  \lim_{n\to\infty}\sum_{k=1}^n b_{n-k+1}a_k=a.
\]

\noindent\textbf{证}\quad 因 $\sum b_n=1$，只需证
\[
  \lim_{n\to\infty}\sum_{k=1}^n b_{n-k+1}(a_k-a)=0.
\]
设 $|a_k-a|\le M$. 由 $\sum b_n$ 的 Cauchy 准则，可使尾和
$\sum_{k=1}^{n-m}b_{n-k+1}<\varepsilon/(2M)$；又由 $a_k\to a$，当 $k$ 充分大时 $|a_k-a|<\varepsilon/2$. 将和式分成前 $n-m$ 项与后 $m$ 项，即得结论.

\noindent\textbf{new 练习}\quad 证明极限
\[
  \lim_{n\to\infty}\sum_{k=1}^{n-2}
  \frac1{2^{n-1-k}}\left[\ln\left(1+\frac1{2^{k+1}-1}\right)\right]=0.
\]

\noindent\textbf{证 I}\quad 直接引用例 5.1.57.

\noindent\textbf{证 II}\quad 由 $\sum 2^{-k}=1$ 的 Cauchy 准则和
$\ln\left(1+\dfrac1{2^{k+1}-1}\right)\to0$，将和式拆成两段 $I_1+I_2$，分别估计为 $<\varepsilon/2$，故极限为 $0$.

\noindent\textbf{单元练习 5.1}

\noindent\textbf{$\star$5.1.1}\quad 设 $k,i,j$ 都是自然数，且 $k=i+j$，试求级数
\[
  \sum_{n=1}^{\infty}\frac1{(kn-i)(kn+j)}
\]
的和. 提示：通项 $=\dfrac1k\left(\dfrac1{kn-i}-\dfrac1{kn+j}\right)$.

\noindent\textbf{$\star$5.1.2}\quad 设 $\{a_n\}$ 为等差数列，$a_{n+1}-a_n=d>0$，$m$ 为一正整数，计算
\[
  S=\sum_{n=1}^{\infty}\frac1{a_n a_{n+1}\cdots a_{n+m}}.
\]
提示：通项
\[
  =\frac1{md}\left(\frac1{a_n\cdots a_{n+m-1}}-\frac1{a_{n+1}\cdots a_{n+m}}\right),
\]
故 $S=\dfrac1{md\,a_1a_2\cdots a_m}$.

\noindent\textbf{$\star$5.1.3}\quad 证明级数
\[
  1+\frac1{\sqrt3}-\frac1{\sqrt2}+\frac1{\sqrt5}+\frac1{\sqrt7}
  -\frac1{\sqrt4}+\frac1{\sqrt9}+\frac1{\sqrt{11}}-\frac1{\sqrt6}+\cdots
\]
发散到 $+\infty$. 提示：
\[
  S_{3n}=\sum_{k=1}^n\left(\frac1{\sqrt{4k-3}}+\frac1{\sqrt{4k-1}}-\frac1{\sqrt{2k}}\right)\to+\infty.
\]

\noindent\textbf{$\star$5.1.4}\quad 证明：当 $p\ge1$ 时，
\[
  \sum_{n=1}^{\infty}\frac1{(n+1)\sqrt[p]{n}}<p.
\]
提示：通项 $\le p[n^{-1/p}-(n+1)^{-1/p}]$.

\noindent\textbf{※5.1.5}\quad 证明：若删去调和级数中所有分母含有数字 9 的项，则新级数收敛，且和小于 80.

\noindent\textbf{$\star$5.1.6}\quad 证明下列级数收敛：
\[
  1)\ \sum_{n=1}^{\infty}\left[\frac1n-\ln\left(1+\frac1n\right)\right];
  \qquad
  2)\ \sum_{n=1}^{\infty}\left[e-\left(1+\frac1{1!}+\cdots+\frac1{n!}\right)\right].
\]

\noindent\textbf{*5.1.7}\quad 设 $a_n=n^{n^\alpha}-1$，讨论级数 $\displaystyle\sum a_n$ 的敛散性. 提示：$\alpha<-1$ 时收敛，$\alpha\ge-1$ 时发散.

\noindent\textbf{$\star$5.1.8}\quad 设正项级数 $\displaystyle\sum a_n$ 收敛，证明级数
\[
  \sum_{n=1}^{\infty}\frac{a_n}{\sqrt{r_{n-1}}+\sqrt{r_n}}
\]
仍收敛，其中 $r_n=\displaystyle\sum_{k=n+1}^{\infty}a_k$. 提示：
\[
  \frac{a_n}{\sqrt{r_{n-1}}+\sqrt{r_n}}=\sqrt{r_{n-1}}-\sqrt{r_n}.
\]

\noindent\textbf{5.1.9}\quad 证明对数判别法：
\[
  \frac{\ln(1/a_n)}{\ln n}\ge1+\alpha\quad(a_n>0)
  \Rightarrow \sum a_n\text{ 收敛},
\]
若 $\dfrac{\ln(1/a_n)}{\ln n}\le1$，则级数发散.

\noindent\textbf{5.1.10}\quad $\{x_n\}$ 是正项单调递增并且有界的序列，证明级数
\[
  \sum_{n=1}^{\infty}\left(1-\frac{x_n}{x_{n+1}}\right)
\]
收敛.（国外赛题）

\noindent\textbf{提示}\quad 部分和
\[
  S_n\le \frac1{x_1}\sum_{i=1}^n(x_{i+1}-x_i).
\]

\noindent\textbf{※5.1.11}\quad （Lobachevskii 判别法）证明：若 $a_n>0,\ a_n\downarrow0$，则
\[
  \sum_{n=1}^{\infty}a_n\quad\text{与}\quad
  \sum_{m=1}^{\infty}p_m2^{-m}\quad
  (p_m=\max\{n\mid a_n\ge2^{-m}\})
\]
同时敛散.

\noindent\textbf{提示}\quad 注意 $p_m=\max\{n\mid a_n\ge2^{-m}\}$，$\{a_n\}$ 单调下降，表明
$\{a_1,a_2,\cdots,a_{p_m}\}$ 的每项都大于等于 $2^{-m}$；
同时 $\{a_{p_m+1},a_{p_m+2},\cdots\}$ 的每项都小于 $2^{-m}$.
由此，只需证明两点：
\[
  1)\ \sum_{k=1}^{\infty}a_k\text{ 与 }
  \sum_{m=2}^{\infty}(p_m-p_{m-1})2^{-m+1}\text{ 同时敛散};
\]
\[
  2)\ \sum_{m=2}^{\infty}(p_m-p_{m-1})2^{-m+1}\text{ 与 }
  \sum_{m=1}^{\infty}p_m2^{-m}\text{ 同时敛散}.
\]

\noindent\textbf{证}\quad 注意级数的敛散性与前面有限项无关. 正项级数的敛散性取决于部分和是否有界，因此正项级数“部分和序列”与其“子列”同时敛散.

$1^\circ$（证 1））i）（证明：
$\sum_{k=1}^{\infty}a_k$ 收敛 $\Rightarrow$
$\sum_{m=2}^{\infty}(p_m-p_{m-1})2^{-m+1}$ 收敛.）根据 $p_m$ 的定义，知
\[
  \sum_{k=p_{m-1}+1}^{p_m}a_k
  \ge \sum_{k=p_{m-1}+1}^{p_m}2^{-m}
  =(p_m-p_{m-1})2^{-m}\quad(m=2,3,\cdots).
\]
由此得
\[
  \sum_{k=1}^{p_N}a_k
  =\sum_{m=2}^N\sum_{k=p_{m-1}+1}^{p_m}a_k
  \ge\sum_{m=2}^N(p_m-p_{m-1})2^{-m}\ge0.
\qquad(1)
\]
又 $\sum_{k=1}^{\infty}a_k$ 收敛，而
$\left\{\sum_{k=1}^{p_N}a_k\right\}_N$ 是
$\left\{S_n\right\}=\left\{\sum_{k=1}^n a_k\right\}$ 的子列，故
$\left\{\sum_{k=1}^{p_N}a_k\right\}_N$ 收敛，由式 (1) 可知
$\sum_{m=2}^{\infty}(p_m-p_{m-1})2^{-m+1}$ 收敛.

ii）（证明：
$\sum_{m=2}^{\infty}(p_m-p_{m-1})2^{-m+1}$ 收敛 $\Rightarrow$
$\sum_{k=1}^{\infty}a_k$ 收敛.）由
\[
  \sum_{k=p_{m-1}+1}^{p_m}a_k
  <\sum_{k=p_{m-1}+1}^{p_m}2^{-m+1}
  =(p_m-p_{m-1})2^{-m+1},
\]
得
\[
  \sum_{k=1}^{p_N}a_k
  =\sum_{m=2}^N\sum_{k=p_{m-1}+1}^{p_m}a_k
  <\sum_{m=2}^N(p_m-p_{m-1})2^{-m+1}.
\]
因此，由 $\sum_{m=2}^{\infty}(p_m-p_{m-1})2^{-m+1}$ 收敛可知
$\left\{\sum_{k=1}^{p_N}a_k\right\}_N$ 收敛，从而
$\sum_{k=1}^{\infty}a_k$ 收敛.

$2^\circ$（证明 2）：$\sum_{m=1}^{\infty}p_m2^{-m}$ 收敛
$\Leftrightarrow\sum_{m=2}^{\infty}(p_m-p_{m-1})2^{-m+1}$ 收敛.）

$(\Rightarrow)$ 由于
\[
  \sum_{m=2}^N(p_m-p_{m-1})2^{-m+1}
  =2\sum_{m=2}^Np_m2^{-m}
    -\sum_{m=2}^Np_{m-1}2^{-(m-1)},
\]
自明.

$(\Leftarrow)$ 若 $\sum_{m=2}^{\infty}(p_m-p_{m-1})2^{-m+1}$ 收敛，则由下式可知
$\sum_{m=2}^{\infty}p_m2^{-m}$ 也收敛：
\begin{equation}
\begin{aligned}
\sum_{m=2}^Np_m2^{-m}
&=\sum_{m=2}^Np_m(2-1)2^{-m}\\
&=\left(\sum_{m=2}^Np_m2^{-m+1}\right)
  -\left(\sum_{m=2}^Np_m2^{-m}\right)\\
&=(p_22^{-1}+p_32^{-2}+\cdots+p_N2^{-N+1})\\
&\quad-(p_22^{-2}+p_32^{-3}+\cdots+p_{N-1}2^{-N+1}+p_N2^{-N})\\
&=p_22^{-1}+\sum_{m=3}^N(p_m-p_{m-1})2^{-m+1}+p_N2^{-N}.
\end{aligned}
\qquad(2)
\end{equation}
注意
\[
\begin{aligned}
p_N2^{-N}
&\le p_N2^{-(N-1)}\\
 \le (p_N-p_{N-1})2^{-(N-1)}+p_{N-1}2^{-(N-1)}\\
&\le (p_N-p_{N-1})2^{-(N-1)}
 +(p_{N-1}-p_{N-2})2^{-(N-2)}+p_{N-2}2^{-(N-2)}\\
&\le\cdots\le\sum_{m=2}^N(p_m-p_{m-1})2^{-(m-1)}+p_12^{-1}.
\end{aligned}
\]
故从式 (2) 可看出：若
$\sum_{m=2}^{\infty}(p_m-p_{m-1})2^{-m+1}$ 收敛，则
$\sum_{m=2}^{\infty}p_m2^{-m}$ 收敛. 证毕.

\noindent\textbf{5.1.12}\quad 设 $0<x_1<\pi,\ x_n=\sin x_{n-1}\ (n=2,3,\cdots)$，证明：级数
$\displaystyle\sum_{n=1}^{\infty}x_n^p$ 当 $p>2$ 时收敛，当 $p\le2$ 时发散.（吉林大学）

\noindent\textbf{提示}\quad 参看例 1.5.19，$x_n^2\sim\dfrac3n$.

\noindent\textbf{*5.1.13}\quad 证明级数
\[
  1+\frac12-\frac13+\frac14+\frac15-\frac16+\cdots
\]
发散.

\noindent\textbf{提示}\quad 可用 Cauchy 准则：
\[
  |S_{6n}-S_{3n}|\ge\frac13\left(\frac1{n+1}+\frac1{n+2}+\cdots+\frac1{2n}\right)\ge\frac16.
\]

\noindent\textbf{$\star$5.1.14}\quad 设 $a_n\ne0\ (n=1,2,\cdots)$ 且
$\lim_{n\to\infty}a_n=a\ (a\ne0)$，求证：级数
\[
  \sum_{n=1}^{\infty}|a_{n+1}-a_n|\quad\text{与}\quad
  \sum_{n=1}^{\infty}\left|\frac1{a_{n+1}}-\frac1{a_n}\right|
\]
同时收敛或同时发散.（上海交通大学）

\noindent\textbf{提示}\quad $\exists m,M:0<m<M$，使得
$m\le |a_n|\le M$，
\[
  \frac1{M^2}|a_{n+1}-a_n|
  \le \left|\frac1{a_{n+1}}-\frac1{a_n}\right|
  \le \frac1{m^2}|a_{n+1}-a_n|.
\]

\noindent\textbf{5.1.15}\quad 设 $\varphi(x)$ 是 $(-\infty,+\infty)$ 上的连续周期函数，周期为 $1$，且
$\displaystyle\int_0^1\varphi(x)\,dx=0$，$f(x)$ 在 $[0,1]$ 上可微，且有连续的一阶导数，
$a_n=\displaystyle\int_0^1f(x)\varphi(nx)\,dx,\ n=1,2,\cdots$. 证明：级数
$\displaystyle\sum_{n=1}^{\infty}a_n^2$ 收敛.（华东师范大学）

\noindent\textbf{提示}\quad 可令 $\Phi(x)=\displaystyle\int_0^x\varphi(t)\,dt$. 用分部积分，
$a_n=-\dfrac1n\displaystyle\int_0^1\Phi(nx)f'(x)\,dx$，从而
$|a_n^2|\le\dfrac{M}{n^2}$（$M$ 为常数）.

\noindent\textbf{5.1.16}\quad 设 $f(x)$ 于 $[1,+\infty)$ 上可导，$f'(x)$ 单调递增，且
$f(x)\to A(x\to+\infty)$，证明 $\displaystyle\sum_{n=2}^{\infty}f'(n)$ 收敛.

\noindent\textbf{提示}\quad 由 $f(n)-f(n-1)\le f'(n)\le f(n+1)-f(n)$，可知
$f'(n)\nearrow0,\ f'(n)\le0$，$\sum-f'(n)$ 以 $\sum[f(n-1)-f(n)]$ 为优级数.

\noindent\textbf{5.1.17}\quad 设 $a_n>0(n=1,2,\cdots)$ 且
$\displaystyle\sum_{n=1}^{\infty}a_n$ 收敛，$r_n=\displaystyle\sum_{k=n}^{\infty}a_k$. 试证：
\[
  1)\ \sum_{n=1}^{\infty}\frac{a_n}{r_n}\text{ 发散};\qquad
  2)\ \sum_{n=1}^{\infty}\frac{a_n}{\sqrt{r_n}}\text{ 收敛}.
\]

\noindent\textbf{提示}\quad 1）可用 Cauchy 准则，
\[
  \sum_{k=n+1}^{n+p}\frac{a_k}{r_k}
  \ge \frac1{r_{n+1}}\sum_{k=n+1}^{n+p}a_k\to1\ (p\to\infty),
\]
故 $\exists p\in\mathbb N$ 使得 $\displaystyle\sum_{k=n+1}^{n+p}\dfrac{a_k}{r_k}\ge\dfrac12$.

2）可证 $\displaystyle\sum\frac{a_n}{\sqrt{r_n}}$ 以
$2\sum(\sqrt{r_{n-1}}-\sqrt{r_n})$ 为优级数.

\noindent\textbf{再提示}\quad 2）
\[
  \sqrt{r_{n-1}}-\sqrt{r_n}
  =\frac1{2\sqrt{\xi_{n-1}}}(r_{n-1}-r_n)
  =\frac{a_n}{2\sqrt{\xi_{n-1}}}\ge\frac{a_n}{2\sqrt{r_n}}
  \quad(r_n<\xi_{n-1}<r_{n-1}).
\]
又因 $r_n\to0$，
$\sum_{i=1}^n(\sqrt{r_{n-1}}-\sqrt{r_n})=\sqrt{r_0}-\sqrt{r_n}\to\sqrt{r_0}$，故
$\sum\dfrac{a_n}{\sqrt{r_n}}$ 收敛.

\noindent\textbf{5.1.18}\quad 设 $f(x)$ 是在 $(-\infty,+\infty)$ 内的可微函数，且满足：
\[
  1)\ f(x)>0;\qquad 2)\ |f'(x)|\le m|f(x)|,\quad 0<m<1.
\]
任取 $a_0$，定义 $a_n=\ln f(a_{n-1}),\ n=1,2,\cdots$，证明：级数
\[
  \sum_{n=1}^{\infty}(a_n-a_{n-1})
\]
绝对收敛.（西安电子科技大学）

\noindent\textbf{提示}\quad
\[
  |a_n-a_{n-1}|
  =|\ln f(a_{n-1})-\ln f(a_{n-2})|
  =\left|\frac{f'(\xi_n)}{f(\xi_n)}(a_{n-1}-a_{n-2})\right|,
\]
\[
  \frac{|a_n-a_{n-1}|}{|a_{n-1}-a_{n-2}|}
  \le \left|\frac{f'(\xi_n)}{f(\xi_n)}\right|\le m<1,
\]
故 $\sum|a_n-a_{n-1}|$ 收敛.

\noindent\textbf{*5.1.19}\quad 设 $\displaystyle\sum_{n=1}^{\infty}a_n$ 收敛，$0<p_n\nearrow+\infty$，试证：
\[
  \lim_{n\to\infty}\frac{p_1a_1+p_2a_2+\cdots+p_na_n}{p_n}=0.
\]

\noindent\textbf{提示}\quad 记 $S_n=\displaystyle\sum_{k=1}^n a_k$，$a_n=S_n-S_{n-1}$，代入变形可知
\[
\begin{aligned}
\text{原式}
&=\lim_{n\to\infty}
\frac{S_1(p_1-p_2)+S_2(p_2-p_3)+\cdots+S_{n-1}(p_{n-1}-p_n)+S_np_n}{p_n}\\
&\stackrel{\text{Stolz 公式}}{=}
\lim_{n\to\infty}\frac{S_{n-1}(p_{n-1}-p_n)}{p_n-p_{n-1}}+S
=S-S=0.
\end{aligned}
\]

\noindent\textbf{※5.1.20}\quad 设 $a_n>0,\ \displaystyle\sum_{n=1}^{\infty}a_n$ 收敛，$na_n$ 单调，证明：
\[
  \lim_{n\to\infty}na_n\ln n=0.
\]

\noindent\textbf{提示}\quad 因 $a_n>0,\ \sum_{n=1}^{\infty}a_n$ 收敛，$na_n$ 单调，必为单减（否则：
$a_n\ge\dfrac{a_1}{n}>0,\ \sum_{n=1}^{\infty}a_n$ 发散）.

\noindent\textbf{再提示}\quad （本习题是例 4.5.27 的离散形式，证法类似.）设 $n$ 是任一正整数，$[\sqrt n]$ 表示“$\sqrt n$ 的整数部分”. 因 $|na_n|\downarrow$，有
\[
\begin{aligned}
\sum_{k=[\sqrt n]}^n a_k
&=\sum_{k=[\sqrt n]}^n ka_k\cdot\frac1k
\ge na_n\sum_{k=[\sqrt n]}^n\frac1k\\
&\ge na_n\sum_{k=[\sqrt n]}^{n-1}\int_k^{k+1}\frac1x\,dx
\ge na_n\int_{[\sqrt n]}^n\frac{dx}{x}\\
&=na_n\ln\frac{n}{[\sqrt n]}
\ge na_n\ln\sqrt n=\frac12na_n\ln n.
\end{aligned}
\]
因 $\sum_{n=1}^{\infty}a_n$ 收敛，根据 Cauchy 准则，$\forall\varepsilon>0$，$\exists N\in\mathbb N$，当 $n>N$ 时，
$\displaystyle\sum_{k=n}^{n+p}a_k<\dfrac{\varepsilon}{2}$. 因此只要取 $n$ 充分大，就有 $[\sqrt n]>N$，
进而 $0<na_n\ln n\le2\displaystyle\sum_{k=[\sqrt n]}^na_k<\varepsilon$. 说明
$\lim_{n\to\infty}na_n\ln n=0$.

\noindent\textbf{注}\quad 下面的证法十分简捷，正确吗？

“利用 Cauchy 积分判别法，知级数
$\displaystyle\sum_{n=2}^{\infty}\frac1{n\ln n}$ 发散. 现知正项级数
$\displaystyle\sum_{n=1}^{\infty}a_n$ 收敛，因此 $a_n$ 是
$\dfrac1{n\ln n}$ 的高阶无穷小量（当 $n\to\infty$ 时），故
$\displaystyle\lim_{n\to\infty}na_n\ln n=\lim_{n\to\infty}\dfrac{a_n}{1/(n\ln n)}=0$.”

回答：此证法不对. 因为：“两正项级数：
$\displaystyle\sum_{n=1}^{\infty}a_n$ 收敛，
$\displaystyle\sum_{n=1}^{\infty}b_n$ 发散（$\lim_{n\to\infty}b_n=0$），并不能断言
$\{a_n\}$ 是比 $\{b_n\}$ 高阶的无穷小量”. 例如：
\[
  \sum_{n=1}^{\infty}a_n=\sum_{n=1}^{\infty}\frac1{n^2}\text{（收敛）},\quad
  \sum_{n=1}^{\infty}b_n=1+\frac12+\frac1{3^2}+\frac14+\frac1{5^2}+\frac16+\frac1{7^2}+\cdots\text{（发散）},
\]
（因 $\displaystyle\sum_{k=1}^{\infty}b_{2k}
=\sum_{k=1}^{\infty}\dfrac1{2k}=\dfrac12\sum_{k=1}^{\infty}\dfrac1k=+\infty$，而
$\sum_{n=1}^{\infty}b_n$（正项级数）项数比 $\sum_{k=1}^{\infty}b_{2k}$ 多，故
$\sum_{n=1}^{\infty}b_n$ 发散）. 但是
\[
  \frac{a_n}{b_n}=
  \begin{cases}
  1, & n\text{ 为奇数},\\
  \dfrac1n, & n\text{ 为偶数},
  \end{cases}
\]
极限 $\displaystyle\lim_{n\to\infty}\frac{a_n}{b_n}$ 不存在.

\noindent\textbf{※5.1.21}\quad 设数 $a>0$，$\{p_n\}$ 是一个数列，并且
$p_n>0,\ p_{n+1}\ge p_n$，证明：级数
\[
  \sum_{n=1}^{\infty}\frac{p_n-p_{n-1}}{p_np_{n-1}^{\,a}}
\]
收敛.（国外赛题）

\noindent\textbf{提示}\quad 将通项 $\dfrac{p_n-p_{n-1}}{p_np_{n-1}^{\,a}}$ 放大为
$\dfrac1{p_1^{a-1}}\dfrac{p_n-p_{n-1}}{p_np_{n-1}}$（当 $a\ge1$ 时）或
$\dfrac1a\dfrac{p_n^a-p_{n-1}^a}{p_n^ap_{n-1}^a}$（当 $a<1$ 时），然后用连锁消去法（见例 5.1.3 至例 5.1.5）证明级数收敛.

\noindent\textbf{证}\quad $1^\circ$ 当 $a\ge1$ 时，因
$\dfrac1{p_{k-1}^{a-1}}\le\dfrac1{p_1^{a-1}}\ (k\ge2)$，
\[
\begin{aligned}
0<\sum_{k=2}^n\frac{p_k-p_{k-1}}{p_kp_{k-1}^{\,a}}
&=\sum_{k=2}^n\frac1{p_{k-1}^{a-1}}\cdot\frac{p_k-p_{k-1}}{p_kp_{k-1}}\\
&\le \frac1{p_1^{a-1}}\sum_{k=2}^n\frac{p_k-p_{k-1}}{p_kp_{k-1}}\\
&=\frac1{p_1^{a-1}}\sum_{k=2}^n\left(\frac1{p_{k-1}}-\frac1{p_k}\right)
=\frac1{p_1^{a-1}}\left(\frac1{p_1}-\frac1{p_n}\right).
\end{aligned}
\]
因为 $\dfrac1{p_n}\downarrow$，有下界 $0$，故级数收敛.

$2^\circ$ 当 $a<1$ 时，
\[
  0<\sum_{k=2}^n\frac{p_k-p_{k-1}}{p_kp_{k-1}^{\,a}}
  =\sum_{k=2}^n\left(\frac{p_k-p_{k-1}}{p_kp_{k-1}^{\,a}}\cdot
  \frac{p_k^a-p_{k-1}^a}{p_k^a-p_{k-1}^a}\right)
  \cdot\frac{p_k^a-p_{k-1}^a}{p_k^ap_{k-1}^{\,a}}.
\qquad(1)
\]
应用 Lagrange 定理：$\exists\xi:p_{k-1}<\xi<p_k$，使得
$p_k^a-p_{k-1}^a=a\xi^{a-1}(p_k-p_{k-1})$，得
\[
  \frac1{p_k^{1-a}}\cdot\frac{p_k-p_{k-1}}{p_k^a-p_{k-1}^a}
  =\frac1{p_k^{1-a}}\cdot
  \frac{p_k-p_{k-1}}{a\xi^{a-1}(p_k-p_{k-1})}
  \le\frac1{p_k^{1-a}}\cdot\frac1{ap_k^{a-1}}=\frac1a.
\qquad(2)
\]
由式 (1) 得
\[
  0<\sum_{k=2}^n\frac{p_k-p_{k-1}}{p_kp_{k-1}^{\,a}}
  \le\frac1a\sum_{k=2}^n\frac{p_k^a-p_{k-1}^a}{p_k^ap_{k-1}^{\,a}}
  =\frac1a\left(\frac1{p_1^a}-\frac1{p_n^a}\right),
\]
故级数收敛. 证毕.

\noindent\textbf{5.1.22}\quad 举出一个收敛级数 $\displaystyle\sum_{n=1}^{\infty}a_n$ 的例子，使级数
$\displaystyle\sum_{n=1}^{\infty}a_n\ln n$ 发散.

\noindent\textbf{提示}\quad 例如 $a_n=\dfrac1{n\ln n(\ln\ln n)^2}$.

\noindent\textbf{再提示}\quad 用积分判别法易证该级数 $\displaystyle\sum_{n=1}^{\infty}a_n$ 收敛. 又因 $n$ 充分大时，
$\dfrac1{n(\ln\ln n)^2}>\dfrac1{n\ln n}$，而
$\displaystyle\sum_{n=1}^{\infty}\dfrac1{n\ln n}$ 发散，故
\[
  \sum_{n=1}^{\infty}a_n\ln n
  =\sum_{n=1}^{\infty}\frac1{n(\ln\ln n)^2}
\]
发散.

\noindent\textbf{5.1.23}\quad 序列 $\{b_n\}(n=1,2,\cdots)$ 具有下列性质：
$b_n>0,\ \lim_{n\to\infty}b_n=+\infty$. 作出序列 $\{a_n\}$，使
\[
  a_n\ge0,\qquad \sum_{n=1}^{\infty}a_n<\infty,\qquad
  \sum_{n=1}^{\infty}a_nb_n=+\infty.
\]
（国外赛题）

\noindent\textbf{提示}\quad $\forall k\in\mathbb N$，因为
$\lim_{n\to\infty}b_n=+\infty$，$\exists n_k\in\mathbb N$ 使
$b_{n_k}>k\ (k=1,2,\cdots)$. 顺次可使 $b_{n_k}>b_{n_{k-1}}$. 令
\[
  a_n=
  \begin{cases}
  \dfrac1{k^2}, & n=n_k,\\
  0, & n\ne n_k
  \end{cases}
\]
即可.

\noindent\textbf{5.1.24}\quad 设 $\{n_k\}$ 是自然数列 $\{n\}$ 的子列，试证：
\[
  1)\ \text{当 }n_k-n_{k-1}\ge k\text{ 时，}\sum_{k=1}^{\infty}\frac1{n_k}\text{ 收敛};
\]
\[
  2)\ \text{当 }n_k-n_{k-1}\le g\text{（常数）时，}\sum_{k=1}^{\infty}\frac1{n_k}\text{ 发散};
\]
\[
  ※3)\ \text{当 }n_k-n_{k-1}\ge k^r\ (r>0)\text{ 时，}\sum_{k=1}^{\infty}\frac1{n_k}\text{ 收敛}.
\]

\noindent\textbf{提示}\quad 1）$n_k\ge k+(k-1)+\cdots+1=\dfrac{k}{2}(1+k)$，
$0<\dfrac1{n_k}<\dfrac2{k^2}$.
\[
  2)\ \sum_{k=1}^{\infty}\frac1{n_k}
  \ge\sum_{k=1}^{\infty}\frac1{(k-1)g+n_1}=+\infty.
\]

\noindent\textbf{$\star$5.1.25}\quad 对函数
$\displaystyle\zeta(s)=\sum_{n=1}^{\infty}\frac1{n^s}\ (s>1)$，证明：
\[
  \zeta(s)=s\int_1^{+\infty}\frac{[x]}{x^{s+1}}\,dx,
\]
其中 $[x]$ 为 $x$ 的整数部分.（西北师范大学）

\noindent\textbf{提示}\quad
\[
  s\int_1^{+\infty}\frac{[x]}{x^{s+1}}\,dx
  =\sum_{n=1}^{\infty}s\int_n^{n+1}\frac{n}{x^{s+1}}\,dx
  =\sum_{n=1}^{\infty}\left[\frac n{n^s}-\frac n{(n+1)^s}\right].
\]

\noindent\textbf{再提示}\quad 上式
\[
  =\sum_{n=1}^{\infty}\frac1{n^{s-1}}
   -\sum_{n=1}^{\infty}\frac1{(n+1)^{s-1}}
   +\sum_{n=1}^{\infty}\frac1{(n+1)^s}
  =1+\sum_{n=2}^{\infty}\frac1{n^s}
  =\sum_{n=1}^{\infty}\frac1{n^s}.
\]

\noindent\textbf{$\star$5.1.26}\quad 1）求证：当 $s>0$ 时，
$\displaystyle\int_1^{+\infty}\frac{x-[x]}{x^{s+1}}\,dx$ 收敛；

2）求证：当 $s>1$ 时，
\[
  \int_1^{+\infty}\frac{x-[x]}{x^{s+1}}\,dx
  =\frac1{s-1}-\frac1s\sum_{n=1}^{\infty}\frac1{n^s}.
\]
$[x]$ 表示 $x$ 的整数部分.（北京航空航天大学）

\noindent\textbf{提示}\quad 利用上题结果.

\noindent\textbf{5.1.27}\quad 求
\[
  \lim_{t\to+\infty}\left(\frac1t+\frac{2t}{t^2+1^2}+\frac{2t}{t^2+2^2}+\cdots+\frac{2t}{t^2+n^2}+\cdots\right).
\]
《$\pi$》

\noindent\textbf{提示}\quad 可参看例 5.1.55 和例 5.1.56.

\noindent\textbf{$\star$5.1.28}\quad 设 $k>0,\ a>0$，证明：
\[
  1)\ \int_a^{+\infty}\frac{\sin 2n\pi x}{x^k}\,dx\text{ 收敛};\qquad
  2)\ \sum_{n=1}^{\infty}\frac1n\left|\int_a^{+\infty}\frac{\sin 2n\pi x}{x^k}\,dx\right|\text{ 收敛}.
\]
（上海交通大学）

\noindent\textbf{提示}\quad 1）可用反常积分的 Dirichlet 判别法；2）可先用第二中值定理.

\noindent\textbf{再提示}\quad 1）当 $x\to+\infty$ 时，$\dfrac1{x^k}\downarrow0$，而
\[
  \left|\int_a^A\sin2n\pi x\,dx\right|\le\frac1{n\pi}\quad(\forall A>a),
\]
故原积分收敛.

2）
\[
\begin{aligned}
\left|\int_a^A\frac{\sin2n\pi x}{x^k}\,dx\right|
&=\left|\frac1{a^k}\int_a^\xi\sin2n\pi x\,dx+
\frac1{A^k}\int_\xi^A\sin2n\pi x\,dx\right|\\
&\le \frac1{n\pi a^k}+\frac1{n\pi A^k}.
\end{aligned}
\]
不等式两边同时取极限 $(A\to+\infty)$，知
\[
  \sum_{n=1}^{\infty}\frac1n\left|\int_a^{+\infty}\frac{\sin2n\pi x}{x^k}\,dx\right|
  \le \frac1{\pi a^k}\sum_{n=1}^{\infty}\frac1{n^2}<+\infty.
\]

\noindent\textbf{*5.1.29}\quad 证明：
\[
  \lim_{n\to\infty}\left\{\sum_{k=2}^n\frac1{k\ln k}-\ln\ln n\right\}
\]
存在（有限）.（北京师范大学）

\noindent\textbf{提示}\quad 可参看例 5.1.54.

\noindent\textbf{new $\star$5.1.30}\quad 已知 $a_n>0$，级数
$\displaystyle\sum_{n=1}^{\infty}\frac1{a_n}$ 发散. 试证：级数
$\displaystyle\sum_{n=1}^{\infty}\frac1{a_n+1}$ 亦发散.（中国科学院）

\noindent\textbf{提示}\quad （反证法）若 $\displaystyle\sum_{n=1}^{\infty}\frac1{a_n+1}$ 收敛，则 $n\to\infty$ 时，
$\dfrac1{a_n+1}\to0$，即 $a_n+1\to+\infty$，得 $a_n\to+\infty$，故
\[
  \frac{\dfrac1{a_n}}{\dfrac1{a_n+1}}=\frac{a_n+1}{a_n}=1+\frac1{a_n}\to1.
\]
因此，$\forall\varepsilon>0,\ \exists N>0$，当 $n>N$ 时：
$0<\dfrac1{a_n}<(1+\varepsilon)\dfrac1{a_n+1}$. 于是
$\displaystyle\sum_{n=1}^{\infty}\frac1{a_n}$ 也收敛. 与题设矛盾.

\noindent\textbf{new *5.1.31}\quad 设 $u(x)$ 在 $[0,+\infty)$ 上有连续导数，且
\[
  \int_0^{+\infty}\left(|u(x)|^2+|u'(x)|^2\right)\,dx=M<+\infty.
\]
试证：1）$\exists\{x_n\}_{n=1}^{\infty}\subset[0,+\infty)$ 使得
$\{x_n\}_{n=1}^{\infty}\to+\infty,\ \lim_{n\to\infty}u(x_n)=0$；

2）$\exists C>0$，使得
\[
  \sup_{x\in[0,+\infty)}u(x)\le C\left[\int_0^{+\infty}\left(|u(x)|^2+|u'(x)|^2\right)\,dx\right]^{1/2}.
\]
（华中师范大学）

\noindent\textbf{提示}\quad 1）$\forall k,\ \exists x_k\in[k,k+1]$，使得
\[
  \int_k^{k+1}u^2(x)\,dx
  \stackrel{\text{积分中值定理}}{=}u^2(x_k),
\]
故
\[
  \sum_{k=0}^n|u(x_k)|^2
  =\sum_{k=0}^n\int_k^{k+1}|u(x)|^2\,dx
  \le\int_0^{+\infty}|u(x)|^2\,dx\le M.
\]
因此 $\sum_{k=0}^{\infty}|u(x_k)|^2$ 收敛，$\lim_{k\to+\infty}|u(x_k)|=0$.

2）（利用均值不等式）
\[
  |u(x)u'(x)|\le\frac{|u(x)|^2+|u'(x)|^2}{2}.
\qquad(1)
\]
\begin{equation}
\begin{aligned}
|u^2(x_2)-u^2(x_1)|
&=\left|\int_{x_1}^{x_2}d(u^2(x))\right|
 =\left|\int_{x_1}^{x_2}2u(x)u'(x)\,dx\right|\\
&\le\int_{x_1}^{x_2}\left(|u(x)|^2+|u'(x)|^2\right)\,dx.
\end{aligned}
\qquad(2)
\end{equation}

\noindent\textbf{再提示}\quad 2）因题设：
$\displaystyle\int_0^{+\infty}\left(|u(x)|^2+|u'(x)|^2\right)\,dx$ 收敛，根据 Cauchy 准则，
$\forall\varepsilon>0,\ \exists A>0$，当 $x_2>x_1>A$ 时有
\[
  \int_{x_1}^{x_2}\left(|u(x)|^2+|u'(x)|^2\right)\,dx<\varepsilon.
\]
因此
\[
  |u^2(x_2)-u^2(x_1)|\stackrel{(2)}{\le}
  \int_{x_1}^{x_2}\left(|u(x)|^2+|u'(x)|^2\right)\,dx<\varepsilon.
\]
再由 Cauchy 准则的充分性，即知 $\lim_{x\to+\infty}u^2(x)$ 存在. 加之 $u(x)$ 连续，故
$u^2(x)$ 在 $[0,+\infty)$ 上有界，$u(x)$ 同样如此，即：$\exists C>0$，可使得
\[
  \sup_{0\le x<+\infty}|u(x)|
  \le C M^{1/2}
  =C\left[\int_0^{+\infty}\left(|u(x)|^2+|u'(x)|^2\right)\,dx\right]^{1/2}.
\]

\noindent\textbf{new **5.1.32}\quad 设 $a_1=a>0,\ a_2=b>0$，
\[
  a_{n+2}=2+\frac1{a_{n+1}^2}+\frac1{a_n^2}.
\qquad(1)
\]
试证：$\{a_n\}$ 收敛.（华东师范大学）

\noindent\textbf{提示}\quad $1^\circ$ 由式 (1) 知：$\forall k\ge3$，有 $a_k>2$. 因此当 $n\ge5$ 时，
$2<a_n\le\dfrac52$. 故 $\forall m,n\ge5$，由
\[
  |a_m-a_n|<\frac52-2=\frac12
\qquad(2)
\]
及
\[
  0<\frac{a_{n-1}+a_{n-3}}{a_{n-1}^2a_{n-3}^2}<\frac5{2^4}
\qquad(3)
\]
知（当 $n\ge8$ 时）
\begin{equation}
\begin{aligned}
|a_n-a_{n-1}|
&=\left|\frac1{a_{n-1}^2}+\frac1{a_{n-2}^2}
 -\left(\frac1{a_{n-2}^2}+\frac1{a_{n-3}^2}\right)\right|\\
&=\left|\frac1{a_{n-1}^2}-\frac1{a_{n-3}^2}\right|
=\frac{a_{n-1}+a_{n-3}}{a_{n-1}^2a_{n-3}^2}|a_{n-1}-a_{n-3}|
<\frac5{2^5}.
\end{aligned}
\qquad(4)
\end{equation}
当 $n\ge11$ 时，
\[
  |a_n-a_{n-1}|=\frac{a_{n-1}+a_{n-3}}{a_{n-1}^2a_{n-3}^2}|a_{n-1}-a_{n-3}|
  \stackrel{\text{式}(3)}{\le}\frac5{2^4}\left(|a_{n-1}-a_{n-2}|+|a_{n-2}-a_{n-3}|\right).
\]
因此有迭代关系：
\[
  |a_n-a_{n-1}|\le\frac5{2^3}\max\{|a_{n-1}-a_{n-2}|,\ |a_{n-2}-a_{n-3}|\}.
\]
反复迭代下去，当 $n=2k+8$ 或 $n=2k+9$ 时，至少迭代 $k$ 次，有
\[
  |a_n-a_{n-1}|\le\left(\frac58\right)^k
  \max\{|a_8-a_7|,\ |a_7-a_6|\}\stackrel{\text{记}}{=}aq^k
  \quad\left(q=\frac58,\ a>0\right).
\]
于是
\[
  0\le\sum_{k=8}^n(a_k-a_{k-1})
  \le\sum_{k=8}^n|a_k-a_{k-1}|
  \le a\sum_{k=8}^nq^k\quad\left(q=\frac58<1\right).
\]
因 $\sum_{k=8}^{\infty}q^k$ 收敛，上式表明
$\sum_{k=8}^{\infty}(a_k-a_{k-1})$ 绝对收敛，可知
$\{a_n\}=\left\{a_7+\sum_{k=8}^n(a_k-a_{k-1})\right\}$ 收敛. 证毕.

\section*{* §5.2\quad 函数项级数}

所谓函数项级数 $\displaystyle\sum_{n=1}^{\infty}u_n(x)$ 在某区间 $I$ 上收敛，是指它逐点收敛. 意即：对每一固定点
$x\in I$，作为数项级数，$\displaystyle\sum_{n=1}^{\infty}u_n(x)$ 总是收敛的. 因此对收敛性，可用上一节数项级数的各种判别法进行判断. 本节的任务，主要讨论一致收敛性的判断及其应用.

\noindent\textbf{导读}\quad 函数项级数（及序列）一致收敛问题是数学院系学生的难点、重点，也是各类考试的热点，非数学院系的学生从略.

\subsection*{一、一致收敛性的判断}

证明一致收敛性一般有如下几种方法：a）利用定义；b）利用 Cauchy 准则；c）利用常用的几个判别法；d）利用一致有界与等度连续. 下面我们对这些分别进行介绍和讨论.

\noindent\textbf{a.\quad 利用定义证明一致收敛性}

\noindent\textbf{要点}\quad $1^\circ\ \varepsilon-N$ 方法.

i）要用定义证明 $\displaystyle\sum_{n=1}^{\infty}u_n(x)$ 在区间 $I$ 上一致收敛，应首先设法求出和函数
\[
  S(x)=\sum_{n=1}^{\infty}u_n(x),
\]
写出部分和 $S_n(x)=\displaystyle\sum_{k=1}^n u_k(x)$，然后对任意给定的 $\varepsilon>0$，找出与 $x$ 无关的
$N=N(\varepsilon)$，使得 $n>N$ 时有 $|S(x)-S_n(x)|<\varepsilon$.

ii）$S_n(x)\not\rightrightarrows S(x)$（$n\to\infty$ 时关于 $x\in I$），等价于：
$\exists\varepsilon_0>0,\ \forall N>0,\ \exists n>N,\ \exists x_N\in I$ 使得
$|S(x_N)-S_n(x_N)|\ge\varepsilon_0$.

亦等价于：$\exists\varepsilon_0>0,\ \exists\{x_n\}\subset I$，使得
$|S(x_n)-S_n(x_n)|\ge\varepsilon_0$.

特别来讲，若发现有 $x_0\in I$，或为 $I$ 的端点，使得 $x\to x_0$ 时有
$S(x)-S_n(x)\not\to0$（对充分大的 $n$ 成立），则在区间 $I$ 上
$S_n(x)\not\rightrightarrows S(x)$（当 $n\to\infty$ 时）.

$2^\circ$ “放大法”：若 $\forall n,\ \exists\alpha_n>0$，使得
$|S(x)-S_n(x)|\le\alpha_n\ (\forall x\in I)$，且 $n\to\infty$ 时 $\alpha_n\to0$，则当 $n\to\infty$ 时，
$S_n(x)\rightrightarrows S(x)$（于 $I$ 上）.

$3^\circ$ 确界法. 当 $n\to\infty$ 时，$S_n(x)\rightrightarrows S(x)$ 等价于
\[
  \lim_{n\to\infty}\sup_{x\in I}|S(x)-S_n(x)|=0.
\]

\noindent\textbf{$\varepsilon-N$ 方法}

\noindent\textbf{$\star$例 5.2.1}\quad $f(x)$ 是 $(-\infty,+\infty)$ 上的连续函数，
$f_n(x)=\displaystyle\sum_{k=0}^{n-1}\dfrac1n f\left(x+\dfrac{k}{n}\right)$，证明：函数序列
$\{f_n(x)\}(n=1,2,3,\cdots)$ 在任何有限区间上一致收敛.（首都师范大学，北京师范大学）

\noindent\textbf{分析}\quad 我们看到
$f_n(x)=\displaystyle\sum_{k=0}^{n-1}\dfrac1n f\left(x+\dfrac{k}{n}\right)$ 正好是积分
$\displaystyle\int_0^1f(x+t)\,dt$ 的一个积分和，因为 $f(x)$ 连续，该积分有意义，故当
$n\to\infty$ 时，$f_n(x)\to\displaystyle\int_0^1f(x+t)\,dt$.

设 $[a,b]$ 是任意一个有限区间，要证明当 $n\to\infty$ 时，
$f_n(x)\rightrightarrows\displaystyle\int_0^1f(x+t)\,dt$ 于 $[a,b]$ 上，即对任一
$\varepsilon>0$，要找 $N>0$，使得当 $n>N$ 时，
\[
  \left|f_n(x)-\int_0^1f(x+t)\,dt\right|<\varepsilon\quad(\forall x\in[a,b]).
\qquad(1)
\]
因为
\[
  f_n(x)=\sum_{k=0}^{n-1}\frac1n f\left(x+\frac{k}{n}\right)
  =\sum_{k=0}^{n-1}\int_{k/n}^{(k+1)/n}f\left(x+\frac{k}{n}\right)\,dt,
\]
\[
  \int_0^1f(x+t)\,dt
  =\sum_{k=0}^{n-1}\int_{k/n}^{(k+1)/n}f(x+t)\,dt,
\]
所以
\begin{equation}
\begin{aligned}
\left|f_n(x)-\int_0^1f(x+t)\,dt\right|
&=\left|\sum_{k=0}^{n-1}\int_{k/n}^{(k+1)/n}f\left(x+\frac{k}{n}\right)\,dt
 -\sum_{k=0}^{n-1}\int_{k/n}^{(k+1)/n}f(x+t)\,dt\right|\\
&=\left|\sum_{k=0}^{n-1}\int_{k/n}^{(k+1)/n}
\left(f\left(x+\frac{k}{n}\right)-f(x+t)\right)\,dt\right|\\
&\le\sum_{k=0}^{n-1}\int_{k/n}^{(k+1)/n}
\left|f\left(x+\frac{k}{n}\right)-f(x+t)\right|\,dt.
\end{aligned}
\qquad(2)
\end{equation}
故要式 (1) 成立，只需式 (2) 右端的
\[
  \sum_{k=0}^{n-1}\int_{k/n}^{(k+1)/n}
  \left|f\left(x+\frac{k}{n}\right)-f(x+t)\right|\,dt<\varepsilon
  \quad(\forall x\in[a,b]).
\qquad(3)
\]
为此只要能使
\[
  \left|f\left(x+\frac{k}{n}\right)-f(x+t)\right|<\varepsilon,
\qquad(4)
\]
则式 (3) 自然成立. 注意到 $t\in\left[\dfrac{k}{n},\dfrac{k+1}{n}\right]$，因此点
$\left(x+\dfrac{k}{n}\right)$ 与点 $(x+t)$ 的距离
\[
  \left|\left(x+\frac{k}{n}\right)-(x+t)\right|
  =\left|\frac{k}{n}-t\right|<\frac1n.
\qquad(5)
\]
利用 Cantor 定理，$f$ 在 $[a,b+1]$ 上一致连续. 所以 $\forall\varepsilon>0,\ \exists\delta>0$，当
$x',x''\in[a,b+1]$，$|x'-x''|<\delta$ 时，便有
$|f(x')-f(x'')|<\varepsilon$. 故取 $N=\dfrac1\delta$，当 $n>N$ 时有
$\dfrac1n<\dfrac1N=\delta$，于是由式 (5) 推得式 (4)，从而推得式 (3) 成立. 问题获证.

\noindent\textbf{new 练习}\quad 设 $f(x)$ 在 $[0,1]$ 上连续，试证：
\[
  \sum_{i=1}^n\frac{x}{n}f\left(i\frac{x}{n}\right)\rightrightarrows\int_0^xf(t)\,dt
\qquad(1)
\]
（当 $n\to\infty$ 时，关于 $x\in[0,1]$ 一致）.

\noindent\textbf{提示}\quad 因
\[
  \int_0^xf(t)\,dt=\sum_{i=1}^n\int_{(i-1)x/n}^{ix/n}f(t)\,dt
  =\sum_{i=1}^n f(\xi_i)\frac{x}{n},\quad (i-1)\frac{x}{n}\le\xi_i\le i\frac{x}{n},
\]
故
\begin{equation}
\begin{aligned}
\left|\sum_{i=1}^n\frac{x}{n}f\left(i\frac{x}{n}\right)-\int_0^xf(t)\,dt\right|
&=\left|\frac{x}{n}\sum_{i=1}^n\left(f\left(i\frac{x}{n}\right)-f(\xi_i)\right)\right|\\
&\le\frac{x}{n}\sum_{i=1}^n\left|f\left(i\frac{x}{n}\right)-f(\xi_i)\right|
\le \frac{x}{n}\sum_{i=1}^n\omega_i^{[0,x]/n}.
\end{aligned}
\qquad(2)
\end{equation}
其中 $\omega_i^{[0,x]/n}$ 表示当 $[0,x]$ $n$ 等分时，$f$ 在第 $i$ 个小区间上的振幅.

\noindent\textbf{证}\quad 根据 $f$ 在 $[0,1]$ 上可积的充要条件（§4.2 的定理 1）：$\forall\varepsilon>0,\ \exists\delta>0$，使得
$\lambda=\max_{1\le i\le n}|\Delta x_i|<\delta$ 时，
$0<\sum_{i=1}^n\omega_i\Delta x_i<\varepsilon$. 取 $N:\dfrac1N<\delta$，当
$n>N$ 时，$\forall x\in(0,1]$，在 $[0,x]$ 上用 $[0,x]$ 的 $n$ 等分点，在
$(x,1]$ 上用 $[0,1]$ 的 $n$ 等分点，两者组成 $[0,1]$ 的一个分割，且有
$\lambda=\max_{1\le i\le n}|\Delta x_i|\le\dfrac1n<\dfrac1N<\delta$，于是
\[
  \left|\sum_{i=1}^n\frac{x}{n}f\left(i\frac{x}{n}\right)-\int_0^xf(t)\,dt\right|
  \stackrel{\text{式}(2)}{\le}
  \frac{x}{n}\sum_{i=1}^n\omega_i^{[0,x]/n}
  <\sum_{i=1}^n\frac1n\omega_i<\varepsilon.
\qquad(3)
\]
总之：$\forall\varepsilon>0,\ \exists N>0$，当 $n>N$ 时，$\forall x\in(0,1]$，有式 (3) 成立，式 (1) 获证.

\noindent\textbf{$\star$例 5.2.2}\quad 设函数 $f(x)$ 在 $(-\infty,+\infty)$ 上有连续的导函数
$f'(x)$，$f_n(x)=e^n[f(x+e^{-n})-f(x)]\ (n=1,2,\cdots)$，证明：
$\{f_n(x)\}(n=1,2,\cdots)$ 在任一有限开区间 $(a,b)$ 内一致收敛于 $f'(x)$.（福建师范大学）

\noindent\textbf{证}\quad （利用微分中值定理.）
\[
  |f_n(x)-f'(x)|
  =\left|\frac{f(x+e^{-n})-f(x)}{e^{-n}}-f'(x)\right|
  =|f'(\xi)-f'(x)|\quad(x<\xi\le x+e^{-n}).
\]
因 $f'(x)$ 在 $[a,b+1]$ 上一致连续，$\forall\varepsilon>0,\ \exists\delta>0(\delta<1)$，当
$x_1,x_2\in[a,b+1]$，$|x_1-x_2|<\delta$ 时，有
$|f'(x_1)-f'(x_2)|<\varepsilon$. 取 $N=\ln\dfrac1\delta$，则 $n>N$ 时（此时
$0<e^{-n}<\delta$），有
\[
  |f_n(x)-f'(x)|=|f'(\xi)-f'(x)|<\varepsilon.
\]
故当 $n\to\infty$ 时，$f_n(x)\rightrightarrows f'(x)$ 于 $(a,b)$ 上.

\noindent\textbf{例 5.2.3}\quad 设 $f(x)$ 在 $[0,1]$ 上连续，$g_n(x)$ 为阶梯函数
\[
  g_n(x)=\sum_{k=1}^n f\left(\frac{k}{n}\right)\left(\chi_{k/n}(x)-\chi_{(k-1)/n}(x)\right),
\]
其中
\[
  \chi_{i/n}(x)=
  \begin{cases}
  1, & 0\le x<\dfrac{i}{n},\\
  0, & \dfrac{i}{n}\le x\le1
  \end{cases}
  \quad(i=1,2,\cdots,n).
\]
试证：当 $n\to\infty$ 时，$g_n(x)\rightrightarrows f(x)$ 于 $[0,1)$ 上.

\noindent\textbf{注}\quad 按定义，$[0,1]$ 上每项
\[
  f\left(\frac{k}{n}\right)\left(\chi_{k/n}(x)-\chi_{(k-1)/n}(x)\right)
  =
  \begin{cases}
  f\left(\dfrac{k}{n}\right), & x\in\left[\dfrac{k-1}{n},\dfrac{k}{n}\right),\\
  0, & \text{其他}.
  \end{cases}
\]
因此 $f(x)$ 与 $g_n(x)$ 的图形如图 5.2.1 所示，曲线是 $f(x)$ 的图形，水平线段是
$g_n(x)$ 的图形，$f(x)$ 在 $x=\dfrac{k}{n}$ 处的值作为
$g_n(x)$ 在整个小区间 $\left[\dfrac{k-1}{n},\dfrac{k}{n}\right)$ 上的值.

\begin{center}
  \includegraphics[width=0.58\textwidth]{5-2-1.png}
\end{center}

\noindent\textbf{证}\quad 因 $\forall x\in[0,1)$，$\exists k\in\{1,2,\cdots,n\}$，使得
$x\in\left[\dfrac{k-1}{n},\dfrac{k}{n}\right)$. 于是
$g_n(x)=f\left(\dfrac{k}{n}\right)$，
\[
  |f(x)-g_n(x)|=\left|f(x)-f\left(\frac{k}{n}\right)\right|
  \quad(k=1,2,\cdots,n).
\qquad(1)
\]
又因 $f(x)$ 在 $[0,1]$ 上连续，所以在 $[0,1]$ 上一致连续，$\forall\varepsilon>0,\ \exists\delta>0$（与 $x$ 无关），当
$|x_1-x_2|<\delta(x_1,x_2\in[0,1])$ 时，有
\[
  |f(x_1)-f(x_2)|<\varepsilon.
\qquad(2)
\]
令 $N=\dfrac1\delta$，则 $n>N$ 时，
$\left|x-\dfrac{k}{n}\right|\le\dfrac1n<\delta$，从而利用式 (1) 和式 (2) 有
\[
  |f(x)-g_n(x)|=\left|f(x)-f\left(\frac{k}{n}\right)\right|<\varepsilon
  \quad(\forall x\in[0,1)).
\]
此即表明当 $n\to\infty$ 时，$g_n(x)\rightrightarrows f(x)$ 于 $[0,1)$ 上.

\noindent\textbf{$\star$例 5.2.4（Heine 定理的推广）}\quad 试证 $x\to a$ 时
$f(x,y)\rightrightarrows\varphi(y)$（关于 $y\in I$）的充要条件是：
$\forall\{x_n\}\to a\ (x_n\ne a)$，有
$f(x_n,y)\rightrightarrows\varphi(y)$（关于 $y\in I$）（$n\to\infty$）.（郑州大学）

\noindent\textbf{证}\quad $1^\circ$ 必要性.（对 $\{x_n\}\to a(x_n\ne a)$，要证明 $n\to\infty$ 时
$f(x_n,y)\rightrightarrows\varphi(y)$（关于 $y\in I$）. 即要 $\forall\varepsilon>0$，找
$N>0$ 使得 $n>N$ 时，有
$|f(x_n,y)-\varphi(y)|<\varepsilon(\forall y\in I)$.）因已知
$x\to a$ 时 $f(x,y)\rightrightarrows\varphi(y)$（关于 $y\in I$），所以
$\forall\varepsilon>0,\ \exists\delta>0$，当
$0<|x-a|<\delta$ 时，有
$|f(x,y)-\varphi(y)|<\varepsilon(\forall y\in I)$. 既然 $x_n\ne a$ 且
$x_n\to a$（当 $n\to\infty$ 时），所以对此 $\delta>0$，$\exists N>0$，当 $n>N$ 时，有
$0<|x_n-a|<\delta$，从而
$|f(x_n,y)-\varphi(y)|<\varepsilon(\forall y\in I)$. 此即
$f(x_n,y)\rightrightarrows\varphi(y)(y\in I,n\to\infty)$.

$2^\circ$ 充分性. 假设 $x\to a$ 时 $f(x,y)\not\rightrightarrows\varphi(y)$（关于 $y\in I$），则
$\exists\varepsilon_0>0$，使得 $\forall\dfrac1n>0$，
$\exists x_n\left(0<|x_n-a|<\dfrac1n\right)$ 及 $y_n\in I$ 满足
$|f(x_n,y_n)-\varphi(y_n)|\ge\varepsilon_0$. 如此我们得到
$\{x_n\}\to a,\ x_n\ne a$，但
$f(x_n,y)\not\rightrightarrows\varphi(y)$（于 $I$ 上），与已知条件矛盾.

\noindent\textbf{$\star$例 5.2.5}\quad 证明：若
$\displaystyle\sum_{n=1}^{\infty}u_n(x)$ 在区间 $I$ 上收敛，则
$\displaystyle\sum_{n=1}^{\infty}u_n(x)$ 在 $I$ 上一致收敛的充要条件是
$\forall\{x_n\}\subset I$，有
\[
  \lim_{n\to\infty}r_n(x_n)=0
  \quad\left(\text{其中 }r_n(x)=\sum_{k=n+1}^{\infty}u_k(x)\text{ 为级数余和}\right).
\]

\noindent\textbf{证法}\quad 与上题类似，留给读者，并请写出函数序列的相应结果.

\noindent\textbf{注}\quad 本题的充分性的否定形式为：若 $\exists\{x_n\}\subset I$ 使得
$r_n(x_n)\not\to0(n\to\infty)$，则
$\displaystyle\sum_{n=1}^{\infty}u_n(x)$ 在 $I$ 上非一致收敛.

作为该例的一个应用，可见例 5.2.27 之证 II.

\noindent\textbf{放大法}\quad 如前所述，放大法在于把级数的余和
$r_n(x)=S(x)-S_n(x)$ 的绝对值进行适当放大，使得在区间 $I$ 上
$|r_n(x)|=|S(x)-S_n(x)|\le\alpha_n$（$\alpha_n$ 与 $x$ 无关），且
$\alpha_n\to0(n\to\infty)$，则该级数在所论区间上一致收敛.

实现放大有很多技巧，下面各例分别是通过已知的不等式，求极值，利用已知的余项估计、递推放大等典型方法来实现的. 如下例是利用 Cauchy 不等式进行放大.

\noindent\textbf{例 5.2.6}\quad 若 $f_n(x)$ 在 $[a,b]$ 上可积，$n=1,2,\cdots$，且
$f(x)$ 与 $g(x)$ 在 $[a,b]$ 上都可积，
\[
  \lim_{n\to\infty}\int_a^b|f_n(x)-f(x)|^2\,dx=0,
\]
设
\[
  h(x)=\int_a^x f(t)g(t)\,dt,\qquad
  h_n(x)=\int_a^x f_n(t)g(t)\,dt,
\]
则在 $[a,b]$ 上 $h_n(x)\rightrightarrows h(x)$.（东北师范大学）

\noindent\textbf{证}\quad
\[
\begin{aligned}
|h(x)-h_n(x)|
&=\left|\int_a^x f(t)g(t)\,dt-\int_a^x f_n(t)g(t)\,dt\right|\\
&=\left|\int_a^x(f(t)-f_n(t))g(t)\,dt\right|\\
&\le\int_a^x|f(t)-f_n(t)|\,|g(t)|\,dt\\
&\le\left(\int_a^x|f(t)-f_n(t)|^2\,dt\right)^{1/2}
     \left(\int_a^x|g(t)|^2\,dt\right)^{1/2}\\
&\le\left(\int_a^b|f(t)-f_n(t)|^2\,dt\right)^{1/2}
     \left(\int_a^b|g(t)|^2\,dt\right)^{1/2}\to0\quad(n\to\infty),
\end{aligned}
\]
所以当 $n\to\infty$ 时，$h_n(x)\rightrightarrows h(x)$ 于 $[a,b]$ 上.

下例是通过求极大值得到放大.

\noindent\textbf{例 5.2.7}\quad 给定函数序列：
\[
  f_n(x)=\frac{x(\ln n)^\alpha}{n^x}\quad(n=2,3,4,\cdots).
\]
试问当 $\alpha$ 取何值时，$\{f_n(x)\}$ 在 $[0,+\infty)$ 上一致收敛.（广西大学）

\noindent\textbf{解}\quad
\[
  f_n'(x)=\frac{(\ln n)^{\alpha+1}}{n^x}\left(\frac1{\ln n}-x\right).
\]
可见，当 $x<\dfrac1{\ln n}$ 时，$f_n(x)\nearrow$；当
$x>\dfrac1{\ln n}$ 时，$f_n(x)\searrow$，函数 $f_n(x)$ 在
$x=\dfrac1{\ln n}$ 处取极大值. 注意极限函数
$f(x)\equiv\lim_{n\to\infty}f_n(x)=0$. 故
\[
\begin{aligned}
\sup_{x\in(0,+\infty)}|f(x)-f_n(x)|
&=\max_{x\in(0,+\infty)}|f_n(x)|
 =f_n\left(\frac1{\ln n}\right)\\
&=\frac{(\ln n)^{\alpha-1}}{n^{1/\ln n}}
 =\frac1e(\ln n)^{\alpha-1}
\begin{cases}
\not\to0, & \alpha\ge1,\\
\to0, & \alpha<1
\end{cases}
\quad(n\to\infty).
\end{aligned}
\]
（这里 $n^{1/\ln n}=(e^{\ln n})^{1/\ln n}=e$.）所以当且仅当 $\alpha<1$ 时，
$\{f_n(x)\}$ 在 $[0,+\infty)$ 上一致收敛.

注意，放大法一般来说只是一致收敛的充分条件. 但如本例用求极大值的方法，得到
$\alpha_n=\max\{|S(x)-S_n(x)|\}$，则 $\alpha_n\to0(n\to\infty)$ 的条件不仅是充分的而且是必要的.

下两例是利用级数的余和估计：

对于 Leibniz 级数，级数余和
$r_n=S-S_n=\displaystyle\sum_{k=n+1}^{\infty}(-1)^ka_k$，有估计式
$|r_n|\le a_{n+1}$.

\noindent\textbf{例 5.2.8}\quad 试证：
\[
  \sum_{n=1}^{\infty}\frac{n(-1)^n}{n^2+x^2}
\]
在 $(-\infty,+\infty)$ 内一致收敛.

\noindent\textbf{证}\quad 设函数 $f(y)=\dfrac{y}{y^2+x^2}$，则
$f'(y)=\dfrac{x^2-y^2}{(y^2+x^2)^2}$. 可见 $\forall x\in(-\infty,+\infty)$，当 $n$ 充分大时，级数通项的绝对值
$\dfrac{n}{n^2+x^2}\searrow0(n\to\infty)$，故该级数为 Leibniz 级数. 因而
\[
  |r_n(x)|\le\frac{n+1}{(n+1)^2+x^2}\le\frac1{n+1}\to0\quad(n\to\infty).
\]
所以 $\displaystyle\sum_{n=1}^{\infty}\dfrac{n(-1)^n}{n^2+x^2}$ 在
$(-\infty,+\infty)$ 内一致收敛.

\noindent\textbf{例 5.2.9}\quad 讨论
\[
  \sum_{n=1}^{\infty}\frac{nx}{(1+x)(1+2x)\cdots(1+nx)}
\]
在 $(0,a)$ 与 $(a,+\infty)$ 内的一致收敛性.

\noindent\textbf{解}\quad 用 D'Alembert 判别法，容易知道该级数在 $(0,+\infty)$ 内处处收敛.
\[
\begin{aligned}
r_n(x)
&=\sum_{k=n+1}^{\infty}\frac{kx}{(1+x)(1+2x)\cdots(1+kx)}\\
&=\sum_{k=n+1}^{\infty}\frac{kx+1-1}{(1+x)(1+2x)\cdots(1+kx)}\\
&=\sum_{k=n+1}^{\infty}\left\{
\frac1{(1+x)(1+2x)\cdots[1+(k-1)x]}
-\frac1{(1+x)(1+2x)\cdots(1+kx)}
\right\}\\
&=\frac1{(1+x)(1+2x)\cdots(1+nx)}.
\end{aligned}
\]
\[
  \sup_{x\in(0,a)}|r_n(x)|=1\not\to0,\qquad
  \sup_{x\in(a,+\infty)}|r_n(x)|
  =\frac1{(1+a)(1+2a)\cdots(1+na)}\to0\quad(n\to\infty).
\]
所以此级数在 $(0,a)$ 内非一致收敛，在 $(a,+\infty)$ 内一致收敛.

有些函数序列是用递推形式给出的，这时可考虑用递推的方式进行放大. 如

\noindent\textbf{$\star$例 5.2.10}\quad 设 $f_1(x)$ 在 $[a,b]$ 上正常可积，
$f_{n+1}(x)=\displaystyle\int_a^x f_n(t)\,dt,\ n=1,2,\cdots$. 证明：函数序列
$\{f_n(x)\}$ 在 $[a,b]$ 上一致收敛于零.（吉林大学）

\noindent\textbf{证}\quad 因为 $f_1(x)$ 在 $[a,b]$ 上正常可积，故在 $[a,b]$ 上有界，即
$\exists M>0$，使得 $|f_1(x)|\le M(\forall x\in[a,b])$. 从而
\[
  |f_2(x)|\le\int_a^x|f_1(t)|\,dt\le M(x-a),
\]
\[
  |f_3(x)|\le\int_a^x|f_2(t)|\,dt
  \le M\int_a^x(t-a)\,dt=\frac{M(x-a)^2}{2!}.
\]
一般来说，若对 $n$ 有
\[
  |f_n(x)|\le\frac{M(x-a)^{n-1}}{(n-1)!},
\]
则
\[
  |f_{n+1}(x)|\le\int_a^x|f_n(t)|\,dt
  \le\frac{M}{(n-1)!}\int_a^x(t-a)^{n-1}\,dt
  =\frac{M(x-a)^n}{n!}.
\]
所以
\[
  |f_n(x)|\le\frac{M(b-a)^{n-1}}{(n-1)!}\to0\quad(n\to\infty),
\]
故 $f_n(x)\rightrightarrows0$（当 $n\to\infty$ 时）（关于 $x\in[a,b]$）.

\noindent\textbf{例 5.2.11}\quad 证明：若 $K(x,t)$ 在
$D=[a\le x\le b,\ a\le t\le b]$ 上连续，$u_0(x)$ 在 $[a,b]$ 上连续，且对任意
$x\in[a,b]$，令
\[
  u_n(x)=\int_a^x K(x,t)u_{n-1}(t)\,dt,\quad n=1,2,\cdots,
\]
则函数序列 $\{u_n(x)\}$ 在 $[a,b]$ 上一致收敛.（东北师范大学）

\noindent\textbf{提示}\quad 在有界闭区域上连续的二元函数必有界，即 $\exists M>0$，当
$a\le x\le b,\ a\le t\le b$ 时，恒有 $|K(x,t)|\le M$.

\noindent\textbf{例 5.2.12}\quad 假设

1）$f(x)$ 在 $(-\infty,+\infty)$ 内连续；

2）$x\ne0$ 时有 $|f(x)|<|x|$；

3）$f_1(x)=f(x),\ f_2(x)=f(f_1(x)),\cdots,\ f_n(x)=f(f_{n-1}(x)),\cdots$.

试证：$f_n(x)$ 在 $[-A,A]$ 上一致收敛（其中 $A$ 为正常数）.（南京大学，吉林大学）

\noindent\textbf{证}\quad 因 $x\ne0$ 时有 $0\le |f(x)|<|x|$，故令 $x\to0$，取极限（已知
$f(x)$ 连续）得 $0\le f(0)\le0,\ f(0)=0$，从而由条件 2），在
$[-A,A](A>0)$ 上恒有 $|f(x)|\le |x|$.

由此，$\forall\varepsilon>0$（不妨设 $\varepsilon<A$），当
$x\in[-\varepsilon,\varepsilon]$ 时，有 $|f(x)|\le |x|\le\varepsilon$；在
$[-A,-\varepsilon]\cup[\varepsilon,A]$ 上，$\left|\dfrac{f(x)}{x}\right|<1$ 连续，且有最大值
$q:0<q<1$. 于是，$|f(x)|\le q|x|\le qA$. 总之在 $[-A,A]$ 上恒有
\[
  |f(x)|\le\max\{\varepsilon,qA\}.
\]
$\forall x\in[-A,A]$，若 $|f(x)|\le\varepsilon$，则
$|f_2(x)|=|f(f(x))|\le |f(x)|\le\varepsilon$.
若 $|f(x)|\in[\varepsilon,A]$，则
$|f_2(x)|=|f(f(x))|\le q|f(x)|\le q^2A$.

所以，总有
\[
  |f_2(x)|\le\max\{\varepsilon,q^2A\}.
\]
同理，由 $|f_{n-1}(x)|\le\max\{\varepsilon,q^{n-1}A\}$ 可推出
$|f_n(x)|\le\max\{\varepsilon,q^nA\}$. 故此式对一切 $n$ 成立. 由于当
$n\to\infty$ 时，$q^nA\to0$，对 $\varepsilon>0$，$\exists N>0$，$n>N$ 时，
$q^nA<\varepsilon$，所以 $|f_n(x)|\le\max\{\varepsilon,q^nA\}=\varepsilon$. 证毕.

放大法也要注意根据具体情况，作灵活的处理.

\noindent\textbf{例 5.2.13}\quad 设 $\alpha_n>0,\ \lim_{n\to\infty}\alpha_n=0$，
$|u_n(x)|\le\alpha_n(x\in I)$ 且 $u_i(x)u_j(x)=0(i\ne j)(\forall x\in I)$. 试证
$\displaystyle\sum_{n=1}^{\infty}u_n(x)$ 在 $I$ 上一致收敛（这里 $I$ 为任意区间）.

\noindent\textbf{证}\quad $\forall x\in I$，$|u_n(x)|$ 至多只有一项不为 $0$. 因此（若用
$S(x)$ 表示和函数，$S_n(x)$ 表示部分和）
\[
  |S(x)-S_n(x)|\le \sup_{k>n}|u_k(x)|\le \sup_{k>n}\alpha_k\to0\quad(n\to\infty).
\]
故 $\displaystyle\sum_{n=1}^{\infty}u_n(x)$ 在 $I$ 上一致收敛.

\noindent\textbf{b.\quad 利用 Cauchy 准则判断一致收敛性}

\noindent\textbf{要点}\quad 用 Cauchy 准则判断级数（或函数序列）是否一致收敛完全取决于充分后的“片段”是否能一致地任意小，而无须求出和函数（或极限函数）. 这一点比用定义法优越. 如
$\displaystyle\sum_{n=1}^{\infty}u_n(x)$ 在区间 $I$ 上一致收敛的充要条件是：
$\forall\varepsilon>0,\ \exists N>0$，当 $n>N$ 时，
\[
  \left|\sum_{k=n+1}^{n+p}u_k(x)\right|<\varepsilon
  \quad(\forall x\in I,\ \forall p\in\mathbb N).
\]
$\displaystyle\sum_{n=1}^{\infty}u_n(x)$ 在 $I$ 上非一致收敛的充要条件是：
$\exists\varepsilon_0>0,\ \forall N>0,\ \exists n>N,\ \exists x\in I,\ \exists p\in\mathbb N$，使得
\[
  \left|\sum_{k=n+1}^{n+p}u_k(x)\right|\ge\varepsilon_0.
\]

特别，若通项 $u_n(x)\not\rightrightarrows0(n\to\infty)$ 关于 $x\in I$，则
$\displaystyle\sum_{n=1}^{\infty}u_n(x)$ 非一致收敛.

可见：“当 $n\to\infty$ 时，$u_n(x)\rightrightarrows0$（关于 $x\in I$）”是
$\displaystyle\sum_{n=1}^{\infty}u_n(x)$ 在 $I$ 中一致收敛的必要条件. 因此，若找到
$\{x_n\}\subset I$ 使得 $\{u_n(x_n)\}\not\to0(n\to\infty)$，则当 $n\to\infty$ 时
$u_n(x)\not\rightrightarrows0$（关于 $x\in I$）. 从而
$\displaystyle\sum_{n=1}^{\infty}u_n(x)$ 在 $I$ 中非一致收敛.

\noindent\textbf{$\star$例 5.2.14}\quad 设 $\{u_n(x)\}$ 为 $[a,b]$ 上的可导函数序列，且在
$[a,b]$ 上有
\[
  \left|\sum_{k=1}^n u_k'(x)\right|\le C,
\qquad(1)
\]
$C$ 是不依赖于 $x$ 和 $n$ 的正常数. 证明：若
$\displaystyle\sum_{n=1}^{\infty}u_n(x)$ 在 $[a,b]$ 上收敛，则必为一致收敛.（华东师范大学）

\noindent\textbf{证 I}\quad $1^\circ$ $\displaystyle\sum_{n=1}^{\infty}u_n(x)$ 在 $[a,b]$ 上收敛，所以，
$\forall x_0\in[a,b],\ \forall\varepsilon>0,\ \exists N=N(\varepsilon,x_0)>0$，当
$n>N$ 时，有
\[
  \left|\sum_{k=n+1}^{n+p}u_k(x_0)\right|<\frac{\varepsilon}{2}
  \quad(\forall p\in\mathbb N).
\qquad(2)
\]
故
\[
\left|\sum_{k=n+1}^{n+p}u_k(x)\right|
\le
\left|\sum_{k=n+1}^{n+p}u_k(x)-\sum_{k=n+1}^{n+p}u_k(x_0)\right|
\left|\sum_{k=n+1}^{n+p}u_k(x_0)\right|.
\]
此式右端第一项中，对函数 $\displaystyle\sum_{k=n+1}^{n+p}u_k(x)$ 应用微分中值定理：
\[
  \left|\sum_{k=n+1}^{n+p}u_k(x)\right|
  \le \left|\sum_{k=n+1}^{n+p}u_k'(\xi)(x-x_0)\right|
  +\left|\sum_{k=n+1}^{n+p}u_k(x_0)\right|
  <2C|x-x_0|+\frac{\varepsilon}{2}
\]
（$\xi$ 在 $x$ 与 $x_0$ 之间）.

取 $\delta=\dfrac{\varepsilon}{4C}$，则 $|x-x_0|<\delta$ 时，有
\[
  \left|\sum_{k=n+1}^{n+p}u_k(x)\right|<\varepsilon
  \quad(\forall p\in\mathbb N,\ \forall x\in(x_0-\delta,x_0+\delta)).
\qquad(3)
\]
至此，虽然式 (3) 不在整个区间 $[a,b]$ 上同时成立，但在 $x_0$ 的邻域
$(x_0-\delta,x_0+\delta)$ 成立.

$2^\circ$（注意到 $x_0\in[a,b]$ 是任意的. 于是，对 $[a,b]$ 上每一点，都采用上述步骤.）
$\forall x_0\in[a,b]$，$\exists N(\varepsilon,x_0)>0$，当 $n>N$，
$x\in(x_0-\delta,x_0+\delta)$ 时有
\[
  \left|\sum_{k=n+1}^{n+p}u_k(x)\right|<\varepsilon
  \quad(\forall p\in\mathbb N).
\]
如此，$\{(x_0-\delta,x_0+\delta)\mid x_0\in[a,b]\}$ 组成了 $[a,b]$ 上的一个开覆盖，由有限覆盖定理，其中存在有限子覆盖. 不妨设之为
$\{(x_i-\delta,x_i+\delta)\}_{i=1}^r$. 令
$N=\max\{N(\varepsilon,x_i)\}$，则当 $n>N$ 时，$\forall x\in[a,b]$，有
\[
  \left|\sum_{k=n+1}^{n+p}u_k(x)\right|<\varepsilon
  \quad(\forall p\in\mathbb N).
\]

\noindent\textbf{证 II}\quad $\forall\varepsilon>0$，取 $m$ 充分大，将 $[a,b]$ $m$ 等分，使分得的每个小区间长度
$\delta<\dfrac{\varepsilon}{4C}$. 顺次以 $x_1,x_2,\cdots,x_m$ 表示各小区间的中点. 因
$\displaystyle\sum_{n=1}^{\infty}u_n(x)$ 收敛，对 $\varepsilon>0$，$\exists N_i=N(\varepsilon,x_i)$，当
$n>N_i$ 时有
\[
  \left|\sum_{k=n+1}^{n+p}u_k(x_i)\right|<\frac{\varepsilon}{2}
  \quad(\forall p\in\mathbb N).
\qquad(4)
\]
令 $N=\max\{N_1,N_2,\cdots,N_m\}$，则 $\forall x\in[a,b]$（不妨设 $x$ 位于第 $i$ 个小区间，
$i\in\{1,2,\cdots,m\}$），有
\[
\begin{aligned}
\left|\sum_{k=n+1}^{n+p}u_k(x)\right|
&=\left|\sum_{k=n+1}^{n+p}u_k(x_i)
 +\int_{x_i}^x\left(\sum_{k=n+1}^{n+p}u_k(t)\right)'dt\right|\\
&\le\left|\sum_{k=n+1}^{n+p}u_k(x_i)\right|
 +\int_{x_i}^x\left|\sum_{k=n+1}^{n+p}u_k'(t)\right|dt\\
&<\frac{\varepsilon}{2}+2C|x-x_i|\le\frac{\varepsilon}{2}+\frac{\varepsilon}{2}=\varepsilon.
\end{aligned}
\]

上面我们看到使用 Cauchy 准则证明一致收敛，十分重要的问题是将“片段”
$\displaystyle\sum_{k=n+1}^{n+p}u_k(x)$ 进行变形. 作为这种变形的一个重要方法是 Abel 变换.

\noindent\textbf{*例 5.2.15}\quad 设函数序列 $f_0(x),f_1(x),\cdots$ 在区间 $I$ 上有定义，且满足
\[
  \text{i） }|f_0(x)|\le M;\qquad
  \text{ii） }\sum_{n=0}^m|f_n(x)-f_{n+1}(x)|\le M,\quad m=0,1,2,\cdots,
\]
其中 $M$ 是常数. 试证：如果级数 $\displaystyle\sum_{n=0}^{\infty}b_n$ 收敛，则级数
$\displaystyle\sum_{n=0}^{\infty}b_nf_n(x)$ 必在区间 $I$ 上一致收敛.（吉林大学）

\noindent\textbf{证}\quad 因 $\displaystyle\sum_{n=0}^{\infty}b_n$ 收敛，故
$\forall\varepsilon>0,\ \exists N>0$，当 $n>N$ 时，
\[
  \left|\sum_{k=n+1}^{n+p}b_k\right|<\varepsilon\quad(\forall p\in\mathbb N).
\]
记 $S_i=\displaystyle\sum_{k=n+1}^{n+i}b_k$，于是
$|S_i|<\varepsilon(i=1,2,\cdots)$. 故
\[
\begin{aligned}
\left|\sum_{k=n+1}^{n+p}b_kf_k(x)\right|
&\stackrel{\text{仿用 Abel 变换}}{=}
\left|S_1f_{n+1}+(S_2-S_1)f_{n+2}+\cdots+(S_p-S_{p-1})f_{n+p}\right|\\
&=\left|S_1(f_{n+1}-f_{n+2})+\cdots+S_{p-1}(f_{n+p-1}-f_{n+p})+S_pf_{n+p}\right|\\
&\le |S_1|\,|f_{n+1}-f_{n+2}|+\cdots+|S_{p-1}|\,|f_{n+p-1}-f_{n+p}|+|S_p|\,|f_{n+p}|\\
&\le\varepsilon\left(\sum_{k=n+1}^{n+p-1}|f_k-f_{k+1}|+|f_{n+p}|\right).
\end{aligned}
\]
因为
\[
  \sum_{k=n+1}^{n+p-1}|f_k-f_{k+1}|
  \le\sum_{k=0}^{n+p-1}|f_k-f_{k+1}|\le M
\]
（条件 ii）），故
\[
  |f_{n+p}|=|f_0-f_0+f_1-f_1+\cdots+f_{n+p-1}-f_{n+p-1}+f_{n+p}|
  \le |f_0|+|f_0-f_1|+\cdots+|f_{n+p-1}-f_{n+p}|\le2M.
\]
所以，
\[
  \left|\sum_{k=n+1}^{n+p}b_kf_k(x)\right|<\varepsilon\cdot3M\quad(\forall p\in\mathbb N).
\]
故 $\displaystyle\sum_{k=0}^{\infty}b_kf_k(x)$ 在 $I$ 上一致收敛.

\noindent\textbf{用 Cauchy 准则证明非一致收敛}

\noindent\textbf{$\star$例 5.2.16}\quad 求证：级数
\[
  \frac{\sin x}{1}+\frac{\sin2x}{2}+\cdots+\frac{\sin nx}{n}+\cdots
\]
在 $x=0$ 的邻域内非一致收敛.

\noindent\textbf{分析}\quad 我们的目标是证明每个标号 $n$ 之后均有“片段”$\ge\varepsilon$（某个事先指定的正数）. 片段
$\displaystyle\sum_{k=n+1}^{n+p}\dfrac{\sin kx}{k}$ 的麻烦在于每项有因子 $\sin kx$，否则
$\displaystyle\sum_{k=n+1}^{n+p}\dfrac1k$ 是调和级数 $\displaystyle\sum_{n=1}^{\infty}\dfrac1n$ 的片段，例如取
$p=n$，则片段 $\displaystyle\sum_{k=n+1}^{2n}\dfrac1k\ge n\cdot\dfrac1{2n}=\dfrac12$.

另一方面，我们看到函数 $\sin x$ 在 $\left[\dfrac{\pi}{4},\dfrac{\pi}{2}\right]$ 上恒大于等于
$\sin\dfrac{\pi}{4}$，因此我们只要保持使
$kx\in\left[\dfrac{\pi}{4},\dfrac{\pi}{2}\right]$，那么
\[
  \sum_{k=n+1}^{2n}\frac{\sin kx}{k}
  \ge \sin\frac{\pi}{4}\sum_{k=n+1}^{2n}\frac1k
  >\frac12\sin\frac{\pi}{4}.
\]
如此，我们想到取 $x=x_n=\dfrac{\pi}{4n}$，从而 $\forall n\in\mathbb N$ 有
\[
  \left|\sum_{k=n+1}^{2n}\frac1k\sin kx_n\right|
  =\left|\sum_{k=n+1}^{2n}\frac{\sin\left(k\cdot\dfrac{\pi}{4n}\right)}{k}\right|.
\]
这是因为 $n+1\le k\le2n$，所以
\[
  \frac{\pi}{4}<(n+1)\frac{\pi}{4n}\le k\frac{\pi}{4n}\le2n\frac{\pi}{4n}=\frac{\pi}{2},
\]
即 $k\cdot\dfrac{\pi}{4n}\in\left[\dfrac{\pi}{4},\dfrac{\pi}{2}\right]$. 故对
$\varepsilon_0=\dfrac14\sqrt2$，$\forall n$，有片段
\[
  \left|\sum_{k=n+1}^{2n}\frac{\sin kx_n}{k}\right|
  \ge \sum_{k=n+1}^{2n}\frac1k\sin\frac{\pi}{4}
  >\frac12\sin\frac{\pi}{4}=\frac14\sqrt2=\varepsilon_0,
\]
级数非一致收敛. 简洁的证明请读者写出.

\noindent\textbf{例 5.2.17}\quad 证明：
\[
  \sum_{n=1}^{\infty}\frac1n\left[e^x-\left(1+\frac{x}{n}\right)^n\right]
\]
在 $(0,+\infty)$ 上非一致收敛.（北京大学）

\noindent\textbf{证}\quad 通项
$\dfrac1n\left[e^x-\left(1+\dfrac{x}{n}\right)^n\right]\to0$（当 $n\to\infty$ 时），
$x\in(0,+\infty)$. 事实上，$\forall n\in\mathbb N$，当 $x\to+\infty$ 时，易知
\[
  \frac1n\left[e^x-\left(1+\frac{x}{n}\right)^n\right]\to+\infty.
\]
所以该级数在 $(0,+\infty)$ 上非一致收敛.

\noindent\textbf{new 练习 1}\quad 试证：函数项级数
\[
  \sum_{n=1}^{\infty}\frac{n^{x+2}}{(1+nx)^n}
\]
在 $(1,+\infty)$ 内收敛，在 $(1,+\infty)$ 内非一致收敛，但和函数在 $(1,+\infty)$ 内连续.（武汉大学）

\noindent\textbf{提示}\quad 1）
\[
  0<\frac{n^{x+2}}{(1+nx)^n}
  <\frac{n^{x+2}}{(nx)^n}
  =\frac{n^2}{x^n}\le\frac{n^2}{\alpha^n}\quad(x\ge\alpha>1),
\]
而 $\sqrt[n]{\dfrac{n^2}{\alpha^n}}\to\dfrac1\alpha<1$. 因此，级数在
$(1,+\infty)$ 内逐点收敛，内闭一致收敛，且和函数连续；

2）令 $x_n=\dfrac1n+1(n\to\infty)$，
\[
  \frac{n^{x_n+2}}{(1+nx_n)^n}
  =\frac{n^2}{\left[\left(1+\frac2n\right)^{n/2}\right]^2}\to+\infty,
\]
说明级数在 $(1,+\infty)$ 内非一致收敛.（事实上：取 $\varepsilon_0=\dfrac12$，当 $n$ 充分大时，
$\dfrac{n^{x_n+2}}{(1+nx_n)^n}>1$，于是（当 $p>1$ 时）
\[
  \sum_{k=n}^{n+p}\frac{k^{x_k+2}}{(1+kx_k)^k}
  >\frac{n^{x_n+2}}{(1+nx_n)^n}>\frac12=\varepsilon_0.
\]

\noindent\textbf{new 练习 2}\quad 讨论级数
\[
  \sum_{n=1}^{\infty}\frac{\sqrt{n+1}-\sqrt n}{n^x}
\]
的收敛性与一致收敛性.（浙江大学）

\noindent\textbf{解}\quad $1^\circ$（在 $\left(\dfrac12,+\infty\right)$ 内收敛，且内闭一致收敛.）因
$\forall\alpha>\dfrac12$，当 $x\ge\alpha$ 时，
\[
  \left|\frac{\sqrt{n+1}-\sqrt n}{n^x}\right|
  =\frac1{n^x(\sqrt{n+1}+\sqrt n)}
  \le\frac1{n^{x+1/2}}\le\frac1{n^{\alpha+1/2}}\quad(n\ge1).
\]
而 $\displaystyle\sum_{n=1}^{\infty}\dfrac1{n^{\alpha+1/2}}$ 收敛，故级数在
$[\alpha,+\infty)$ 上一致收敛，又因 $\alpha>\dfrac12$ 任意，所以在
$\left(\dfrac12,+\infty\right)$ 内逐点收敛，且内闭一致收敛.

$2^\circ$（在 $\left(\dfrac12,+\infty\right)$ 内非一致收敛.）我们已知
$\displaystyle\sum_{n=1}^{\infty}\dfrac1n$ 发散，故 $\forall n\in\mathbb N$，当
$m>n,\ m$ 充分大时，有
$\displaystyle\sum_{k=n+1}^m\dfrac1k>2$. 故当
$x\to\left(\dfrac12\right)^+$ 时，
\[
  \left|\sum_{k=n+1}^m\frac{\sqrt{k+1}-\sqrt k}{k^x}\right|
  =\sum_{k=n+1}^m\frac1{k^x(\sqrt{k+1}+\sqrt k)}
  \ge \sum_{k=n+1}^m\frac1{k^{x+1/2}}
  \xrightarrow{x\to(1/2)^+}\sum_{k=n+1}^m\frac1k>2.
\]
可见若取 $\varepsilon_0=1$，则 $\forall n\in\mathbb N$，总 $\exists m>n,\ x>\dfrac12$（且与
$\dfrac12$ 充分接近），使得
\[
  \left|\sum_{k=n+1}^m\frac{\sqrt{k+1}-\sqrt k}{k^x}\right|\ge\varepsilon_0=1.
\]
这表明原级数在 $\left(\dfrac12,+\infty\right)$ 内非一致收敛.

\noindent\textbf{c.\quad 利用常用的判别法证明一致收敛性}

这里所说的常用判别法指：$M$ 判别法（即 Weierstrass 判别法），Abel 判别法、Dirichlet 判别法及 Dini 判别法. 下面分别加以讨论.

\noindent\textbf{1）$M$ 判别法}

\noindent\textbf{要点}\quad 根据 $M$ 判别法，要证明
$\displaystyle\sum_{n=1}^{\infty}u_n(x)$ 在区间 $I$ 上一致收敛，只要找到收敛的优级数. 即：将通项
$u_n(x)$ 放大，使
\[
  |u_n(x)|\le M_n,\quad \forall x\in I.
\]
若级数 $\displaystyle\sum_{n=1}^{\infty}M_n$ 收敛，则
$\displaystyle\sum_{n=1}^{\infty}u_n(x)$ 在 $I$ 上一致收敛（
$\displaystyle\sum M_n$ 称为它的优级数）.

求优级数的方法，除有些可以用观察法之外，通常还可用如下方法：求 $u_n(x)$ 在区间
$I$ 上的最大值；利用已知的不等式；用 Taylor 公式、微分中值定理等各种方法变形再放大.

\noindent\textbf{利用 $u_n(x)$ 的最大值进行放大}

\noindent\textbf{$\star$例 5.2.18}\quad 证明
\[
  \sum_{n=1}^{\infty}x^n(1-x)^2
\]
在 $[0,1]$ 上一致收敛.（安徽大学）

\noindent\textbf{证}\quad 对通项 $u_n(x)=x^n(1-x)^2$ 求导，令
\[
  u_n'(x)=nx^{n-1}(1-x)^2-2x^n(1-x)=0,
\]
得出全部驻点 $x=0,1,\dfrac{n}{n+2}$. 因
$u_n\left(\dfrac{n}{n+2}\right)>u_n(0)=u_n(1)=0$，所以
$u_n\left(\dfrac{n}{n+2}\right)$ 为 $u_n(x)$ 在 $[0,1]$ 上的最大值. 如此
\[
  x^n(1-x)^2
  \le\left(\frac{n}{n+2}\right)^n\left(1-\frac{n}{n+2}\right)^2
  \le\left(1-\frac{n}{n+2}\right)^2
  =\left(\frac2{n+2}\right)^2\le\frac4{n^2}.
\]
因 $\displaystyle\sum_{n=1}^{\infty}\dfrac4{n^2}$ 收敛，所以
$\displaystyle\sum_{n=1}^{\infty}x^n(1-x)^2$ 在 $[0,1]$ 上一致收敛.

\noindent\textbf{new 练习}\quad 试证函数项级数
\[
  \sum_{n=1}^{\infty}x^3e^{-nx^2}
\]
在 $[0,+\infty)$ 上一致收敛.（北京大学）

\noindent\textbf{提示}\quad 求通项的最大值可得
$\displaystyle\sum_{n=1}^{\infty}x^3e^{-nx^2}$ 的优级数：
\[
  \sum_{n=1}^{\infty}\left(\frac3{2n}\right)^{3/2}e^{-3/2}.
\]

\noindent\textbf{利用已知的不等式进行放大}

\noindent\textbf{例 5.2.19}\quad 证明
\[
  \sum_{n=1}^{\infty}\arctan\frac{2x}{x^2+n^3}
\]
在 $(-\infty,+\infty)$ 内一致收敛.

\noindent\textbf{提示}\quad
\[
  \left|\arctan\frac{2x}{x^2+n^3}\right|
  \le2\frac{|x|}{x^2+n^3}
  =\frac{\sqrt{x^2n^3}}{(x^2+n^3)/2}\cdot\frac1{n^{3/2}}<\frac1{n^{3/2}}.
\]

\noindent\textbf{$\star$例 5.2.20}\quad 证明函数项级数
\[
  \sum_{n=1}^{\infty}\frac1n\left[e^x-\left(1+\frac{x}{n}\right)^n\right]
\]
在任意有限区间 $[a,b]$ 上一致收敛.（北京大学）

\noindent\textbf{提示}\quad 可利用例 3.4.5 中的不等式，如设 $|a|,|b|\le M$，则当 $n$ 充分大时，
\[
  \frac1n\left[e^x-\left(1+\frac{x}{n}\right)^n\right]
  \le \frac{x^2}{n^2}e^x\le \frac{M^2}{n^2}e^M.
\]

\noindent\textbf{利用 Taylor 公式等进行变形后放大}

\noindent\textbf{$\star$例 5.2.21}\quad 设一元函数 $f$ 在 $x=0$ 的邻域里有二阶连续导数，
$f(0)=0,\,0<f'(0)<1$. 函数 $f_n$ 是 $f$ 的 $n$ 次复合，证明级数
$\displaystyle\sum_{n=1}^{\infty}f_n(x)$ 在 $x=0$ 的邻域里一致收敛.（中国科学技术大学）

\noindent\textbf{分析}\quad 因为 $f$ 在 $x=0$ 的邻域内有二阶连续导数，当 $\delta>0$ 充分小时，在
$[-\delta,\delta]$ 上 $f''(x)$ 连续，且
\[
  f(x)=f(0)+f'(0)x+\frac1{2!}f''(\xi)x^2
  =f'(0)x+\frac12 f''(\xi)x^2\quad(|\xi|<|x|\le\delta).
\]
既然 $f''(x)$ 在 $[-\delta,\delta]$ 上连续，所以 $\exists M>0$，使得
$|f''(x)|\le M$. 于是
\[
  |f(x)|\le |x|\left(f'(0)+\frac12M\delta\right)=q\cdot |x|.
\qquad(1)
\]
（这里记 $f'(0)+\dfrac12M\delta=q$.）重复使用得
\[
  |f_2(x)|=|f(f(x))|\le q|f(x)|\le q\cdot q|x|=q^2|x|<q^2\delta,
\]
\[
  |f_n(x)|\le q|f_{n-1}(x)|\le\cdots\le q^n|x|\le q^n\delta.
\qquad(2)
\]
为使级数 $\displaystyle\sum_{n=1}^{\infty}q^n\delta$ 收敛，必须使得正数
$q=f'(0)+\dfrac12M\delta<1$. 但 $0<f'(0)<1$，故只要从式（1）开始把邻域进一步缩小，取
$\delta_1<\delta$，用 $\delta_1$ 代替 $\delta$ 使得
$\dfrac12M\delta_1<1-f'(0)$（即
$\delta_1<\dfrac2M(1-f'(0))$），则
$q_1=f'(0)+\dfrac12M\delta_1<1$. 从而式（2）成为
$|f_n(x)|\le q_1^n\delta_1$（当 $|x|<\delta_1<\delta$ 时），可知
$\displaystyle\sum_{n=1}^{\infty}f_n(x)$ 在 $x=0$ 的邻域里一致收敛（简洁的证明请读者写出）.

\noindent\textbf{例 5.2.22}\quad 证明 $\displaystyle\sum_{n=1}^{\infty}x^2e^{-nx}$ 在
$(0,+\infty)$ 内一致收敛.（西南大学）

\noindent\textbf{提示}\quad
\[
  x^2e^{-nx}
  =\frac{x^2}{1+nx+\dfrac{n^2x^2}{2}+\cdots}
  <\frac{x^2}{\dfrac{n^2x^2}{2}}=\frac2{n^2}.
\]

\noindent\textbf{2）Abel 判别法与 Dirichlet 判别法}

\noindent\textbf{要点}\quad 根据 Abel 判别法与 Dirichlet 判别法，要证明级数
$\displaystyle\sum_{n=1}^{\infty}u_n(x)$ 一致收敛，关键是将通项写成两个因子相乘，使之符合判别法的条件. 即令
$u_n(x)=a_n(x)\cdot b_n(x)$，

若：

i）$\displaystyle\sum_{n=1}^{\infty}a_n(x)$ 在 $I$ 上一致收敛；

ii）$|b_n(x)|$ 一致有界，且对每个固定的 $x$，$|b_n(x)|$ 关于 $n$ 单调，

则由 Abel 判别法，$\displaystyle\sum_{n=1}^{\infty}u_n(x)$ 在 $I$ 上一致收敛.

若：

i）$\displaystyle\sum_{k=1}^{n}a_k(x)$ 关于 $x$ 与 $n$ 一致有界；

ii）$b_n(x)$ 对每个固定的 $x$，关于 $n$ 单调，且 $n\to\infty$ 时 $b_n(x)\rightrightarrows0$（于 $I$ 上），

则由 Dirichlet 判别法，$\displaystyle\sum_{n=1}^{\infty}u_n(x)$ 在 $I$ 上一致收敛.

\noindent\textbf{例 5.2.23}\quad 假设 $b>0$，$a_1,a_2,\cdots$ 均为常数，且级数
$\displaystyle\sum_{n=1}^{\infty}a_n$ 收敛，试证：
\[
  \sum_{n=1}^{\infty}a_n\frac1{n!}\int_0^x t^ne^{-t}\,dt
\]
在 $[0,b]$ 上一致收敛.

\noindent\textbf{证}\quad $1^\circ$ $\displaystyle\sum_{n=1}^{\infty}a_n$ 收敛，自然关于 $x$ 一致收敛.

$2^\circ$
\[
  0\le\frac1{n!}\int_0^x t^ne^{-t}\,dt
  \le\frac1{n!}\int_0^{+\infty}t^ne^{-t}\,dt=1
\]
（$\forall n\in\mathbb N,\ \forall x\in[0,b]$），即
$\dfrac1{n!}\displaystyle\int_0^x t^ne^{-t}\,dt$ 一致有界.

$3^\circ$ 当 $n>b$ 时，$\forall x\in[0,b]$，
\[
  \frac1{(n+1)!}\int_0^x t^{n+1}e^{-t}\,dt
  =\frac1{n!}\int_0^x\frac{t}{n+1}\cdot t^ne^{-t}\,dt
  \le\frac1{n!}\int_0^x t^ne^{-t}\,dt,
\]
即 $\dfrac1{n!}\displaystyle\int_0^x t^ne^{-t}\,dt$ 关于 $n$ 单调.

根据 Abel 判别法，$\displaystyle\sum_{n=1}^{\infty}a_n\frac1{n!}\int_0^x t^ne^{-t}\,dt$
在 $[0,b]$ 上一致收敛.

\noindent\textbf{$\star$例 5.2.24}\quad 证明：级数
$\displaystyle\sum_{n=1}^{\infty}(-1)^n\dfrac{e^{x^2}+\sqrt n}{n^{3/2}}$
在任何有限区间 $[a,b]$ 上一致收敛，但在任何一点 $x_0$ 处不绝对收敛.（四川大学）

第二结论根据
\[
  \sum_{n=1}^{\infty}\frac{e^{x^2}+\sqrt n}{n^{3/2}}
  =\sum_{n=1}^{\infty}\frac{e^{x^2}}{n^{3/2}}+\sum_{n=1}^{\infty}\frac1n
\]
明显成立，这里只证第一结论.

\noindent\textbf{证 I}\quad $1^\circ$
\[
  \left|\sum_{k=1}^{n}(-1)^k\right|\le2\quad(n=1,2,\cdots).
\]

$2^\circ$
\[
  \frac{e^{x^2}+\sqrt n}{n^{3/2}}
  =\frac{e^{x^2}}{n^{3/2}}+\frac1n
  \ge\frac{e^{x^2}}{(n+1)^{3/2}}+\frac1{n+1}\quad(n=1,2,\cdots),
\]
即 $\dfrac{e^{x^2}+\sqrt n}{n^{3/2}}$ 关于 $n$ 单调下降.

$3^\circ$ 当 $x\in[a,b]$ 时，
\[
  \left|\frac{e^{x^2}+\sqrt n}{n^{3/2}}\right|
  \le\frac{e^{c^2}+\sqrt n}{n^{3/2}}\to0\quad(n\to\infty)
\]
（其中 $c=\max\{|a|,|b|\}$），因此
$\dfrac{e^{x^2}+\sqrt n}{n^{3/2}}\rightrightarrows0$（当 $n\to\infty$ 时）.

根据 Dirichlet 判别法，$\displaystyle\sum_{n=1}^{\infty}(-1)^n\dfrac{e^{x^2}+\sqrt n}{n^{3/2}}$
在 $[a,b]$ 上一致收敛.

\noindent\textbf{证 II}\quad
\[
  \sum_{n=1}^{\infty}(-1)^n\frac{e^{x^2}+\sqrt n}{n^{3/2}}
  =\sum_{n=1}^{\infty}(-1)^n\frac{e^{x^2}}{n^{3/2}}
  +\sum_{n=1}^{\infty}\frac{(-1)^n}{n}.
\]
因为级数 $\displaystyle\sum_{n=1}^{\infty}(-1)^n\dfrac1{n^{3/2}}$ 与
$\displaystyle\sum_{n=1}^{\infty}\dfrac{(-1)^n}{n}$ 皆为 Leibniz 级数，故收敛，自然关于
$x$ 一致收敛. 又因为 $e^{x^2}$ 在 $[a,b]$ 上为有界函数，一致收敛级数各项同乘某有界函数后仍一致收敛. 所以，
$\displaystyle\sum_{n=1}^{\infty}(-1)^n\dfrac{e^{x^2}}{n^{3/2}}$ 一致收敛. 进而两个一致收敛级数的和级数
$\displaystyle\sum_{n=1}^{\infty}(-1)^n\dfrac{e^{x^2}+\sqrt n}{n^{3/2}}$ 也一致收敛.

\noindent\textbf{证 III}\quad $\forall x\in[a,b]$，
$\displaystyle\sum_{n=1}^{\infty}(-1)^n\dfrac{e^{x^2}+\sqrt n}{n^{3/2}}$ 是 Leibniz 级数，故收敛，且余和的绝对值：
\[
  |r_n(x)|=\left|\frac{e^{x^2}+\sqrt{n+1}}{(n+1)^{3/2}}\right|
  \le\frac{e^{c^2}}{(n+1)^{3/2}}+\frac1{n+1}\to0\quad(n\to\infty)
\]
（其中 $c=\max\{|a|,|b|\}$）. 故 $r_n(x)\rightrightarrows0$（$n\to\infty$）于
$[a,b]$. 因此级数在 $[a,b]$ 上一致收敛.

\noindent\textbf{注}\quad 本例说明一致收敛不意味着绝对收敛.

\noindent\textbf{例 5.2.25}\quad 判断级数
$\displaystyle\sum_{n=1}^{\infty}\frac{(-1)^{n-1}x^2}{(1+x^2)^n}$ 在 $(-\infty,+\infty)$ 内的一致收敛性.

\noindent\textbf{提示}\quad 考虑
\[
  \sum_{n=1}^{\infty}\frac{(-1)^{n-1}}{n}\cdot\frac{nx^2}{(1+x^2)^n}.
\]
Abel 判别法与 Dirichlet 判别法有时可连环使用. 如

\noindent\textbf{$\star$例 5.2.26}\quad 证明：
\[
  \sum_{n=1}^{\infty}(1-x)\frac{x^n}{1-x^{2n}}\sin nx
\]
在 $\left(\dfrac12,1\right)$ 内一致收敛.

\noindent\textbf{证}\quad 因
\[
  \sum_{n=1}^{\infty}(1-x)\frac{x^n}{1-x^{2n}}\sin nx
  =\sum_{n=1}^{\infty}\frac1{1+x^n}\cdot\frac{(1-x)x^n}{1-x^n}\sin nx,
\]
其中 $\dfrac1{1+x^n}$ 关于 $n$ 单调，且一致有界：
$\left|\dfrac1{1+x^n}\right|\le1$. 根据 Abel 判别法，要证明该级数一致收敛，只要证明级数
$\displaystyle\sum_{n=1}^{\infty}\dfrac{(1-x)x^n}{1-x^n}\sin nx$
在 $\left(\dfrac12,1\right)$ 内一致收敛. 但是
\[
\begin{aligned}
1^\circ\quad
\left|\sum_{k=1}^{n}\sin kx\right|
&=\left|\frac1{2\sin\dfrac{x}{2}}\sum_{k=1}^{n}2\sin\frac{x}{2}\sin kx\right|\\
&=\left|\frac1{2\sin\dfrac{x}{2}}\sum_{k=1}^{n}\left[\cos\left(k-\frac12\right)x-\cos\left(k+\frac12\right)x\right]\right|\\
&=\left|\frac{\cos\dfrac12x-\cos\left(nx+\dfrac{x}{2}\right)}{2\sin\dfrac{x}{2}}\right|\\
&\le\frac1{\sin\dfrac{x}{2}}\le\frac1{\sin\dfrac14}\quad
\left(x\in\left(\frac12,1\right),\ n=1,2,\cdots\right).
\end{aligned}
\qquad(1)
\]
（此即表明，$\displaystyle\sum_{k=1}^{n}\sin kx$ 在 $\left(\dfrac12,1\right)$ 内一致有界.）

$2^\circ$
\[
  \frac{(1-x)x^n}{1-x^n}
  =\frac{x^n}{1+x+x^2+\cdots+x^{n-1}},\quad x\in\left(\frac12,1\right),
\]
可见它关于 $n$ 单调. 又
\[
  0\le\frac{x^n}{1+x+\cdots+x^{n-1}}
  \le\frac{x^n}{nx^{n-1}}<\frac1n\to0
  \quad\left(x\in\left(\frac12,1\right)\right),
\]
即在 $\left(\dfrac12,1\right)$ 上，
$\dfrac{(1-x)x^n}{1-x^n}\rightrightarrows0$（当 $n\to\infty$ 时）. 于是根据 Dirichlet 判别法，级数
$\displaystyle\sum_{n=1}^{\infty}\dfrac{(1-x)x^n}{1-x^n}\sin nx$
在 $\left(\dfrac12,1\right)$ 内一致收敛. 证毕.

\noindent\textbf{注}\quad 若将区间改为 $[0,1)$，上面证法中，式（1）不再成立. 但假如我们将 $[0,1)$ 分为两段，在
$\left[0,\dfrac12\right]$ 上，
\[
  \left|(1-x)\frac{x^n}{1-x^{2n}}\sin nx\right|
  \le\frac{\dfrac1{2^n}}{1-\dfrac1{2^{2n}}},
\]
利用 $M$ 判别法容易证明该级数在 $\left[0,\dfrac12\right]$ 上一致收敛；在
$\left(\dfrac12,1\right)$ 内，应用本题的结果. 于是可知原级数在 $[0,1)$ 上一致收敛.

\noindent\textbf{3）Dini 定理及其应用}

\noindent\textbf{$\star$例 5.2.27（Dini 定理）}\quad 设 $u_n(x)\ge0$ 在 $[a,b]$ 上连续，
$n=1,2,\cdots$. 又 $\displaystyle\sum_{n=1}^{\infty}u_n(x)$ 在 $[a,b]$ 上收敛于连续函数
$f(x)$，则 $\displaystyle\sum_{n=1}^{\infty}u_n(x)$ 在 $[a,b]$ 上一致收敛于 $f(x)$.（东北师范大学，南京大学）

\noindent\textbf{证 I}\quad （从正面证明.）（已知 $u_n(x)\ge0$（$n=1,2,\cdots$），
$\displaystyle\sum_{n=1}^{\infty}u_n(x)=f(x)$，因此 $\forall x\in[a,b]$，
$S_n(x)=\displaystyle\sum_{k=1}^{n}u_k(x)\nearrow f(x)$（$n\to\infty$），
$r_n(x)=f(x)-S_n(x)\searrow0$. 我们要证明
$S_n(x)\rightrightarrows f(x)$（$n\to\infty$）在 $[a,b]$ 上，只需证明
$r_n(x)\rightrightarrows0$（$n\to\infty$）于 $[a,b]$. 因此问题在于对任意 $\varepsilon>0$，找到
$N>0$，使得当 $n>N$ 时，$|r_n(x)|<\varepsilon$（$\forall x\in[a,b]$）. 我们的做法是先在局部里找到
$N$，然后在整体上找出 $N$.）

已知 $r_n(x)\searrow0$，所以 $\forall n\in\mathbb N,\ \forall x\in[a,b]$，有
$r_n(x)\ge0$，且 $\forall x_0\in[a,b],\ \forall\varepsilon>0$，
$\exists N(x_0,\varepsilon)>0$，当 $n\ge N(x_0,\varepsilon)$ 时，有
$0\le r_n(x_0)<\varepsilon$. 将 $n$ 固定，令 $n=N_0=N(x_0,\varepsilon)$，因为
$r_n(x)=f(x)-S_n(x)$ 在 $[a,b]$ 上连续. 既然 $r_n(x_0)<\varepsilon$，所以
$\exists\delta_0>0$，当 $x\in(x_0-\delta_0,x_0+\delta_0)$ 时，$r_n(x)<\varepsilon$，从而当
$n>N_0$ 时有 $r_n(x)<\varepsilon$，即 $|r_n(x)|<\varepsilon$
（$x\in(x_0-\delta_0,x_0+\delta_0)$）.

（至此，我们虽然在整个 $[a,b]$ 上未能找到公共的 $N>0$，使得 $n>N$ 时，
$|r_n(x)|<\varepsilon$ 在 $[a,b]$ 上同时成立，但我们已能在任一点 $x_0\in[a,b]$ 的某个邻域里找到一个相应的
$N_0$，使得 $n>N_0$ 时，在此邻域里恒有 $|r_n(x)|<\varepsilon$. 剩下的问题是从局部转化为整体，这就要用到有限覆盖定理.）

如上所述，对每个点 $x_\lambda\in[a,b]$，可找到相应的邻域
$(x_\lambda-\delta_\lambda,x_\lambda+\delta_\lambda)$ 及对应的 $N_\lambda$，使得
$n>N_\lambda$ 时，对 $x\in(x_\lambda-\delta_\lambda,x_\lambda+\delta_\lambda)$ 恒有
$|r_n(x)|<\varepsilon$. 如此，
$\{(x_\lambda-\delta_\lambda,x_\lambda+\delta_\lambda)\mid x_\lambda\in[a,b]\}$
构成了 $[a,b]$ 的一个开覆盖，从而必存在有限子覆盖. 不妨记它们为
$\{(x_1-\delta_1,x_1+\delta_1),\cdots,(x_r-\delta_r,x_r+\delta_r)\}$，于是
$\forall x\in[a,b]$，总 $\exists i\in\{1,2,\cdots,r\}$ 使得
$x\in(x_i-\delta_i,x_i+\delta_i)$，取 $N=\max\{N_1,N_2,\cdots,N_r\}$，那么当
$n>N$ 时，恒有 $|r_n(x)|<\varepsilon$. 证毕.

\noindent\textbf{证 II}\quad （反证法）若在 $[a,b]$ 上非一致收敛，则 $\exists\varepsilon_0>0$，使得
$\forall N>0,\ \exists n>N,\ \exists x\in[a,b]$ 使得 $|r_n(x)|\ge\varepsilon_0$. 取
$N=1$，知 $\exists n_1>1,\ \exists x_1\in[a,b]$，使得
$|r_{n_1}(x_1)|\ge\varepsilon_0$；令 $N=n_1$，知
$\exists n_2>n_1,\ \exists x_2\in[a,b]$，使得
$|r_{n_2}(x_2)|\ge\varepsilon_0$；如此下去，我们得到
$\{n_k\}$ 的子列 $n_1<n_2<\cdots<n_k<\cdots$，使得
\[
  |r_{n_k}(x_k)|\ge\varepsilon_0\quad(k=1,2,\cdots).
\qquad(1)
\]
利用致密性原理（即 Bolzano--Weierstrass 定理），在有界数列 $\{x_k\}$ 里，存在收敛子列
$\{x_{k_j}\}\to x_0\in[a,b]$（当 $j\to\infty$ 时）. 因
$|r_n(x)|\searrow$（关于 $n$），所以 $\forall m\in\mathbb N$，当 $n_{k_j}>m$ 时，有
\[
  |r_m(x_{k_j})|\ge |r_{n_{k_j}}(x_{k_j})|\ge\varepsilon_0.
\]
由于 $r_m(x)=f(x)-S_m(x)$ 连续，所以当 $j\to\infty$ 时，在
$|r_m(x_{k_j})|\ge\varepsilon_0$ 里取极限，知
$|r_m(x_0)|\ge\varepsilon_0$（$\forall m\in\mathbb N$），与
$\displaystyle\sum_{n=1}^{\infty}u_n(x)$ 在 $[a,b]$ 上收敛矛盾. 证毕.

\noindent\textbf{注}\quad Dini 定理条件 $u_n(x)\ge0$（$n=1,2,\cdots$）改变为：“固定 $x$ 时，各
$u_n(x)$ 保持同号（当 $x$ 变化时 $u_n(x)$ 可以变号）”时，结论仍然成立. 此时 $x$ 固定，令
$n\to\infty$，仍有 $|r_n(x)|\searrow0$；上述证明适当修改后，依然有效.

\noindent\textbf{$\star$例 5.2.28}\quad 在区间 $[0,1]$ 上：

1）证明函数序列 $\left(1+\dfrac{x}{n}\right)^n$（$n=1,2,\cdots$）一致收敛；

2）证明函数序列
\[
  f_n(x)=\frac1{e^{x/n}+\left(1+\dfrac{x}{n}\right)^n}\quad(n=1,2,\cdots)
\]
一致收敛；

3）求出极限
\[
  \lim_{n\to\infty}\int_0^1
  \frac{dx}{e^{x/n}+\left(1+\dfrac{x}{n}\right)^n}.
\]
（武汉大学）

\noindent\textbf{证}\quad 1）\textbf{证 I}\quad 当 $n\to\infty$ 时，
$\left(1+\dfrac{x}{n}\right)^n\nearrow e^x$，且全部在 $[0,1]$ 上连续，故由 Dini 定理知
$\left(1+\dfrac{x}{n}\right)^n\rightrightarrows e^x$.

\noindent\textbf{证 II}\quad 由
\[
  \left[e^x-\left(1+\frac{x}{n}\right)^n\right]'_x
  =e^x-\left(1+\frac{x}{n}\right)^{n-1}>0,
\]
知 $e^x-\left(1+\dfrac{x}{n}\right)^n\nearrow$（关于 $x$），故当 $x\in[0,1]$ 时，
\[
  0\le e^x-\left(1+\frac{x}{n}\right)^n
  \le e-\left(1+\frac1n\right)^n\to0\quad(n\to\infty).
\]
所以在 $[0,1]$ 上，$\left(1+\dfrac{x}{n}\right)^n\rightrightarrows e^x$（$n\to\infty$）.

2）由于当 $x\in[0,1]$ 时，
\[
\begin{aligned}
  \left|\frac1{1+e^x}
  -\frac1{e^{x/n}+\left(1+\dfrac{x}{n}\right)^n}\right|
  &=
  \left|\frac{e^{x/n}+\left(1+\dfrac{x}{n}\right)^n-1-e^x}
  {(1+e^x)\left[e^{x/n}+\left(1+\dfrac{x}{n}\right)^n\right]}\right|\\
  &\le \left|\left(1+\frac{x}{n}\right)^n-e^x\right|+|e^{x/n}-1|\\
  &=e^x-\left(1+\frac{x}{n}\right)^n+e^{x/n}-1\to0\quad(n\to\infty),
\end{aligned}
\]
因此，在 $[0,1]$ 上，
\[
  \frac1{e^{x/n}+\left(1+\dfrac{x}{n}\right)^n}
  \rightrightarrows\frac1{1+e^x}\quad(n\to\infty).
\]

3）由 2）知，可在积分号下取极限：
\[
  \lim_{n\to\infty}\int_0^1
  \frac{dx}{e^{x/n}+\left(1+\dfrac{x}{n}\right)^n}
  =\int_0^1\frac1{1+e^x}\,dx
  =\int_0^1\frac{d(e^x)}{e^x(1+e^x)}
  =1+\ln\frac2{1+e}.
\]

\noindent\textbf{d.\quad 利用一致有界与等度连续证明一致收敛性}

\noindent\textbf{$\star$例 5.2.29}\quad 设 $\{f_n(x)\}$ 是区间 $(a,b)$ 内的连续函数序列，并且对任一
$x_0\in(a,b)$，$\{f_n(x_0)\}$ 都是有界的. 证明：$\{f_n(x)\}$ 在 $(a,b)$ 的某一非空子区间上一致有界.（北京大学，云南大学）

\noindent\textbf{证}\quad （反证法）若 $\{f_n(x)\}$ 在 $(a,b)$ 内任何非空子区间上都不一致有界，那么在
$(a,b)$ 内就可找到一个区间套
$\Delta_1\supset\Delta_2\supset\cdots\supset\Delta_n\supset\cdots$，使得在
$\Delta_n$ 上恒有 $f_n(x)>n$，如此在区间套的公共点 $x_0$ 上，$\{f_n(x_0)\}$ 无界. 与已知条件矛盾.

假设 $\{f_n(x)\}$ 在任何（非空）子区间上都不一致有界，则
$\exists x_1\in(a,b)$ 及 $n_1\in\mathbb N$，使得 $f_{n_1}(x_1)>1$. 又因
$f_{n_1}$ 连续，根据保号性，在含 $x_1$ 的某个闭子区间 $\Delta_1\subset(a,b)$ 上，恒有
$f_{n_1}(x)>1$.

$\{f_n(x)\}$ 在 $\Delta_1$ 上仍不一致有界，所以
$\exists x_2\in\Delta_1$ 及 $n_2\in\mathbb N$，使得 $f_{n_2}(x_2)>2$. 根据连续保号性，
$\exists$ 闭子区间 $\Delta_2\subset\Delta_1$，使得 $\Delta_2$ 上面有
$f_{n_2}(x)>2$. 如此下去，使得一串闭区间
$\Delta_1\supset\Delta_2\supset\cdots\supset\Delta_k\supset\cdots$，在
$\Delta_k$ 上恒有 $f_{n_k}(x)>k$.

利用区间套定理，$\exists x_0\in\Delta_k$（$k=1,2,\cdots$）及 $n_k\in\mathbb N$，使得
$f_{n_k}(x_0)>k$（$k=1,2,\cdots$）. 故 $\{f_n(x)\}$ 在 $x_0\in(a,b)$ 处无界，与已知条件矛盾.

\noindent\textbf{例 5.2.30}\quad 设 $\{f_n(x)\}$ 在区间 $[0,1]$ 上一致有界，试证存在一个子列，其在
$[0,1]$ 的一切有理点上收敛.

\noindent\textbf{证}\quad 我们知道 $[0,1]$ 的全体有理点可以排成一个数列 $\{a_n\}$（如
$\{a_n\}=\{0,1,\dfrac12,\dfrac13,\dfrac23,\dfrac14,\dfrac34,\dfrac15,\dfrac25,\dfrac35,\dfrac45,\dfrac16,\cdots\}$）.

因 $\{f_n(x)\}$ 一致有界，故 $\{f_n(a_1)\}$ 是有界数列. 由致密性原理，其中存在收敛子列，为了便于叙述，记此收敛子列为
$\{f_{1n}(a_1)\}$，于是 $\{f_{1n}(x)\}\subset\{f_n(x)\}$ 在 $x=a_1$ 处收敛. 同理，因
$\{f_{1n}(a_2)\}$ 是有界数列，又必存在收敛子列 $\{f_{2n}(a_2)\}$. 即
$\{f_{2n}(x)\}\subset\{f_{1n}(x)\}$，$\{f_{2n}(x)\}$ 在 $x=a_1,a_2$ 处都收敛. 如此进行下去，不断地在子列里取子列，使
$\{f_{kn}(x)\}$ 在 $a_1,a_2,\cdots,a_k$ 处收敛，于是得到一串子列：
\[
\begin{matrix}
  f_{11}(x),& f_{12}(x),& f_{13}(x),& \cdots,& f_{1n}(x),& \cdots\\
  f_{21}(x),& f_{22}(x),& f_{23}(x),& \cdots,& f_{2n}(x),& \cdots\\
  f_{31}(x),& f_{32}(x),& f_{33}(x),& \cdots,& f_{3n}(x),& \cdots\\
  \cdots\\
  f_{n1}(x),& f_{n2}(x),& f_{n3}(x),& \cdots,& f_{nn}(x),& \cdots\\
  \cdots
\end{matrix}
\qquad(1)
\]
最后用上表的对角线元素组成一个子列 $\{f_{nn}(x)\}$，即
$f_{11}(x),f_{22}(x),f_{33}(x),\cdots,f_{nn}(x),\cdots$. 易知此级数在点
$a_i$（$i=1,2,\cdots$）处收敛. 事实上，$\forall a_i$（$i=1,2,\cdots$），已知
（1）中第 $i$ 个子列在 $a_i$ 处收敛，而
$f_{ii}(x),f_{i+1,i+1}(x),\cdots$ 是第 $i$ 个子列的子列. 故
$\{f_{nn}(x)\}$ 在 $a_i$ 点收敛. 由此知 $\{f_{nn}(x)\}$ 在
$\{a_1,a_2,\cdots,a_n,\cdots\}$ 上收敛.

\noindent\textbf{等度连续}

\noindent\textbf{定义}\quad 设 $\mathcal M$ 是区间 $I$ 上定义的函数族，所谓族 $\mathcal M$ 上的函数在 $I$ 上等度连续，是指：
$\forall\varepsilon>0,\ \exists\delta>0$，当 $x_1,x_2\in I$ 且
$|x_1-x_2|<\delta$ 时，有
\[
  |f(x_1)-f(x_2)|<\varepsilon\quad(\forall f\in\mathcal M).
\]
特别，$I$ 上定义的函数序列 $\{f_n(x)\}$ 在 $I$ 上等度连续，是指：
$\forall\varepsilon>0,\ \exists\delta>0$，当 $x_1,x_2\in I$ 且
$|x_1-x_2|<\delta$ 时，有
\[
  |f_n(x_1)-f_n(x_2)|<\varepsilon\quad(\forall n\in\mathbb N).
\]
显然，若 $\mathcal M$ 是有限族（即由有限个函数组成），且 $I$ 为有界闭区间，那么
$\mathcal M$ 中每个函数连续，就必然等度连续. 下面会看到，若 $\mathcal M$ 为无穷族，
$\mathcal M$ 中每个成员连续，$\mathcal M$ 不见得是等度连续的.

\noindent\textbf{例 5.2.31}\quad 证明：若序列 $\{f_n(x)\}$ 在 $I$ 上等度连续，且
$\displaystyle\lim_{n\to\infty}f_n(x)=f(x)$（$x\in I$），那么 $f(x)$ 在 $I$ 上一致连续.

\noindent\textbf{证}\quad 因 $\{f_n(x)\}$ 等度连续，所以 $\forall\varepsilon>0,\ \exists\delta>0$，当
$x_1,x_2\in I,\ |x_1-x_2|<\delta$ 时，有
$|f_n(x_1)-f_n(x_2)|<\dfrac\varepsilon2$. 令 $n\to\infty$，取极限可得
$|f(x_1)-f(x_2)|\le\dfrac\varepsilon2<\varepsilon$. 此即表明 $f(x)$ 在 $I$ 上一致连续.

该例结果表明，若 $\displaystyle\lim_{n\to\infty}f_n(x)=f(x)$，而 $f(x)$ 不一致连续，则
$\{f_n(x)\}$ 不可能等度连续. 例如：$\{x^n\}$ 在 $[0,1)$ 上正是如此.

\noindent\textbf{例 5.2.32}\quad 若 $\mathcal M$ 是区间 $I$ 上定义的函数族，
$\forall f\in\mathcal M$ 皆在 $I$ 上可微，且
$\{|f'(x)|\mid f\in\mathcal M\}$ 在 $I$ 上一致有界，那么 $\mathcal M$ 在 $I$ 上等度连续.

\noindent\textbf{证}\quad 因 $\{|f'(x)|\mid f\in\mathcal M\}$ 一致有界，故
$\exists M>0$ 使得 $|f'(x)|\le M$（$\forall x\in I,\ \forall f\in\mathcal M$）. 于是
$\forall\varepsilon>0$，取 $\delta=\dfrac\varepsilon M$，则当
$x_1,x_2\in I$ 且 $|x_1-x_2|<\delta$ 时，恒有
\[
  |f(x_1)-f(x_2)|=|f'(\xi)(x_1-x_2)|
  \le M|x_1-x_2|<M\cdot\frac\varepsilon M=\varepsilon\quad(\forall f\in\mathcal M),
\]
即 $\mathcal M$ 在 $I$ 上等度连续.

\noindent\textbf{$\star$例 5.2.33}\quad 设函数序列 $\{f_n(x)\}$ 在区间 $[a,b]$ 上为等度连续的. 试证：若在
$[a,b]$ 上 $f_n(x)\to f(x)$（$n\to\infty$），则
$f_n(x)\rightrightarrows f(x)$ 在 $[a,b]$ 上（$n\to\infty$）.

\noindent\textbf{证 I}\quad （$\forall\varepsilon>0$，先对每个 $x_\lambda\in[a,b]$ 找
$\delta_\lambda>0$，使得当 $|x-x_\lambda|<\delta_\lambda$ 时，
$|f(x)-f(x_\lambda)|<\varepsilon$，然后应用有限覆盖定理.）

$\forall x_\lambda\in[a,b]$，因 $\displaystyle\lim_{n\to\infty}f_n(x_\lambda)=f(x_\lambda)$，所以
$\forall\varepsilon>0,\ \exists N_\lambda=N(\varepsilon,x_\lambda)$，当
$n>N_\lambda$ 时有 $|f_n(x_\lambda)-f(x_\lambda)|<\dfrac\varepsilon3$. 由于
$\{f_n\}$ 等度连续，从而 $f$ 连续（见例 5.2.31）. 故对此 $\varepsilon>0$，
$\exists\delta_\lambda>0$，当 $|x-x_\lambda|<\delta_\lambda$ 时，有
\[
  |f_n(x)-f_n(x_\lambda)|<\frac\varepsilon3,\qquad
  |f(x)-f(x_\lambda)|<\frac\varepsilon3.
\]
于是
\[
\begin{aligned}
  |f_n(x)-f(x)|
  &\le |f_n(x)-f_n(x_\lambda)|
    +|f_n(x_\lambda)-f(x_\lambda)|+|f(x_\lambda)-f(x)|\\
  &<\frac\varepsilon3+\frac\varepsilon3+\frac\varepsilon3=\varepsilon.
\end{aligned}
\]
这时 $\{(x_\lambda-\delta_\lambda,x_\lambda+\delta_\lambda)\mid x_\lambda\in[a,b]\}$
组成 $[a,b]$ 的一个开覆盖. 根据有限覆盖定理，其中存在有限子覆盖，记之为
$\{(x_i-\delta_i,x_i+\delta_i)\mid i=1,2,\cdots,r\}$. 取
$N=\max\{N_1,N_2,\cdots,N_r\}$，则 $n>N$ 时，$\forall x\in[a,b]$，
$\exists i\in\{1,2,\cdots,r\}$ 使得 $x\in(x_i-\delta_i,x_i+\delta_i)$，从而有
$|f_n(x)-f(x)|<\varepsilon$. 这就证明了：在 $[a,b]$ 上
$f_n(x)\rightrightarrows f(x)$（$n\to\infty$）.

\noindent\textbf{证 II}\quad 由 $\{f_n(x)\}$ 在 $[a,b]$ 上等度连续，知 $f(x)$ 在
$[a,b]$ 上一致连续. 因而 $\forall\varepsilon>0,\ \exists\delta>0$，当
$x',x''\in[a,b]$，$|x'-x''|<\delta$ 时，有
\[
  |f_n(x')-f_n(x'')|<\frac\varepsilon3,\qquad
  |f(x')-f(x'')|<\frac\varepsilon3.
\]
令将 $[a,b]$ $k$ 等分，使每个小区间的长度小于 $\delta$（这是可以办到的，只要令
$\dfrac{b-a}{k}<\delta$，即 $k>\dfrac{b-a}{\delta}$ 便可）. 记 $k$ 等分的各分点为
$a=a_0<a_1<\cdots<a_k=b$，因为 $f_n(a_i)\to f(a_i)$（$n\to\infty$），所以对上述
$\varepsilon>0$，$\exists N_i>0$ 使得 $n>N_i$ 时，有
$|f_n(a_i)-f(a_i)|<\dfrac\varepsilon3$（$i=1,2,\cdots,k$）. 令
$N=\max\{N_1,N_2,\cdots,N_k\}$，则当 $n>N$ 时，$\forall x\in[a,b]$，
$\exists a_i$（$i\in\{1,2,\cdots,k\}$）使得 $|a_i-x|<\delta$，
\[
\begin{aligned}
  |f_n(x)-f(x)|
  &\le |f_n(x)-f_n(a_i)|+|f_n(a_i)-f(a_i)|+|f(a_i)-f(x)|\\
  &<\frac\varepsilon3+\frac\varepsilon3+\frac\varepsilon3=\varepsilon.
\end{aligned}
\]
这就证明了：在 $[a,b]$ 上 $f_n(x)\rightrightarrows f(x)$（$n\to\infty$）.

\noindent\textbf{注}\quad 从证 II 中容易看出，条件“$f_n(x)\to f(x)$（$n\to\infty$）”只需在
$[a,b]$ 上某个稠密子集 $\{a_i\}$ 上成立就够了.

\noindent\textbf{new 练习}\quad 设 $\{f_n(x)\}$ 是 $[a,b]$ 上定义的函数序列，满足：

1）$\forall x_0\in[a,b]$，$\{f_n(x_0)\}$ 有界；
2）$\{f_n(x)\}$ 在 $[a,b]$ 上等度连续.

试证：存在子列 $\{f_{n_k}(x)\}$ 在 $[a,b]$ 上一致收敛.（南开大学）

\noindent\textbf{提示}\quad 设 $\{a_n\}$ 是 $[a,b]$ 上全体有理数组成的数列. 根据条件 1），对于
$a_i$，数列 $\{f_n(a_i)\}$ 有界，可应用致密性原理，采用例 5.2.30 中的方法，从子列里找子列，最后取对角线元素，可作出
$\{f_n(x)\}$ 的子列 $\{f_{n_k}(x)\}$，使之在 $\{a_n\}$ 的每点都收敛.

\noindent\textbf{证}\quad （利用等度连续，证明 $\{f_{n_k}(x)\}$ 在无理点上也收敛.）因
$\{f_n(x)\}$ 等度连续：$\forall\varepsilon>0,\ \exists\delta>0$，$x',x''\in[a,b]$，
$|x'-x''|<\delta$，有
\[
  |f_n(x')-f_n(x'')|<\frac\varepsilon3\quad(\forall n>0).
\]
因此，对于上面作出的子列也有
\[
  |f_{n_k}(x')-f_{n_k}(x'')|<\frac\varepsilon3\quad(\forall k>0).
\]
设 $x_0$ 是 $[a,b]$ 上任一无理点，根据有理点的稠密性，可取有理点 $x_r$，使得
$|x_r-x_0|<\delta$. 于是
\[
  |f_{n_k}(x_r)-f_{n_k}(x_0)|<\frac\varepsilon3,\qquad
  |f_{n_{k+p}}(x_r)-f_{n_{k+p}}(x_0)|<\frac\varepsilon3.
\]
又因 $\{f_{n_k}(x)\}$ 在 $x_r$ 收敛，（利用 Cauchy 准则必要性）$\exists N>0$，当
$k>N$ 时，$|f_{n_k}(x_r)-f_{n_{k+p}}(x_r)|<\dfrac\varepsilon3$（$\forall p\in\mathbb N$）. 所以
\[
\begin{aligned}
  |f_{n_k}(x_0)-f_{n_{k+p}}(x_0)|
  &\le |f_{n_k}(x_0)-f_{n_k}(x_r)|
    +|f_{n_k}(x_r)-f_{n_{k+p}}(x_r)|\\
  &\quad +|f_{n_{k+p}}(x_r)-f_{n_{k+p}}(x_0)|\\
  &<\frac\varepsilon3+\frac\varepsilon3+\frac\varepsilon3=\varepsilon\quad(\forall p\in\mathbb N).
\end{aligned}
\]
（根据 Cauchy 准则充分性）$\{f_{n_k}(x_0)\}$ 收敛. 由无理点 $x_0\in[a,b]$ 的任意性，
$\{f_{n_k}(x)\}$ 在无理点上都收敛.

以上证明了：$\{f_{n_k}(x)\}$ 在 $[a,b]$ 上逐点收敛，再加上等度连续性，根据例 5.2.33（或重复证法），可知
$\{f_{n_k}(x)\}$ 在 $[a,b]$ 上一致收敛.

作为例 5.2.33 的逆命题：

\noindent\textbf{$\star$例 5.2.34}\quad 设 $f_n(x)$（$n=1,2,\cdots$）在闭区间 $[a,b]$ 上连续，且当
$n\to\infty$ 时，$f_n(x)\rightrightarrows f(x)$ 于 $[a,b]$，试证 $\{f_n(x)\}$ 等度连续.（郑州大学）

\noindent\textbf{证}\quad 因 $[a,b]$ 上 $f_n(x)$ 连续，且 $f_n(x)\rightrightarrows f(x)$（当
$n\to\infty$ 时），故 $f(x)$ 连续，从而在 $[a,b]$ 上一致连续：
$\forall\varepsilon>0,\ \exists\delta_0>0$，当
$x',x''\in[a,b],\ |x'-x''|<\delta_0$ 时，有
$|f(x')-f(x'')|<\dfrac\varepsilon3$. 由一致收敛性，对此 $\varepsilon>0$，
$\exists N>0$，当 $n>N$ 时，有 $|f_n(x)-f(x)|<\dfrac\varepsilon3$
（$\forall x\in[a,b]$）. 如此
\[
\begin{aligned}
  |f_n(x')-f_n(x'')|
  &\le |f_n(x')-f(x')|+|f(x')-f(x'')|+|f(x'')-f_n(x'')|\\
  &<\frac\varepsilon3+\frac\varepsilon3+\frac\varepsilon3=\varepsilon.
\end{aligned}
\]
对剩下的 $f_1(x),f_2(x),\cdots,f_N(x)$ 应用 Cantor 定理知，对每个
$f_i$（$i=1,2,\cdots,N$），$\exists\delta_i>0$（$i=1,2,\cdots,N$），当
$x',x''\in[a,b],\ |x'-x''|<\delta_i$ 时，有 $|f_i(x')-f_i(x'')|<\varepsilon$. 如此令
$\delta=\min\{\delta_0,\delta_1,\cdots,\delta_N\}$，当
$\forall x',x''\in[a,b],\ |x'-x''|<\delta$ 时，对一切 $n\in\mathbb N$，恒有
$|f_n(x')-f_n(x'')|<\varepsilon$. 证毕.

\noindent\textbf{例 5.2.35}\quad 可微函数序列 $\{f_n(x)\}$ 在 $[a,b]$ 上收敛，且
$\exists M>0$，使 $|f_n'(x)|\le M$（$n=1,2,\cdots,\ x\in[a,b]$），试证
$\{f_n(x)\}$ 在 $[a,b]$ 上一致收敛.（上海交通大学）

\noindent\textbf{提示}\quad 可利用例 5.2.32 和例 5.2.33，也可直接证明.

\noindent\textbf{练习}\quad 设在 $[a,b]$ 上 $f_n(x)$ 连续，且 $f_n'(x)$ 一致有界，
$f_n(x)$ 逐点收敛于极限函数 $f(x)$. 试证函数 $f(x)$ 在 $[a,b]$ 上连续.（北京大学）

\noindent\textbf{提示}\quad $\exists M>0$：
$|f_n'(x)|\le M$（$\forall x\in[a,b],\ \forall n\in\mathbb N$）
$\Rightarrow |f_n(x')-f_n(x'')|\le M|x'-x''|$
$\Rightarrow |f(x')-f(x'')|\le M|x'-x''|$.

\noindent\textbf{例 5.2.36}\quad 设 $\mathcal M$ 是定义在 $[a,b]$ 上的连续函数族，一致有界且等度连续，试证在
$\mathcal M$ 中存在函数序列 $f_1(x),f_2(x),\cdots,f_n(x),\cdots$ 在
$[a,b]$ 上一致收敛.（吉林大学）

\noindent\textbf{提示}\quad 参看例 5.2.30 及例 5.2.33.

\noindent\textbf{$\star$例 5.2.37}\quad 设函数序列 $\{f_n(x)\}$ 与 $\{g_n(x)\}$ 在区间 $I$ 上一致收敛，而且对每个
$n=1,2,\cdots$，$f_n(x)$ 与 $g_n(x)$ 在 $I$ 上有界（界可随 $n$ 而异），证明
$\{f_n(x)g_n(x)\}$ 在 $I$ 上亦一致收敛.（华东师范大学，北京大学）

\noindent\textbf{分析}\quad 不妨设 $f_n(x)\rightrightarrows f(x)$，
$g_n(x)\rightrightarrows g(x)$（$n\to\infty$，关于 $x\in I$）. 要证明
$f_n(x)g_n(x)\rightrightarrows f(x)g(x)$，利用不等式
\[
\begin{aligned}
  |f(x)g(x)-f_n(x)g_n(x)|
  &\le |f(x)g(x)-f(x)g_n(x)|
    +|f(x)g_n(x)-f_n(x)g_n(x)|\\
  &=|f(x)|\,|g(x)-g_n(x)|+|g_n(x)|\,|f(x)-f_n(x)|.
\end{aligned}
\qquad(1)
\]
可见只要证明：$\exists M>0$，使得
\[
  |f(x)|\le M,\qquad |g_n(x)|\le M\quad(\forall x\in I,\ \forall n\in\mathbb N).
\qquad(2)
\]
因为有了这样的 $M$，则
\[
  \text{式（1）}\le M|g(x)-g_n(x)|+M|f(x)-f_n(x)|.
\qquad(3)
\]
于是根据 $g_n(x)\rightrightarrows g(x)$，$f_n(x)\rightrightarrows f(x)$，可知
$\forall\varepsilon>0,\ \exists N>0$ 使得当 $n>N$ 时，
\[
  |g(x)-g_n(x)|<\frac1{2M}\varepsilon,\qquad
  |f(x)-f_n(x)|<\frac1{2M}\varepsilon,
\]
式（3）$\le M\dfrac{\varepsilon}{2M}+M\dfrac{\varepsilon}{2M}=\varepsilon$.

下面设法证明，存在满足式（2）的 $M$.

事实上，因 $f_n(x)\rightrightarrows f(x)$，所以对于 $\varepsilon=1$，
$\exists n_1$ 使得
\[
  |f(x)-f_{n_1}(x)|<\varepsilon=1\quad(\forall x\in I).
\]
\[
  |f(x)|\le |f(x)-f_{n_1}(x)|+|f_{n_1}(x)|\le1+|f_{n_1}(x)|,
\]
因 $f_{n_1}(x)$ 有界，$\exists M_1>0$，使得
$|f_{n_1}(x)|\le M_1$（$\forall x\in I$）. 故
$|f(x)|\le1+M_1$（$\forall x\in I$）.

同理，由 $g_n(x)\rightrightarrows g(x)$ 及每个 $g_n(x)$ 有界，可知
$\exists M_2>0$ 使得 $|g(x)|\le1+M_2$（$\forall x\in I$）. 又由
$g_n(x)\rightrightarrows g(x)$，可知对 $\varepsilon=1$，
$\exists N>0$，当 $n>N$ 时，$|g_n(x)-g(x)|<1$. 所以
\[
  |g_n(x)|\le |g_n(x)-g(x)|+|g(x)|\le2+M_2\quad(\forall x\in I).
\]
$g_1(x),g_2(x),\cdots,g_N(x)$ 分别有界，因此 $\exists G_i$，使得
$|g_i(x)|\le G_i$（$i=1,2,\cdots,N$）. 最后，令
$M=\max\{1+M_1,2+M_2,G_1,\cdots,G_N\}$，则
\[
  |f(x)|\le M,\qquad |g_n(x)|\le M\quad(\forall x\in I,\ \forall n\in\mathbb N).
\]
问题得证.

\subsection*{二、一致收敛级数的性质}

\noindent\textbf{a.\quad 关于逐项取极限}

\noindent\textbf{例 5.2.38（逐项取极限定理）}\quad 设级数
$\displaystyle\sum_{n=1}^{\infty}u_n(x)$ 在 $x_0$ 的某个空心邻域
$U_0(x_0)=\{x\mid0<|x-x_0|<\delta\}$ 里一致收敛，
$\displaystyle\lim_{x\to x_0}u_n(x)=c_n$，则
$\displaystyle\sum_{n=1}^{\infty}c_n$ 收敛，且
\[
  \lim_{x\to x_0}\sum_{n=1}^{\infty}u_n(x)
  =\sum_{n=1}^{\infty}\lim_{x\to x_0}u_n(x)
  =\sum_{n=1}^{\infty}c_n.
\qquad(1)
\]
（西安电子科技大学）

\noindent\textbf{证 I}\quad $1^\circ$ 因 $\displaystyle\sum_{n=1}^{\infty}u_n(x)$ 在
$U_0(x_0)$ 内一致收敛，所以 $\forall\varepsilon>0,\ \exists N>0$，当 $n>N$ 时，
$\forall p\in\mathbb N$，有
\[
  \left|\sum_{k=n+1}^{n+p}u_k(x)\right|<\varepsilon,\quad \forall x\in U_0(x_0).
\]
令 $x\to x_0$，取极限得
$\left|\displaystyle\sum_{k=n+1}^{n+p}c_k\right|\le\varepsilon$. 由 Cauchy 准则，级数
$\displaystyle\sum_{n=1}^{\infty}c_n=c$ 收敛（$c$ 为某个常数）.

$2^\circ$ 由 $S(x)\equiv\displaystyle\sum_{n=1}^{\infty}u_n(x)$ 一致收敛及
$c=\displaystyle\sum_{n=1}^{\infty}c_n$ 的收敛性，易知 $\forall\varepsilon>0$，
$\exists n\in\mathbb N$，使得
\[
  |S(x)-S_n(x)|<\frac\varepsilon3,\qquad
  \left|c-\sum_{k=1}^{n}c_k\right|<\frac\varepsilon3.
\]
其中 $S_n(x)=\displaystyle\sum_{k=1}^{n}u_k(x)$. 将 $n$ 固定，因
$S_n(x)=\displaystyle\sum_{k=1}^{n}u_k(x)\to\displaystyle\sum_{k=1}^{n}c_k$（当
$x\to x_0$ 时），故对 $\varepsilon>0$，$\exists\delta>0$，当
$|x-x_0|<\delta$ 时，
$\left|S_n(x)-\displaystyle\sum_{k=1}^{n}c_k\right|<\dfrac\varepsilon3$. 从而
\[
  |S(x)-c|
  \le |S(x)-S_n(x)|
  +\left|S_n(x)-\sum_{k=1}^{n}c_k\right|
  +\left|\sum_{k=1}^{n}c_k-c\right|
  <\frac\varepsilon3+\frac\varepsilon3+\frac\varepsilon3=\varepsilon.
\]
即式（1）成立.

\noindent\textbf{证 II}\quad 关于 $\displaystyle\sum_{n=1}^{\infty}c_n$ 的收敛性同上证明，现证明可逐项取极限.
在 $x_0$ 处补充定义，令 $u_n(x_0)=c_n$（$n=1,2,\cdots$），则 $u_n(x)$ 在
$x_0$ 处连续，并且 $\displaystyle\sum_{n=1}^{\infty}u_n(x)$ 在 $x_0$ 的邻域
$(x_0-\delta,x_0+\delta)$ 内一致收敛. 从而由和函数连续定理，可知
$\displaystyle\sum_{n=1}^{\infty}u_n(x)$ 在 $x_0$ 处连续，从而
\[
  \lim_{x\to x_0}\sum_{n=1}^{\infty}u_n(x)
  =\sum_{n=1}^{\infty}u_n(x_0)
  =\sum_{n=1}^{\infty}c_n
  =\sum_{n=1}^{\infty}\lim_{x\to x_0}u_n(x).
\]

\noindent\textbf{例 5.2.39}\quad 假定函数 $u_n(x)$ 在区间 $(0,1)$ 内单调增加，并且
$u_n(x)\ge0,\ n=1,2,\cdots$. 又假定 $\displaystyle\sum_{n=1}^{\infty}u_n(x)$ 在
$(0,1)$ 内逐点收敛，并且有上界，试证
$\displaystyle\sum_{n=1}^{\infty}u_n(x)$ 在 $(0,1)$ 内一致收敛，并且
\[
  \lim_{x\to1^-}\sum_{n=1}^{\infty}u_n(x)
  =\sum_{n=1}^{\infty}\lim_{x\to1^-}u_n(x).
\qquad(1)
\]
（厦门大学）

\noindent\textbf{证}\quad 根据上例（逐项取极限的定理），要证明式（1），只要证明
$\displaystyle\sum_{n=1}^{\infty}u_n(x)$ 在 $(0,1)$ 内一致收敛，并且极限
$\displaystyle\lim_{x\to1^-}u_n(x)$ 存在.

$1^\circ$ 先来证明 $\displaystyle\lim_{x\to1^-}u_n(x)$ 存在. 因已知
$u_n(x)\ge0$，又 $S(x)\equiv\displaystyle\sum_{n=1}^{\infty}u_n(x)$ 收敛并且有上界，所以存在
$M\ge0$，使得 $u_n(x)\le S(x)\le M,\ \forall x\in(0,1)$. 而 $u_n(x)$ 单调增加，故
$\displaystyle\lim_{x\to1^-}u_n(x)$ 存在. 另外，若记
$\displaystyle\lim_{x\to1^-}u_n(x)=u_n(1)$，则在 $(0,1)$ 上有
\[
  0\le u_n(x)\le u_n(1)\quad(n=1,2,\cdots).
\qquad(2)
\]

$2^\circ$ 剩下只要证明 $\displaystyle\sum_{n=1}^{\infty}u_n(x)$ 在 $(0,1)$ 内一致收敛. 式（2）表明只要证明级数
$\displaystyle\sum_{n=1}^{\infty}u_n(1)$ 收敛即可. 因为
$\displaystyle\sum_{k=1}^{n}u_k(x)\le S(x)\le M$，令 $x\to1^-$ 取极限，得
$\displaystyle\sum_{k=1}^{n}u_k(1)\le M$（$n=1,2,\cdots$）. 因
$u_k(1)\ge0$，故 $\displaystyle\sum_{n=1}^{\infty}u_n(1)$ 收敛. 从而根据
$M$ 判别法，$\displaystyle\sum_{n=1}^{\infty}u_n(x)$ 在 $(0,1)$ 上一致收敛. 证毕.

\noindent\textbf{b.\quad 和函数的连续性}

\noindent\textbf{要点}\quad 和函数连续的定理常以如下三种形式叙述：

\noindent\textbf{定理 1}\quad 若 $u_n(x)$ 在区间 $I$ 上连续（$n=1,2,\cdots$），
$\displaystyle\sum_{n=1}^{\infty}u_n(x)$ 在 $I$ 上一致收敛，则
$S(x)\equiv\displaystyle\sum_{n=1}^{\infty}u_n(x)$ 在 $I$ 上连续.

\noindent\textbf{定理 2}\quad 若 $u_n(x)$ 在 $x=x_0$ 处连续（$n=1,2,\cdots$），
$\displaystyle\sum_{n=1}^{\infty}u_n(x)$ 在 $x_0$ 的某个邻域里一致收敛，则
$S(x)\equiv\displaystyle\sum_{n=1}^{\infty}u_n(x)$ 在 $x=x_0$ 处连续.

\noindent\textbf{定理 3}\quad 若 $u_n(x)$ 在 $(a,b)$ 内连续（$n=1,2,\cdots$），
$\displaystyle\sum_{n=1}^{\infty}u_n(x)$ 在 $(a,b)$ 内内闭一致收敛（意指在
$(a,b)$ 内的任一闭区间上分别一致收敛），则
$S(x)\equiv\displaystyle\sum_{n=1}^{\infty}u_n(x)$ 在 $(a,b)$ 内连续.

下面我们将看到，以上三定理在使用时各有好处. 如

\noindent\textbf{例 5.2.40}\quad 证明：
$\displaystyle\sum_{n=-\infty}^{+\infty}\dfrac1{(n-x)^2}$ 当 $x\ne$ 整数时收敛，周期为
1，并且当 $x\ne$ 整数时和函数连续.

\noindent\textbf{证}\quad 因
\[
  \sum_{n=-\infty}^{+\infty}\frac1{(n-x)^2}
  =\sum_{n=0}^{+\infty}\frac1{(n-x)^2}
  +\sum_{n=1}^{+\infty}\frac1{(-n-x)^2}.
\qquad(1)
\]
当 $x\ne$ 整数，$n\to+\infty$ 时，
$\dfrac1{(n-x)^2}\sim\dfrac1{n^2}$，
$\dfrac1{(-n-x)^2}\sim\dfrac1{n^2}$，故级数（1）收敛. 又
\[
\begin{aligned}
  f(x+1)
  &\equiv \sum_{n=-\infty}^{+\infty}\frac1{[n-(x+1)]^2}
  =\sum_{n=-\infty}^{+\infty}\frac1{(n-1-x)^2}\quad(\text{令 }n-1=k)\\
  &=\sum_{k=-\infty}^{+\infty}\frac1{(k-x)^2}=f(x)\quad(x\ne\text{整数}),
\end{aligned}
\]
所以，和函数 $f(x)$ 以 1 为周期，其连续性只需在 $(0,1)$ 内证明. 由于 $(0,1)$ 内，
\[
  \left|\frac1{(n-x)^2}\right|\le\frac1{(n-1)^2},\qquad
  \left|\frac1{(-n-x)^2}\right|=\frac1{(n+x)^2}\le\frac1{n^2},
\]
所以级数（1）在 $(0,1)$ 内一致收敛，$f(x)$ 在 $(0,1)$ 内连续（定理 1）. 证毕.

\noindent\textbf{例 5.2.41}\quad 设 $u_n(x)$ 在 $[a,b]$ 上连续（$n=1,2,\cdots$），
$\displaystyle\sum_{n=1}^{\infty}u_n(x)$ 在 $(a,b)$ 内一致收敛，求证
$\displaystyle\sum_{n=1}^{\infty}u_n(x)$ 在 $[a,b]$ 上一致连续.（北京师范大学）

\noindent\textbf{证}\quad 由于 $u_n(x)$ 在 $[a,b]$ 上连续，
$\displaystyle\lim_{x\to a^+}u_n(x)=u_n(a)$，
$\displaystyle\lim_{x\to b^-}u_n(x)=u_n(b)$，又因
$\displaystyle\sum_{n=1}^{\infty}u_n(x)$ 在 $(a,b)$ 内一致收敛，利用逐项取极限定理（或重复例 5.2.38 的证明），可知级数在
$x=a,\ x=b$ 处收敛，于是，$\displaystyle\sum_{n=1}^{\infty}u_n(x)$ 在
$[a,b]$ 上一致收敛. 根据和函数连续性定理（定理 1）知，
$\displaystyle\sum_{n=1}^{\infty}u_n(x)$ 在 $[a,b]$ 上连续；再由 Cantor 定理，
$\displaystyle\sum_{n=1}^{\infty}u_n(x)$ 在 $[a,b]$ 上一致连续.

\noindent\textbf{new $\star$练习}\quad 设函数项级数
$\displaystyle\sum_{n=1}^{\infty}u_n(x)$ 在开区间 $(a,b)$ 内满足：

1）对每个 $n\ge1$，$u_n(x)$ 一致连续；
2）$\displaystyle\sum_{n=1}^{\infty}u_n(x)$ 一致收敛于 $S(x)$.

试证：$S(x)$ 在 $(a,b)$ 内一致连续.（华中科技大学）

\noindent\textbf{证}\quad 由于 $u_n(x)$ 在 $(a,b)$ 上一致连续，（再根据例 2.2.6）知
$u_n(x)$ 在两端点有内侧极限，故能补充 $u_n(x)$ 在端点的定义，可使 $u_n(x)$ 在
$[a,b]$ 上连续，（根据 Cantor 定理）知 $u_n(x)$ 在 $[a,b]$ 上一致连续.（借助三角不等式）可得到
$\displaystyle\sum_{n=1}^{\infty}u_n(x)$ 在端点收敛. 因此，
$\displaystyle\sum_{n=1}^{\infty}u_n(x)$ 在 $[a,b]$ 上一致收敛. 根据一致收敛保持连续性定理，可知和函数在
$[a,b]$ 上连续.（利用 Cantor 定理）从而和函数在 $[a,b]$ 上一致连续，即
$S(x)$ 在 $(a,b)$ 内一致连续.

\noindent\textbf{注}\quad 本题表明：任意有限区间 $I$（不论是否为闭）一致收敛能保持一致连续性.

\noindent\textbf{$\star$例 5.2.42}\quad 设 $\{x_n\}$ 是 $(0,1)$ 内的一个序列：
$0<x_n<1$，且 $x_i\ne x_j$（$i\ne j$），试讨论函数
\[
  f(x)=\sum_{n=1}^{\infty}\frac{\operatorname{sgn}(x-x_n)}{2^n}
\]
在 $(0,1)$ 内的连续性，其中
\[
  \operatorname{sgn}x=
  \begin{cases}
    1,& x>0,\\
    0,& x=0,\\
    -1,& x<0.
  \end{cases}
\]
（北京大学）

\noindent\textbf{解}\quad $1^\circ$ 因
\[
  \left|\frac{\operatorname{sgn}(x-x_n)}{2^n}\right|\le\frac1{2^n},
  \quad \forall x\in(0,1),
\]
且 $\displaystyle\sum_{n=1}^{\infty}\dfrac1{2^n}$ 收敛，所以
$\displaystyle\sum_{n=1}^{\infty}\dfrac{\operatorname{sgn}(x-x_n)}{2^n}$ 在
$(0,1)$ 内一致收敛.

$2^\circ$ 设 $x_0\ne x_n$（$n=1,2,\cdots$）为 $(0,1)$ 内任意一点，则通项
$u_n(x)=\dfrac{\operatorname{sgn}(x-x_n)}{2^n}$ 在 $x=x_0$ 处连续，由
$1^\circ$ 并应用和函数连续定理 2，知 $f(x)$ 在 $x=x_0$ 处连续.

$3^\circ$ 设 $x_k$ 是 $\{x_n\}$ 中任意一点，因
\[
  f(x)=\sum_{n\ne k}\frac{\operatorname{sgn}(x-x_n)}{2^n}
  +\frac{\operatorname{sgn}(x-x_k)}{2^k},
\]
右边第一项在 $x=x_k$ 处连续，第二项在 $x=x_k$ 处间断，因此
$f(x)$ 在 $x=x_k$ 处间断.

\noindent\textbf{注}\quad $\{x_n\}$ 可以在 $(0,1)$ 内稠密，因此在证明 $x\ne x_n$ 处连续时，无法用定理 1，只能用定理 2.

\noindent\textbf{$\star$例 5.2.43}\quad 证明：
$f(x)=\displaystyle\sum_{n=1}^{\infty}\left(x+\dfrac1n\right)^n$ 在
$(-1,1)$ 内连续；

\noindent\textbf{证}\quad $\forall q:0<q<1$，考虑内闭区间 $[-q,q]\subset(-1,1)$. 因
\[
  \left|\left(x+\frac1n\right)^n\right|
  \le\left(|x|+\frac1n\right)^n
  \le\left(q+\frac1n\right)^n
  \quad(\forall x\in[-q,q]),
\]
且 $\displaystyle\sum_{n=1}^{\infty}\left(q+\dfrac1n\right)^n$ 收敛（因为
$\sqrt[n]{\left(q+\dfrac1n\right)^n}\to q<1$（$n\to\infty$）），所以
$\displaystyle\sum_{n=1}^{\infty}\left(x+\dfrac1n\right)^n$ 在 $[-q,q]$ 上一致收敛. 由
$q$ 的任意性，由定理 3 可知
$f(x)=\displaystyle\sum_{n=1}^{\infty}\left(x+\dfrac1n\right)^n$ 在 $(-1,1)$ 内连续.

\noindent\textbf{注}\quad 该级数在 $(-1,1)$ 内收敛，但非一致收敛. 因此本例不能直接使用定理 1. 和函数连续性定理还可用于推断非一致收敛. 如

\noindent\textbf{$\star$例 5.2.44}\quad 证明：
$\displaystyle\sum_{n=1}^{\infty}\dfrac{x}{(1+x^2)^n}$ 在 $(0,+\infty)$ 内非一致收敛.

\noindent\textbf{证}\quad 当 $x=0$ 时，级数和为 0；当 $x\ne0$ 时，级数是等比级数. 所以
\[
  S(x)=\sum_{n=1}^{\infty}\frac{x}{(1+x^2)^n}
  =
  \begin{cases}
    0,& \text{当 }x=0\text{ 时},\\
    \dfrac1x,& \text{当 }x\ne0\text{ 时}.
  \end{cases}
\]
$S(x)$ 在 $x=0$ 处间断，因此该级数在 $[0,+\infty)$ 上不一致收敛（定理 1），进而在
$(0,+\infty)$ 内也不一致收敛（因为：假若在 $(0,+\infty)$ 内一致收敛，加之级数在
$x=0$ 处收敛，便可推知级数在 $[0,+\infty)$ 上一致收敛，矛盾）.

\noindent\textbf{c.\quad 和函数的可微性与逐项求导}

\noindent\textbf{要点}\quad 若要证明 $\displaystyle\sum_{n=1}^{\infty}u_n(x)$ 在区间 $I$ 上可微，且可逐项求导，即在
$I$ 上，
\[
  \left(\sum_{n=1}^{\infty}u_n(x)\right)'
  =\sum_{n=1}^{\infty}u_n'(x),
\]
只要证明如下三条即可：

1）级数 $\displaystyle\sum_{n=1}^{\infty}u_n(x)$ 在 $I$ 上收敛（或者只要验证在
$I$ 上至少有一个收敛点）；

2）$u_n(x)$ 在 $I$ 上有连续导数（$n=1,2,\cdots$）；

3）$\displaystyle\sum_{n=1}^{\infty}u_n'(x)$ 在 $I$ 上一致收敛（或在 $I$ 的任一内闭区间上一致收敛）.

（对于函数序列，有类似叙述.）

\noindent\textbf{$\star$例 5.2.45}\quad 证明：
$f(x)=\displaystyle\sum_{n=1}^{\infty}ne^{-nx}$ 在 $(0,+\infty)$ 内收敛，但不一致收敛，而和函数在
$(0,+\infty)$ 内无穷次可微.

\noindent\textbf{证 I}\quad $1^\circ$ $\forall x\in(0,+\infty)$，因
$n^2\cdot ne^{-nx}\to0$（$n\to\infty$），所以
$\displaystyle\sum_{n=1}^{\infty}ne^{-nx}$ 收敛.

$2^\circ$ $\forall n\in\mathbb N$，因 $ne^{-nx}\to n\ne0$（$x\to0$），所以在
$(0,+\infty)$ 上，级数通项 $ne^{-nx}\nrightarrow0$（$n\to\infty$）. 从而
$\displaystyle\sum_{n=1}^{\infty}ne^{-nx}$ 在 $(0,+\infty)$ 内非一致收敛.

$3^\circ$ 因为级数在 $(0,+\infty)$ 内收敛，
$(ne^{-nx})'=-n^2e^{-nx}$ 连续，
$-\displaystyle\sum_{n=1}^{\infty}n^2e^{-nx}$ 在 $(0,+\infty)$ 内内闭一致收敛
（因 $\forall\varepsilon>0$，有
$0<n^2e^{-nx}\le n^2e^{-n\varepsilon}$（$x\in[\varepsilon,+\infty)$），而级数
$\displaystyle\sum_{n=1}^{\infty}n^2e^{-n\varepsilon}$ 收敛），故 $f(x)$ 可微，且
\[
  f'(x)=\left(\sum_{n=1}^{\infty}u_n(x)\right)'
  =\sum_{n=1}^{\infty}u_n'(x),\quad x\in(0,+\infty).
\]
一般而言，若已有
\[
  f^{(k)}(x)=(-1)^k\sum_{n=1}^{\infty}n^ke^{-nx},
\qquad(1)
\]
则通过类似方法可证 $f^{(k)}(x)$ 可导，且可逐项求导得到. 这就证明了 $f(x)$ 任意次可微.
式（1）对任意 $k\in\mathbb N$ 成立.

\noindent\textbf{证 II}\quad 级数在 $(0,+\infty)$ 内收敛，但不一致收敛，证法同证 I. 现证
$f(x)$ 无穷次可微.
\[
  f(x)=e^{-x}+2e^{-2x}+3e^{-3x}+\cdots+ne^{-nx}+(n+1)e^{-(n+1)x}+\cdots,
\]
\[
  e^{-x}f(x)=e^{-2x}+2e^{-3x}+\cdots+(n-1)e^{-nx}+ne^{-(n+1)x}+\cdots.
\]
所以
\[
  (1-e^{-x})f(x)=f(x)-e^{-x}f(x)
  =\sum_{n=1}^{\infty}e^{-nx}=\frac{e^{-x}}{1-e^{-x}}.
\]
故 $f(x)$ 在 $(0,+\infty)$ 内任意次可微.

\noindent\textbf{new $\star$练习}\quad 设
$f(x)=\displaystyle\sum_{n=1}^{\infty}ne^{-n}\cos nx,\ x\in[0,2\pi]$. 试证：在
$[0,2\pi]$ 上：

1）$f(x)$ 连续；2）$f(x)$ 有连续导数；3）
$\displaystyle\max_{0\le x\le2\pi}|f(x)|=\dfrac{e}{(1-e)^2}$.（华东师范大学）

\noindent\textbf{提示}\quad 1）（当 $x>1$ 时）$(xe^{-x})'=(1-x)e^{-x}<0$，故
$xe^{-x}\searrow$，能用 Cauchy 积分判别法. 由 $\displaystyle\int_1^{+\infty}xe^{-x}\,dx$
收敛，知 $\displaystyle\sum_{n=1}^{\infty}ne^{-n}$ 收敛. 又因
$ne^{-n}|\cos nx|\le ne^{-n}$，由 $M$ 判别法知，
$\displaystyle\sum_{n=1}^{\infty}ne^{-n}\cos nx$ 一致收敛.

2）类似地，可验证满足逐项求导的条件，知 $f(x)$ 有连续导数.

3）由
\[
  |f(x)|\le \sum_{n=1}^{\infty}ne^{-n}|\cos nx|
  \le \sum_{n=1}^{\infty}ne^{-n}|\cos(n\cdot0)|=f(0)\quad(\forall x\in[0,2\pi]),
\]
知
\[
  \max_{0\le x\le2\pi}|f(x)|=f(0)
  =\sum_{n=1}^{\infty}ne^{-n}
  =x\left(\sum_{n=1}^{\infty}x^n\right)'\bigg|_{x=e^{-1}}
  =x\left(\frac{x}{1-x}\right)'\bigg|_{x=e^{-1}}
  =\frac{e}{(1-e)^2}.
\]

\noindent\textbf{例 5.2.46}\quad 证明 Riemann $\zeta$ 函数
$\zeta(x)=\displaystyle\sum_{n=1}^{\infty}\dfrac1{n^x}$ 在
$(1,+\infty)$ 内连续，并有各阶连续导数.（北京大学）

\noindent\textbf{提示}\quad 可在 $(1,+\infty)$ 的内闭区间上应用数学归纳法进行证明.

\noindent\textbf{new 练习}\quad 试证：函数
$\zeta(x)=\displaystyle\sum_{n=1}^{\infty}\dfrac1{n^x}$ 在
$(1,+\infty)$ 内连续，但非一致连续.（浙江大学）

\noindent\textbf{提示}\quad “若在区间 $(1,+\infty)$ 上能找到收敛序列 $\{x_n\}$，使得
$\zeta(x_n)\to+\infty$，则 $\zeta(x)$ 在 $(1,+\infty)$ 上非一致连续”. 因为：一方面，当
$n,m$ 充分大时，$|x_n-x_m|$ 可以任意小；但另一方面，利用
$\zeta(x_n)\to+\infty$，$\forall m>n>0$，固定 $n$，将 $m$ 取得足够大时，可使得
$|\zeta(x_n)-\zeta(x_m)|$ 任意变大，与一致连续相矛盾.

\noindent\textbf{证}\quad 令 $s_n=\displaystyle\sum_{k=1}^{n}\dfrac1k$，则
$\displaystyle\lim_{n\to\infty}s_n=+\infty$（例 1.2.10）.

因
\[
  \lim_{x\to1^+}\sum_{k=1}^{n}\frac1{k^x}
  =\sum_{k=1}^{n}\frac1k=s_n,
\]
故可取 $x_n\in\left(1,1+\dfrac1n\right)$ 使得
$s_n-1<\displaystyle\sum_{k=1}^{n}\dfrac1{k^{x_n}}<s_n+1$. 如此，一方面，
$\{x_n\}$ 收敛；另一方面，
\[
  \zeta(x_n)=\sum_{k=1}^{\infty}\frac1{k^{x_n}}
  >\sum_{k=1}^{n}\frac1{k^{x_n}}>s_n-1\to+\infty\quad(n\to\infty).
\]
如提示所述，这表明函数 $\zeta(x)$ 在 $(1,+\infty)$ 内非一致连续.

\noindent\textbf{$\star$例 5.2.47}\quad 设 $f(x)$ 在 $(-\infty,+\infty)$ 内有任意阶导数，级数
\[
\cdots+f^{(n)}(x)+\cdots+f''(x)+f'(x)+f(x)
+\int_0^x f(t_1)\,dt_1
+\int_0^xdt_2\int_0^{t_2}f(t_1)\,dt_1+\cdots
\]
\[
  +\int_0^xdt_n\int_0^{t_n}dt_{n-1}\cdots\int_0^{t_2}f(t_1)\,dt_1+\cdots
\]
按两个方向在 $(-\infty,+\infty)$ 内一致收敛. 试求级数的和函数 $F(x)$.（同济大学）

\noindent\textbf{解}\quad $f(x)$ 有各阶导数，自然各阶导数都连续；该级数逐项求导之后的级数仍是它自己，因而一致收敛，满足逐项求导三条件，所以
$\dfrac{dF}{dx}=F(x)$，得 $\dfrac{dF(x)}{F(x)}=dx$. 两边同时积分得
$\ln F(x)=x+C$，即
\[
  F(x)=C_1e^x\quad(\text{其中 }C_1=e^C\text{ 为常数}),
\]
令 $x=0$，知 $C_1=f(0)+f'(0)+\cdots+f^{(n)}(0)+\cdots$.

\noindent\textbf{注}\quad 逐项求导定理中的条件，只是充分的，有时条件不满足，还可利用导数定义证明和函数的可微性.

\noindent\textbf{*例 5.2.48}\quad 设 $\{a_n\}$ 为区间 $[0,1]$ 上全体有理数组成的数列，
$u_n(x)$ 在 $[0,1]$ 上一致地满足 Lipschitz 条件，即：
$\exists L>0$，使得 $\forall x_1,x_2\in[0,1]$，有
\[
  |u_n(x_1)-u_n(x_2)|\le L|x_1-x_2|\quad(\forall n\in\mathbb N).
\qquad(1)
\]
又设 $u_n(0)=0$，$u_n(x)$ 只在 $a_n$ 处无导数（$n=1,2,\cdots$），在
$[0,1]$ 其他地方有导数. 试证：
$f(x)=\displaystyle\sum_{n=1}^{\infty}\dfrac{u_n(x)}{2^n}$ 在
$[0,1]$ 上连续；在 $[0,1]$ 的有理点上不可导，无理点上可导，且
\[
  f'(x)=\sum_{n=1}^{\infty}\frac{u_n'(x)}{2^n}.
\]

\noindent\textbf{注意}\quad 由于有理点的稠密性，本题不能使用逐项求导的定理.

\noindent\textbf{证}\quad $1^\circ$ 由式（1）可知 $u_n(x)$ 在 $[0,1]$ 上连续且等度连续，现证它们一致有界. 事实上，$\forall x\in[0,1]$，
\[
  |u_n(x)|\le |u_n(x)-u_n(0)|+|u_n(0)|\le L|x-0|+0\le L
  \quad(\forall n\in\mathbb N).
\]
$2^\circ$ 由于
\[
  \left|\frac{u_n(x)}{2^n}\right|\le\frac L{2^n},
\]
且 $\displaystyle\sum_{n=1}^{\infty}\dfrac L{2^n}$ 收敛，所以级数
$\displaystyle\sum_{n=1}^{\infty}\dfrac{u_n(x)}{2^n}$ 在 $[0,1]$ 上一致收敛，从而和函数
$f(x)$ 在 $[0,1]$ 上连续.

$3^\circ$ 设 $x_0\in(0,1)$ 为任一无理数，来求 $f'(x_0)$：
\[
  \frac{f(x_0+h)-f(x_0)}{h}
  =\sum_{n=1}^{\infty}\frac{u_n(x_0+h)-u_n(x_0)}{h2^n},
\]
右端是关于自变量 $h$ 的函数项级数，取 $\delta>0$ 充分小，使得
$(x_0-\delta,x_0+\delta)\subset(0,1)$. 由已知条件式（1），有
\[
  \left|\frac{u_n(x_0+h)-u_n(x_0)}{h\cdot2^n}\right|
  \le\frac L{2^n},\quad0<|h|<\delta.
\]
而 $\displaystyle\sum_{n=1}^{\infty}\dfrac L{2^n}$ 收敛，所以
$\displaystyle\sum_{n=1}^{\infty}\dfrac{u_n(x_0+h)-u_n(x_0)}{h\cdot2^n}$
关于 $h$ 在 $(x_0-\delta,x_0+\delta)$ 内一致收敛. 于是可利用逐项取极限的定理，
\[
  f'(x_0)=\lim_{h\to0}\frac{f(x_0+h)-f(x_0)}h
  =\sum_{n=1}^{\infty}\lim_{h\to0}\frac{u_n(x_0+h)-u_n(x_0)}h\cdot\frac1{2^n}
  =\sum_{n=1}^{\infty}\frac1{2^n}u_n'(x_0).
\]

$4^\circ$ 设 $x_0=a_k$（有理点），这时
\[
  f(x)=\sum_{n=1}^{\infty}\frac{u_n(x)}{2^n}
  =\sum_{n\ne k}\frac{u_n(x)}{2^n}+\frac{u_k(x)}{2^k}.
\]
用 $3^\circ$ 中的方法，可知右边第一项在 $x_0=a_k$ 处可导，但根据已知条件，第二项在
$a_k$ 处不可导，故 $f(x)$ 在 $a_k$ 处不可导（$k=1,2,\cdots$）. 即
$f(x)$ 在 $[0,1]$ 的有理点上无导数.

\noindent\textbf{※例 5.2.49}\quad 试构造一个函数，使之在 $(-\infty,+\infty)$ 内处处连续，但处处不可微.

\noindent\textbf{方法}\quad 用一串锯齿波（如图 5.2.2 中 $u_n(x)$）进行叠加，这些波的振幅按
$n$ 以等比数列的方式无限变小，使级数
$f(x)=\displaystyle\sum_{n=1}^{\infty}u_n(x)$ 一致收敛，从而 $f(x)$ 连续. 另一方面，让这些波的周期随
$n$ 而无限变小（即振动得越来越快，无限变快），使得和波 $f(x)$ 的图像无限“粗糙”，$f(x)$ 处处不可导.

\begin{center}
  \includegraphics[width=0.62\textwidth]{5-2-2.png}
\end{center}

\noindent\textbf{解}\quad 在每个区间 $[m,m+1]$ 上（$m=0,\pm1,\pm2,\cdots$）令
\[
  u_0(x)=
  \begin{cases}
    x-m,& x\in\left[m,m+\dfrac12\right),\\
    m+1-x,& x\in\left[m+\dfrac12,m+1\right],
  \end{cases}
\]
\[
  u_1(x)=\frac14u_0(4x),\quad\cdots,\quad
  u_n(x)=\frac1{4^n}u_0(4^nx),\quad\cdots.
\]
下面我们来证明
\[
  f(x)=\sum_{n=1}^{\infty}u_n(x)
\qquad(1)
\]
在 $(-\infty,+\infty)$ 内连续，但处处不可导. 事实上：

$1^\circ$ 因为 $|u_0(x)|\le1$，
\[
  |u_n(x)|=\left|\frac1{4^n}u_0(4^nx)\right|\le\frac1{4^n},
  \quad\forall x\in(-\infty,+\infty),
\]
$\displaystyle\sum_{n=1}^{\infty}\dfrac1{4^n}$ 收敛，所以
$f(x)=\displaystyle\sum_{n=1}^{\infty}u_n(x)$ 在
$(-\infty,+\infty)$ 内一致收敛，$f(x)$ 处处连续.

$2^\circ$ $\forall x_0\in(-\infty,+\infty)$，我们取
\[
  x_n=x_0\pm\frac1{4^n}.
\qquad(2)
\]
若我们证明了极限
$\displaystyle\lim_{n\to\infty}\dfrac{f(x_n)-f(x_0)}{x_n-x_0}$ 不存在，则
$f(x)$ 在 $x_0$ 处不可导. 从而由 $x_0$ 的任意性知 $f(x)$ 处处不可导. 由式（1），
\[
  \frac{f(x_n)-f(x_0)}{x_n-x_0}
  =\frac1{x_n-x_0}\left(\sum_{k=1}^{\infty}u_k(x_n)-\sum_{k=1}^{\infty}u_k(x_0)\right)
  =\sum_{k=1}^{\infty}\frac{u_k(x_n)-u_k(x_0)}{x_n-x_0}.
\]
注意 $|x_n-x_0|=\dfrac1{4^n}$. 对于 $k\ge n$ 的 $u_k$ 而言，$\dfrac1{4^n}$ 是
$u_k$ 的周期 $\dfrac1{4^k}$ 的整数倍，因此
\[
  u_k(x_n)=u_k(x_0)\quad(\text{当 }k\ge n\text{ 时}).
\]
故
\[
  \frac{f(x_n)-f(x_0)}{x_n-x_0}
  =\sum_{k=0}^{n-1}\frac{u_k(x_n)-u_k(x_0)}{x_n-x_0}.
\qquad(3)
\]
至此式（2）中的 $\pm$ 号尚未选定，为使式（3）无极限，我们规定（2）中的符号如此选取：使
$x_n$ 与 $x_0$ 位于锯齿波 $u_{n-1}(x)$ 呈直线的同一半波区间上. 于是
$\dfrac{u_{n-1}(x_n)-u_{n-1}(x_0)}{x_n-x_0}$ 等于 1 或 $-1$. 因
$u_{n-2}(x)$ 的一个周期包含 $u_{n-1}(x)$ 的四个周期，故 $x_n$ 与 $x_0$ 位于
$u_{n-2}(x)$ 的同一直线段区间上，
$\dfrac{u_{n-2}(x_n)-u_{n-2}(x_0)}{x_n-x_0}=1$ 或 $-1$. 同理，式（3）中每一项皆为 1 或
$-1$. 如此选定之后，式（3）是这种级数 $\displaystyle\sum_{n=1}^{\infty}a_n$ 的部分和；其中
$a_n=1$ 或 $-1$. 由于 $a_n\nrightarrow0$（当 $n\to\infty$ 时），故
$\displaystyle\sum_{n=1}^{\infty}a_n$ 发散，从而式（3）当 $n\to\infty$ 时无极限. 证毕.

\noindent\textbf{d.\quad 逐项积分与积分号下取极限}

\noindent\textbf{要点}\quad 1）若 $u_n(x)$ 在 $[a,b]$ 上可积（$n=1,2,\cdots$），级数
$\displaystyle\sum_{k=1}^{\infty}u_k(x)$ 在 $[a,b]$ 上一致收敛，则可逐项积分，即
\[
  \int_a^b\sum_{k=1}^{\infty}u_k(x)\,dx
  =\sum_{k=1}^{\infty}\int_a^b u_k(x)\,dx.
\tag{A}
\]

若未知该级数在 $[a,b]$ 上一致收敛，能否进行逐项积分，取决于是否有
\[
  \lim_{n\to\infty}\int_a^b R_n(x)\,dx=0,
\tag{B}
\]
其中 $R_n(x)$ 是 $\displaystyle\sum_{k=1}^{\infty}u_k(x)$ 前 $n$ 项和的余项.（因为
\[
\begin{aligned}
\text{式（A）}\Longleftrightarrow{}&
\lim_{n\to\infty}\left|\int_a^b\sum_{k=1}^{\infty}u_k(x)\,dx
-\sum_{k=1}^n\int_a^b u_k(x)\,dx\right|=0\\
\Longleftrightarrow{}&
\lim_{n\to\infty}\left|\int_a^b\left(\sum_{k=1}^{\infty}u_k(x)-\sum_{k=1}^n u_k(x)\right)\,dx\right|=0\\
\Longleftrightarrow{}&
\lim_{n\to\infty}\int_a^b R_n(x)\,dx=0\quad(\text{式（B）}.)
\end{aligned}
\]

2）函数项级数 $\displaystyle\sum_{k=1}^{\infty}u_k(x)$ 在区间 $I$ 上有意义，并且在 $I$ 上一致收敛，若极限
$\displaystyle\lim_{x\to x_0}u_n(x)$ 都存在，则可逐项取极限，即
\[
  \lim_{x\to x_0}\sum_{k=1}^{\infty}u_k(x)
  =\sum_{k=1}^{\infty}\lim_{x\to x_0}u_k(x).
\tag{C}
\]

\noindent\textbf{※例 5.2.50}\quad 设 $h(x),f_n'(x)$ 在 $[a,b]$ 上连续，$n=1,2,\cdots$，又对 $[a,b]$ 中任意的
$x_1,x_2$ 和正整数 $n$ 有
$|f_n(x_1)-f_n(x_2)|\le\dfrac Mn|x_1-x_2|$，其中 $M>0$ 为常数，求证
\[
  \lim_{n\to\infty}\int_a^b h(x)f_n'(x)\,dx=0.
\]
（南京大学）

\noindent\textbf{分析}\quad 要证明 $\displaystyle\lim_{n\to\infty}\int_a^b h(x)f_n'(x)\,dx=0$，关键问题在于证明
$f_n'(x)\rightrightarrows0$（当 $n\to\infty$ 时，关于 $x\in[a,b]$）. 因为 $h(x)$ 在 $[a,b]$ 上连续，所以
$h(x)$ 在 $[a,b]$ 上有界. 若 $f_n'(x)\rightrightarrows0$，便有
$h(x)f_n'(x)\rightrightarrows0$，从而可在积分号下取极限，得出欲求的结果. 下面只证
$f_n'(x)\rightrightarrows0$.

\noindent\textbf{证}\quad 因 $f_n'(x)$ 在 $[a,b]$ 上连续，所以一致连续，$\forall\,\dfrac1n>0$，$\exists\,\delta>0$，当
$|x_1-x_2|<\delta$ 时，有 $|f_n'(x_1)-f_n'(x_2)|<\dfrac1n$. 取 $m$ 充分大，使得
$\dfrac{b-a}{m}<\delta$，将 $[a,b]$ $m$ 等分：
$a=x_0<x_1<\cdots<x_m=b$，利用已知条件
\[
  |f_n(x_i)-f_n(x_{i-1})|\le\frac Mn|x_i-x_{i-1}|.
\]
由微分中值定理，$\exists\,\xi_i\in(x_{i-1},x_i)$，使得
$|f_n'(\xi_i)|\le\dfrac Mn$. 于是 $\forall x\in[a,b]$，$x$ 必属于某个小区间
$[x_{i-1},x_i]$，所以
\[
  |f_n'(x)|\le |f_n'(x)-f_n'(\xi_i)|+|f_n'(\xi_i)|
  <\frac1n+\frac Mn=\frac{1+M}{n},
\]
故 $f_n'(x)\rightrightarrows0\ (n\to\infty)$.

\noindent\textbf{例 5.2.51}\quad 设 $g(x)$ 及 $f_n(x)\ge0\ (n=1,2,\cdots)$ 在 $[a,b]$ 上有界可积，且
$\forall c\in(a,b)$，当 $n\to\infty$ 时 $f_n(x)\to0$ 于 $[c,b]$ 上；
$\displaystyle\lim_{n\to\infty}\int_a^b f_n(x)\,dx=1$，
$\displaystyle\lim_{x\to a+}g(x)=A$. 试证
\[
  \lim_{n\to\infty}\int_a^b g(x)f_n(x)\,dx=A.
\]

\noindent\textbf{分析}\quad 已知 $\displaystyle\lim_{n\to\infty}\int_a^b f_n(x)\,dx=1$，得
$\displaystyle\lim_{n\to\infty}\int_a^b Af_n(x)\,dx=A$. 要证
$\displaystyle\lim_{n\to\infty}\int_a^b g(x)f_n(x)\,dx=A$，只需证明
\[
  \lim_{n\to\infty}\int_a^b g(x)f_n(x)\,dx
  =\lim_{n\to\infty}\int_a^b Af_n(x)\,dx
\]
或
\[
  \lim_{n\to\infty}\int_a^b (A-g(x))f_n(x)\,dx=0.
\]
将积分拆成两项：
\[
  \int_a^b(A-g(x))f_n(x)\,dx
  =\int_a^{a+\delta}(A-g(x))f_n(x)\,dx
  +\int_{a+\delta}^b(A-g(x))f_n(x)\,dx.
\]
因 $\displaystyle\lim_{x\to a+}g(x)=A$，可见 $\delta>0$ 取得充分小时第一项能任意小，再将 $\delta$ 固定，因
$f_n(x)\to0$（当 $n\to\infty$ 时）于 $[a+\delta,b]$ 上，所以 $n$ 充分大时第二项能任意小.

\noindent\textbf{证}\quad 由 $\displaystyle\lim_{x\to a+}g(x)=A$ 知：$\forall\varepsilon>0$，$\exists\delta>0$（$\delta<b-a$），使得当
$a<x<a+\delta$ 时，$|A-g(x)|<\dfrac\varepsilon4$. 因此
\[
\begin{aligned}
\left|\int_a^b(A-g(x))f_n(x)\,dx\right|
\le{}&\int_a^{a+\delta}|A-g(x)|f_n(x)\,dx
+\int_{a+\delta}^b(|A|+|g(x)|)f_n(x)\,dx\\
\le{}&\frac\varepsilon4\int_a^{a+\delta}f_n(x)\,dx
+(|A|+M)\int_{a+\delta}^b f_n(x)\,dx
\end{aligned}
\]
（这里 $M$ 表示 $g(x)$ 的界，即 $|g(x)|\le M$ 于 $[a,b]$ 上）. 因
$\displaystyle\lim_{n\to\infty}\int_a^b f_n(x)\,dx=1$，所以
$\exists N_1>0$，$n>N_1$ 时，
$0\le\displaystyle\int_a^{a+\delta}f_n(x)\,dx\le\int_a^b f_n(x)\,dx<2$.

又因在 $[a+\delta,b]$ 上，$f_n(x)\rightrightarrows0$（当 $n\to\infty$ 时），所以
$\exists N_2>0$，当 $n>N_2$ 时，
$0\le f_n(x)<\dfrac{\varepsilon}{2(|A|+M)(b-a-\delta)}$. 取
$N=\max\{N_1,N_2\}$，则 $n>N$ 时，有
\[
  \left|\int_a^b(A-g(x))f_n(x)\,dx\right|
  <\frac\varepsilon2+\frac\varepsilon2=\varepsilon.
\]
证毕.

\noindent\textbf{※练习 1}\quad 求证等式
\[
  \int_0^1\frac1{x^x}\,dx=\sum_{n=1}^{\infty}\frac1{n^n}.
\]
（中国科学院）

\noindent\textbf{提示}\quad 1）
\[
  \int_0^1\frac1{x^x}\,dx
  =\int_0^1 e^{-x\ln x}\,dx
  =\int_0^1\sum_{n=0}^{\infty}\frac{(-x\ln x)^n}{n!}\,dx
  =\sum_{n=0}^{\infty}\int_0^1\frac{(-x\ln x)^n}{n!}\,dx.
\qquad(1)
\]
2）
\[
  \sum_{n=0}^{\infty}\int_0^1\frac{(-x\ln x)^n}{n!}\,dx
  =\sum_{n=0}^{\infty}\frac{(-1)^n}{n!}\int_0^1x^n\ln^n x\,dx
  =\sum_{n=1}^{\infty}\frac1{n^n}.
\qquad(2)
\]

\noindent\textbf{证}\quad $1^\circ$（证明可逐项求积分.）因为在 $(0,1)$ 上，$f(x)=x\ln x<0$，$f(1)=0$，应用 L'Hospital 法则，易知
$f(0+)=0$. 因此只要令 $f(0)=0$，则 $f(x)$ 在 $[0,1]$ 上连续，从而有界（记为 $M$）.（或由
$f'(x)=\ln x+1=0$ 得 $(0,1)$ 内唯一可疑点 $x=e^{-1}$，而区间端点上 $f(0)=f(1)=0$. 所以
$|f(x)|\le |f(e^{-1})|\overset{\text{记}}=M$，于 $[0,1]$ 上.）

由此知：
\[
  \left|\frac{(-x\ln x)^n}{n!}\right|
  =\frac1{n!}(-f(x))^n\le\frac1{n!}M^n.
\]
而级数 $\displaystyle\sum_{n=1}^{\infty}\frac1{n!}M^n$ 收敛，故
$\displaystyle\int_0^1\sum_{n=0}^{\infty}\frac{(-x\ln x)^n}{n!}\,dx$
在 $[0,1]$ 上一致收敛，可逐项求积分. 等式（1）成立.

$2^\circ$（利用分部积分.）
\[
\begin{aligned}
\int_0^1(x\ln x)^n\,dx
={}&\int_0^1x^n\ln^n x\,dx
=\frac1{n+1}\int_0^1\ln^n x\,d(x^{n+1})\\
={}&\left.\frac{x^{n+1}\ln^n x}{n+1}\right|_0^1
-\frac n{n+1}\int_0^1x^n\ln^{n-1}x\,dx\\
={}&\frac{(-1)\cdot n}{n+1}\int_0^1x^n\ln^{n-1}x\,dx\\
={}&\cdots=\frac{(-1)^n n!}{(n+1)^n}\int_0^1x^n\,dx
=\frac{(-1)^n n!}{(n+1)^{n+1}},
\end{aligned}
\]
说明式（2）成立. 于是
\[
  \int_0^1\frac1{x^x}\,dx
  =\sum_{n=0}^{\infty}\int_0^1\frac1{n!}(-x\ln x)^n\,dx
  =\sum_{n=0}^{\infty}\frac{(-1)^n}{n!}\cdot\frac{(-1)^n n!}{(n+1)^{n+1}}
  =\sum_{n=1}^{\infty}\frac1{n^n}.
\]

\noindent\textbf{※练习 2}\quad 设 $\{f_n(x)\}$ 是定义在 $[-1,1]$ 上的连续函数序列，在原点附近一致有界，若
$\forall\delta:1>\delta>0$，$f_n(x)$ 在 $[-1,-\delta]\cup[\delta,1]$ 上一致收敛于零，且函数
$g(x)$ 在 $[-1,1]$ 上连续，$g(0)=0$，求证：
\[
  \lim_{n\to\infty}\int_{-1}^1 f_n(x)g(x)\,dx=g(0).
\qquad(1)
\]

\noindent\textbf{提示}\quad 1）在原点附近 $g$ 能任意小，$|f_n|$ 一致有界，故 $f_ng$ 在原点邻域内的积分可任意小.

2）其他地方 $g$ 保持有界而 $|f_n|$ 一致趋向零，故其他地方 $f_ng$ 的积分也可任意小.

\noindent\textbf{证}\quad $1^\circ$ $\{f_n(x)\}$ 在原点附近一致有界，意思即：$\exists M>0$ 及 $\delta_1>0$ 使得
\[
  |f_n(x)|\le M\quad(\forall x\in(-\delta_1,\delta_1),\ \forall n\in\mathbb N).
\qquad(2)
\]
又 $g(x)$ 连续，且 $g(0)=0$，故 $\forall\varepsilon>0(\varepsilon<1)$，$\exists\delta_2>0$（$\delta_2<\dfrac12$）使得当
$|x|<\delta_2$ 时有
\[
  |g(x)|=|g(x)-g(0)|<\frac{\varepsilon}{4M\delta_2}.
\qquad(3)
\]
取 $\delta=\min\{\delta_1,\delta_2\}$，则
\[
\left|\int_{-1}^1 f_n(x)g(x)\,dx\right|
\le\left|\int_{-\delta}^{\delta}f_n(x)g(x)\,dx\right|
+\left|\left(\int_{-1}^{-\delta}+\int_{\delta}^1\right)f_n(x)g(x)\,dx\right|
\overset{\text{记}}=I+J.
\qquad(4)
\]
利用式（2）和（3）：
\[
  I=\left|\int_{-\delta}^{\delta}f_n(x)g(x)\,dx\right|
  \le\int_{-\delta}^{\delta}|f_n(x)||g(x)|\,dx
  \le M\cdot\frac{\varepsilon}{4M\delta}\cdot2\delta
  =\frac\varepsilon2.
\qquad(5)
\]
$2^\circ$ 函数 $g(x)$ 在 $[-1,1]$ 连续，故有界，即存在 $M_1>0$，使得 $|g(x)|<M_1$. 于是
\[
  J\le\left(\int_{-1}^{-\delta}+\int_{\delta}^1\right)|f_n(x)||g(x)|\,dx
  \le M_1\left(\int_{-1}^{-\delta}+\int_{\delta}^1\right)|f_n(x)|\,dx.
\qquad(6)
\]
根据题设知在 $[-1,-\delta]\cup[\delta,1]$ 上，$\forall\varepsilon>0(\varepsilon<1)$，$\exists N>0$ 使得
$n>N$ 时有 $|f_n(x)|<\dfrac{\varepsilon}{4M_1(1-\delta)}$. 故式（6）可写成
\[
  |J|\le\frac{\varepsilon}{4M_1(1-\delta)}\cdot M_1\left(\int_{-1}^{-\delta}+\int_{\delta}^1\right)dx
  =\frac\varepsilon2.
\qquad(7)
\]
$3^\circ$ 最后得
\[
  \left|\int_{-1}^1 f_n(x)g(x)\,dx\right|
  \overset{\text{式（4）}}{\le} I+J
  \overset{\text{式（5）,（7）}}{\le}\frac\varepsilon2+\frac\varepsilon2=\varepsilon.
\]
此式表明式（1）成立，命题获证.

\noindent\textbf{※练习 3}\quad 设 $\{|f_n(x)|\}$ 是定义在 $[-1,1]$ 上的连续函数序列，若

i）$\displaystyle\lim_{n\to\infty}\int_{-1}^1 f_n(x)\,dx=1$；ii）$\forall\delta:1>\delta>0$，$f_n(x)$ 在
$[-1,-\delta]\cup[\delta,1]$ 上一致收敛于零.

求证：对 $[-1,1]$ 上任意连续函数 $g(x)$，
\[
  \lim_{n\to\infty}\int_{-1}^1 f_n(x)g(x)\,dx=g(0).
\qquad(1)
\]
（华东师范大学）

\noindent\textbf{提示}\quad 1）利用拟合法. 由 $\displaystyle\lim_{n\to\infty}\int_{-1}^1 f_n(x)\,dx=1$ 出发，可将问题转化为：求证
\[
  \lim_{n\to\infty}\int_{-1}^1 f_n(x)(g(x)-g(0))\,dx=0.
\qquad(2)
\]

2）
\[
\begin{aligned}
\left|\int_{-1}^1 f_n(x)(g(x)-g(0))\,dx\right|
\le{}&\left|\int_{-\delta}^{\delta}f_n(x)(g(x)-g(0))\,dx\right|\\
&+\left|\left(\int_{-1}^{-\delta}+\int_{\delta}^1\right)f_n(x)(g(x)-g(0))\,dx\right|
\overset{\text{记}}=I+J.
\end{aligned}
\qquad(3)
\]

3）在 $[-1,-\delta]\cup[\delta,1]$ 上，$g(x)-g(0)$ 有界，$f_n(x)$ 一致收敛于零，因此第二项
$J\to0\ (n\to\infty)$.

\noindent\textbf{分析}\quad 证明式（3）第一项 $I$ 的极限为 0 是本题的关键. 上题（练习 2）在这里靠的是
$f_n$ 一致有界，另一因子趋向零；例 5.2.51 在这里靠的是 $f_n\ge0$. 如果这里也有条件 $f_n\ge0$，那么
$\forall\varepsilon>0$，让 $|g(x)-g(0)|\le\varepsilon$，则
\[
  I=\left|\int_{-\delta}^{\delta}f_n(x)(g(x)-g(0))\,dx\right|
  \le\int_{-\delta}^{\delta}f_n(x)|g(x)-g(0)|\,dx
  \le\varepsilon\int_{-\delta}^{\delta}f_n(x)\,dx
  \le\varepsilon\int_{-1}^1 f_n(x)\,dx\le\varepsilon,
\]
或者利用第一积分中值定理，$\exists\,\xi\in(-\delta,\delta)$，使得
\[
  I=\left|\int_{-\delta}^{\delta}f_n(x)(g(x)-g(0))\,dx\right|
  \le |g(\xi)-g(0)|\left|\int_{-\delta}^{\delta}f_n(x)\,dx\right|
  \le\varepsilon\int_{-\delta}^{\delta}f_n(x)\,dx
  \le\varepsilon\int_{-1}^1 f_n(x)\,dx\le\varepsilon,
\]
则问题也可以解决. 这里撤销了传统条件“$f_n\ge0$”，是本题的特色和创新. 由于取消 $f_n\ge0$ 后 $f_n$ 可以无穷多次变号，我们不可随便使用第一积分中值定理.

\noindent\textbf{证}\quad 下面我们将证明虽然在 $[0,1]$ 上 $f_n$ 可能变号，但当 $\delta>0$ 且充分小时，$f_n(x)$ 在
$[-\delta,\delta]$ 上会保持恒正，因而可在 $[-\delta,\delta]$ 上应用第一积分中值定理. 剩下问题就可按练习 2 的方法同样处理.

利用条件 i）和条件 ii），在 $[-1,-\delta]\cup[\delta,1]$ 上，$\forall\varepsilon>0(\varepsilon\le\dfrac12)$，$\exists N>0$，当
$n>N$ 时有
\[
  |f_n(x)|<\frac{\varepsilon}{2(1-\delta)}.
\]
于是
\[
  \int_{-\delta}^{\delta}f_n(x)\,dx
  =\int_{-1}^1f_n(x)\,dx
  -\left|\left(\int_{-1}^{-\delta}+\int_{\delta}^1\right)f_n(x)\,dx\right|
  \ge1-\varepsilon\ge\frac12.
\]
（设 $n>N$）在 $(-\delta,\delta)$ 上，当 $x\to0$ 时，若 $f_n(x)$ 由正变为负，则每次变负后，迟早会返回正，且值
$\ge\dfrac12$（不然一直为负时 $\displaystyle\int_{-\delta}^{\delta}f_n(x)\,dx$ 不可能 $\ge\dfrac12$（$\delta>0$ 充分小））. 进而
$f_n(x)$ 在 $(-\delta,\delta)$ 上不可能无穷次变号，否则 $f_n(x)$ 的值就会上、下穿过 $[0,1/2]$ 无穷多遍，于是，每个
$y\in[0,1/2]$ 都是函数图像 $y=f_n(x)$ 的极限点（极限点不唯一），极限 $\displaystyle\lim_{x\to0}f_n(x)$ 不存在，与
$|f_n(x)|$ 的连续性矛盾. 因此，动点在 $(-\delta,\delta)$ 上趋向原点时最多只有有限次变负，最后一次变负返回后，将一直为正. 故让
$(-\delta,\delta)$ 充分缩小，就能保证 $f_n(x)$ 在更小的 $(-\delta,\delta)$ 上恒为正.

通过上面的分析，我们知道当 $x\to0$ 时，虽然 $f_n(x)$ 开始可以变号，但当 $\delta>0$ 充分小时，就会在
$[-\delta,\delta]$ 上保持恒正，就可以应用第一积分中值定理. 剩下的工作与上题（练习 2）一样.

\noindent\textbf{※例 5.2.52}\quad 试证级数
$\displaystyle\sum_{n=1}^{\infty}x^{2n}\ln x$ 在 $(0,1)$ 内不一致收敛，但在 $[0,1]$ 上可逐项积分.

\noindent\textbf{证}\quad $1^\circ$ 当 $x=1$ 时，级数通项
$u_n(1)=x^{2n}\ln x\big|_{x=1}=0$；当 $0<x<1$ 时，
$\displaystyle\sum_{n=1}^{\infty}x^{2n}\ln x$ 为等比级数，所以和
\[
  S(x)=
  \begin{cases}
    \dfrac{x^2}{1-x^2}\ln x,&0<x<1,\\
    0,&x=1.
  \end{cases}
\]
可见
\[
  S(1-)=\lim_{x\to1-}\frac{x^2\ln[1-(1-x)]}{(1+x)(1-x)}
  =\frac12\ne S(1),
\]
故该级数非一致收敛（根据和函数连续定理）.

$2^\circ$（证明能逐项积分.）因
\[
  R_n(x)=\sum_{k=n+1}^{\infty}x^{2k}\ln x
  =\frac{x^{2n+2}}{1-x^2}\ln x
  =\frac{x^2\ln x}{1-x^2}\cdot x^{2n},
\]
其中 $\displaystyle\lim_{x\to0+}\frac{x^2\ln x}{1-x^2}$ 及
$\displaystyle\lim_{x\to1-}\frac{x^2\ln x}{1-x^2}$ 都有有限极限，且
$\dfrac{x^2\ln x}{1-x^2}$ 在 $(0,1)$ 内连续，所以
$\dfrac{x^2\ln x}{1-x^2}$ 在 $(0,1)$ 内有界，即 $\exists M>0$，使得
$\left|\dfrac{x^2\ln x}{1-x^2}\right|\le M$. 故
$|R_n(x)|\le Mx^{2n}$，
\[
  \left|\int_0^1 R_n(x)\,dx\right|
  \le\int_0^1|R_n(x)|\,dx
  \le M\int_0^1x^{2n}\,dx
  =\frac{M}{2n+1}\to0\quad(\text{当 }n\to\infty\text{ 时}).
\]
此即表明 $\displaystyle\lim_{n\to\infty}\int_0^1 R_n(x)\,dx=0$. 级数可以逐项积分.

\noindent\textbf{e.\quad 和函数的其他性质（综合性问题）}

\noindent\textbf{例 5.2.53}\quad 设 $f_n(x)\ (n=1,2,\cdots)$ 在 $[0,1]$ 上连续，并且
\[
  f_n(x)\ge f_{n+1}(x),\quad \forall x\in[0,1],\ n=1,2,\cdots.
\]
若 $f_n(x)$ 在 $[0,1]$ 收敛于 $f(x)$，试证明 $f(x)$ 在 $[0,1]$ 上达到最大值.（北京大学）

\noindent\textbf{证}\quad 因为 $f_n(x)\downarrow f(x)$，所以
\[
  f(x)\le f_n(x)\quad(\forall n\in\mathbb N,\ \forall x\in[0,1]).
\qquad(1)
\]
因 $f_1(x)$ 在 $[0,1]$ 上连续，所以有上界 $M$，故
$f(x)\le f_1(x)\le M\ (\forall x\in[0,1])$. 因此，
$\mu=\sup\limits_{[0,1]}f(x)$ 存在.

假若 $f(x)$ 在 $[0,1]$ 上不达上确界 $\mu$，则由确界定义，可知 $\exists\{x_n\}\subset[0,1]$，使得
$\displaystyle\lim_{n\to\infty}f(x_n)=\mu$. 利用致密性原理，在有界序列 $\{x_n\}$ 里，必存在收敛的子列
$\{x_{n_k}\}\to x_0\in[0,1]$（当 $k\to\infty$ 时）. 我们证明 $f(x_0)=\mu$. 不然的话，因 $\mu$ 为上确界，
$f(x_0)<\mu$，从而 $\exists\mu_1$ 使得
\[
  f(x_0)<\mu_1<\mu.
\qquad(2)
\]
因为 $\displaystyle\lim_{n\to\infty}f_n(x_0)=f(x_0)$，所以 $\exists n_1$ 使得
$f_{n_1}(x_0)<\mu_1$. 又因为 $f_{n_1}$ 连续，所以 $\exists\delta>0$，当
$x\in(x_0-\delta,x_0+\delta)\cap[0,1]$ 时，有 $f_{n_1}(x)<\mu_1$. 由于
$x_{n_k}\to x_0$（当 $k\to\infty$），对 $\delta>0$，$\exists K>0$，当 $k>K$ 时，
$|x_0-x_{n_k}|<\delta$，于是 $f_{n_1}(x_{n_k})<\mu_1$. 联系式（1），知
$f(x_{n_k})<\mu_1$，故
$\mu=\displaystyle\lim_{k\to\infty}f(x_{n_k})\le\mu_1$. 与式（2）矛盾. 证毕.

\noindent\textbf{例 5.2.54}\quad 设 $\{f_n(x)\}$ 是在 $(-\infty,+\infty)$ 上定义并且连续的函数序列，试构造一个函数
$f(x)$，使之在 $(-\infty,+\infty)$ 上有界、连续，当且仅当对所有 $n:f_n(x)=0$ 时
$f(x)=0$.

\noindent\textbf{提示}\quad 令
\[
  f(x)=\sum_{n=1}^{\infty}a_n\frac{|f_n(x)|}{1+|f_n(x)|}
\]
（其中 $\displaystyle\sum_{n=1}^{\infty}a_n$ 为任意一个收敛的正项级数）.

\noindent\textbf{※例 5.2.55}\quad 设 $f_n(x)$ 在 $[a,b]$ 上可积 $(n=1,2,\cdots)$，且在 $[a,b]$ 上
$f_n(x)\rightrightarrows f(x)$（当 $n\to\infty$ 时），试证 $f(x)$ 在 $[a,b]$ 上可积.（天津大学）

\noindent\textbf{分析}\quad $f$ 可积的充分必要条件是：$\forall\varepsilon>0$，$\exists [a,b]$ 的一个分割 $T$，使得
$\displaystyle\sum_{i=0}^{k-1}\omega_i'\Delta x_i<\varepsilon$. 已知 $f_n$ 可积，所以对 $f_n$ 可找到这种分割 $T$.
又因 $n$ 充分大时，$f_n$ 可以任意逼近 $f$（关于 $x\in[a,b]$ 一致）. 因此，取充分大的 $n$，对 $f_n$ 取分割
$T$，然后以此分割作为 $f$ 的分割即可.

\noindent\textbf{证}\quad 因当 $n\to\infty$ 时，$f_n(x)\rightrightarrows f(x)$（关于 $x\in[a,b]$），故
$\forall\varepsilon>0$，$\exists N>0$，当 $n>N$ 时，对于一切 $x\in[a,b]$，恒有
\[
  |f(x)-f_n(x)|<\frac{\varepsilon}{4(b-a)}.
\]
又因 $f_n$ 在 $[a,b]$ 上可积，对此 $\varepsilon>0$，$\exists [a,b]$ 的一个分割
$T:a=x_0<x_1<\cdots<x_k=b$，使得
\[
  \sum_{i=0}^{k-1}\omega_i\Delta x_i<\frac\varepsilon2,\quad
  \text{其中 }\omega_i=\sup_{x_i',x_i''\in[x_{i-1},x_i]}|f_n(x_i')-f_n(x_i'')|.
\]
但是 $\forall x_i',x_i''\in[x_{i-1},x_i]$，有
\[
\begin{aligned}
|f(x_i')-f(x_i'')|
\le{}& |f(x_i')-f_n(x_i')|+|f_n(x_i')-f_n(x_i'')|
+|f_n(x_i'')-f(x_i'')|\\
\le{}&\frac{\varepsilon}{4(b-a)}+\omega_i+\frac{\varepsilon}{4(b-a)}
=\frac{\varepsilon}{2(b-a)}+\omega_i.
\end{aligned}
\]
所以
\[
  \omega_i'\equiv\sup_{x_i',x_i''\in[x_{i-1},x_i]}|f(x_i')-f(x_i'')|
  \le\frac{\varepsilon}{2(b-a)}+\omega_i.
\]
因此
\[
  \sum_{i=0}^{k-1}\omega_i'\Delta x_i
  \le\frac{\varepsilon}{2(b-a)}\sum_{i=0}^{k-1}\Delta x_i
  +\sum_{i=0}^{k-1}\omega_i\Delta x_i
  \le\frac\varepsilon2+\frac\varepsilon2=\varepsilon.
\]
证毕.

\subsection*{三、一致逼近问题}

\noindent\textbf{定理（Weierstrass 多项式逼近定理）}\quad 有界闭区间 $[a,b]$ 上的连续函数 $f(x)$ 可用多项式一致逼近，即存在多项式序列
$\{P_n(x)\}$（其中 $P_n(x)$ 是 $n$ 次多项式），使得 $n\to\infty$ 时
$P_n(x)\rightrightarrows f(x)$（在 $[a,b]$ 上关于 $x$ 一致），亦即：
$\forall\varepsilon>0$，$\exists N>0$，使得当 $n>N$ 时，有
\[
  |f(x)-P_n(x)|<\varepsilon\quad(\forall x\in[a,b]).
\]

\noindent\textbf{例 5.2.56}\quad 设 $f(x)$ 在 $[0,1]$ 上连续，且
$\displaystyle\int_0^1f(x)x^n\,dx=0\ (n=0,1,2,\cdots)$，证：
$f(x)\equiv0$.（中国科学技术大学）

\noindent\textbf{证 I}\quad 由 $f$ 在 $[0,1]$ 上连续，可得 $f(x)$ 在 $[0,1]$ 上必有界：$\exists M>0$，
$|f(x)|\le M\ (\forall x\in[0,1])$. 根据多项式逼近定理，$\forall\varepsilon_n>0$，
$\exists P_n(x)\equiv\displaystyle\sum_{k=0}^n a_kx^k$（多项式），使得 $\forall x\in[0,1]$，恒有
$|f(x)-P_n(x)|<\dfrac{\varepsilon_n}{M}$.

因 $\displaystyle\int_0^1f(x)x^n\,dx=0\ (n=0,1,2,\cdots)$，知
$\displaystyle\int_0^1f(x)P_n(x)\,dx=0$. 故
\[
\begin{aligned}
0\le\int_0^1f^2(x)\,dx
={}&\left|\int_0^1 f(x)(f(x)-P_n(x))\,dx+\int_0^1f(x)P_n(x)\,dx\right|\\
\le{}&\int_0^1|f(x)(f(x)-P_n(x))|\,dx
\le M\cdot\frac{\varepsilon_n}{M}=\varepsilon_n.
\end{aligned}
\]
再令 $\varepsilon_n\to0$，可知 $\displaystyle\int_0^1f^2(x)\,dx=0$，$f(x)\equiv0$.

\noindent\textbf{证 II}\quad 根据多项式逼近定理，存在 $P_n(x)\rightrightarrows f(x)$（当 $n\to\infty$ 时，关于
$x\in[0,1]$ 一致）. 因 $f$ 有界，故 $f(x)P_n(x)\rightrightarrows f^2(x)$（当 $n\to\infty$ 时，关于
$x\in[0,1]$ 一致）. 从而可在积分号下取极限：
\[
  0\le\int_0^1 f^2(x)\,dx
  =\int_0^1\lim_{n\to\infty}f(x)P_n(x)\,dx
  =\lim_{n\to\infty}\int_0^1f(x)P_n(x)\,dx=0,
\]
得 $\displaystyle\int_0^1f^2(x)\,dx=0$，$f(x)\equiv0$.

\noindent\textbf{例 5.2.57}\quad 若实系数多项式序列 $\{P_n(x)\}$ 在 $\mathbb R$ 上一致收敛于函数 $f(x)$，试证：
$f(x)$ 必是多项式函数.（华东师范大学，北京大学）

\noindent\textbf{证}\quad 因已知 $n\to\infty$ 时，$P_n(x)\rightrightarrows f(x)$（关于 $x\in\mathbb R$ 一致）. 由 Cauchy 准则：
$\forall\varepsilon>0$，$\exists N>0$，当 $m>n>N$ 时，有
\[
  |P_m(x)-P_n(x)|<\varepsilon\quad(\forall x\in\mathbb R).
\qquad(1)
\]
由此能断言：当 $m>n>N$ 时，$P_m(x)$，$P_n(x)$ 只可能是同一多项式，最多只是常数项不同（不然的话，当
$x\to+\infty$ 时，$|P_m(x)-P_n(x)|\to+\infty$，与式（1）矛盾）. 因此，$P_n(x)$ 可记作
\[
  P_n(x)=P(x)+a_n\quad(a_n\text{ 是常数项}).
\qquad(2)
\]
代入式（1），得 $\forall\varepsilon>0$，$\exists N>0$，当 $m>n>N$ 时有
$|a_m-a_n|<\varepsilon$. 由 Cauchy 准则知 $\displaystyle\lim_{n\to\infty}a_n\overset{\text{存在}}=c$.
于是在式（2）里令 $n\to\infty$，取极限得 $f(x)=P(x)+c$（$f$ 为多项式）. 证毕.

\noindent\textbf{类题}\quad 设 $I$ 是无穷区间，$f(x)$ 在 $I$ 上连续（不是多项式），试证：不存在多项式序列
$\{P_n(x)\}$ 在 $I$ 上一致收敛于 $f(x)$.（北京大学，中国科学技术大学）

\noindent\textbf{证}\quad （有了上题，本题自明，也可直接证明.）

若有多项式序列 $P_n(x)\rightrightarrows f(x)$（当 $n\to\infty$ 时，关于 $x\in I$），那么（根据 Cauchy 准则）
$\forall\varepsilon>0$，$\exists N>0$，当 $m>n>N$ 时，
\[
  |P_n(x)-P_m(x)|<\varepsilon\quad(\forall x\in I).
\qquad(1)
\]
令 $x\to+\infty$，仍保持成立，说明：（当 $m>n>N$ 时）$P_n(x)-P_m(x)$ 中只含 $x$ 的零次项，换句话说：
对于 $P_n(x)$ 和 $P_m(x)$，如果不看常数项，则 $\{P_n(x)\}_{n>N}$ 统统是同一个多项式，（当
$n\to\infty$ 时）$P_n(x)\to f(x)$，说明 $f(x)$ 只能是一多项式. 矛盾. 证毕.

\noindent\textbf{※例 5.2.58}\quad 设 $f(x)$ 在 $I$ 上有连续导数，
$F_n(x)=n\left[f\left(x+\dfrac1n\right)-f(x)\right]$.

1）证明：若 $I$ 是有界闭区间，则 $\{F_n(x)\}$ 在 $I$ 上一致收敛；

2）如果 $I$ 是有界开区间，问：$\{F_n(x)\}$ 是否仍在 $I$ 上一致收敛.（北京大学）

\noindent\textbf{证}\quad 根据导数定义，明显有
\[
  F_n(x)=n\left[f\left(x+\frac1n\right)-f(x)\right]\to f'(x)\quad(n\to\infty).
\]
要证在 $I$ 上 $\{F_n(x)\}$ 一致收敛，即要证明：$\forall\varepsilon>0$，$\exists N>0$，使得 $n>N$ 时有
\[
  |F_n(x)-f'(x)|<\varepsilon\quad(\forall x\in I).
\qquad(1)
\]
1）当 $I$ 是有界闭区间时：若 $f'(x)$ 连续，根据 Cantor 定理，可知 $f'(x)$ 一致连续，即：
$\forall\varepsilon>0$，$\exists\delta>0$，$\forall x_1,x_2\in I$，当 $|x_1-x_2|<\delta$ 时，有
\[
  |f'(x_2)-f'(x_1)|<\varepsilon.
\qquad(2)
\]
如此，$\forall x\in I$，只要取 $N>\dfrac1\delta$，则当 $n>N$ 时，利用 Lagrange 定理，
$\exists\xi_n:x<\xi_n<x+\dfrac1n$，使得
\[
  f\left(x+\frac1n\right)-f(x)=f'(\xi_n)\frac1n,
\]
$|\xi_n-x|<\dfrac1n<\dfrac1N<\delta$. 由式（2）知 $|f'(\xi_n)-f'(x)|<\varepsilon$. 故
\[
  |F_n(x)-f'(x)|
  =\left|n\left[f\left(x+\frac1n\right)-f(x)\right]-f'(x)\right|
  <|f'(\xi_n)-f'(x)|<\varepsilon,
\]
表明式（1）成立. 所以当 $n\to\infty$ 时，$F_n(x)\rightrightarrows f'(x)$（在 $I$ 上关于 $x$ 一致）. 故
$\{F_n(x)\}$ 在 $I$ 上一致收敛.

2）否！（作反例，使得 $f$ 和 $f'$ 是在局部无限变陡的函数（如：$\dfrac1x,\ln x,\cot x,\cdots$）.）

反例：令 $f(x)=\dfrac1x\ (\forall x\in(0,1))$，
\[
  \lim_{n\to\infty}F_n(x)
  =\lim_{n\to\infty}n\left[f\left(x+\frac1n\right)-f(x)\right]
  =f'(x)=\left(\frac1x\right)'=-\frac1{x^2}.
\]
若取 $\varepsilon_0=\dfrac12$，$\forall n>1$，令 $x_n=\dfrac1n$，则
\[
  f'(x_n)=-\frac1{x_n^2}=-n^2,\quad
  F_n(x_n)=n\left[f\left(x_n+\frac1n\right)-f(x_n)\right]
  =n\left(\frac n2-n\right)=-\frac{n^2}{2}.
\]
因此
\[
  |F_n(x_n)-f'(x_n)|
  =\left|-\frac{n^2}{2}+n^2\right|
  =\frac{n^2}{2}>\varepsilon_0.
\]
说明式（2）不成立. 此反例说明：当有界闭区间改为有界开区间时，原结论不再成立.

\section*{单元练习 5.2}

\noindent\textbf{5.2.1}\quad 1）设

i）$f_n(x)$ 在 $[a,b]$ 上连续，$n=1,2,\cdots$；

ii）$\{|f_n(x)|\}$ 在 $[a,b]$ 上一致收敛于 $f(x)$；

iii）在 $[a,b]$ 上 $f_n(x)\le f(x)$，$n=1,2,\cdots$；

试证 $e^{f_n(x)}$ 在 $[a,b]$ 上一致收敛于 $e^{f(x)}$；

2）若将 1）中条件 iii）去掉，则 $\{e^{f_n(x)}\}$ 是否还一致收敛，试证明你的结论.（河北师范大学）

\noindent\textbf{提示}\quad 2）$\exists M>0$，使 $|f(x)|\le M$，从而 $n$ 充分大时，
\[
  |f_n(x)|\le |f(x)|+|f_n(x)-f(x)|\le M+1,
\]
\[
  0\le |e^{f(x)}-e^{f_n(x)}|\le e^{M+1}|f_n(x)-f(x)|.
\]
或者用 $e^x$ 在 $[-M-1,M+1]$ 上一致连续. 由 $\varepsilon>0$ 找 $\delta>0$，再根据
$f_n\rightrightarrows f$，由 $\delta>0$ 找 $N_0$.

\noindent\textbf{5.2.2}\quad 设
\[
  f_n(x)=\sum_{k=1}^n\frac1n\cos\left(x+\frac kn\right),\quad n=1,2,\cdots,
\]
证明在 $(-\infty,+\infty)$ 上 $\{|f_n(x)|\}$ 一致收敛.（兰州大学）

\noindent\textbf{提示}\quad 可参看例 5.2.1.

\noindent\textbf{※5.2.3}\quad 设
\[
  f_n(x)=\frac{\displaystyle\int_0^x(1-t^2)^n\,dt}{\displaystyle\int_0^1(1-t^2)^n\,dt},\quad
  g_n(x)=\int_0^x f_n(t)\,dt,
\]
试证：

1）当 $n\to\infty$ 时，
\[
  f_n(x)\to
  \begin{cases}
    -1,&-1\le x\le-\varepsilon,\\
    1,&\varepsilon\le x\le1
  \end{cases}
  \quad(0<\varepsilon<1);
\]

2）$g_n(x)\rightrightarrows |x|$ 关于 $x\in[-1,1]$，当 $n\to\infty$ 时.

\noindent\textbf{提示}\quad 1）可参看例 4.1.6.

2）注意 $|x|=\displaystyle\int_0^x\operatorname{sgn}x\,dx$，然后用分段法（见例 4.1.4）证明
$\displaystyle\int_0^x(f_n(x)-\operatorname{sgn}x)\,dx\to0\ (n\to\infty)$，关于 $x\in[-1,1]$ 一致.

\noindent\textbf{※5.2.4}\quad 试证级数 $\displaystyle\sum_{n=1}^{\infty}(-1)^n(1-x)x^n$ 在区间 $[0,1]$ 上绝对收敛、一致收敛，但不是绝对一致收敛（指各项取绝对值之后，仍一致收敛）.

\noindent\textbf{提示}\quad 在 $[0,1]$ 上，
$\displaystyle\sum_{k=0}^{n-1}(1-x)x^k\to S(x)=
\begin{cases}
0,&\text{当 }x=1\text{ 时},\\
1,&\text{当 }0\le x<1\text{ 时},
\end{cases}$
可见原级数绝对收敛，$S(x)$ 不连续，违反保证连续性，故原级数非绝对一致收敛. 原级数为 Leibniz 级数，
\[
  |S(x)-S_{n-1}(x)|=|r_{n-1}(x)|\le(1-x)x^n
  \le\max_{x\in[0,1]}(1-x)x^n
  =\frac1{n+1}\left(\frac n{n+1}\right)^n
  \le\frac1{n+1}\to0\quad(n\to\infty),
\]
故原级数在 $[0,1]$ 上一致收敛.

\noindent\textbf{5.2.5}\quad 判断级数
$\displaystyle\sum_{n=1}^{\infty}\frac{(-1)^n}{x+n}$ 在 $0<x<+\infty$ 内是否一致收敛.

\noindent\textbf{提示}\quad $|r_n|\le\dfrac1{n+1}\to0\ (n\to\infty)$.

\noindent\textbf{5.2.6}\quad 讨论级数 $\displaystyle\sum_{n=1}^{\infty}xe^{-(n-1)x}$ 关于 $0\le x\le1$ 是否一致收敛？（复旦大学）

\noindent\textbf{提示}\quad
\[
  r_n(x)=\frac{xe^{-nx}}{1-e^{-x}}\to0(n\to\infty),
\]
但 $\displaystyle\sup_{x\in[0,1]}r_n(x)=1\ne0$，$[0,1]$ 上非一致收敛. $\forall\,0<\varepsilon<1$，当
$x\in[\varepsilon,1]$ 时，
$0\le r_n(x)\le\dfrac{e^{-n\varepsilon}}{1-e^{-\varepsilon}}\to0(n\to\infty)$，故级数在 $[\varepsilon,1]$ 上一致收敛.

\noindent\textbf{※5.2.7}\quad 讨论级数
$\displaystyle\sum_{n=1}^{\infty}\frac{n^2}{\left(x+\dfrac1n\right)^n}$ 的收敛性和一致收敛性（$x\ge0$）.（华东师范大学）

\noindent\textbf{提示}\quad 1）当 $x\in(1,+\infty)$ 时，因 $\forall c>1$，$x>c$ 时通项
$u_n(x)\le\dfrac{n^2}{c^n}$，而 $\displaystyle\sum_{n=1}^{\infty}\dfrac{n^2}{c^n}$ 收敛，故原级数在
$(1,+\infty)$ 内收敛且内闭一致收敛. 在 $(1,+\infty)$ 内，令 $x_n=1+\dfrac1n$，通项
$u_n(x_n)=\dfrac{n^2}{\left(1+\dfrac2n\right)^n}\nrightarrow0$，故原级数在 $(1,+\infty)$ 内非一致收敛.
2）当 $x\le1$ 时，通项
$u_n(x)\ge\dfrac{n^2}{\left(1+\dfrac1n\right)^n}\to+\infty$，原级数发散.

\noindent\textbf{5.2.8}\quad 讨论级数
$\displaystyle\sum_{n=1}^{\infty}\frac{x^{2n}}{1+x^{2n+1}}\ (x\ge0)$ 的一致收敛性.（南京大学）

\noindent\textbf{提示}\quad 当 $|x|\ge1$ 时，通项不趋 0，发散. $\forall q:0<q<1$，当 $0\le x\le q$ 时，
$0\le$ 通项 $<q^{2n}$，级数在 $[0,q]$ 上一致收敛；但在 $(0,1)$ 上非一致收敛，因通项
$\to\dfrac12\ne0$（当 $x\to1^-$ 时），与 Cauchy 准则矛盾.

\noindent\textbf{※5.2.9}\quad 设函数项级数 $\displaystyle\sum_{n=1}^{\infty}u_n(x)$ 在 $[a,b]$ 上收敛，试证：若对任何
$x\in[a,b]$，$\exists\delta_x>0,G_x>0$，使对任意 $y\in(x-\delta_x,x+\delta_x)\cap[a,b]$ 与自然数 $n$，都有
\[
  \left|\sum_{k=1}^n u_k'(y)\right|<G_x,
\]
则 $\displaystyle\sum_{k=1}^{\infty}u_k(x)$ 在 $[a,b]$ 上
一致收敛.（北京师范大学）

\noindent\textbf{提示}\quad 用有限覆盖定理 $\Rightarrow$ 部分和的导数一致有界 $\Rightarrow$ 等度连续 $\Rightarrow$ 一致收敛.

\noindent\textbf{※5.2.10}\quad 设 $b_1\ge b_2\ge\cdots\ge b_n\ge\cdots\ge0$，试证：级数
$\displaystyle\sum_{n=1}^{\infty}b_n\sin nx$ 在任意区间上一致收敛的充要条件是 $n\to\infty$ 时
$nb_n\to0$.

\noindent\textbf{证}\quad $1^\circ$（必要性）已知
$\displaystyle\sum_{k=1}^{\infty}b_k\sin kx$ 在 $(-\infty,+\infty)$ 上一致收敛，（应用 Cauchy 准则）
$\forall\varepsilon>0$，$\exists N>0$，$\forall n>N$，有
\[
  \left|\sum_{k=n+1}^{2n}b_k\sin kx\right|<\frac{\sqrt2}{4}\varepsilon,\quad
  x\in(-\infty,+\infty).
\]
对此 $n$，取 $x=\dfrac{\pi}{4n}$，则
\[
  \frac{\sqrt2}{4}\varepsilon
  >\sum_{k=n+1}^{2n}b_k\sin\left(k\frac{\pi}{4n}\right)
  >nb_{2n}\sin\frac\pi4=2nb_{2n}\frac{\sqrt2}{4}.
\]
可见 $\displaystyle\lim_{n\to\infty}2nb_{2n}=0$，同理，有
$\displaystyle\lim_{n\to\infty}(2n+1)b_{2n+1}=0$，故
$\displaystyle\lim_{n\to\infty}nb_n=0$ 获证.

$2^\circ$（充分性）因为 $\sin kx$ 是以 $2\pi$ 为周期的奇函数，只要证明
$\displaystyle\sum_{k=1}^{\infty}b_k\sin kx$ 在 $[0,\pi]$ 上一致收敛，则它在任意区间上都一致收敛. 又由
$b_n>0$ 单调递减，$\displaystyle\lim_{n\to\infty}nb_n=0$，可知 $\{b_n\}$ 单调下降趋向 0（关于
$x$ 一致）. 而已知
$\left|\displaystyle\sum_{k=1}^n\sin kx\right|\le\dfrac1{\left|\sin\dfrac x2\right|}$
（见例 5.2.26 式（1）），因此，当 $\dfrac\pi2\le x<\pi$ 时，有
\[
  \left|\sum_{k=1}^n\sin kx\right|
  \le\frac1{\left|\sin\dfrac x2\right|}\le\sqrt2.
\]
利用 Dirichlet 判别法，知
$\displaystyle\sum_{k=1}^{\infty}b_k\sin kx$ 在 $\left[\dfrac\pi2,\pi\right]$ 上一致收敛.

（剩下只需证明：原级数 $\displaystyle\sum_{k=n}^{\infty}b_k\sin kx$ 在
$\left[0,\dfrac\pi2\right)$ 内一致收敛.）根据 Cauchy 准则，只要证明：
$\forall\varepsilon>0$，$\exists N>0$，$\forall m>n>N$，$\forall x\in\left[0,\dfrac\pi2\right)$，有
\[
  \left|\sum_{k=n+1}^{m}b_k\sin kx\right|<\varepsilon.
\qquad(1)
\]
为此，记
\[
  \beta_n=\sup\{kb_k\mid k\ge n\}.
\qquad(2)
\]
则由 $\displaystyle\lim_{n\to\infty}nb_n=0$ 知 $\displaystyle\lim_{n\to\infty}\beta_n=0$. 因此，
$\forall\varepsilon>0$，$\exists N>0$，当 $n>N$ 时，有
$\beta_n<\dfrac{\varepsilon}{2\pi}$，即
\[
  \pi\beta_n<\frac\varepsilon2<\varepsilon.
\qquad(3)
\]

（下面用 Cauchy 准则来证原级数一致收敛.）为此，将 $\left[0,\dfrac\pi2\right)$ 分成三个区间，证明式（3）里的
$N$ 能作为收敛准则所需要的 $N$，分别适用于三区间. $\forall m>n>N$，令
\[
  \left[0,\frac\pi2\right)=\left[0,\frac\pi m\right)\cup
  \left[\frac\pi m,\frac\pi n\right)\cup
  \left[\frac\pi n,\frac\pi2\right),
\]
则

$1^\circ$ 当 $x\in\left[0,\dfrac\pi m\right)$ 时，$0\le mx<\pi$，
\[
  \left|\sum_{k=n+1}^m b_k\sin kx\right|
  \le\sum_{k=n+1}^m b_kkx
  \overset{\text{式（2）}}{\le}(m-n)x\beta_n
  <mx\beta_n<\pi\beta_n
  \overset{\text{式（3）}}{<}\frac\varepsilon2<\varepsilon.
\qquad(4)
\]

下面要用到如下的三角不等式：当 $x\in\left(0,\dfrac\pi2\right]$ 时，$\forall m>n>N\ge1$，有
\[
  \left|\sum_{k=n+1}^m\sin kx\right|\le\frac1{\sin\dfrac x2}.
\qquad(5)
\]
（其推导可参照例 5.2.26 的式（1）.）

$2^\circ$ 当 $x\in\left[\dfrac\pi n,\dfrac\pi2\right)$ 时，$\forall m>n>N$，上面式（5）成立. 因
$b_n>0$ 单调递减，可利用 Abel 变换（见例 5.1.35，这里 $M=\dfrac1{\sin\dfrac x2}$）
\[
\begin{aligned}
\left|\sum_{k=n+1}^m b_k\sin kx\right|
&\overset{\text{Abel变换}}{\le}\frac{2b_{n+1}}{\sin\dfrac x2}
\quad\left(\text{应用 }\sin x>\frac2\pi x\ (0<x<\frac\pi2)\right)\\
&\le 2b_{n+1}\frac1{\dfrac2\pi\cdot\dfrac x2}
=\frac{2b_{n+1}\pi}{x}
\le2(n+1)b_{n+1}\le\pi\beta_n<\frac\varepsilon2.
\end{aligned}
\qquad(6)
\]

$3^\circ$ 当 $x\in\left[\dfrac\pi m,\dfrac\pi n\right)=
\displaystyle\bigcup_{i=n}^{m-1}\left[\dfrac\pi{i+1},\dfrac\pi i\right)$，例如，
$x\in\left[\dfrac\pi{i+1},\dfrac\pi i\right)\ (i\in\{n,\cdots,m-1\})$（此时
$i<\dfrac\pi x\le i+1$）时，
\[
  \left|\sum_{k=n+1}^{m}b_k\sin kx\right|
  \le\left|\sum_{k=n+1}^{i}b_k\sin kx\right|
  +\left|\sum_{k=i+1}^{m}b_k\sin kx\right|
  \overset{\text{记}}=I_1+I_2<\varepsilon,
\]
其中
\[
  I_1=\left|\sum_{k=n+1}^{i}b_k\sin kx\right|
  \overset{\text{式（4）}}{<}\frac\varepsilon2,\quad
  I_2=\left|\sum_{k=i+1}^{m}b_k\sin kx\right|
  \overset{\text{式（6）}}{<}\frac\varepsilon2.
\]
因此，式（1）获证，级数一致收敛. 证毕.

\noindent\textbf{5.2.11}\quad 试证级数
$\displaystyle\sum_{n=1}^{\infty}\frac{x+n(-1)^n}{x^2+n^2}$ 在 $(-\infty,+\infty)$ 内闭一致收敛（即在任何内闭区间
$[a,b]\subset(-\infty,+\infty)$ 上一致收敛）.

\noindent\textbf{提示}\quad 拆成两个级数之和：
$\displaystyle\sum_{n=1}^{\infty}\frac{x}{x^2+n^2}$，$\forall A>0$，$[-A,A]$ 上以
$\displaystyle\sum_{n=1}^{\infty}\frac A{n^2}$ 为优级数；
$\displaystyle\sum_{n=1}^{\infty}\frac{(-1)^nn}{x^2+n^2}$ 是 Leibniz 级数，$|R_n|<\dfrac1n\to0$.

\noindent\textbf{※5.2.12}\quad 指出级数 $\displaystyle\sum_{n=1}^{\infty}\frac{e^{-nx}}n$ 的收敛区间和一致收敛区间，并证明之.（兰州大学）

\noindent\textbf{提示}\quad 当且仅当 $(0,+\infty)$ 内收敛；$\forall a>0$，$[a,+\infty]$ 上一致收敛（Dirichlet）.

\noindent\textbf{5.2.13}\quad 指出 $\displaystyle\lim_{n\to\infty}\frac{x^2}{n}\ln\left(\frac{x^2}{n}+1\right)$ 的收敛与一致收敛的范围.（兰州大学）

\noindent\textbf{提示}\quad $\mathbb R$ 上收敛，且内闭一致收敛（$\forall A>0$，$[-A,A]$ 上一致收敛）.

\noindent\textbf{※5.2.14}\quad 设 $f(x)$ 在 $[0,1]$ 上连续，
\[
  f_1(x)=f(x),\quad f_{n+1}(x)=\int_x^1 f_n(t)\,dt,\quad \forall x\in[0,1],\ n=1,2,\cdots.
\]
求证：$\displaystyle\sum_{n=1}^{\infty}f_n(x)$ 在 $[0,1]$ 上一致收敛.（北京航空航天大学）

\noindent\textbf{提示}\quad $\exists M>0$，使 $[0,1]$ 上 $|f(x)|\le M$.
$|f_n(x)|\le\dfrac{M(1-x)^{n-1}}{(n-1)!}\le\dfrac M{(n-1)!}$，故
$\displaystyle\sum_{n=1}^{\infty}\dfrac M{(n-1)!}$ 为优级数.

\noindent\textbf{5.2.15}\quad 1）证明函数序列
$\left\{\left(1+\dfrac xn\right)^n\right\}\ (n=1,2,\cdots)$ 在 $x\in[0,1]$ 上对 $n$ 单调增大；

2）证明 $\displaystyle\sum_{n=1}^{\infty}\frac{(-1)^n(n+x)^n}{n^{n+1}}$ 在 $[0,1]$ 上一致收敛.（南京航空航天大学）

\noindent\textbf{提示}\quad 1）可用平均值不等式
\[
  \sqrt[n+1]{\left(1+\frac xn\right)^n\cdot1}
  \le\frac{n\left(1+\dfrac xn\right)+1}{n+1}.
\]

2）通项为 $\dfrac{(-1)^n}{n}\left(1+\dfrac xn\right)^n$，可用 Abel 判别法.

\noindent\textbf{5.2.16}\quad 试证：若级数 $\displaystyle\sum_{n=0}^{\infty}a_n$ 收敛，则 Dirichlet 级数
$\displaystyle\sum_{n=1}^{\infty}\frac{a_n}{n^x}$ 在 $[0,+\infty)$ 上一致收敛.（陕西师范大学）

\noindent\textbf{提示}\quad 可用 Abel 判别法.

\noindent\textbf{※5.2.17}\quad 证明级数
$\displaystyle\sum_{n=1}^{\infty}\frac{(-1)^{[\sqrt n]}}{\sqrt n(n+x)}$ 在 $0\le x<+\infty$ 上一致收敛.

\noindent\textbf{提示}\quad 可参看例 5.1.32，
\[
  \sum_{n=1}^{\infty}\frac{(-1)^{[\sqrt n]}}n\cdot\frac1{\sqrt{1+\dfrac xn}}\quad(\text{Abel}).
\]

\noindent\textbf{5.2.18}\quad 试证：$\forall\alpha:0<\alpha<\dfrac\pi2$，函数项级数
$\displaystyle\sum_{n=1}^{\infty}x^n\left(1-\dfrac{2x}{\pi}\right)^n\tan^n x$ 在 $[0,\alpha]$ 上一致收敛. 若记其和函数为
$S(x)$，试证
\[
  \lim_{x\to(\frac\pi2)^-}S(x)=+\infty.
\]
（北京师范大学）

\noindent\textbf{提示}\quad 原级数
$\displaystyle=\sum_{n=1}^{\infty}q^n=\dfrac q{1-q}$（其中
$0<q=x\left(1-\dfrac{2x}{\pi}\right)\tan x
=\dfrac2\pi\cdot x\left(\dfrac\pi2-x\right)\tan x<1$）.

\noindent\textbf{再提示}\quad $1^\circ$ 在 $\left[0,\dfrac\pi4\right]$ 上，函数
$y=x\left(\dfrac\pi2-x\right)$ 和 $y=\tan x$ 都严增，所以 $q(x)$ 也严增；在
$\left[\dfrac\pi4,\dfrac\pi2\right)$ 上，函数 $y=x$ 和
$y=\left(1-\dfrac{2x}{\pi}\right)\tan x$ 都严增，所以 $q(x)$ 也严增. 故在
$\left[0,\dfrac\pi2\right)$ 上 $q(x)$ 严增，且 $q(0)=0$. 于是
$q(x)>0,\ x\in\left(0,\dfrac\pi2\right)$.

$2^\circ$
\[
\begin{aligned}
q\left(\frac\pi2-0\right)
&=\lim_{x\to(\frac\pi2)^-}x\left(1-\frac{2x}{\pi}\right)\tan x\\
&=\lim_{x\to(\frac\pi2)^-}\frac2\pi\cdot x\left(\frac\pi2-x\right)\tan x
\end{aligned}
\]
令 $t=\dfrac\pi2-x$，
\[
  =\lim_{t\to0+}\frac2\pi\left(\frac\pi2-t\right)\cos t\frac t{\sin t}=1.
\]
可见在 $\left[0,\dfrac\pi2\right)$ 上，$0<q=q(x)<1$，从而原级数在
$\left[0,\dfrac\pi2\right)$ 上逐点收敛.

$3^\circ$ 用 $M$ 判别法易证：$\forall\alpha\in\left(0,\dfrac\pi2\right)$，原级数在 $[0,\alpha]$ 上一致收敛.

\noindent\textbf{5.2.19}\quad 证明：
$\displaystyle\sum_{n=1}^{\infty}(-1)^n\dfrac{x^{n^2}+n}{n^2}$ 在任何有穷区间上一致收敛，而在任何一点都不绝对收敛.（华中科技大学）

\noindent\textbf{提示}\quad 可参考例 5.2.24.

\noindent\textbf{5.2.20}\quad 讨论级数 $\displaystyle\sum_{n=1}^{\infty}x^n(\ln x)^2$ 在 $[0,1]$ 区间上的一致收敛性.（北京大学）

\noindent\textbf{提示}\quad 设 $x=0$ 时通项 $=0$，在 $[0,1]$ 上可用 Dini 定理（见例 5.2.27）.

\noindent\textbf{※5.2.21}\quad 设 $g(x)$ 和函数序列 $\{f_n(x)\}\ (n=1,2,\cdots)$ 在区间 $[a,b]$ 上连续，且对任一
$x\in[a,b]$，$\displaystyle\lim_{n\to\infty}f_n(x)=g(x)$，问能否断定 $f_n(x)$ 在 $[a,b]$ 上一致收敛于 $g(x)$？
论证你的结论.（兰州大学）

\noindent\textbf{提示}\quad 例如可在 $[0,1]$ 上考虑 $f_n(x)=x^n(1-x^n)$.

\noindent\textbf{※5.2.22}\quad 证明：
$\displaystyle\sum_{n=1}^{\infty}\left[nxe^{-nx}-(n+1)xe^{-(n+1)x}\right]$ 在 $[0,+\infty)$ 内收敛，但对任何
$A>0$，级数在 $[0,A]$ 上均不一致收敛；再证：上述级数在 $[0,+\infty)$ 内定义了一个连续函数，问级数在
$[0,A]\ (A>0)$ 上可否逐项积分？（南京大学）

\noindent\textbf{提示}\quad 在 $(0,+\infty)$ 上和 $S(x)=xe^{-x}$，
$|r_n(1/n)|=e^{-1}\ne0\ (n\to\infty)$，在 $[0,A]$ 上级数非一致收敛；但
$\forall A>0$，在 $[0,A]$ 上仍能逐项积分. 积分值之和亦等于 $-Ae^{-A}-e^{-A}+1$.

\noindent\textbf{※5.2.23}\quad 在 $(0,1)$ 内任取一数列 $\{a_n\}$（各项互不相同），作级数
$\displaystyle\sum_{k=1}^{\infty}\frac{|x-a_k|}{2^k}$. 证明：

1）该级数在 $(0,1)$ 内定义一个连续函数 $f(x)$；

2）$f(x)$ 在 $x=a_k\ (k=1,2,\cdots)$ 处不可微，而在 $(0,1)$ 内其他点处均可微.（南京大学）

\noindent\textbf{提示}\quad 可参看例 5.2.42 和例 5.2.48.

\noindent\textbf{※5.2.24}\quad 试作 $[0,1]$ 上的连续函数序列 $\{f_n(x)\}$，使之逐点收敛于连续函数 $f(x)$，但
\[
  \lim_{n\to\infty}\int_0^1 f_n(x)\,dx\ne\int_0^1f(x)\,dx.
\]
（安徽大学）

\noindent\textbf{提示}\quad 可考虑 $f_n(x)=nx^n(1-x^n)$ 于 $[0,1]$ 上.

\noindent\textbf{※5.2.25}\quad $k$ 取何值时，

1）$f_n(x)=n^kxe^{-nx}\ (n=1,2,\cdots)$ 在 $[0,1]$ 上收敛；

2）$f_n(x)$ 在 $[0,1]$ 上一致收敛；

3）$\displaystyle\lim_{n\to\infty}\int_0^1 f_n(x)\,dx$ 可在积分号下取极限？

\noindent\textbf{提示}\quad $\forall x\ge0$，$f_n(x)\to0\ (n\to\infty)$，
\[
  |f_n(x)-0|\le\max_{x\ge0}f_n(x)
  =xn^ke^{-nx}\bigg|_{x=\frac1n}
  =n^{k-1}e^{-1}\to
  \begin{cases}
    0,&k<1,\\
    e^{-1},&k=1,\\
    +\infty,&k>1.
  \end{cases}
\]

\noindent\textbf{5.2.26}\quad 证明级数
$\displaystyle\sum_{n=1}^{\infty}\frac{\ln(1+nx)}{nx^n}$ 在 $(1,+\infty)$ 上连续.（西北大学）

\noindent\textbf{提示}\quad $\forall x>1$，可取 $a,b>1$，使 $x\in(a,b)$；再证级数在 $[a,b]$ 上一致收敛.

\noindent\textbf{※5.2.27}\quad 设
\[
  y_{n+1}(x)=\psi(x)+\varphi(y_n(x))\quad(x\in\mathbb R),
\qquad(1)
\]
其中 $\psi(x)$ 是连续有界函数，$y_0(x)=y_0$，$\psi(x_0)=y_0-\varphi(y_0)$. $\varphi$ 满足 Lipschitz 条件：
\[
  |\varphi(y')-\varphi(y'')|\le\alpha|y'-y''|\quad(0<\alpha<1).
\qquad(2)
\]
试证：

1）$\{y_n(x)\}$ 在 $\mathbb R$ 上一致收敛；

2）记 $y(x)=\displaystyle\lim_{n\to\infty}y_n(x)$，则 $y(x)$ 连续，且 $y(x_0)=y_0$；

3）若 $\psi(x)$ 一致连续，则 $y(x)$ 也一致连续.（武汉大学）

\noindent\textbf{证}\quad 1）由式（1），（2）知
\[
  |y_{n+1}(x)-y_n(x)|
  =|\varphi(y_n(x))-\varphi(y_{n-1}(x))|
  \le\alpha|y_n(x)-y_{n-1}(x)|\quad(0<\alpha<1),
\qquad(3)
\]
这表明（1）为压缩映像. $y(x)=\displaystyle\lim_{n\to\infty}y_n(x)$ 存在，且
\[
  y(x)=\sum_{k=0}^{\infty}(y_{k+1}(x)-y_k(x))+y_0.
\qquad(4)
\]
反复用式（3）递推，知：
\[
  |y_{n+1}(x)-y_n(x)|\le\alpha^n|y_1(x)-y_0|\le\alpha^n M.
\]
（因 $y_1(x)-y_0=\psi(x)+\varphi(y_0)-y_0$ 连续有界，记
$M:|y_1(x)-y_0|\le M$.）对 $0<\alpha<1$，
$\displaystyle\sum\alpha^nM$ 收敛，故（4）一致收敛，$y_n(x)\rightrightarrows y(x)$ 于
$\mathbb R$ 上.

2）（1）式令 $n\to\infty$，取极限得
\[
  y(x)=\psi(x)+\varphi(y(x)).
\qquad(5)
\]
于是有 $y(x_0)=\psi(x_0)+\varphi(y(x_0))$. 而已知
$y_0=\psi(x_0)+\varphi(y_0)$，相减得
\[
  |y(x_0)-y_0|=|\varphi(y(x_0))-\varphi(y_0)|
  \le\alpha|y(x_0)-y_0|\quad(0<\alpha<1),
\]
故 $y(x_0)=y_0$. 根据保连续性定理知 $y(x)$ 连续.

3）由（5）式知 $\forall x',x''\in\mathbb R$，有
\[
\begin{aligned}
|y(x')-y(x'')|
&=|\psi(x')-\psi(x'')+\varphi(y(x'))-\varphi(y(x''))|\\
&\le|\psi(x')-\psi(x'')|+\alpha|y(x')-y(x'')|,
\end{aligned}
\]
即
\[
  |y(x')-y(x'')|\le\frac1{1-\alpha}|\psi(x')-\psi(x'')|.
\]
故由 $\psi$ 一致连续，可直接推得 $y(x)$ 一致连续.

\noindent\textbf{5.2.28}\quad 设 $f(x)$ 在 $(-\infty,+\infty)$ 上有任意阶导数 $f^{(n)}(x)$，且任意区间 $[a,b]$ 上
$f^{(n)}(x)\rightrightarrows\phi(x)$（当 $n\to\infty$ 时），求证：
$\phi(x)=ce^x$（其中 $c$ 为常数）.（北京大学）

\noindent\textbf{提示}\quad
\[
  \frac{d\phi(x)}{dx}
  =\lim_{n\to\infty}\frac{df^{(n)}(x)}{dx}
  =\phi(x).
\]

\noindent\textbf{5.2.29}\quad 设
$f(x)=\displaystyle\sum_{n=1}^{\infty}(-1)^{n+1}\dfrac{e^{-nx}}n$，求 1）$f$ 的连续范围；2）$f$ 的可导范围.（北京大学）

\noindent\textbf{提示}\quad 1）在 $x\ge0$ 内闭一致收敛（Abel）；当 $x<0$ 时，通项 $\nrightarrow0$.

2）逐项求导后，$\forall a>0$，$[a,+\infty)$ 上有优级数 $\displaystyle\sum e^{-an}$.

\noindent\textbf{※5.2.30}\quad 设 $f(x)=\displaystyle\sum_{n=0}^{\infty}2^{-2n}\cos 2^nx$，求
$\displaystyle\lim_{x\to0+}x^{-1}(f(x)-f(0))$.（北京师范大学）

\noindent\textbf{提示}\quad $\forall\varepsilon:0<\varepsilon<\dfrac\pi2$，在 $[0,\varepsilon]$ 上级数收敛，且可逐项求导，
\[
  \lim_{x\to0}x^{-1}(f(x)-f(0))=f'(0)
  =\left.\sum_{n=0}^{\infty}(2^{-2n}\cos2^nx)'\right|_{x=0}=0.
\]

\section*{§5.3\quad 幂级数}

\noindent\textbf{导读}\quad 幂级数不仅具有重大的理论意义，而且有着广泛的实用价值. 因而理工科各专业对该内容极为重视，既是教学重点也是各类考试的重点.

\subsection*{一、幂级数的收敛半径与收敛范围}

\noindent\textbf{a.\quad 公式法}

\noindent\textbf{要点}\quad $\displaystyle\sum_{n=1}^{\infty}a_nx^n$ 的收敛半径 $R$ 可按如下公式计算：

i）
\[
  R=\frac1{\displaystyle\varlimsup_{n\to\infty}\sqrt[n]{|a_n|}};
\tag{A}
\]

ii）特别若 $\displaystyle\lim_{n\to\infty}\sqrt[n]{|a_n|}$ 存在或为 $+\infty$，则
\[
  R=\frac1{\displaystyle\lim_{n\to\infty}\sqrt[n]{|a_n|}};
\tag{B}
\]

iii）若 $\displaystyle\lim_{n\to\infty}\left|\frac{a_n}{a_{n+1}}\right|$ 存在或为 $+\infty$，则
\[
  R=\lim_{n\to\infty}\left|\frac{a_n}{a_{n+1}}\right|.
\tag{C}
\]
（约定：在式（A）和（B）中，当分母 $=0$ 时，$R=+\infty$；当分母 $=+\infty$ 时，$R=0$.）

必须注意的是，在求收敛区间时，务必要检验区间端点的敛散性. 对于广义幂级数，可以化为一般幂级数处理，如

\noindent\textbf{例 5.3.1}\quad 求级数
$\displaystyle\sum_{n=1}^{\infty}\left(\sin\frac1{3n}\right)(x^2+x+1)^n$ 的收敛区间.

\noindent\textbf{解}\quad （这是广义幂级数，令 $t=x^2+x+1$，即化为 $t$ 的幂级数
$\displaystyle\sum_{n=1}^{\infty}\left(\sin\dfrac1{3n}\right)t^n$.）利用公式（C），
\[
  R=\lim_{n\to\infty}\frac{\sin\dfrac1{3n}}{\sin\dfrac1{3(n+1)}}
  \overset{\text{等价代换}}=
  \lim_{n\to\infty}\frac{\dfrac1{3n}}{\dfrac1{3(n+1)}}=1.
\]
当 $t=1$ 时，级数 $\displaystyle\sum_{n=1}^{\infty}\sin\dfrac1{3n}$ 发散；当 $t=-1$ 时，级数
$\displaystyle\sum_{n=1}^{\infty}(-1)^n\sin\dfrac1{3n}$ 收敛. 故原级数当且仅当
$-1\le x^2+x+1<1$ 时收敛. 解不等式知收敛区间为 $(-1,0)$.

\noindent\textbf{※例 5.3.2}\quad 讨论级数
\[
  1+\frac1{2x\sqrt2}+\frac1{4x^2\sqrt3}+\frac1{8x^3\sqrt4}
  +\frac1{16x^4\sqrt5}+\cdots+\frac1{2^nx^n\sqrt{n+1}}+\cdots
\]
的收敛性，求出它的收敛区域与一致收敛区域.（北京师范大学）

\noindent\textbf{解}\quad 该级数
\[
  =\sum_{n=0}^{\infty}\frac1{\sqrt{n+1}}\left(\frac1{2x}\right)^n
  \overset{\text{令 }t=\frac1{2x}}{=}
  \sum_{n=0}^{\infty}\frac{t^n}{\sqrt{n+1}}.
\]
对 $t$ 而言，
\[
  R=\lim_{n\to\infty}\frac{a_{n-1}}{a_n}
  =\lim_{n\to\infty}\frac{\sqrt{n+1}}{\sqrt n}=1.
\]
当 $t=1$ 时，级数 $\displaystyle\sum_{n=1}^{\infty}\frac1{\sqrt{n+1}}$ 发散；当 $t=-1$ 时，级数
$\displaystyle\sum_{n=1}^{\infty}\frac{(-1)^n}{\sqrt{n+1}}$ 收敛. 故原级数在
\[
  \left\{x\mid -1\le\frac1{2x}<1\right\}
  =\left\{x\mid x>\frac12\text{ 或 }x\le-\frac12\right\}
\]
内收敛，在 $\left(-\infty,-\dfrac12\right]$ 及 $\left(\dfrac12,+\infty\right)$ 的内闭区间上一致收敛.

\noindent\textbf{例 5.3.3}\quad 求级数
\[
  \sum_{n=1}^{\infty}\frac{1^n+2^n+\cdots+50^n}{n^2}
  \left(\frac{1-x}{1+x}\right)^n
\]
的收敛区间.

\noindent\textbf{解}\quad （利用公式（B））. 由于
\[
  1\le\sqrt[n]{\left(\frac1{50}\right)^n+\left(\frac2{50}\right)^n+\cdots+
  \left(\frac{49}{50}\right)^n+1}\le\sqrt[n]{50}\to1
\]
及 $\sqrt[n]{n}\to1\ (n\to\infty)$. 因此
\[
  \lim_{n\to\infty}\sqrt[n]{\frac{1^n+2^n+\cdots+50^n}{n^2}}
  =\lim_{n\to\infty}\frac{50}{(\sqrt[n]{n})^2}
  \sqrt[n]{\left(\frac1{50}\right)^n+\left(\frac2{50}\right)^n+\cdots+1}=50.
\]
又因 $\displaystyle\sum_{n=1}^{\infty}\frac{1^n+2^n+\cdots+50^n}{n^2}\left(\frac1{50}\right)^n$ 收敛，故
$\displaystyle\sum_{n=1}^{\infty}\frac{1^n+2^n+\cdots+50^n}{n^2}t^n$ 在
$\left[-\dfrac1{50},\dfrac1{50}\right]$ 上收敛. 解不等式
$-\dfrac1{50}\le\dfrac{1-x}{1+x}\le\dfrac1{50}$，得原级数收敛区间为
$\left[\dfrac{49}{51},\dfrac{51}{49}\right]$.

\noindent\textbf{※例 5.3.4}\quad 设 $a_n\ge0$，$\displaystyle\sum_{n=1}^{\infty}a_n$ 收敛，
$b_m=\displaystyle\sum_{n=1}^{\infty}\left(1+\dfrac1{n^m}\right)^n a_n$. 试证：级数
$\displaystyle\sum_{m=1}^{\infty}b_mx^m$ 的收敛半径 $R$ 满足不等式：
$\dfrac1e\le R\le1$.

\noindent\textbf{证}\quad 因为
\[
  \sqrt[m]{\sum_{n=1}^{\infty}a_n}
  \le\sqrt[m]{\sum_{n=1}^{\infty}\left(1+\frac1{n^m}\right)^n a_n}
  =\sqrt[m]{b_m}
  \le\sqrt[m]{e^m}\sqrt[m]{\sum_{n=1}^{\infty}a_n}
  =e\sqrt[m]{\sum_{n=1}^{\infty}a_n},
\]
（注意到 $\displaystyle\lim_{m\to\infty}\sqrt[m]{\sum_{n=1}^{\infty}a_n}=1$）所以
$1\le\displaystyle\lim_{m\to\infty}\sqrt[m]{b_m}\le e$. 从而对于收敛半径 $R$，有
$\dfrac1e\le R\le1$（公式（A））.

不要以为讨论端点是件容易的事. 请看

\noindent\textbf{**例 5.3.5}\quad 求级数
\[
  \sum_{n=0}^{\infty}\frac{\left(1+2\cos\dfrac{n\pi}{4}\right)^n}{n\ln n}x^n
\]
的收敛范围.

\noindent\textbf{解}\quad
\[
\begin{aligned}
R^{-1}
&=\lim_{n\to\infty}\sqrt[n]{\frac{\left(1+2\cos\dfrac{n\pi}{4}\right)^n}{n\ln n}}
=\lim_{n\to\infty}\frac{1+2\cos\dfrac{n\pi}{4}}{\sqrt[n]{n\ln n}}\\
&=\lim_{k\to\infty}\frac3{\sqrt[8k]{8k\ln(8k)}}
=\lim_{n\to\infty}\frac3{\sqrt[n]{n\ln n}}.
\end{aligned}
\]
注意到
\[
  1\le\sqrt[n]{n\ln n}\le\sqrt[n]{n^2}=(\sqrt[n]{n})^2\to1\quad(\text{当 }n\to\infty\text{ 时}),
\]
所以 $R=\dfrac13$. 即在 $\left(-\dfrac13,\dfrac13\right)$ 内原级数收敛.（下面考虑区间端点.）

因为该级数之通项趋向零，因此将相邻 8 项逐次地括在一起，组成的级数与原级数同时敛散，且和值不变. 故原级数可视为 8 个级数之和，即
\[
\begin{aligned}
\text{原级数}
&=\sum_{k=1}^{\infty}\sum_{r=1}^{8}
\frac{\left[1+2\cos\dfrac{(8k+r)\pi}{4}\right]^{8k+r}}
{(8k+r)\ln(8k+r)}x^{8k+r}\\
&=\sum_{r=1}^{8}\sum_{k=1}^{\infty}
\frac{\left(1+2\cos\dfrac{8k+r}{4}\pi\right)^{8k+r}}
{(8k+r)\ln(8k+r)}x^{8k+r}.
\end{aligned}
\]
我们不难证明，不论 $x=\dfrac13$ 或 $x=-\dfrac13$，其中第 8 个级数（$r=8$）发散，其余 7 个级数都收敛，因而级数在端点处不收敛. 故收敛范围仍为
$\left(-\dfrac13,\dfrac13\right)$. 事实上，例如 $x=\dfrac13$，这时第 8 个级数为
\[
  \sum_{k=1}^{\infty}\frac1{(8k+8)\ln(8k+8)}
  =\sum_{k=2}^{\infty}\frac1{8k\ln(8k)},
\]
其通项与 $\displaystyle\sum\dfrac1{k\ln k}$ 的通项同阶，故此级数发散. 其余 7 个级数为
\[
  \sum_{k=1}^{\infty}\frac1{(8k+r)\ln(8k+r)}
  \left(\frac{1+2\cos\dfrac{8k+r}{4}\pi}{3}\right)^{8k+r}
  \quad(r=1,2,\cdots,7),
\qquad(1)
\]
利用 Dirichlet 判别法，易知它们都收敛.（因为
$\dfrac1{(8k+r)\ln(8k+r)}\downarrow0\ (k\to\infty)$，且
\[
  \left|\sum_{k=1}^{n}\left(\frac{1+2\cos\dfrac{8k+r}{4}\pi}{3}\right)^{8k+r}\right|
  \le\sum_{k=1}^{\infty}\left(\frac{1+\sqrt2}{3}\right)^k
  =\frac{1+\sqrt2}{2-\sqrt2}\quad(\forall n\in\mathbb N)
\]
（部分和有界），故式（1）中 7 个级数都收敛.）

对于 $x=-\dfrac13$ 的情况，类似可证.

\noindent\textbf{b.\quad 缺项幂级数的收敛范围}

\noindent\textbf{要点}\quad 对缺项幂级数（如下例），我们可通过补项，或利用上节一般函数项级数的方法处理.

\noindent\textbf{例 5.3.6}\quad 求级数 $\displaystyle\sum_{n=1}^{\infty}n^{n^2}x^{n^3}$ 的收敛范围.

\noindent\textbf{解 I}\quad 根据级数添加若干值为 0 的项，不影响级数的敛散性及和的值，我们可令
\[
  a_k=
  \begin{cases}
    n^{n^2},&\text{当 }k=n^3\text{ 时},\\
    0,&\text{当 }k\ne n^3\text{ 时}
  \end{cases}
  \quad(k=1,2,\cdots).
\]
如此
\[
  \sum_{n=1}^{\infty}n^{n^2}x^{n^3}
  =\sum_{k=1}^{\infty}a_kx^k,
\]
这时
\[
  R^{-1}=\varlimsup_{k\to\infty}\sqrt[k]{|a_k|}
  =\lim_{n\to\infty}\sqrt[n^3]{n^{n^2}}
  =\lim_{n\to\infty}\sqrt[n]{n}=1.
\]
又因 $x=\pm1$ 时，原级数之通项 $(\pm1)^{n^3}n^{n^2}\nrightarrow0\ (n\to\infty)$，故该级数收敛范围为
$(-1,1)$.

\noindent\textbf{解 II}\quad 将原级数看作函数项级数，用根式判别法. 因
\[
  \lim_{n\to\infty}\sqrt[n]{|n^{n^2}x^{n^3}|}=
  \begin{cases}
    +\infty,&\text{当 }|x|\ge1\text{ 时},\\
    0,&\text{当 }|x|<1\text{ 时},
  \end{cases}
\]
所以收敛范围为 $(-1,1)$. $|x|<1$ 时的极限可由如下不等式看出：
\[
  0\le\sqrt[n]{|n^{n^2}x^{n^3}|}
  =\left(\frac{n}{|1/x|^n}\right)^n
  \le\left(\frac12\right)^n\quad(\text{当 }n\text{ 充分大时成立}).
\]

\noindent\textbf{例 5.3.7}\quad 求 $\displaystyle\sum_{n=1}^{\infty}\frac{x^{n^2}}{2^n}$ 的收敛范围.

\noindent\textbf{解}\quad 将原级数看作函数项级数，
\[
  \lim_{n\to\infty}\sqrt[n]{|u_n|}
  =\lim_{n\to\infty}\sqrt[n]{\frac{|x^{n^2}|}{2^n}}
  =\lim_{n\to\infty}\frac{|x|^n}{2}
  =
  \begin{cases}
    0,&|x|<1,\\
    \dfrac12,&x=\pm1,\\
    +\infty,&|x|>1.
  \end{cases}
\]
按根式判别法的极限形式，可知原级数的收敛范围为 $[-1,1]$.

\noindent\textbf{c.\quad 利用收敛半径求极限}

\noindent\textbf{※例 5.3.8}\quad 设数列 $\{a_n\}$ 满足条件
$\displaystyle\varlimsup_{n\to\infty}\sqrt[n]{|a_n|}=1$，记其部分和为
$S_n=\displaystyle\sum_{k=0}^n a_k$，试证
\[
  \varlimsup_{n\to\infty}\sqrt[n]{|S_n|}=1.
\]

\noindent\textbf{分析}\quad 我们的问题等价于已知级数 $\displaystyle\sum_{n=0}^{\infty}a_nx^n$ 的收敛半径
$R=1$，求证级数
\[
  \sum_{n=0}^{\infty}S_nx^n
  =\sum_{n=0}^{\infty}(a_0+a_1+\cdots+a_n)x^n
\]
的收敛半径 $R_1=1$. 换句话说要证明 $R_1=R$. 我们看到

$1^\circ$ 根据级数乘法知，当 $|x|<1$ 时，
\[
  \sum_{n=0}^{\infty}x^n\cdot\sum_{n=0}^{\infty}a_nx^n
  =\sum_{n=0}^{\infty}(a_0+a_1+\cdots+a_n)x^n
  =\sum_{n=0}^{\infty}S_nx^n.
\qquad(1)
\]
可见 $\displaystyle\sum_{n=0}^{\infty}S_nx^n$ 的收敛半径 $R_1\ge1=R$.

$2^\circ$ 若 $\displaystyle\sum_{n=0}^{\infty}S_nx^n$ 收敛，则有
$(1-x)\displaystyle\sum_{n=0}^{\infty}S_nx^n=\displaystyle\sum_{n=0}^{\infty}a_nx^n$. 可见
$\displaystyle\sum_{n=0}^{\infty}a_nx^n$ 也必收敛，故 $R\ge R_1$. 总之，有 $R_1=R=1$.

\noindent\textbf{※例 5.3.9}\quad $C_n^k$ 表示 $n$ 个元素取 $k$ 个的组合数，试证：
\[
  \varlimsup_{n\to\infty}\sqrt[n]{|a_n|}=\frac{\sqrt5+1}{2},
\]
其中
\[
  a_n=
  \begin{cases}
    C_{3^m}^{k},&n=3^m+k\ (m=1,2,\cdots,\ k=0,1,2,\cdots,3^m),\\
    0,&\text{其他}.
  \end{cases}
\]

\noindent\textbf{证}\quad 问题等价于证明级数
$\displaystyle\sum_{n=1}^{\infty}a_nx^n$ 的收敛半径为 $\dfrac{\sqrt5-1}{2}$. 当 $x>0$ 时，各项为正，故任意加括号不影响收敛性.
将级数第 $3^m$ 项至 $2\cdot3^m$ 项括在一起组成新级数：
\[
\begin{aligned}
&\sum_{m=1}^{\infty}\left(C_{3^m}^0x^{3^m}+C_{3^m}^1x^{3^m+1}
\cdots+C_{3^m}^{3^m}x^{3^m+3^m}\right)\\
={}&\sum_{m=1}^{\infty}\left[x^{3^m}(1+x)^{3^m}\right]
=\sum_{m=1}^{\infty}[x(1+x)]^{3^m}.
\end{aligned}
\]
这是一函数项级数.
\[
  \lim_{m\to\infty}\sqrt[m]{|x(1+x)|^{3^m}}=
  \begin{cases}
    0,&|x(1+x)|<1,\\
    1,&|x(1+x)|=1,\\
    +\infty,&|x(1+x)|>1.
  \end{cases}
\]
注意不等式 $|x(1+x)|<1$ 等价于
$-\dfrac{\sqrt5+1}{2}<x<\dfrac{\sqrt5-1}{2}$. 由此知在 $x>0$ 的情况下：级数
$\displaystyle\sum_{n=1}^{\infty}a_nx^n$ 当 $x<\dfrac{\sqrt5-1}{2}$ 时收敛，当
$x>\dfrac{\sqrt5-1}{2}$ 时发散. 因而 $\displaystyle\sum_{n=1}^{\infty}a_nx^n$ 的收敛半径为
$\dfrac{\sqrt5-1}{2}$. 故
$\displaystyle\varlimsup_{n\to\infty}\sqrt[n]{|a_n|}=\dfrac{\sqrt5+1}{2}$.

\noindent\textbf{※例 5.3.10}\quad 设 $a_n\ge0$，$\displaystyle\sum_{n=1}^{\infty}a_n$ 发散，
$\displaystyle\lim_{n\to\infty}\dfrac{a_n}{a_1+a_2+\cdots+a_n}=0$. 试证：
$\displaystyle\varlimsup_{n\to\infty}\sqrt[n]{a_n}=1$.

\noindent\textbf{证 I}\quad 因 $\displaystyle\sum_{n=1}^{\infty}a_n$ 发散，所以
$\displaystyle\sum_{n=1}^{\infty}a_nx^n$ 的收敛半径 $R\le1$，
$\displaystyle\varlimsup_{n\to\infty}\sqrt[n]{a_n}\ge1$. 剩下只要证明相反的不等式.

记 $A_n=a_1+a_2+\cdots+a_n$，则由已知条件知，当 $n\to\infty$ 时，
\[
  1-\frac{A_{n-1}}{A_n}
  =\frac{A_n-A_{n-1}}{A_n}
  =\frac{a_n}{a_1+a_2+\cdots+a_n}\to0.
\]
因此，$\dfrac{A_{n-1}}{A_n}\to1$，进而
$\sqrt[n]{a_1+a_2+\cdots+a_n}=\sqrt[n]{A_n}\to1$.

又因 $\displaystyle\lim_{n\to\infty}\dfrac{a_n}{a_1+a_2+\cdots+a_n}=0$，所以 $\exists N>0$，当 $n>N$ 时，
$\dfrac{a_n}{a_1+a_2+\cdots+a_n}<1$，

\[
  a_n<a_1+a_2+\cdots+a_n.
\]
故
\[
  \varlimsup_{n\to\infty}\sqrt[n]{a_n}
  \le\lim_{n\to\infty}\sqrt[n]{a_1+\cdots+a_n}
  =\lim_{n\to\infty}\sqrt[n]{A_n}=1.
\]

\noindent\textbf{证 II}\quad 同上，易证 $\displaystyle\sum_{n=1}^{\infty}a_nx^n$ 的收敛半径
$R\le1$，$\displaystyle\lim_{n\to\infty}\frac{A_{n-1}}{A_n}=1$. 从而 $\displaystyle\sum_{n=1}^{\infty}A_nx^n=
\sum_{n=1}^{\infty}(a_1+a_2+\cdots+a_n)x^n$ 的收敛半径为 1. 当 $|x|<1$ 时，
\[
  (1-x)\sum_{n=1}^{\infty}(a_1+a_2+\cdots+a_n)x^n
  =\sum_{n=1}^{\infty}a_nx^n,
\]
因此 $R\ge1$. 总之 $R=1$. 证毕.

\subsection*{二、初等函数展开为幂级数}

\noindent\textbf{要点}\quad 将初等函数展开为幂级数，通常方法：

1）通过变形、转换、利用已知的展开式；

2）利用逐项积分或逐项微分法；

3）利用待定系数法；

4）计算指定点的各阶导数，然后利用 Taylor 级数；

5）利用级数的运算（加、减、乘、复合）.

\noindent\textbf{通过变形、变换，利用已知的展开式}
\[
  \sin x=\sum_{n=0}^{\infty}(-1)^n\frac{x^{2n+1}}{(2n+1)!},\quad\text{于 }\mathbb R\text{ 上},
\]
\[
  \cos x=\sum_{n=0}^{\infty}(-1)^n\frac{x^{2n}}{(2n)!},\quad\text{于 }\mathbb R\text{ 上},
\]
\[
  e^x=\sum_{n=0}^{\infty}\frac{x^n}{n!},\quad\text{于 }\mathbb R\text{ 上},
\]
\[
  \ln(1+x)=\sum_{n=1}^{\infty}(-1)^{n-1}\frac{x^n}{n},\quad\text{于 }(-1,1]\text{ 上},
\]
\[
  (1+x)^\alpha=1+\sum_{n=1}^{\infty}
  \frac{\alpha(\alpha-1)\cdots(\alpha-n+1)}{n!}x^n,\quad\text{于 }(-1,1)\text{ 内}.
\]

\noindent\(\star\)\textbf{例 5.3.11}\quad 把下列函数展开成 $x$ 的幂级数，并说明收敛范围：
\[
  1)\ f(x)=\frac1{(1+x)(1+x^2)(1+x^4)};\qquad
  2)\ \varphi(x)=\sin^3x.\quad(\text{武汉大学})
\]

\noindent\textbf{解}\quad 1）
\[
\begin{aligned}
f(x)
&=\frac{1-x}{(1-x)(1+x)(1+x^2)(1+x^4)}
=\frac{1-x}{1-x^8}\\
&=\sum_{n=0}^{\infty}x^{8n}-\sum_{n=0}^{\infty}x^{8n+1}
=1-x+x^8-x^9+\cdots+x^{8n}-x^{8n+1}+\cdots\quad(|x|<1).
\end{aligned}
\]
2）
\[
\begin{aligned}
\sin^3x
&=\frac34\sin x-\frac14\sin3x\\
&=\frac34\sum_{n=0}^{\infty}\frac{(-1)^nx^{2n+1}}{(2n+1)!}
-\frac14\sum_{n=0}^{\infty}\frac{(-1)^n(3x)^{2n+1}}{(2n+1)!}\\
&=\frac34\sum_{n=0}^{\infty}\frac{(-1)^n}{(2n+1)!}\left(1-3^{2n}\right)x^{2n+1},
\quad x\in(-\infty,+\infty).
\end{aligned}
\]

\noindent\textbf{例 5.3.12}\quad 设 $x>0$，求证：
\[
  \ln x=2\left[\frac{x-1}{x+1}+\frac13\left(\frac{x-1}{x+1}\right)^3
  +\frac15\left(\frac{x-1}{x+1}\right)^5+\cdots\right].
\]

\noindent\textbf{提示}\quad 变量替换法：令 $\dfrac{x-1}{x+1}=t$，即
$x=\dfrac{1+t}{1-t}$，从而
$\ln x=\ln\dfrac{1+t}{1-t}=\ln(1+t)-\ln(1-t)$，展开再回到原来变量 $x$.

\noindent\textbf{利用逐项积分或逐项微分法}

值得注意的是：逐项积分与逐项微分，常常只能在区间内部进行，但这并不等于说，所得的展开式一定不会在端点上成立. 如

\noindent\textbf{例 5.3.13}\quad 试求 $f(x)=\arctan\dfrac{2x}{2-x^2}$ 的幂级数展开式.

\noindent\textbf{解}\quad
\[
\begin{aligned}
f(x)
&=\int_0^x f'(t)\,dt
=\int_0^x\left(\arctan\frac{2t}{2-t^2}\right)'\,dt
=\int_0^x\left(1+\frac{t^2}{2}\right)\frac1{1+\left(\dfrac{t^2}{2}\right)^2}\,dt\\
&=\int_0^x\left(1+\frac{t^2}{2}\right)
\sum_{n=0}^{\infty}(-1)^n\left(\frac{t^4}{4}\right)^n\,dt
\quad(\text{此步要求 }|t|<\sqrt2)\\
&=\int_0^x\left[1+\frac{t^2}{2}-\left(\frac{t^2}{2}\right)^2
-\left(\frac{t^2}{2}\right)^3+\left(\frac{t^2}{2}\right)^4
+\left(\frac{t^2}{2}\right)^5-\cdots\right]dt\\
&=\int_0^x\sum_{n=0}^{\infty}(-1)^{\left[\frac n2\right]}
\left(\frac{t^2}{2}\right)^n\,dt\quad(\text{逐项积分})\\
&=\sum_{n=0}^{\infty}(-1)^{\left[\frac n2\right]}
\frac{x^{2n+1}}{2^n(2n+1)}\quad(|x|<\sqrt2).
\end{aligned}
\]

（讨论端点的情况.）当 $x=\sqrt2$ 时，级数
\[
\begin{aligned}
\sum_{n=0}^{\infty}(-1)^{\left[\frac n2\right]}
\frac{(\sqrt2)^{2n+1}}{2^n(2n+1)}
&=\sum_{n=0}^{\infty}(-1)^{\left[\frac n2\right]}\frac{\sqrt2}{2n+1}\\
&=\sqrt2\left(1+\frac13-\frac15-\frac17+\frac19+\frac1{11}-\cdots\right)\\
&=\sqrt2\left[\left(\frac11+\frac13\right)-\left(\frac15+\frac17\right)
+\left(\frac19+\frac1{11}\right)-\cdots\right]\\
&=\sqrt2\left(\sum_{n=0}^{\infty}(-1)^n\frac1{4n+1}
+\sum_{n=0}^{\infty}(-1)^n\frac1{4n+3}\right).
\end{aligned}
\]
可见当 $x=\sqrt2$ 时该级数收敛.

同理，当 $x=-\sqrt2$ 时级数也收敛.

又因 $f(x)=\arctan\dfrac{2x}{2-x^2}$ 在 $x=\pm\sqrt2$ 处连续，所以上面展开式在
$[-\sqrt2,\sqrt2]$ 上成立（利用 Abel 定理）.

\noindent\textbf{利用待定系数法}

\noindent\textbf{例 5.3.14}\quad 求 $\dfrac{x\sin\alpha}{1-2x\cos\alpha+x^2}$（$|x|<1$）的幂级数展开式.

\noindent\textbf{提示}\quad 用待定系数法. 设
\[
  \frac{x\sin\alpha}{1-2x\cos\alpha+x^2}
  =\sum_{n=0}^{\infty}a_nx^n,
\]
则
\[
\begin{aligned}
x\sin\alpha
&=(1-2x\cos\alpha+x^2)\sum_{n=0}^{\infty}a_nx^n\\
&=a_0+a_1x+a_2x^2+a_3x^3+\cdots
-(2a_0\cos\alpha)x-(2a_1\cos\alpha)x^2\\
&\quad-(2a_2\cos\alpha)x^3+\cdots+a_0x^2+a_1x^3+\cdots.
\end{aligned}
\]
比较等式两边同次幂的系数，得 $a_0=0,a_1=\sin\alpha,a_2=\sin2\alpha,\cdots,a_n=\sin n\alpha,\cdots$，这里用到三角恒等式
$\sin(n+1)\alpha=2\sin n\alpha\cdot\cos\alpha-\sin(n-1)\alpha\ (n=2,3,\cdots)$.

\noindent\textbf{利用 Taylor 级数}

\noindent\textbf{*例 5.3.15}\quad 求 $\dfrac{\ln(x+\sqrt{1+x^2})}{\sqrt{1+x^2}}$ 的幂级数展开式.

\noindent\textbf{解 I}\quad （计算 $\displaystyle\sum_{n=0}^{\infty}\frac{f^{(n)}(0)}{n!}x^n$.）设
$y=\dfrac{\ln(x+\sqrt{1+x^2})}{\sqrt{1+x^2}}$，因此
$y'=\dfrac1{1+x^2}\left[1-x\dfrac{\ln(x+\sqrt{1+x^2})}{\sqrt{1+x^2}}\right]$，即
\[
  (1+x^2)y'=1-xy. \qquad(1)
\]
由此两边同时求 $n$ 阶导数，得
\[
  (1+x^2)y^{(n+1)}+(2n+1)xy^{(n)}+n^2y^{(n-1)}=0. \qquad(2)
\]
令 $x=0$，得
\[
  y_0^{(n+1)}=-n^2y_0^{(n-1)}. \qquad(3)
\]
（下标“0”表示在 $x=0$ 处的值）. 在 (1) 中令 $x=0$，得 $y_0'=1$.

对 (1) 两边求导一次，得
\[
  2xy'+(1+x^2)y''=-y-xy'.
\]
令 $x=0$，知 $y_0''=-y_0=0$. 将 $y_0'=1,y_0''=0$ 代入递推公式 (3) 中，得
$y_0^{(2n)}=0\ (n=1,2,\cdots)$，
\[
  y_0^{(2n+1)}=(-1)^n[(2n)!!]^2y_0'=(-1)^n[(2n)!!]^2,
\]
故
\[
  \frac{\ln(x+\sqrt{1+x^2})}{\sqrt{1+x^2}}
  \sim\sum_{n=0}^{\infty}(-1)^n\frac{[(2n)!!]^2}{(2n+1)!}x^{2n+1}
  =\sum_{n=0}^{\infty}(-1)^n\frac{(2n)!!}{(2n+1)!!}x^{2n+1}. \qquad(4)
\]

容易证明右端的级数收敛半径 $R=1$. 利用逐项微分法，可验证该级数的和函数 $y=S(x)$ 是式 (1) 给定的微分方程的解，且 $S(0)=0$，而函数
$y=\dfrac{\ln(x+\sqrt{1+x^2})}{\sqrt{1+x^2}}$ 也是如此. 根据解的唯一性，可知式 (4) 中的“$\sim$”可改写成“$=$”（当 $|x|<1$ 时）. 又因 (4) 中级数当 $x=1$ 时收敛（根据 Leibniz 定理），而函数
$y=\dfrac{\ln(x+\sqrt{1+x^2})}{\sqrt{1+x^2}}$ 在 $x=1$ 处也连续，所以按 Abel 第二定理，式 (4) 中的“$\sim$”改为“$=$”对 $x=1$ 也成立. 同理，对 $x=-1$，等式也成立. 实际上，不难验证：$x=\pm1$ 确确实实都满足等式. 因此，式 (4) 成立的范围是 $[-1,1]$.

\noindent\textbf{解 II}\quad （待定系数法）令
$y=\displaystyle\sum_{n=0}^{\infty}a_nx^n$，代入 (1) 得
\[
  (1+x^2)\sum_{n=1}^{\infty}a_nnx^{n-1}
  =1-x\sum_{n=0}^{\infty}a_nx^n,
\]
即
\[
  a_1+2a_2x+\sum_{n=2}^{\infty}[a_{n+1}(n+1)+a_{n-1}(n-1)]x^n
  =1-a_0x-\sum_{n=2}^{\infty}a_{n-1}x^n.
\]
比较系数得
\[
  a_1=1,\ 2a_2=-a_0,\ 3a_3+a_1=-a_1,\ 4a_4+2a_2=-a_2,\cdots. \qquad(5)
\]
由于
\[
  y\big|_{x=0}
  =\left.\frac{\ln(x+\sqrt{1+x^2})}{\sqrt{1+x^2}}\right|_{x=0}=0,
\]
所以 $a_0=0$. 于是由递推关系 (5) 可得 $a_{2n}=0\ (n=1,2,\cdots)$，
\[
  a_1=1,\ a_3=-\frac23,\ a_5=\frac23\cdot\frac45,\cdots,\ 
  a_{2n+1}=(-1)^n\frac{(2n)!!}{(2n+1)!!},\cdots.
\]
所以
\[
  y=\sum_{n=0}^{\infty}a_nx^n
  =\sum_{n=0}^{\infty}(-1)^n\frac{(2n)!!}{(2n+1)!!}x^{2n+1}.
\]
为方程 (1) 适合条件 $y|_{x=0}=0$ 的唯一解. 故得
\[
  \frac{\ln(x+\sqrt{1+x^2})}{\sqrt{1+x^2}}
  =\sum_{n=0}^{\infty}(-1)^n\frac{(2n)!!}{(2n+1)!!}x^{2n+1}.
\]
重复解 I 中相应的内容，可知此式在 $[-1,1]$ 上成立.

\noindent\textbf{利用级数的运算}

\noindent\textbf{例 5.3.16}\quad 求 $f(x)=\ln^2(1-x)$ 的幂级数展开式.

\noindent\textbf{解 I}\quad $\displaystyle\ln(1-x)=-\sum_{n=1}^{\infty}\frac{x^n}{n}$ 在 $[-1,1)$ 上内闭一致收敛，故 $[-1,1)$ 上可用级数乘法：
\[
\begin{aligned}
f(x)
&=\left(-x-\frac{x^2}{2}-\frac{x^3}{3}-\cdots\right)^2\\
&=\sum_{n=1}^{\infty}\left(\frac11\cdot\frac1n+\frac12\cdot\frac1{n-1}
+\frac13\cdot\frac1{n-2}+\cdots+\frac1n\cdot\frac11\right)x^{n+1}\\
&=\sum_{n=1}^{\infty}\left[\sum_{k=1}^n\frac1{k(n+1-k)}\right]x^{n+1}\\
&=\sum_{n=1}^{\infty}\frac1{n+1}\left[\sum_{k=1}^n
\frac{k+[(n+1)-k]}{k(n+1-k)}\right]x^{n+1}\\
&=\sum_{n=1}^{\infty}\frac1{n+1}\left[\sum_{k=1}^n
\left(\frac1{n+1-k}+\frac1k\right)\right]x^{n+1}\\
&=2\sum_{n=1}^{\infty}\frac1{n+1}\left(\sum_{k=1}^n\frac1k\right)x^{n+1}\\
&=2\sum_{n=1}^{\infty}\left(1+\frac12+\frac13+\cdots+\frac1n\right)\cdot\frac{x^{n+1}}{n+1}.
\end{aligned}
\]
上面的展开式在 $[-1,1)$ 内成立.

\noindent\textbf{解 II}\quad 先求导数
$f'(x)=-\dfrac{2\ln(1-x)}{1-x}$ 的展开式.

\noindent\textbf{例 5.3.17}\quad 求 $f(x)=\dfrac1e(1+x)^{1/x}$ 按 $x$ 的幂的展开式至三次项.

\noindent\textbf{解}\quad （利用幂级数的复合.）
\[
\begin{aligned}
f(x)
&=\frac1e(1+x)^{1/x}
=e^{\frac1x\ln(1+x)-1}
=e^{-\frac x2+\frac{x^2}3-\frac{x^3}4+\cdots}\\
&=1+\left(-\frac x2+\frac{x^2}3-\frac{x^3}4+\cdots\right)
+\frac12\left(-\frac x2+\frac{x^2}3-\cdots\right)^2
+\frac16\left(-\frac x2+\cdots\right)^3\\
&=1-\frac12x+\frac{11}{24}x^2-\frac7{16}x^3+\cdots\quad(|x|<1).
\end{aligned}
\]

\subsection*{三、求和问题}

\noindent\textbf{a.\quad 利用逐项求导与逐项求积分}

\noindent\textbf{要点}\quad 利用逐项求导或逐项积分，将级数化为已知的展开式求和.

\noindent\textbf{**}\(\star\)\textbf{例 5.3.18}\quad 设 $|x|<1$，试求级数
$\displaystyle\sum_{n=0}^{\infty}\frac{(-1)^nx^{n+2}}{(n+1)(n+2)}$ 之和.（安徽大学，中国科学技术大学）

\noindent\textbf{解 I}\quad 原式
\[
\begin{aligned}
&=\sum_{n=0}^{\infty}\frac{(-1)^nx^{n+2}}{(n+1)(n+2)}
=\sum_{n=0}^{\infty}(-1)^n\left(\frac1{n+1}-\frac1{n+2}\right)x^{n+2}\\
&=\sum_{n=0}^{\infty}\frac{(-1)^nx^{n+2}}{n+1}
+\sum_{n=0}^{\infty}\frac{(-1)^{n+1}x^{n+2}}{n+2}\\
&=x\sum_{k=1}^{\infty}\frac{(-1)^{k-1}x^k}{k}
+\sum_{k=2}^{\infty}\frac{(-1)^{k-1}x^k}{k}\\
&=x\ln(1+x)+\ln(1+x)-x\quad(|x|<1).
\end{aligned}
\]

\noindent\textbf{解 II}\quad 令
$f(x)=\displaystyle\sum_{n=0}^{\infty}\dfrac{(-1)^nx^{n+2}}{(n+1)(n+2)}$，则
\[
  f(0)=f'(0)=0,\qquad f''(x)=\sum_{n=0}^{\infty}(-1)^nx^n=\frac1{1+x}.
\]
故
\[
  f'(x)=\int_0^x\frac{dt}{1+t}=\ln(1+x),
\]
从而
\[
  f(x)=\int_0^x\ln(1+t)\,dt=x\ln(1+x)+\ln(1+x)-x.
\]

\noindent\textbf{练习}\quad 级数的求和：
$\displaystyle\sum_{n=1}^{\infty}\frac{(-1)^{n-1}(n+2)}{n(n+1)}$.（南开大学）\hfill《$3\ln2-1$》

\noindent\textbf{提示}\quad （拆分法）$\displaystyle\frac{n+2}{n(n+1)}=\frac2n-\frac1{n+1}$，再用
$\ln(1+x)$ 的 Taylor 展式或用逐项积分、逐项微分的方法.

\noindent\textbf{再提示}\quad
\[
\begin{aligned}
\sum_{n=1}^{\infty}\frac{(-1)^{n-1}(n+2)}{n(n+1)}
&=2\sum_{n=1}^{\infty}\frac{(-1)^{n-1}}n
-\sum_{n=1}^{\infty}\frac{(-1)^{n-1}}{n+1}\\
&=2\ln2+\sum_{k=2}^{\infty}\frac{(-1)^{k-1}}k
=3\ln2-1.
\end{aligned}
\]
或令
\[
  f(x)=\sum_{n=1}^{\infty}\frac{(-1)^{n-1}(n+2)}{n(n+1)}x^{n+1},
\]
则
\[
  \int_0^x f(t)\,dt=x^2\ln(1+x)+x\ln(1+x)-x^2.
\]
从而
\[
  \text{原式}=\left[x^2\ln(1+x)+x\ln(1+x)-x^2\right]\bigg|_{x=1}=3\ln2-1.
\]

\noindent\textbf{例 5.3.19}\quad 证明对任一正整数 $k$，
$\displaystyle\sum_{n=1}^{\infty}\frac{n^k}{n!}$ 是 $e$ 的整数倍.（北京理工大学）

\noindent\textbf{分析}\quad 要证明 $\displaystyle\sum_{n=1}^{\infty}\frac{n^k}{n!}$ 是 $e$ 的整数倍，关键在于求级数的和，为此我们考虑对应的幂级数
$f(x)=\displaystyle\sum_{n=1}^{\infty}\frac{n^k}{n!}x^n$. 若能求出 $f(x)$，令 $x=1$ 就可得原级数的和（此幂级数收敛半径 $=+\infty$）. 求 $f(x)$ 的困难在于通项的分子里有因子 $n^k$（否则
$\displaystyle\sum_{n=1}^{\infty}\frac1{n!}x^n=e^x-1$）.

为了消去 $n^k$，将此级数先乘 $\dfrac1x$ 再逐项积分，这么做一次，分子就消了一个因子 $n$；反复做 $k$ 次，便可消去 $n^k$.

\noindent\textbf{证}\quad 记
\[
  f(x)=\sum_{n=1}^{\infty}\frac{n^k}{n!}x^n
  \qquad\left(R=\lim_{n\to\infty}\left|\frac{a_n}{a_{n+1}}\right|=+\infty\right),
\]
则
\[
\begin{aligned}
\int_0^x\frac1t f(t)\,dt
&=\int_0^x\left(\sum_{n=1}^{\infty}\frac{n^k}{n!}t^{n-1}\right)dt
=\sum_{n=1}^{\infty}\frac{n^k}{n!}\int_0^x t^{n-1}\,dt\\
&=\sum_{n=1}^{\infty}\frac{n^{k-1}}{n!}x^n.
\end{aligned}
\]
反复这样做 $k$ 次，得
\[
  \int_0^x\frac1{t_{k-1}}dt_{k-1}\int_0^{t_{k-1}}\frac1{t_{k-2}}dt_{k-2}
  \cdots\int_0^{t_1}\frac1t f(t)\,dt
  =\sum_{n=1}^{\infty}\frac1{n!}x^n=e^x-1.
\]
于是
\[
  f(x)=\underbrace{(\cdots(((e^x-1)'x)'x)'x\cdots)'x}_{k\text{ 次}}=p_k(x)e^x,
\]
其中 $p_k(x)$ 为 $x$ 的整数系数 $k$ 次多项式. 由此知
\[
  \sum_{n=1}^{\infty}\frac{n^k}{n!}=f(1)=p_k(1)e\quad(p_k(1)\text{ 为整数}),
\]
即为 $e$ 的整数倍.

（说明：若进行到第 $k-1$ 次时为 $p_{k-1}(x)e^x$，其中 $p_{k-1}$ 为整数系数 $k-1$ 次多项式，则第 $k$ 次应为
$(p_{k-1}(x)e^x)'x=(p'_{k-1}(x)x+p_{k-1}(x)x)e^x=p_k(x)e^x$，其中 $p_k(x)$ 为整数系数 $k$ 次多项式，故上式对一切 $k$ 成立.）

\noindent\textbf{b.\quad 方程式法}

\noindent\textbf{要点}\quad 设法证明级数的和满足某个方程式，然后求此方程的解.

\noindent\(\star\)\textbf{例 5.3.20}\quad 试求下列幂级数的和函数：
\[
  S(x)=1+x+\frac{x^2}{2}+\frac{x^3}{1\cdot3}+\frac{x^4}{2\cdot4}
  +\frac{x^5}{1\cdot3\cdot5}+\frac{x^6}{2\cdot4\cdot6}+\cdots.
\]
（广西大学）

\noindent\textbf{提示}\quad 收敛半径 $=+\infty$，逐项微分可知
\[
  S'(x)=1+xS(x),
\]
且 $S(0)=1$. 解此微分方程知
\[
  S(x)=e^{x^2/2}\left(\int_0^x e^{-t^2/2}\,dt+1\right).
\]

\noindent\(\star\)\textbf{例 5.3.21}\quad 证明：若函数 $f(x)$ 在 $[0,1]$ 上连续，令
\[
  f_0(x)=f(x),\quad f_{n+1}(x)=\int_x^1 f_n(y)\,dy
  \quad(x\in[0,1],\ n=0,1,2,\cdots),
\]
则 $\displaystyle\sum_{n=1}^{\infty}f_n(x)$ 在 $[0,1]$ 上一致收敛于
$\varphi(x)=\displaystyle\int_x^1 e^{y-x}f(y)\,dy$.（东北师范大学）

\noindent\textbf{证}\quad $1^\circ$ （先证明该级数一致收敛.）

因 $f(x)$ 在 $[0,1]$ 上连续，所以有界，即 $\exists M>0$，使 $|f(x)|\le M$ 于 $[0,1]$ 上，由此知
\[
  |f_1(x)|=\left|\int_x^1 f_0(y)\,dy\right|
  =\left|\int_x^1 f(y)\,dy\right|\le M(1-x),
\]
\[
  |f_2(x)|=\left|\int_x^1 f_1(y)\,dy\right|
  \le M\int_x^1(1-y)\,dy=M\frac{(1-x)^2}{2!},
\]
用数学归纳法易证
\[
  |f_n(x)|\le M\frac{(1-x)^n}{n!}\quad(\forall n=1,2,3,\cdots).
\]
但 $\displaystyle\sum M\dfrac{(1-x)^n}{n!}=Me^{1-x}$ 在全数轴上成立，在 $[0,1]$ 上一致收敛，所以
$\displaystyle\sum_{n=1}^{\infty}f_n(x)$ 在 $[0,1]$ 上绝对一致收敛.

$2^\circ$ （证明和满足微分方程.）记原级数之和为
\[
  \varphi(x)=\int_x^1 f(t)\,dt+\int_x^1dt_1\int_{t_1}^1 f(t_2)\,dt_2+\cdots. \qquad(1)
\]
此式两端同时加 $f(x)$，再同时在 $[x,1]$ 上取积分得
\[
  \int_x^1 f(t)\,dt+\int_x^1\varphi(t)\,dt=\varphi(x). \qquad(2)
\]
由此求导得
\[
  \varphi'(x)+\varphi(x)+f(x)=0. \qquad(3)
\]
从式 (2) 看出
\[
  \varphi(1)=0. \qquad(4)
\]
在条件 (4) 下求解微分方程 (3) 可得
$\displaystyle\varphi(x)=\int_x^1 f(y)e^{y-x}\,dy$.

未学过微分方程的读者可这样来求解：设 $\varphi(x)=u(x)e^{-x}$，则代入式 (3) 得
$u'(x)=-f(x)e^x$，所以
\[
  u(x)=-\int_0^x f(t)e^t\,dt+C. \qquad(5)
\]
根据 (4) 应有 $u(1)=0$，故知 $C=\displaystyle\int_0^1 f(t)e^t\,dt$. 代入 (5) 得
\[
  u(x)=-\int_0^x f(t)e^t\,dt+\int_0^1 f(t)e^t\,dt
  =\int_x^1 f(t)e^t\,dt=\int_x^1 f(y)e^y\,dy.
\]
因此
\[
  \varphi(x)=e^{-x}\int_x^1 f(y)e^y\,dy
  =\int_x^1 f(y)e^{y-x}\,dy.
\]

\noindent\textbf{c.\quad 利用 Abel 第二定理计算数项级数的和}

\noindent\textbf{要点}\quad 根据 Abel 第二定理，若要计算某收敛的数项级数
$\displaystyle\sum_{n=1}^{\infty}a_n$ 的和，我们只要设法求出幂级数
$\displaystyle\sum_{n=1}^{\infty}a_nx^n$ 在 $(-1,1)$ 内的和函数 $S(x)$，然后令 $x\to1^-$，取极限，则
\[
  \sum_{n=1}^{\infty}a_n=\lim_{x\to1^-}S(x).
\]
而和函数 $S(x)$ 可以通过逐项积分、逐项微分，或方程法等方法求解.

\noindent\(\star\)\textbf{例 5.3.22}\quad 求数项级数
$\displaystyle\frac12-\frac15+\frac18-\frac1{11}+\cdots$ 的和.

\noindent\textbf{解}\quad 级数
\[
  \frac12-\frac15+\frac18-\frac1{11}+\cdots
  =\sum_{k=1}^{\infty}\frac{(-1)^{k-1}}{3k-1},
\]
$\dfrac1{3k-1}\downarrow0$. 根据 Leibniz 定理，该级数收敛. 下面以 $S$ 表示其和. 易知幂级数
$\displaystyle\sum_{k=1}^{\infty}\frac{(-1)^{k-1}}{3k-1}x^{3k-1}$ 的收敛区间为 $(-1,1]$，于是由 Abel 第二定理，
$S=\displaystyle\lim_{x\to1-0}S(x)$. 注意 $S(0)=0$，因此
\[
  S(x)=S(x)-S(0)=\int_0^x S'(t)\,dt\quad(\text{当 }|x|<1\text{ 时}).
\]
利用逐项微分法，
\[
  S'(x)=\sum_{k=1}^{\infty}(-1)^{k-1}x^{3k-2}=\frac{x}{1+x^3}.
\]
因此
\[
\begin{aligned}
S(x)
&=\int_0^x S'(t)\,dt=\int_0^x\frac{t}{1+t^3}\,dt\\
&=-\frac13\ln(1+x)+\frac16\ln(1-x+x^2)
+\frac1{\sqrt3}\arctan\frac2{\sqrt3}\left(x-\frac12\right)
+\frac1{\sqrt3}\arctan\frac1{\sqrt3}.
\end{aligned}
\]
从而
\[
  S=\lim_{x\to1-0}S(x)=\ln\frac1{\sqrt2}+\frac{\pi}{3\sqrt3}
  =\frac{\sqrt3}{9}\pi-\frac13\ln2.
\]

\noindent\textbf{练习 1}\quad 求 $1-\dfrac14+\dfrac17-\dfrac1{10}+\cdots$.

\noindent\textbf{提示}\quad 设
$f(x)=\displaystyle\sum_{n=0}^{\infty}(-1)^n\dfrac{x^{3n+1}}{3n+1}$，则原式
$=\displaystyle\lim_{x\to1^-}f(x)$（因收敛半径为 1）.

\noindent\textbf{解}\quad 设 $f(x)=\displaystyle\sum_{n=0}^{\infty}(-1)^n\dfrac{x^{3n+1}}{3n+1}$，则
\[
  f'(x)=\sum_{n=0}^{\infty}(-1)^nx^{3n}
  =\frac1{1+x^3}
  =\frac13\left(\frac1{1+x}+\frac{2-x}{1-x+x^2}\right)\quad(x\in(-1,1)).
\]
\[
\begin{aligned}
f(x)
&=\frac13\int_0^x\frac1{1+t}\,dt+\frac13\int_0^x\frac{2-t}{1-t+t^2}\,dt\\
&=\frac13\ln(1+x)-\frac16\int_0^x\frac{-3+(-1+2t)}{1-t+t^2}\,dt\\
&=\frac13\ln(1+x)+\frac12\int_0^x\frac1{1-t+t^2}\,dt
-\frac16\int_0^x\frac{d(1-t+t^2)}{1-t+t^2}\\
&=\frac13\ln(1+x)+\frac1{\sqrt3}\arctan\frac{2t-1}{\sqrt3}\bigg|_0^x
-\frac16\ln(1-x+x^2).
\end{aligned}
\]
根据 Abel 第二定理，原式
$=\displaystyle\lim_{x\to1^-}f(x)=\dfrac13\ln2+\dfrac{\pi}{3\sqrt3}$.

\noindent\textbf{练习 2}\quad 已知级数 $\displaystyle\sum_{n=1}^{\infty}\frac{a_n}{n+1}$ 收敛，证明：
\[
  \int_0^1\sum_{n=1}^{\infty}a_nx^n\,dx=\sum_{n=1}^{\infty}\frac{a_n}{n+1}, \qquad(1)
\]
并计算 $1-\dfrac12+\dfrac13-\dfrac14+\cdots$.（华中科技大学）

\noindent\textbf{解}\quad 由 $\displaystyle\sum_{n=1}^{\infty}\frac{a_n}{n+1}$ 收敛可知，
$\dfrac{a_n}{n+1}\to0\ (n\to\infty)$，可得当 $n$ 充分大时，
$\left|\dfrac{a_n}{n+1}\right|\le1$，亦即
$0<\sqrt[n]{\left|\dfrac{a_n}{n+1}\right|}\le1$. 故
$\displaystyle\sum_{n=1}^{\infty}\frac{a_n}{n+1}x^n$ 的收敛半径 $R\ge1$.

在 $(-1,1)$ 内部，$\displaystyle\sum_{n=1}^{\infty}\frac{a_n}{n+1}x^n$ 收敛，记和为 $S(x)$：
\[
\begin{aligned}
S(x)
&=\int_0^x S'(t)\,dt
=\int_0^x\left(\sum_{n=1}^{\infty}\frac{a_n}{n+1}t^{n+1}\right)'\,dt\\
&=\int_0^x\left(\sum_{n=1}^{\infty}a_nt^n\right)dt\quad(\forall x\in(-1,1)).
\end{aligned}
\]
应用 Abel 第二定理，得
\[
  \sum_{n=1}^{\infty}\frac{a_n}{n+1}
  =\lim_{x\to1^-}S(x)
  =\lim_{x\to1^-}\int_0^x\left(\sum_{n=1}^{\infty}a_nt^n\right)dt
  =\int_0^1\left(\sum_{n=1}^{\infty}a_nt^n\right)dt.
\]
式 (1) 获证. 若令 $a_n=(-1)^n\ (\forall n\in\mathbb N)$，则得
\[
  \sum_{n=1}^{\infty}\frac{(-1)^n}{n+1}
  =\int_0^1\left[\sum_{n=1}^{\infty}(-t)^n\right]dt. \qquad(2)
\]
因此
\[
  1-\frac12+\frac13-\frac14+\cdots
  =1+\sum_{n=1}^{\infty}\frac{(-1)^n}{n+1}
  \stackrel{(2)}{=}1+\int_0^1\left[\sum_{n=1}^{\infty}(-t)^n\right]dt
  =1-\int_0^1\frac{t}{1+t}\,dt=\ln2.
\]

\noindent\textbf{*例 5.3.23}\quad 设 $\displaystyle\lim_{n\to\infty}a_n=a\in\mathbb R$.

1）证明：幂级数 $\displaystyle\sum_{n=1}^{\infty}a_nx^n$ 的收敛半径 $R\ge1$；

2）若 $f(x)\equiv\displaystyle\sum_{n=1}^{\infty}a_nx^n$，试证：
\[
  \lim_{x\to1^-}(1-x)f(x)=a, \qquad(1)
\]
\[
  \lim_{x\to1^-}(1-x)\int_0^x\frac{f(t)}{1-t}\,dt=a. \qquad(2)
\]
（中国科学技术大学）

\noindent\textbf{证}\quad 1）$\big||a_n|-|a|\big|\le|a_n-a|\to0$，因此
$\displaystyle\lim_{n\to\infty}|a_n|=|a|$. $\forall\varepsilon>0$，当 $n$ 充分大时，若 $a=0$，则
$0\le\sqrt[n]{|a_n|}\le\sqrt[n]{\varepsilon}$，从而
$\displaystyle\varlimsup_{n\to\infty}\sqrt[n]{|a_n|}\le1$；若 $|a|>0$，对任意给定的
$\varepsilon:0<\varepsilon<|a|$，则
$\sqrt[n]{|a|-\varepsilon}\le\sqrt[n]{|a_n|}\le\sqrt[n]{|a|+\varepsilon}$. 从而
$\displaystyle\lim_{n\to\infty}\sqrt[n]{|a_n|}=1$. 故收敛半径
$R=\dfrac1{\displaystyle\lim_{n\to\infty}\sqrt[n]{|a_n|}}\ge1$.

2）因
\[
  (1-x)f(x)=\sum_{n=1}^{\infty}a_nx^n-\sum_{n=1}^{\infty}a_nx^{n+1}
  \overset{a_0=0}{=}\sum_{n=1}^{\infty}(a_n-a_{n-1})x^n,
\]
由 1）知 $R\ge1$. 于是
\[
\begin{aligned}
\lim_{x\to1^-}(1-x)f(x)
&=\lim_{x\to1^-}\sum_{n=1}^{\infty}(a_n-a_{n-1})x^n\\
&=\sum_{k=1}^{\infty}(a_k-a_{k-1})\quad(R>1\text{ 时用连续性，}R=1\text{ 时用 Abel 定理})
\end{aligned}
\]
\[
  =\lim_{n\to\infty}\sum_{k=1}^n(a_k-a_{k-1})=\lim_{n\to\infty}a_n=a.
\]
式 (1) 获证.

幂级数在收敛区间内，绝对收敛，可以相乘，能逐项积分，所以，$\forall x:0\le x<1$，
\[
\begin{aligned}
F(x)&\overset{\text{记}}{=}\int_0^x\frac1{1-t}\cdot f(t)\,dt
=\int_0^x\left(\sum_{i=1}^{\infty}t^{i-1}\cdot\sum_{n=1}^{\infty}a_nt^n\right)dt\\
&=\int_0^x\sum_{n=1}^{\infty}(a_1+a_2+\cdots+a_n)t^n\,dt\\
&=\sum_{n=1}^{\infty}\frac{a_1+a_2+\cdots+a_n}{n+1}x^{n+1}
\overset{\text{记}}{=}\sum_{n=1}^{\infty}A_nx^{n+1},
\end{aligned}
\]
其中系数
\[
  A_n=\frac{a_1+a_2+\cdots+a_n}{n+1}
  =\frac{a_1+a_2+\cdots+a_n}{n}\cdot\frac{n}{n+1}
  \to a\cdot1=a.
\]
说明 $\displaystyle F(x)=\sum_{n=1}^{\infty}A_nx^{n+1}$ 符合式 (1) 的条件，应用式 (1)，即得式 (2). 证毕.

\noindent\textbf{*例 5.3.24}\quad 求级数
$\displaystyle 1+\sum_{n=1}^{\infty}(-1)^n\frac{(2n-1)!!}{(2n)!!}$ 的值.

\noindent\textbf{解}\quad （考虑 $\displaystyle 1+\sum_{n=1}^{\infty}\frac{(2n-1)!!}{(2n)!!}x^n$ 在 $x=-1$ 的情况.）

$1^\circ$ （原级数收敛.）因 $\dfrac{(2n-1)!!}{(2n)!!}\downarrow$，又
\[
  \sqrt{1\cdot3}<\frac{1+3}{2}=2,\quad
  \sqrt{3\cdot5}<\frac{3+5}{2}=4,\quad\cdots,\quad
  \sqrt{(2n-1)(2n+1)}<2n,
\]
\[
  0<\frac{(2n-1)!!}{(2n)!!}
  =\frac{\sqrt{1\cdot3}\sqrt{3\cdot5}\cdots\sqrt{(2n-1)(2n+1)}}{2\cdot4\cdots2n}\cdot\frac1{\sqrt{2n+1}}
  \le\frac1{\sqrt{2n+1}}\to0\quad(n\to\infty),
\]
根据 Leibniz 定理，原级数收敛.

$2^\circ$ 幂级数 $\displaystyle 1+\sum_{n=1}^{\infty}\frac{(2n-1)!!}{(2n)!!}x^n$ 的收敛半径为 1，这是因为
\[
  R=\lim_{n\to\infty}\frac{a_{n-1}}{a_n}
  =\lim_{n\to\infty}\frac{2n}{2n-1}=1.
\]
由此知该幂级数在 $(-1,1)$ 内收敛.

$3^\circ$ （计算幂级数的和
$\displaystyle f(x)=1+\sum_{n=1}^{\infty}\frac{(2n-1)!!}{(2n)!!}x^n$.）因为
\[
  f(x)=1+\frac12x+\frac{1\cdot3}{2\cdot4}x^2
  +\frac{1\cdot3\cdot5}{2\cdot4\cdot6}x^3+\cdots
  +\frac{1\cdot3\cdot5\cdots(2n-1)}{2\cdot4\cdot6\cdots(2n)}x^n+\cdots,
\]
所以
\[
  f'(x)=\frac12+\frac{1\cdot3}{2\cdot4}\cdot2x
  +\frac{1\cdot3\cdot5}{2\cdot4\cdot6}\cdot3x^2+\cdots
\]
\[
  +\frac{1\cdot3\cdot5\cdots(2n-1)}{2\cdot4\cdot6\cdots(2n)}\cdot nx^{n-1}+\cdots,
\]
\[
  2f'(x)=1+\frac12\cdot3x+\frac{1\cdot3}{2\cdot4}\cdot5x^2+\cdots
  +\frac{1\cdot3\cdot5\cdots(2n-3)}{2\cdot4\cdot6\cdots(2n-2)}(2n-1)x^{n-1}+\cdots,
\]
\[
\begin{aligned}
2xf'(x)&=x+\frac12\cdot3x^2+\cdots
+\frac{1\cdot3\cdot5\cdots(2n-3)}{2\cdot4\cdot6\cdots(2n-2)}(2n-1)x^n+\cdots\\
&=\frac12\cdot2x+\frac{1\cdot3}{2\cdot4}\cdot4x^2+\cdots
+\frac{1\cdot3\cdot5\cdots(2n-3)(2n-1)}{2\cdot4\cdot6\cdots(2n-2)(2n)}\cdot2nx^n+\cdots.
\end{aligned}
\]
于是
\[
  2f'(x)-2xf'(x)=1+\frac12x+\frac{1\cdot3}{2\cdot4}x^2+\cdots
  +\frac{1\cdot3\cdots(2n-1)}{2\cdot4\cdots(2n)}x^n+\cdots=f(x).
\]
由此
\[
  \frac{f'(x)}{f(x)}=\frac1{2(1-x)}.
\]
积分得
\[
  \ln f(x)=\ln\frac1{\sqrt{1-x}}.
\]
因此
\[
  f(x)=\frac1{\sqrt{1-x}}.
\]
即
\[
  1+\sum_{n=1}^{\infty}\frac{(2n-1)!!}{(2n)!!}x^n=\frac1{\sqrt{1-x}}\quad(|x|<1). \qquad(1)
\]

$4^\circ$ 令 $x\to-1$ 取极限，利用 Abel 第二定理，知
\[
  1+\sum_{n=1}^{\infty}(-1)^n\frac{(2n-1)!!}{(2n)!!}
  =\lim_{x\to-1^+}\left[1+\sum_{n=1}^{\infty}\frac{(2n-1)!!}{(2n)!!}x^n\right]
  =\lim_{x\to-1^+}\frac1{\sqrt{1-x}}=\frac{\sqrt2}{2}.
\]

\noindent\textbf{**例 5.3.25}\quad 计算极限
$\displaystyle\lim_{m,n\to\infty}\sum_{i=1}^m\sum_{j=1}^n\frac{(-1)^{i+j}}{i+j}$.（吉林大学）

\noindent\textbf{分析}\quad 为了使和式化简，必须消去分母里的 $i+j$，因此各项分别乘以 $x^{i+j}$，再逐项求导.

\noindent\textbf{解}\quad
\[
\begin{aligned}
S_{m,n}
&\equiv\sum_{i=1}^m\sum_{j=1}^n\frac{(-1)^{i+j}}{i+j}
=\sum_{i=1}^m\sum_{j=1}^n(-1)^{i+j}\frac{x^{i+j}}{i+j}\bigg|_{x=0}^{x=1}\\
&=\sum_{i=1}^m\sum_{j=1}^n(-1)^{i+j}\int_0^1 x^{i+j-1}\,dx
=\int_0^1\left[\sum_{i=1}^m\sum_{j=1}^n(-1)^{i+j}x^{i+j-1}\right]dx\\
&=\int_0^1\sum_{i=1}^m(-x)^i\cdot\sum_{j=1}^n(-1)^jx^{j-1}\,dx\\
&=\int_0^1\frac{-x-(-x)^{m+1}}{1+x}\cdot\frac{-1-(-1)^{n+1}x^n}{1+x}\,dx\\
&=\int_0^1\frac1{(1+x)^2}\left[x+(-1)^{m+1}x^{m+1}
+(-1)^{n+1}x^{n+1}+(-1)^{m+n}x^{m+n+1}\right]dx.
\end{aligned}
\]
注意到
\[
  \lim_{k\to\infty}\int_0^1\frac{x^k}{(1+x)^2}\,dx=0
  \quad\left(0\le\int_0^1\frac{x^k}{(1+x)^2}\,dx\le\int_0^1x^k\,dx=\frac1{k+1}\to0\right),
\]
所以
\[
\begin{aligned}
\lim_{m,n\to\infty}S_{m,n}
&=\int_0^1\frac{x}{(1+x)^2}\,dx
=\int_0^1\frac1{1+x}\,dx-\int_0^1\frac{dx}{(1+x)^2}\\
&=\ln(1+x)\bigg|_0^1+\frac1{1+x}\bigg|_0^1
=\ln2-\frac12.
\end{aligned}
\]

\noindent\textbf{例 5.3.26}\quad 求极限
$\displaystyle\lim_{x\to0}\sum_{n=1}^{\infty}\frac{(-1)^n}{x^2+\dfrac{2^n}{n(n+1)}}$.

\noindent\textbf{提示}\quad
\[
\begin{aligned}
\lim_{x\to0}\sum_{n=1}^{\infty}\frac{(-1)^n}{x^2+\dfrac{2^n}{n(n+1)}}
&\overset{\text{因一致收敛}}{=}
\sum_{n=1}^{\infty}\lim_{x\to0}\frac{(-1)^n}{x^2+\dfrac{2^n}{n(n+1)}}\\
&=\sum_{n=1}^{\infty}(-1)^n\frac{n(n+1)}{2^n}.
\end{aligned}
\]
（关于一致收敛性，可参看例 5.2.8 的证法.）用逐项积分、逐项微分法求级数
$\displaystyle\sum_{n=1}^{\infty}(-1)^n\frac{n(n+1)}{2^n}x^{n-1}$ 的和.

\subsection*{四、幂级数的应用}

\noindent\textbf{a.\quad 计算积分}

\noindent\textbf{要点}\quad 将被积函数展开为幂级数，然后逐项积分.

\noindent\textbf{例 5.3.27}\quad 证明：
\[
  \frac2\pi\int_0^{\pi/2}\frac{d\theta}{\sqrt{1-k^2\sin^2\theta}}
  =1+\frac14k^2+\cdots+
  \left[\frac{(2n-1)!!}{(2n)!!}\right]^2k^{2n}+\cdots\quad(|k|<1).
\]

\noindent\textbf{提示}\quad 利用展开式
\[
\begin{aligned}
(1+x)^\alpha
&=\sum_{n=0}^{\infty}C_\alpha^nx^n\\
&=1+\alpha x+\frac{\alpha(\alpha-1)}{1\cdot2}x^2+\cdots
+\frac{\alpha(\alpha-1)\cdots(\alpha-n+1)}{1\cdot2\cdots n}x^n+\cdots\quad(|x|<1).
\end{aligned}
\]
当 $|k|<1$ 时，
\[
  (1-k^2\sin^2\theta)^{-1/2}
  =\sum_{n=0}^{\infty}\frac{(2n-1)!!}{(2n)!!}k^{2n}\sin^{2n}\theta,
\]
关于 $\theta$ 在 $\left[0,\dfrac\pi2\right]$ 上一致收敛. 再逐项积分，并利用 Wallis 公式
\[
  \int_0^{\pi/2}\sin^{2n}\theta\,d\theta
  =\frac{(2n-1)!!}{(2n)!!}\frac\pi2
\]
即得.

\noindent\(\star\)\textbf{例 5.3.28}\quad 计算积分
$\displaystyle\int_0^1\frac{\ln x}{1-x^2}\,dx$.

\noindent\textbf{解 I}\quad
\[
  \int_0^1\frac{\ln x}{1-x^2}\,dx
  =\int_0^1\frac{1-x^2+x^2}{1-x^2}\ln x\,dx
  =\int_0^1\ln x\,dx+\int_0^1\frac{x^2}{1-x^2}\ln x\,dx.
\]
因 $\displaystyle\int_0^1\ln x\,dx=-1$ 及
$\dfrac{x^2}{1-x^2}\ln x=\displaystyle\sum_{n=1}^{\infty}x^{2n}\ln x$，故
\[
  \text{上式}=-1+\int_0^1\sum_{n=1}^{\infty}x^{2n}\ln x\,dx.
\]
重复例 5.2.52 的证明，可知级数 $\displaystyle\sum_{n=1}^{\infty}x^{2n}\ln x$ 虽然在 $[0,1]$ 上不一致收敛，但仍可以在 $[0,1]$ 上逐项积分，因此
\[
\begin{aligned}
\text{上式}
&=-1+\sum_{n=1}^{\infty}\int_0^1x^{2n}\ln x\,dx
=-1-\sum_{n=1}^{\infty}\frac1{(2n+1)^2}
=-\sum_{n=0}^{\infty}\frac1{(2n+1)^2}\\
&=-\left[\sum_{n=0}^{\infty}\frac1{(2n+1)^2}
+\sum_{n=1}^{\infty}\frac1{(2n)^2}\right]
+\frac1{2^2}\sum_{n=1}^{\infty}\frac1{n^2}\\
&=-\sum_{n=1}^{\infty}\frac1{n^2}
+\frac1{2^2}\sum_{n=1}^{\infty}\frac1{n^2}
=-\frac{\pi^2}{6}+\frac{\pi^2}{24}
=-\frac{\pi^2}{8}.
\end{aligned}
\]

\noindent\textbf{解 II}\quad （利用内闭一致收敛，逐项积分，再取极限.）
$\forall\alpha,\beta:0<\alpha<\dfrac12<\beta<1$，因级数
$\displaystyle\sum_{n=0}^{\infty}x^{2n}=\dfrac1{1-x^2}$ 在 $[\alpha,\beta]$ 上一致收敛，可知级数
$\displaystyle\sum_{n=0}^{\infty}x^{2n}\ln x=\dfrac{\ln x}{1-x^2}$ 亦在 $[\alpha,\beta]$ 上一致收敛. 因此
\[
  \int_\alpha^\beta\frac{\ln x}{1-x^2}\,dx
  =\sum_{n=0}^{\infty}\int_\alpha^\beta x^{2n}\ln x\,dx. \qquad(1)
\]
但
\[
  \left|\int_\alpha^\beta x^{2n}\ln x\,dx\right|
  \le\left|\int_0^1x^{2n}\ln x\,dx\right|
  =\frac1{(2n+1)^2}.
\]
所以级数 (1) 关于 $\alpha,\beta\in(0,1)$ 一致收敛. 故
\[
\begin{aligned}
\int_0^1\frac{\ln x}{1-x^2}\,dx
&=\lim_{\substack{\alpha\to0^+\\ \beta\to1^-}}\int_\alpha^\beta\frac{\ln x}{1-x^2}\,dx
=\lim_{\substack{\alpha\to0^+\\ \beta\to1^-}}\sum_{n=0}^{\infty}\int_\alpha^\beta x^{2n}\ln x\,dx\\
&=\sum_{n=0}^{\infty}\lim_{\substack{\alpha\to0^+\\ \beta\to1^-}}\int_\alpha^\beta x^{2n}\ln x\,dx
=\sum_{n=0}^{\infty}\int_0^1x^{2n}\ln x\,dx\\
&=\sum_{n=0}^{\infty}\frac{-1}{(2n+1)^2}
=-\frac{\pi^2}{8}.
\end{aligned}
\]

\noindent\(\star\)\textbf{例 5.3.29}\quad 若
$f(x)=\displaystyle\sum_{n=0}^{+\infty}a_nx^n\ (a_n>0,\ n=0,1,2,\cdots)$ 的收敛半径为 $+\infty$，且
$\displaystyle\sum_{n=0}^{\infty}a_nn!$ 收敛，则
$\displaystyle\int_0^{+\infty}e^{-x}f(x)\,dx$ 也收敛，且
$\displaystyle\int_0^{+\infty}e^{-x}f(x)\,dx=\sum_{n=0}^{+\infty}a_nn!$.（东北师范大学）

\noindent\textbf{证}\quad
\[
\begin{aligned}
\int_0^{+\infty}e^{-x}f(x)\,dx
&=\int_0^{+\infty}\left(e^{-x}\sum_{n=0}^{+\infty}a_nx^n\right)dx
=\lim_{A\to+\infty}\int_0^A\left[\sum_{n=0}^{+\infty}(a_nx^ne^{-x})\right]dx\\
&=\lim_{A\to+\infty}\sum_{n=0}^{+\infty}a_n\left(\int_0^A x^ne^{-x}\,dx\right)\qquad(1)\\
&=\sum_{n=0}^{+\infty}a_n\lim_{A\to+\infty}\int_0^A x^ne^{-x}\,dx\qquad(2)\\
&=\sum_{n=0}^{+\infty}a_n\int_0^{+\infty}x^ne^{-x}\,dx
=\sum_{n=0}^{+\infty}a_n\cdot n!.\qquad(3)
\end{aligned}
\]
这里等式 (1) 成立是因为
\[
  |a_nx^ne^{-x}|=a_nx^n\left(1+x+\cdots+\frac{x^n}{n!}+\cdots\right)^{-1}
  <\frac{a_nx^n}{x^n/n!}=a_nn!\quad(\forall x\ge0),
\]
且 $\displaystyle\sum_{n=0}^{\infty}a_nn!$ 收敛，故
$\displaystyle\sum_{n=0}^{\infty}a_nx^ne^{-x}$ 在 $[0,A]$ 上一致收敛，可逐项积分. 等式 (2) 成立是因为
\[
  \left|a_n\int_0^A x^ne^{-x}\,dx\right|
  =a_n\int_0^A x^ne^{-x}\,dx
  \le a_n\int_0^{+\infty}x^ne^{-x}\,dx=a_nn!\quad(\forall A\ge0),
\]
且已知 $\displaystyle\sum_{n=0}^{\infty}a_nn!$ 收敛，因此
$\displaystyle\sum_{n=0}^{\infty}a_n\int_0^A x^ne^{-x}\,dx$ 关于 $A$ 在 $[0,+\infty)$ 上一致收敛，故可逐项求极限.

至于 $\displaystyle\int_0^{+\infty}x^ne^{-x}\,dx=n!$，可反复使用分部积分法 $n$ 次得到，或利用 Euler 积分：
\[
  \int_0^{+\infty}x^ne^{-x}\,dx=\Gamma(n+1)=n!.
\]

\noindent\textbf{*例 5.3.30}\quad 试利用已知的公式
\[
  \int_{-1}^1\frac{dx}{(\alpha-x)\sqrt{1-x^2}}=\frac{\pi}{\sqrt{\alpha^2-1}}
  \quad(\text{当 }\alpha>1\text{ 时}) \qquad(1)
\]
证明：
\[
  \int_{-1}^1\frac{x^{2n}}{\sqrt{1-x^2}}\,dx
  =\frac{(2n-1)!!}{(2n)!!}\pi\quad(n=1,2,\cdots).
\]

\noindent\textbf{方法}\quad 将式 (1) 两端同时按 $\alpha^{-1}$ 的幂次展开，然后比较同次幂的系数.

\noindent\textbf{解}\quad $1^\circ$ 式 (1) 左端
\[
\begin{aligned}
\int_{-1}^1\frac{dx}{(\alpha-x)\sqrt{1-x^2}}
&=\frac1\alpha\int_{-1}^1\frac{dx}{\left(1-\dfrac{x}{\alpha}\right)\sqrt{1-x^2}}\\
&=\frac1\alpha\int_{-1}^1\left[\sum_{n=0}^{\infty}\left(\frac{x}{\alpha}\right)^n\cdot\frac1{\sqrt{1-x^2}}\right]dx\\
&=\frac1\alpha\sum_{n=0}^{\infty}\frac1{\alpha^n}\int_{-1}^1\frac{x^n}{\sqrt{1-x^2}}\,dx\qquad(2)\\
&=\sum_{n=0}^{\infty}\frac1{\alpha^{2n+1}}\int_{-1}^1\frac{x^{2n}}{\sqrt{1-x^2}}\,dx\quad(\text{奇次项积分为零}).
\end{aligned}
\]

$2^\circ$ 式 (1) 右端
\[
  \frac{\pi}{\sqrt{\alpha^2-1}}
  =\frac{\pi}{\alpha}\left(1-\frac1{\alpha^2}\right)^{-1/2}
  =\frac{\pi}{\alpha}\sum_{n=0}^{\infty}C_{-1/2}^n\left(-\frac1{\alpha^2}\right)^n.
\]
\[
  =\sum_{n=0}^{\infty}\left[\frac{(2n-1)!!}{(2n)!!}\right]\frac{\pi}{\alpha^{2n+1}}.
\]
总之，
\[
  \sum_{n=0}^{\infty}\frac1{\alpha^{2n+1}}\int_{-1}^1\frac{x^{2n}}{\sqrt{1-x^2}}\,dx
  =\sum_{n=0}^{\infty}\frac1{\alpha^{2n+1}}\left[\frac{(2n-1)!!}{(2n)!!}\pi\right].
\]
比较 $\alpha$ 同次幂的系数，便得欲证的等式.

现补充证明等式 (2) 的合理性. 因级数
$\displaystyle\sum_{n=0}^{\infty}\left(\dfrac{x}{\alpha}\right)^n\dfrac1{\sqrt{1-x^2}}$
关于 $x$ 在 $(-1,1)$ 内闭一致收敛，所以 $\forall[a,b]\subset(-1,1)$，在 $[a,b]$ 上可逐项取积分：
\[
  \int_a^b\sum_{n=0}^{\infty}\frac1{\alpha^n}\frac{x^n}{\sqrt{1-x^2}}\,dx
  =\sum_{n=0}^{\infty}\frac1{\alpha^n}\int_a^b\frac{x^n}{\sqrt{1-x^2}}\,dx.
\]
又 $\displaystyle\sum_{n=0}^{\infty}\frac1{\alpha^n}$ 收敛（自然关于 $a,b\in[-1,1]$ 一致），
\[
  \left|\int_a^b\frac{x^n}{\sqrt{1-x^2}}\,dx\right|
  \le\int_{-1}^1\frac{|x|^n}{\sqrt{1-x^2}}\,dx
  \le\int_{-1}^1\frac{dx}{\sqrt{1-x^2}}=\pi
  \quad(a,b\in[-1,1],\ n\in\mathbb N),
\]
即对于 $a,b\in[-1,1]$，$n\in\mathbb N$ 一致有界. 故由 Abel 判别法，级数
$\displaystyle\sum_{n=0}^{\infty}\frac1{\alpha^n}\int_a^b\dfrac{x^n}{\sqrt{1-x^2}}\,dx$
关于 $a,b\in[-1,1]$ 一致收敛，从而可逐项取极限：
\[
\begin{aligned}
\int_{-1}^1\left(\sum_{n=0}^{\infty}\frac1{\alpha^n}\frac{x^n}{\sqrt{1-x^2}}\right)dx
&=\lim_{\substack{a\to(-1)^+\\ b\to1^-}}\int_a^b\left(\sum_{n=0}^{\infty}\frac1{\alpha^n}\frac{x^n}{\sqrt{1-x^2}}\right)dx\\
&=\lim_{\substack{a\to(-1)^+\\ b\to1^-}}\sum_{n=0}^{\infty}\frac1{\alpha^n}\int_a^b\frac{x^n}{\sqrt{1-x^2}}\,dx\\
&=\sum_{n=0}^{\infty}\frac1{\alpha^n}\lim_{\substack{a\to(-1)^+\\ b\to1^-}}\int_a^b\frac{x^n}{\sqrt{1-x^2}}\,dx\\
&=\sum_{n=0}^{\infty}\frac1{\alpha^n}\int_{-1}^1\frac{x^n}{\sqrt{1-x^2}}\,dx.
\end{aligned}
\]
证毕.

\noindent\textbf{b.\quad 证明不等式}

幂级数是表达函数的重要工具，因此也可应用于证明不等式. 如

\noindent\(\star\)\textbf{例 5.3.31}\quad 证明不等式
$e^x+e^{-x}\le2e^{x^2/2},\ x\in(-\infty,+\infty)$.（中山大学）

\noindent\textbf{证}\quad 因
\[
  e^x+e^{-x}=2\operatorname{ch}x=2\sum_{n=0}^{\infty}\frac{x^{2n}}{(2n)!},\qquad
  2e^{x^2/2}=2\sum_{n=0}^{\infty}\frac{x^{2n}}{(2n)!!},
\]
而 $\dfrac{x^{2n}}{(2n)!}\le\dfrac{x^{2n}}{(2n)!!}$，故
$e^x+e^{-x}\le2e^{x^2/2}$.

\noindent\textbf{c.\quad 近似计算}

幂级数常常用于近似计算. 如

\noindent\(\star\)\textbf{例 5.3.32}\quad 求 $\pi$ 的近似值，计算到小数点后第三位（误差不超过 $10^{-3}$）.（安徽大学）

\noindent\textbf{提示}\quad
\[
  \frac\pi6=\arctan x\bigg|_{x=1/\sqrt3}
  =\sum_{n=1}^{\infty}(-1)^{n-1}\frac{x^{2n-1}}{2n-1}\bigg|_{x=1/\sqrt3}.
\]
利用 Leibniz 级数的余和估计：
\[
  |r_n|\le|a_{n+1}|=\frac1{3^{n+1/2}}\cdot\frac1{2n+1}.
\]
要 $|r_n|<\dfrac1{1000}$，只要
\[
  \frac1{3^n}\frac1{2n+1}<\frac1{1000},
\]
由此可知应取的项数 $n\ge5$.

\subsection*{五、综合性问题}

\noindent\(\star\)\textbf{例 5.3.33}\quad 设
$S(x)=\displaystyle\sum_{n=1}^{\infty}e^{-n}\sin n^2x\ (-\infty<x<+\infty)$. 试证：

1）$S(x)$ 在全数轴 $\mathbb R$ 上可任意次微分；

2）$S(x)$ 的 Maclaurin 级数当 $x\ne0$ 时发散.（武汉大学）

\noindent\textbf{证}\quad $\forall k\ge0,\ \forall x\in\mathbb R$，有
\[
  \left|e^{-n}n^{2k}\sin\left(n^2x+k\frac\pi2\right)\right|\le e^{-n}n^{2k},
\]
$\displaystyle\sum_{n=1}^{\infty}e^{-n}n^{2k}$ 收敛. 因此
$\displaystyle\sum_{n=1}^{\infty}e^{-n}n^{2k}\sin\left(n^2x+k\frac\pi2\right)$ 在 $\mathbb R$ 上一致收敛，由此不难用数学归纳法证明结论 1），且
\[
  S^{(k)}(x)=\sum_{n=1}^{\infty}e^{-n}n^{2k}\sin\left(n^2x+k\frac\pi2\right)
  \quad(x\in\mathbb R). \qquad(1)
\]
下面只就结论 2）进行证明. 根据式 (1)，
\[
  S^{(0)}(0)=S(0)=0,\qquad S^{(2k)}(0)=0\quad(k=1,2,\cdots),
\]
\[
  S^{(2k+1)}(0)=(-1)^k\sum_{n=1}^{\infty}e^{-n}n^{2(2k+1)}
  \quad(k=0,1,2,\cdots).
\]
故 $S(x)$ 的 Maclaurin 级数为
\[
  S(x)=\sum_{k=0}^{\infty}\frac{S^{(k)}(0)}{k!}x^k
  =\sum_{k=0}^{\infty}\frac1{(2k+1)!}
  \left[(-1)^k\sum_{n=1}^{\infty}e^{-n}n^{2(2k+1)}\right]x^{2k+1}.
\]
要证明 $x\ne0$ 时发散，若能证明通项 $\nrightarrow0$ 即可.

事实上，$\forall x\ne0$，其通项 $u_k(x)$：
\[
\begin{aligned}
|u_k(x)|
&=\frac1{(2k+1)!}\sum_{n=1}^{\infty}e^{-n}n^{2(2k+1)}|x|^{2k+1}\\
&\ge\frac1{(2k+1)!}e^{-(2k+1)}(2k+1)^{2(2k+1)}|x|^{2k+1}\\
&\ge e^{-(2k+1)}(2k+1)^{2k+1}|x|^{2k+1}.
\end{aligned}
\]
当 $k>\left|\dfrac{e}{2x}\right|$ 时，有 $|x|>\dfrac e{2k}$，因此
\[
  |u_k(x)|\ge e^{-(2k+1)}(2k+1)^{2k+1}\left(\frac e{2k}\right)^{2k+1}\ge1.
\]
证毕.

\noindent\textbf{*例 5.3.34}\quad 设幂级数
$f(x)=a_0+a_1x+\cdots+a_nx^n+\cdots$ 在 $x=1$ 处收敛，试求极限
\[
  \lim_{\nu\to\infty}\frac{\nu^\nu}{2^\nu\nu!}
  \left[f^{(n)}(0)-f^{(n)}\left(\frac{2^\nu\cdot\nu!}{\nu^\nu}\right)\right],
\]
其中 $\nu$ 为自然数.（吉林大学）

\noindent\textbf{证}\quad （利用 Stirling 公式：
$n!=\sqrt{2\pi n}\,n^ne^{-n+\theta_n/(12n)}\ (0<\theta_n<1)$.）
\[
  \frac{2^\nu\cdot\nu!}{\nu^\nu}
  =\frac{2^\nu}{\nu^\nu}\sqrt{2\pi\nu}\,\nu^\nu e^{-\nu+\theta_\nu/(12\nu)}
  =\frac{\sqrt{2\pi\nu}e^{\theta_\nu/(12\nu)}}{\left(e/2\right)^\nu}\to0
  \quad(\nu\to\infty).
\]
所以
\[
\begin{aligned}
\lim_{\nu\to\infty}\frac{\nu^\nu}{2^\nu\cdot\nu!}
\left[f^{(n)}(0)-f^{(n)}\left(\frac{2^\nu\cdot\nu!}{\nu^\nu}\right)\right]
&=\lim_{h\to0}\frac{f^{(n)}(0)-f^{(n)}(h)}h\\
&=-f^{(n+1)}(0)=-a_{n+1}(n+1)!.
\end{aligned}
\]
因为该幂级数在 $(-1,1)$ 内收敛，逐项微分 $n+1$ 次后，令 $x=0$，可知
$f^{(n+1)}(0)=a_{n+1}(n+1)!$.

\noindent\(\star\)\textbf{例 5.3.35}\quad 设 $f(x)$ 在 $(-\infty,+\infty)$ 上无穷次可微并且满足：

1）存在 $M>0$，使得 $|f^{(k)}(x)|\le M\ (\forall x\in(-\infty,+\infty),\ k=0,1,2,\cdots)$；

2）$f\left(\dfrac1{2^n}\right)=0\ (n=1,2,\cdots)$.

证明：在 $(-\infty,+\infty)$ 上 $f(x)\equiv0$.（广西大学）

\noindent\textbf{证}\quad $1^\circ$ 因为各阶导数一致有界，所以 $(-\infty,+\infty)$ 内处处有
\[
  f(x)=\sum_{n=0}^{\infty}\frac{f^{(n)}(0)}{n!}x^n. \qquad(1)
\]
（因
\[
  |R_n(x)|=\left|\frac{f^{(n+1)}(\xi)}{(n+1)!}x^{n+1}\right|
  \le\frac{M}{(n+1)!}|x|^{n+1}\to0\quad(n\to\infty).
\]）

$2^\circ$ 由 $f\left(\dfrac1{2^n}\right)=0\ (n=1,2,\cdots)$ 知
$f(0)=\displaystyle\lim_{n\to\infty}f\left(\dfrac1{2^n}\right)=0$，
\[
  f'(0)=\lim_{n\to\infty}\frac{f\left(\dfrac1{2^n}\right)-f(0)}{\dfrac1{2^n}}=0.
\]
根据 Rolle 定理，$\exists\xi_i^{(1)}:\dfrac12>\xi_1^{(1)}>\dfrac1{2^2}>\xi_2^{(1)}>\dfrac1{2^3}>\cdots>\dfrac1{2^n}>\xi_n^{(1)}>\dfrac1{2^{n+1}}>\cdots$，使得
$f'(\xi_i^{(1)})=0$，从而 $\xi_n^{(1)}\to0\ (n\to\infty)$，
\[
  f''(0)=\lim_{n\to\infty}\frac{f'(\xi_n^{(1)})-f'(0)}{\xi_n^{(1)}}=0.
\]
类似可证，若 $f^{(k)}(x)$ 在
$\xi_1^{(k)}>\xi_2^{(k)}>\cdots>\xi_n^{(k)}>\cdots\ (\xi_n^{(k)}\to0)$ 处有
$f^{(k)}(\xi_n^{(k)})=0$，且 $f^{(k)}(0)=0$，便可推出 $f^{(k+1)}(x)$ 亦然. 于是
$n=0,1,2,\cdots$ 时恒有 $f^{(n)}(0)=0$，从而
\[
  f(x)=\sum_{n=0}^{\infty}\frac{f^{(n)}(0)}{n!}x^n\equiv0
  \quad(-\infty<x<+\infty).
\]

\noindent\textbf{注}\quad 条件 1）（各阶导数一致有界）不可缺少，否则命题可以不成立. 如
\[
  f(x)=
  \begin{cases}
    e^{-x^{-2}}\sin\dfrac\pi x,&x\ne0,\\
    0,&x=0.
  \end{cases}
\]

\noindent\(\star\)\textbf{例 5.3.36}\quad 设
$f(x)=\displaystyle\sum_{n=1}^{\infty}\dfrac{x^n}{n^2}\ (0\le x\le1)$. 求证：当 $0<x<1$ 时，有
\[
  f(x)+f(1-x)+\ln x\ln(1-x)=\frac{\pi^2}{6}. \qquad(1)
\]
（北京航空航天大学）

\noindent\textbf{证}\quad $f(x)=\displaystyle\sum_{n=1}^{\infty}\dfrac{x^n}{n^2}$ 的收敛半径
$R=\displaystyle\lim_{n\to\infty}\sqrt[n]{n^2}=1$. $x=1$ 时，
$f(1)=\displaystyle\sum_{n=1}^{\infty}\dfrac1{n^2}=\dfrac{\pi^2}{6}$. 级数在 $(0,1)$ 内可逐项微分，$f(x)$ 有连续导数. 因此，
\[
\begin{aligned}
[f(x)+f(1-x)+\ln x\ln(1-x)]'
&=f'(x)-f'(1-x)+\frac{\ln(1-x)}x+\frac{\ln x}{x-1}\\
&=\sum_{n=1}^{\infty}\frac{x^{n-1}}n-\sum_{n=1}^{\infty}\frac{(1-x)^{n-1}}n
-\sum_{n=1}^{\infty}\frac{x^{n-1}}n+\sum_{n=1}^{\infty}\frac{(-1)^{n-1}(x-1)^{n-1}}n=0.
\end{aligned}
\]
于是
\[
  f(x)+f(1-x)+\ln x\ln(1-x)=C,\quad x\in(0,1).
\]
令 $x\to0^+$，取极限知
$C=f(1)=\displaystyle\sum_{n=1}^{\infty}\dfrac1{n^2}=\dfrac{\pi^2}{6}$，等式得证.

\noindent\(\star\)\textbf{例 5.3.37}\quad 证明 Tauber 定理：设在 $-1<x<1$ 上，有
$f(x)=\displaystyle\sum_{n=0}^{\infty}a_nx^n$，$\displaystyle\lim_{n\to\infty}na_n=0$. 若
$\displaystyle\lim_{x\to1-0}f(x)=S$，则 $\displaystyle\sum_{n=0}^{\infty}a_n$ 收敛且其和为 $S$.（吉林大学）

\noindent\textbf{分析}\quad 用加一项、减一项的办法，得
\[
\begin{aligned}
\left|\sum_{k=0}^n a_k-S\right|
&=\left|\sum_{k=0}^n a_k-\sum_{k=0}^n a_kx^k-\sum_{k=n+1}^{\infty}a_kx^k
+\sum_{k=0}^{\infty}a_kx^k-S\right|\\
&\le\left|\sum_{k=0}^n a_k(1-x^k)\right|
+\left|\sum_{k=n+1}^{\infty}a_kx^k\right|
+\left|\sum_{k=0}^{\infty}a_kx^k-S\right|.\qquad(1)
\end{aligned}
\]
因此问题在于证明：$n$ 充分大，$x$ 充分接近 1 时右端三项都能任意小.

\noindent\textbf{证}\quad 由 $\displaystyle\lim_{n\to\infty}na_n=0$，知
$\displaystyle\lim_{n\to\infty}n|a_n|=0$，从而
$\displaystyle\lim_{n\to\infty}\dfrac{\sum_{k=0}^n k|a_k|}{n}=0$. 又因
$\displaystyle\lim_{x\to1-0}f(x)=S$，所以
$\left|f\left(1-\dfrac1n\right)-S\right|\to0$（当 $n\to\infty$ 时）. 故
$\forall\varepsilon>0,\ \exists N>0$，使得 $n>N$ 时，有
\[
  0\le\frac{\sum_{k=0}^n k|a_k|}{n}<\frac\varepsilon3,\qquad
  n|a_n|<\frac\varepsilon3,\qquad
  \left|f\left(1-\frac1n\right)-S\right|<\frac\varepsilon3.
\]
利用式 (1)，得
\[
  \left|\sum_{k=0}^n a_k-S\right|
  \le\left|\sum_{k=1}^n a_k(1-x^k)\right|
  +\left|\sum_{k=n+1}^{\infty}a_kx^k\right|
  +\left|\sum_{k=0}^{\infty}a_kx^k-S\right|.
\]
取 $x=1-\dfrac1n$，不等式右边第一项
\[
\begin{aligned}
\left|\sum_{k=0}^n a_k(1-x^k)\right|
&=\left|\sum_{k=1}^n a_k(1-x)(1+x+x^2+\cdots+x^{k-1})\right|\\
&\le\sum_{k=1}^n |a_k|(1-x)k
=\frac{\sum_{k=1}^n k|a_k|}{n}<\frac\varepsilon3;
\end{aligned}
\]
右边第二项
\[
  \left|\sum_{k=n+1}^{\infty}a_kx^k\right|
  \le\frac1n\sum_{k=n+1}^{\infty}k|a_k|x^k
  <\frac\varepsilon{3n}\sum_{k=n+1}^{\infty}x^k
  \le\frac\varepsilon{3n}\cdot\frac1{1-x}
  =\frac\varepsilon3;
\]
右边第三项
\[
  \left|\sum_{k=0}^{\infty}a_kx^k-S\right|
  =\left|f\left(1-\frac1n\right)-S\right|<\frac\varepsilon3.
\]
故 $\left|\displaystyle\sum_{k=0}^n a_k-S\right|<\dfrac\varepsilon3+\dfrac\varepsilon3+\dfrac\varepsilon3=\varepsilon$. 证毕.

\noindent\textbf{*例 5.3.38}\quad 设
$f_n(x)=1+x+\dfrac{x^2}{2!}+\cdots+\dfrac{x^n}{n!}$，其中 $n$ 为自然数. 试证：方程
$f_n(x)\cdot f_{n+1}(x)=0$ 在实数域内有唯一实根.（四川师范大学）

\noindent\textbf{分析}\quad 讨论 $f_n(x)=0$ 的根的情况.

当 $n=0$ 时，$f_0(x)=1$，无实根；

当 $n=1$ 时，$f_1(x)=1+x$，有唯一实根；

当 $n=2$ 时，$f_2(x)=1+x+\dfrac{x^2}{2}=\dfrac12[(x+1)^2+1]>0$，无实根；

当 $n=3$ 时，$f_3(x)=1+x+\dfrac{x^2}{2}+\dfrac{x^3}{6}$，三次方程至少有一实根，而
$f_3'(x)=f_2(x)>0$，$f_3(x)$ 严增，故 $f_3(x)=0$ 只有唯一实根. 不妨记此实根为 $x_3$，因 $f_3$ 严增，可知
\[
  \text{当 }n=4\text{ 时},\quad
  f_4'(x)=f_3(x)
  \begin{cases}
    =0,&x=x_3,\\
    <0,&x<x_3,\\
    >0,&x>x_3,
  \end{cases}
\]
因此 $f_4(x)$ 在 $x=x_3$ 处取极小值也是最小值，故 $\forall x$，
\[
  f_4(x)\ge\min f_4(x)=f_4(x_3)=f_3(x_3)+\frac{x_3^4}{4!}=0+\frac{x_3^4}{4!}>0.
\]
因此 $f_4(x)$ 无实根.

类似地，可写出一般的推理过程：$f_{2n}$ 无实根 $\Rightarrow f_{2n+1}$ 仅有唯一实根
$\Rightarrow f_{2n+2}$ 无实根（留作练习）.

于是，在 $\mathbb R$ 上，$f_n(x)=0$ 无实根（当 $n$ 为偶数时），仅有唯一实根（当 $n$ 为奇数时）. 总之，
$f_n(x)\cdot f_{n+1}(x)=0$ 在 $\mathbb R$ 上仅有唯一实根.

\noindent\textbf{*练习}\quad 设 $a_k>0,\ \displaystyle\sum_{k=0}^{\infty}a_k=1,\ f(x)=\displaystyle\sum_{k=0}^{\infty}a_kx^k-x$. 试证：

1）$f(x)$ 在 $[0,1]$ 上为严格凸函数（见例 3.4.8 前的定理 5）；

2）若 $f'(1)$ 存在，$f(x)$ 在 $[0,1)$ 上有零点的充分必要条件是 $f'(1)>0$.（中南大学）

\noindent\textbf{证}\quad 1）$\displaystyle\sum_{k=0}^{\infty}a_k=1$，说明 $\displaystyle\sum_{k=0}^{\infty}a_kx^k$ 的收敛半径 $R\ge1$.

在 $[0,1)$ 上，$f''(x)=\displaystyle\sum_{k=2}^{\infty}a_kk(k-1)x^{k-2}>0$，可知在 $[0,1)$ 上 $f'(x)$ 严增，且 $f(x)$ 严格凸（§3.4 定理 5）.

2）（必要性）（反证法）假设 $f(x)$ 在 $[0,1)$ 上有零点，但 $f'(1)\le0$. 因为 1° 中已证 $f'(x)$ 严增，那么 $f'(1)\le0$ 意味着：在 $[0,1)$ 上恒有 $f'(x)<0$. 因此，在 $[0,1)$ 上 $f(x)$ 严减. 而
\[
  f(1)=\left(\sum_{k=0}^{\infty}a_kx^k-x\right)_{x=1}
  =\sum_{k=0}^{\infty}a_k-1=0,
\]
说明：在 $[0,1)$ 上 $f(x)$ 严减，$f(1)=0$. 可见在 $[0,1)$ 上恒有 $f(x)>0$，故 $f(x)$ 在 $[0,1)$ 上无零点，矛盾. 因此，$f(x)$ 在 $[0,1)$ 上若有零点，则 $f'(1)>0$.

（充分性）若 $f'(1)>0$，根据导数定义，在 $x=1$ 的左邻域应有 $f(x)<f(1)=0$，故 $\exists\varepsilon>0$，使得 $f(1-\varepsilon)<0$. 但另一方面 $f(0)=a_0>0$，根据连续的介值性，在 $(0,1)$ 里，必有 $f(x)$ 的零点. 证毕.

\noindent\textbf{*例 5.3.39}\quad 设
\[
  p_n(x)=1+x+\frac{x^2}{2!}+\cdots+\frac{x^n}{n!},
\]
$x_m$ 是 $p_{2m+1}(x)=0$ 的实根，求证：$x_m<0$，且 $\displaystyle\lim_{m\to+\infty}x_m=-\infty$.

\noindent\textbf{提示}\quad 1° $\forall m\in\mathbb N,\ x\ge0$ 时，$p_{2m+1}(x)>0$；$x<0$ 且 $|x|$ 充分大时，$p_{2m+1}(x)<0$. 所以 $p_{2m+1}(x)=0$ 的根 $x_m$ 存在（因介值性）. 又 $p'_{2m+1}(x)=p_{2m}(x)>0$，$p_{2m+1}(x)$ 严增，所以根唯一（见上题），$x_m<0$.

2° $\forall x\in(-\infty,0),\ p_n(x)\to e^x>0$（当 $n\to+\infty$ 时），所以 $p_{2m+1}(x)=0$ 的根 $x_m\to-\infty$（当 $m\to+\infty$ 时）. 因为若 $m\to+\infty$ 时 $p_{2m+1}(x)=0$ 的根 $x_m\not\to-\infty$，则 $\exists A>0$，使得 $(-A,0)$ 中含有 $\{x_m\}$ 的一个无穷子列\footnote{任意数列 $\{a_n\}$：$a_n\to-\infty$（$n\to+\infty$ 时）$\Leftrightarrow\forall A>0,\ a_n\notin(-\infty,-A)$ 最多只有有限项.}. 从而存在收敛子列 $x_{m_k}\to x_0$（$x_0$ 为某有限数，$x_0\ge-A$），
\[
  0<e^{-A}=\lim_{k\to+\infty}p_{2m_k+1}(-A)
  \le \lim_{k\to+\infty}p_{2m_k+1}(x_{m_k})=0,
\]
矛盾.

\noindent\textbf{*例 5.3.40}\quad 设 $f(x)$ 是仅有正实根的多项式，$-\dfrac{f'(x)}{f(x)}=\displaystyle\sum_{n=0}^{\infty}c_nx^n$，试证 $\displaystyle\lim_{n\to\infty}\frac1{\sqrt[n]{c_n}}$，$\displaystyle\lim_{n\to\infty}\frac{c_n}{c_{n+1}}$ 存在且都等于 $f(x)$ 的最小根.（北京大学）

\noindent\textbf{分析}\quad 要证明 $\dfrac1{\sqrt[n]{c_n}}$ 与 $\dfrac{c_n}{c_{n+1}}$ 都以 $f$ 的最小根为极限，必须设法求出 $c_n$ 与 $f$ 根的关系.

\noindent\textbf{证}\quad 因为 $f$ 是仅有正根的多项式，所以可设全部的正根为
\[
  0<a_1<a_2<\cdots<a_k. \qquad(1)
\]
于是
\[
  f(x)=A(x-a_1)^{r_1}(x-a_2)^{r_2}\cdots(x-a_k)^{r_k}
\]
（其中 $r_i$ 表示根 $a_i$ 的重数，$i=1,2,\cdots,k$）. 因此
\[
\begin{aligned}
-\frac{f'(x)}{f(x)}
&=-\left(\frac{r_1}{x-a_1}+\cdots+\frac{r_k}{x-a_k}\right)\\
&=\frac{r_1}{a_1}\frac1{1-\dfrac{x}{a_1}}+\cdots+
\frac{r_k}{a_k}\frac1{1-\dfrac{x}{a_k}}\\
&=\frac{r_1}{a_1}\sum_{n=0}^{\infty}\left(\frac{x}{a_1}\right)^n+\cdots+
\frac{r_k}{a_k}\sum_{n=0}^{\infty}\left(\frac{x}{a_k}\right)^n\\
&=\sum_{n=0}^{\infty}\left(\frac{r_1}{a_1^{n+1}}+\cdots+\frac{r_k}{a_k^{n+1}}\right)x^n
=\sum_{n=0}^{\infty}c_nx^n.
\end{aligned}
\]
比较系数得 $c_n=\dfrac{r_1}{a_1^{n+1}}+\cdots+\dfrac{r_k}{a_k^{n+1}}$（展开式的唯一性），所以
\[
  \lim_{n\to\infty}\frac{c_n}{c_{n+1}}
  =\lim_{n\to\infty}\frac{c_{n-1}}{c_n}
  =\lim_{n\to\infty}
  \frac{\dfrac{r_1}{a_1^n}+\cdots+\dfrac{r_k}{a_k^n}}
       {\dfrac{r_1}{a_1^{n+1}}+\cdots+\dfrac{r_k}{a_k^{n+1}}}.
\]
分子分母同乘 $a_1^n$，并注意式 (1)，得
\[
  \lim_{n\to\infty}\frac{c_n}{c_{n+1}}
  =\frac{r_1}{r_1/a_1}=a_1\quad(\text{最小根}).
\]
故
\[
  \lim_{n\to\infty}\frac1{\sqrt[n]{c_n}}
  \ \text{存在，}\quad
  \lim_{n\to\infty}\frac1{\sqrt[n]{c_n}}
  =\lim_{n\to\infty}\frac{c_n}{c_{n+1}}=a_1.
\]

\noindent\textbf{注}\quad 此例说明：仅具有正实根的多项式，其对数导数的幂级数展开式的收敛半径等于其最小根.

（拟合法）

\noindent\textbf{※例 5.3.41}\quad 设 $S_n=\displaystyle\sum_{k=0}^n a_k$，$\sigma_n=\dfrac{S_0+S_1+S_2+\cdots+S_{n-1}}{n}$. 试证：

1）若 $\{\sigma_n\}$ 收敛，则 $a_n=o(n)$（当 $n\to\infty$ 时）；

2）若 $\{\sigma_n\}$ 收敛，则 $f(x)\equiv\displaystyle\sum_{n=0}^{\infty}a_nx^n$ 在 $(-1,1)$ 内绝对收敛，且
\[
  f(x)=(1-x)^2\sum_{n=0}^{\infty}(n+1)\sigma_{n+1}x^n;
\]

3）若 $\displaystyle\lim_{n\to\infty}\sigma_n=S$，则 $\displaystyle\lim_{x\to1-0}f(x)=S$.

\noindent\textbf{证}\quad 1）因为
\[
  \sigma_n=\frac{S_0+S_1+S_2+\cdots+S_{n-1}}{n},
\]
所以 $S_n=(n+1)\sigma_{n+1}-n\sigma_n$，
\[
  \frac{S_n}{n}=\frac{n+1}{n}\sigma_{n+1}-\sigma_n\to0\quad(n\to\infty).
\]
故
\[
  \frac{a_n}{n}=\frac{S_n-S_{n-1}}{n}
  =\frac{S_n}{n}-\frac{n-1}{n}\frac{S_{n-1}}{n-1}\to0\quad(n\to\infty).
\]

2）因为 $\dfrac{a_n}{n}\to0(n\to\infty)$，所以当 $n$ 充分大时，有 $|a_n|<n$，$|a_nx^n|<n|x|^n$，而 $\displaystyle\sum_{n=1}^{\infty}nx^n$ 的收敛半径为 1，故知级数 $f(x)\equiv\displaystyle\sum_{n=0}^{\infty}a_nx^n$ 在 $(-1,1)$ 内绝对收敛. 如此，当 $x\in(-1,1)$ 时，利用级数乘法得
\[
  \frac{f(x)}{1-x}
  =\sum_{n=0}^{\infty}(a_0+a_1+\cdots+a_n)x^n
  =\sum_{n=0}^{\infty}S_nx^n,
\]
\[
  \frac{f(x)}{(1-x)^2}
  =\sum_{n=0}^{\infty}(S_0+S_1+\cdots+S_n)x^n
  =\sum_{n=0}^{\infty}(n+1)\sigma_{n+1}x^n,
\]
即
\[
  f(x)=(1-x)^2\sum_{n=0}^{\infty}(n+1)\sigma_{n+1}x^n. \qquad(1)
\]

3）下面利用拟合法，来证明 $\displaystyle\lim_{x\to1-0}f(x)=S$. 由
\[
  \frac1{1-x}=\sum_{n=0}^{\infty}x^n,\quad
  \frac1{(1-x)^2}=\sum_{n=0}^{\infty}(n+1)x^n,
\]
可得 1 的分解式：
\[
  1=(1-x)^2\sum_{n=0}^{\infty}(n+1)x^n.
\]
两端同乘 $S$，因此 $S$ 可写成与 (1) 类似的形式（此即拟合法的思想）：
\[
  S=(1-x)^2\sum_{n=0}^{\infty}(n+1)Sx^n.
\]
于是
\[
  f(x)-S=(1-x)^2\sum_{n=0}^{\infty}(n+1)(\sigma_{n+1}-S)x^n. \qquad(2)
\]
已知 $\displaystyle\lim_{n\to\infty}\sigma_n=S$，所以 $\forall\varepsilon>0,\ \exists N>0,\ n>N$ 时，$|\sigma_{n+1}-S|<\dfrac\varepsilon2$. 故式 (2) 中 $n>N$ 各项的和（当 $0<x<1$ 时）
\[
  \left|(1-x)^2\sum_{n>N}(n+1)(\sigma_{n+1}-S)x^n\right|
  \le\frac\varepsilon2(1-x)^2\sum_{n>N}(n+1)x^n
  \le\frac\varepsilon2(1-x)^2\sum_{n=0}^{\infty}(n+1)x^n
  =\frac\varepsilon2.
\]
(2) 中前 $N+1$ 项里：$\displaystyle\sum_{n=0}^N(n+1)(\sigma_{n+1}-S)x^n$ 当 $x\to1^-$ 时保持有界，另一因子 $(1-x)^2$ 为无穷小量. 因此 $\exists\delta>0$，当 $0<1-x<\delta$ 时，有
\[
  \left|(1-x)^2\sum_{n=0}^N(n+1)(\sigma_{n+1}-S)x^n\right|<\frac\varepsilon2.
\]
总之
\[
  |f(x)-S|<\frac\varepsilon2+\frac\varepsilon2=\varepsilon.
\]
即
\[
  S=\lim_{x\to1^-}f(x).
\]

\noindent\textbf{\(\star\)例 5.3.42}\quad 设 $f(x)=\dfrac1{1+x-2x^2}$，试证：$\displaystyle\sum_{n=0}^{\infty}\dfrac{n!}{f^{(n)}(0)}$ 绝对收敛.（浙江大学）

\noindent\textbf{证}\quad
\[
\begin{aligned}
f(x)&=\frac1{1+x-2x^2}
=\frac13\left(\frac1{1-x}+\frac2{1+2x}\right)\\
&=\frac13\left[\sum_{n=0}^{\infty}x^n+\sum_{n=0}^{\infty}(-1)^n2^{n+1}x^n\right]\\
&=\sum_{n=0}^{\infty}\frac{1+(-1)^n2^{n+1}}3x^n.
\end{aligned}
\]
另一方面，
\[
  f(x)=\sum_{n=0}^{\infty}\frac{f^{(n)}(0)}{n!}x^n.
\]
由展开式唯一性，知
\[
  \sum_{n=0}^{\infty}\frac{n!}{f^{(n)}(0)}
  =\sum_{n=0}^{\infty}\frac3{1+(-1)^n2^{n+1}}.
\]
因
\[
  \left|\frac3{1+(-1)^n2^{n+1}}\right|\le\frac3{2^n},
\]
且 $\displaystyle\sum_{n=0}^{\infty}\dfrac3{2^n}$ 收敛. 故级数 $\displaystyle\sum_{n=0}^{\infty}\dfrac{n!}{f^{(n)}(0)}$ 收敛，且绝对收敛.

\noindent\textbf{*例 5.3.43}\quad 设 $f(x)=\displaystyle\sum_{n=1}^{\infty}\dfrac1{n^2\ln(1+n)}x^n$，试证：

1）$f(x)$ 在 $[-1,1]$ 上连续；\qquad 2）$f(x)$ 在 $x=-1$ 处可导；

3）$\displaystyle\lim_{x\to1^-}f'(x)=+\infty$；\qquad 4）$f(x)$ 在 $x=1$ 处不可导.（浙江大学）

\noindent\textbf{提示}\quad 1）在 $x=1$ 处，$\displaystyle\sum_{n=1}^{\infty}\dfrac1{n^2\ln(1+n)}$ 有优级数 $\displaystyle\sum_{n=1}^{\infty}\dfrac1{n^2}$；

在 $x=-1$ 处，$\displaystyle\sum_{n=1}^{\infty}\dfrac{(-1)^n}{n^2\ln(1+n)}$ 是 Leibniz 级数. 所以结论 1）成立.

2）逐项求导后的级数在 $x=-1$ 处仍是 Leibniz 级数，收敛，所以级数在 $[-1,0]$ 上一致收敛，能逐项求导，$f(x)$ 在 $x=-1$ 处有右导数.

3）（思路：$\displaystyle f'(x)=\sum_{n=1}^{\infty}\frac1{n\ln(1+n)}x^{n-1}\ge
\sum_{k=1}^n\frac1{k\ln(1+k)}x^{k-1}\xrightarrow{x\to1^-}
\sum_{k=1}^n\frac1{k\ln(1+k)}\xrightarrow{n\to\infty}+\infty$.）

严格证明如下：在 $(0,1)$ 内，$f'(x)=\displaystyle\sum_{n=1}^{\infty}\dfrac1{n\ln(1+n)}x^{n-1}$. 用积分判别法易知
\[
  \sum_{k=2}^{n+1}\frac1{k\ln k}\to+\infty.
\]
而
\[
  \sum_{k=2}^{n+1}\frac1{k\ln k}
  \le\sum_{i=1}^n\frac1{i\ln(1+i)},\quad
  S_n\overset{\text{记}}{=}\sum_{i=1}^n\frac1{i\ln(1+i)}\to+\infty
\]
（当 $n\to\infty$ 时）. 即 $\forall M>0,\ \exists N>0$，当 $n>N$ 时，有 $S_n>M+1,\ S_n-1>M$.

而 $f'_n(x)\overset{\text{表示}}{=}\displaystyle\sum_{k=1}^n\dfrac{x^{k-1}}{k\ln(1+k)}\to S_n=\displaystyle\sum_{k=1}^n\dfrac1{k\ln(1+k)}$（当 $x\to1^-$ 时）. 故对于 $\varepsilon=1$，$\exists\delta>0$，当 $1-\delta<x<1$ 时，$0<S_n-f'_n(x)<\varepsilon=1$. 所以，$\forall M>0,\ \exists\delta>0$，当 $1-\delta<x<1$ 时，$f'(x)>f'_n(x)>S_n-1>M$. 因此，$\displaystyle\lim_{x\to1^-}f'(x)=+\infty$.

4）$f'_-(1)=\displaystyle\lim_{x\to1^-}\frac{f(x)-f(1)}{x-1}
\stackrel{x<\xi<1}{=}
\lim_{x\to1^-}\frac{f'(\xi)(x-1)}{x-1}=+\infty$，说明 $f(x)$ 在 $x=1$ 处不可导.

\noindent\textbf{例 5.3.44}\quad 已知在 $(-1,1)$ 内，幂级数 $\displaystyle\sum_{n=0}^{\infty}a_nx^n=f(x)$. 试证：若存在趋向 0 的序列 $\{x_n\}$：$0\ne x_n\in(-1,1)$，$\displaystyle\lim_{n\to\infty}x_n=0$，使得 $f(x_n)=0$（$n=1,2,\cdots$），则 $\forall x\in(-1,1)$，$f(x)\equiv0$.（南开大学）

\noindent\textbf{证}\quad 1° $a_0=\left.\displaystyle\sum_{n=0}^{\infty}a_nx^n\right|_{x=0}=f(0)=\displaystyle\lim_{n\to\infty}f(x_n)=0$.

2° 因 $f(0)=f(x_n)=0$，根据 Rolle 定理，$\exists\xi_n$，在 0 与 $x_n$ 之间，使得 $f'(\xi_n)=0\ (n\in\mathbb N)$. 对 $f'(x)=\displaystyle\sum_{n=1}^{\infty}na_nx^{n-1}$（此时常数项为 $a_1$），重复 1° 中的推理得
\[
  a_1=\left.\sum_{n=1}^{\infty}na_nx^{n-1}\right|_{x=0}
  =f'(0)=\lim_{n\to\infty}f'(\xi_n)=0.
\]

3° 继续这样做下去，顺次可得 $a_k=0\ (k=2,3,\cdots)$，说明 $\displaystyle\sum_{n=0}^{\infty}a_nx^n$ 是零级数，$f(x)\equiv0$.

\noindent\textbf{注}\quad 1）请将本题与例 5.3.35 进行比较.

2）假设在区间 $(-a,a)$ 内 $f(x)\equiv\displaystyle\sum_{n=0}^{\infty}a_nx^n,\ g(x)\equiv\displaystyle\sum_{n=0}^{\infty}b_nx^n$，存在序列 $\{x_n\}\to0$，使得 $f(x_n)=g(x_n)\ (n=1,2,\cdots)$. 记 $F(x)=f(x)-g(x)=\displaystyle\sum_{n=0}^{\infty}(a_n-b_n)x^n$，有 $F(x_n)=0$，即 $a_n=b_n\ (n=1,2,\cdots)$. 则 $f(x)\equiv g(x)$，$\forall x\in(-a,a)$. 如此，我们得到重要结论：若 $S(x)=\displaystyle\sum_{n=0}^{\infty}a_nx^n$ 的收敛半径 $R>0$，$\{x_n\}$ 是收敛于 0 的点列，则函数 $f:f(x)\equiv S(x)$（$\forall x\in(-R,R)$）$\Leftrightarrow f(x_n)=S(x_n)\ (n=1,2,\cdots)$.

换句话说：$\displaystyle\sum_{n=0}^{\infty}a_nx^n$ 的值能被它（在收敛于 0 的点列上）的值一意确定.

\noindent\textbf{例 5.3.45}\quad 求下列幂级数的收敛域：

1）$\displaystyle\sum_{n=1}^{\infty}
\frac{x^n}{1+\dfrac12+\cdots+\dfrac1n}$；（中国科学院）\hfill 《$[-1,1)$》

2）$\displaystyle\sum_{n=1}^{\infty}
\frac{x^n}{1+\dfrac1{\sqrt2}+\cdots+\dfrac1{\sqrt n}}$；\hfill 《$[-1,1)$》

3）$\displaystyle\sum_{n=2}^{\infty}
\frac{x^n}{\dfrac1{2\ln2}+\dfrac1{3\ln3}+\cdots+\dfrac1{n\ln n}}$. \hfill 《$[-1,1)$》

\noindent\textbf{提示}\quad 1）
\[
  \frac{a_n}{a_{n+1}}
  =\frac{1+\dfrac12+\cdots+\dfrac1{n+1}}
         {1+\dfrac12+\cdots+\dfrac1n}
  =1+\frac{\dfrac1{n+1}}{1+\dfrac12+\cdots+\dfrac1n}.
\]
注意到 $1\le1+\dfrac12+\cdots+\dfrac1n\le n$，因此
\[
  1+\frac{\dfrac1{n+1}}n
  \le\frac{a_n}{a_{n+1}}
  \le1+\frac{\dfrac1{n+1}}1,
\]
故 $R=\displaystyle\lim_{n\to\infty}\dfrac{a_n}{a_{n+1}}=1$，收敛区间为 $(-1,1)$.

（讨论端点.）在端点 $x=-1$ 处，级数为 Leibniz 级数（通项正负交错，绝对值单降，趋向零），因此收敛；

在端点 $x=1$ 处，
\[
  \sum_{k=1}^n\frac1{1+\dfrac12+\cdots+\dfrac1k}
  \ge\sum_{k=1}^n\frac1k\to+\infty\quad(n\to\infty),
\]
级数发散. 故收敛范围是 $[-1,1)$.

2），3）类似.

\noindent\textbf{*例 5.3.46}\quad 设函数 $f:\mathbb N\to\mathbb R$，$a_n(n\in\mathbb N)$ 为实数，且对于充分大的 $x$，有
\[
  f(x)=a_0+\frac{a_1}{x}+\frac{a_2}{x^2}+\cdots+\frac{a_n}{x^n}+\cdots. \qquad(1)
\]
试证：$\displaystyle\sum_{n=1}^{\infty}f(n)$ 收敛的充要条件是 $a_0=a_1=0$.（浙江大学）

\noindent\textbf{证}\quad 根据题意，对充分大的 $x$：$y=\dfrac1x>0$，$\displaystyle\sum_{n=0}^{\infty}a_ny^n$ 收敛，说明 $\dfrac1{\overline{\lim}|a_n|}=r>0$，即
\[
  \overline{\lim}|a_n|=\frac1r<+\infty. \qquad(2)
\]
说明 $\{a_n\}$ 有界，即 $\exists M:|a_n|\le M\ (n=1,2,\cdots)$. 因此
\[
\begin{aligned}
|f(n)-a_0|
&=\left|\sum_{k=1}^{\infty}\frac{a_k}{n^k}\right|
\le M\sum_{k=1}^{\infty}\frac1{n^k}
=M\frac{\dfrac1n}{1-\dfrac1n}\\
&=M\frac1{n-1}\to0\quad(n\to\infty). \qquad(3)
\end{aligned}
\]

1°（证明：若 $\displaystyle\sum_{n=1}^{\infty}f(n)$ 收敛，则必有 $a_0=a_1=0$.）事实上，$\displaystyle\sum_{n=1}^{\infty}f(n)$ 收敛，则 $f(n)\to0\ (n\to\infty)$. 由式 (3) 可知：$|a_0|\le |a_0-f(n)|+|f(n)|\to0\ (n\to\infty)$，故 $a_0=0$.

\[
  \sum_{n=2}^N f(n)=a_1\sum_{n=2}^N\frac1n
  +\sum_{n=2}^N\sum_{k=2}^{\infty}\frac{a_k}{n^k},
\]
其中若 $a_1\ne0$，则 $a_1\displaystyle\sum_{n=2}^N\dfrac1n\to+\infty\ (N\to\infty)$. 而
\[
  \left|\sum_{n=2}^N\sum_{k=2}^{\infty}\frac{a_k}{n^k}\right|
  \le M\sum_{n=2}^N\sum_{k=2}^{\infty}\frac1{n^k}
  =M\sum_{n=2}^N\left(\frac1{n-1}-\frac1n\right)<M,
\]
所以，$\displaystyle\lim_{N\to\infty}\sum_{n=2}^N f(n)=+\infty$，与 $\displaystyle\sum_{n=1}^{\infty}f(n)$ 收敛矛盾. 故只能 $a_1=0$.

2°（证明：若 $a_0=a_1=0$，则 $\displaystyle\sum_{n=1}^{\infty}f(n)$ 收敛.）若 $a_0=a_1=0$，则对充分大的 $N$，
\[
\begin{aligned}
\sum_{n=N+1}^{N+p}|f(n)|
&\le\sum_{n=N+1}^{N+p}\sum_{k=2}^{\infty}\frac{|a_k|}{n^k}
\le\sum_{n=N+1}^{N+p}\sum_{k=2}^{\infty}\frac{M}{n^k}\\
&=M\sum_{n=N+1}^{N+p}\frac{\dfrac1{n^2}}{1-\dfrac1n}
=M\sum_{n=N+1}^{N+p}\left(\frac1{n-1}-\frac1n\right)
\le\frac MN\to0\quad(N\to\infty)\quad(\forall p>0).
\end{aligned}
\]
根据 Cauchy 准则，级数 $\displaystyle\sum_{n=1}^{\infty}f(n)$ 收敛.

\section*{单元练习 5.3}

\noindent\textbf{5.3.1}\quad 对于幂级数 $\displaystyle\sum_{n=1}^{\infty}\frac{2^n\ln n}{n}x^n$，

1）求出收敛半径；

2）讨论在收敛域端点上的收敛性；

3）指出在什么样区间上级数一致收敛.（内蒙古大学）

\noindent\textbf{提示}\quad 在 $\left[-\dfrac12,\dfrac12\right)$ 内收敛，$\forall\alpha:|\alpha|<\dfrac12$，在 $\left[-\dfrac12,\alpha\right]$ 上一致收敛.

\noindent\textbf{5.3.2}\quad 若级数 $\displaystyle\sum_{n=0}^{\infty}a_nx^n$ 有收敛半径 $R_1$，而级数 $\displaystyle\sum_{n=0}^{\infty}b_nx^n$ 有收敛半径 $R_2$，求级数

1）$\displaystyle\sum_{n=0}^{\infty}(a_n+b_n)x^n$；\qquad
2）$\displaystyle\sum_{n=0}^{\infty}a_nb_nx^n$

的收敛半径. \hfill 《1）$R=\min(R_1,R_2)$；2）$R\ge R_1R_2$》

\noindent\textbf{提示}\quad 1）若 $R_1<x<R_2$，则 $\displaystyle\sum_{n=0}^{\infty}(a_n+b_n)x^n$ 发散.

2）若 $x=R_1R_2\theta^2\ (0<\theta<1)$，则 $\displaystyle\sum_{n=0}^{\infty}a_nb_nx^n$ 收敛.

\noindent\textbf{5.3.3}\quad 设 $a_n\ge0$，$\displaystyle\sum_{n=0}^{\infty}a_nx^n$ 的收敛半径为 1，和函数为 $f(x)$，若 $\displaystyle\sum_{n=0}^{\infty}a_n$ 发散，求证 $\displaystyle\lim_{x\to1^-}f(x)=+\infty$.

\noindent\textbf{提示}\quad 可用反证法.

\noindent\textbf{再提示}\quad 否则 $\exists M>0$ 及 $\{x_n\}\to1^-$（当 $n\to\infty$ 时），使得 $f(x_n)=\displaystyle\sum_{k=1}^{\infty}a_kx_n^k\le M$. 于是 $\forall m$：
\[
  0\le\sum_{k=1}^m a_k=\lim_{n\to\infty}\sum_{k=1}^m a_kx_n^k\le M,
\]
与 $\displaystyle\sum_{k=1}^{\infty}a_k$ 发散矛盾.

\noindent\textbf{\(\star\)5.3.4}\quad 证明：$y(x)=\displaystyle\sum_{n=0}^{\infty}\dfrac{x^{4n}}{(4n)!}$ 满足 $y^{(4)}=y$.（中国科学技术大学）

\noindent\textbf{提示}\quad 因 $\displaystyle\sum_{n=0}^{\infty}\dfrac{t^n}{(4n)!}$ 在 $\mathbb R$ 上处处收敛，故 $\displaystyle\sum_{n=0}^{\infty}\dfrac{x^{4n}}{(4n)!}$ 亦然. 逐项求导四次可得 $y^{(4)}=y$.

\noindent\textbf{注}\quad
\[
  \left[\sum_{n=0}^{\infty}\frac{x^{4n}}{(4n)!}\right]'
  =\left[\sum_{n=0}^{\infty}\frac{t^n}{(4n)!}\right]'_{t=x^4}\cdot t'_x
  =\sum_{n=1}^{\infty}\frac{nt^{n-1}}{(4n)!}\bigg|_{t=x^4}\cdot4x^3
  =\sum_{n=1}^{\infty}\frac{x^{4n-1}}{(4n-1)!}
  =\sum_{n=0}^{\infty}\left[\frac{x^{4n}}{(4n)!}\right]',
\]
表明该级数仍可直接进行逐项求导.

\noindent\textbf{\(\star\)5.3.5}\quad 求极限 $\displaystyle\lim_{n\to\infty}\sum_{k=1}^n\frac{k+2}{k!+(k+1)!+(k+2)!}$.（四川师范大学）

\noindent\textbf{提示}\quad
\[
  \frac{k+2}{k!+(k+1)!+(k+2)!}
  =\frac1{(k+1)!}-\frac1{(k+2)!},
\]
原极限
\[
  =\left[\sum_{k=1}^{\infty}\frac{x^{k+1}}{(k+1)!}
  -\sum_{k=1}^{\infty}\frac{x^{k+2}}{(k+2)!}\right]_{x=1}'
  =\left.\frac{x^2}{2}\right|_{x=1}
  =\frac12.
\]

\noindent\textbf{*5.3.6}\quad 设序列 $\{a_n\}_{n=1}^{\infty}$，$\{b_n\}_{n=1}^{\infty}$ 满足：$a_n>0$，级数 $\displaystyle\sum_{n=0}^{\infty}a_nx^n$ 当 $|x|<1$ 时收敛，当 $x=1$ 时发散，又 $\displaystyle\lim_{n\to\infty}\dfrac{b_n}{a_n}=A\ (0\le A<+\infty)$，证明：
\[
  \lim_{x\to1^-}\frac{\displaystyle\sum_{n=0}^{\infty}b_nx^n}
  {\displaystyle\sum_{n=0}^{\infty}a_nx^n}=A.
\]
（南京大学）

\noindent\textbf{证}\quad
\[
  \left|\frac{\displaystyle\sum_{k=1}^n b_kx^k}{\displaystyle\sum_{k=1}^n a_kx^k}-A\right|
  =
  \left|\frac{\displaystyle\sum_{k=1}^n a_k\left(\dfrac{b_k}{a_k}-A\right)x^k}
  {\displaystyle\sum_{k=1}^n a_kx^k}\right|. \qquad(1)
\]
$\forall\varepsilon>0,\ \exists N>0$，当 $n>N$ 时 $\left|\dfrac{b_k}{a_k}-A\right|<\dfrac\varepsilon2$，
\[
  \text{式 }(1)\le
  \frac{\displaystyle\sum_{k=1}^N(|b_k|+a_kA)}
  {\displaystyle\sum_{k=1}^n a_kx^k}
  +\frac\varepsilon2. \qquad(2)
\]
又因 $\displaystyle\sum_{k=1}^{\infty}a_k=+\infty$，$\exists N_1>N$，使得
\[
  \frac{\displaystyle\sum_{k=1}^{N}(|b_k|+a_kA)}
  {\displaystyle\sum_{k=1}^{N_1}a_k}<\frac\varepsilon4.
\]
式 (2) 右端
\[
  <\frac\varepsilon4\cdot
  \frac{\displaystyle\sum_{k=1}^{N_1}a_k}{\displaystyle\sum_{k=1}^{N_1}a_kx^k}\cdot
  \frac{\displaystyle\sum_{k=1}^{N_1}a_kx^k}{\displaystyle\sum_{k=1}^{n}a_kx^k}
  +\frac\varepsilon2. \qquad(3)
\]
因
\[
  \frac{\displaystyle\sum_{k=1}^{N_1}a_k}
  {\displaystyle\sum_{k=1}^{N_1}a_kx^k}\to1\quad(x\to1^-),
\]
故 $\exists\delta>0$，当 $0<x<1-\delta$ 时，有
\[
  \frac{\displaystyle\sum_{k=1}^{N_1}a_k}
  {\displaystyle\sum_{k=1}^{N_1}a_kx^k}<2.
\]
于是当 $n>N_1,\ 0<1-x<\delta$ 时，式 (3)
$<\dfrac\varepsilon4\cdot2\cdot1+\dfrac\varepsilon2=\varepsilon$，证毕.

\noindent\textbf{※5.3.7}\quad 设 $\dfrac{v_n}{v_{n-1}}=a\sqrt{\dfrac{n-1}{n+1}},\ n=2,3,\cdots,\ |a|<1$，
\[
  x_{n+1}=x_n+cv_n^2,\quad n=1,2,\cdots,\ c>0.
\]
求 $\displaystyle\lim_{n\to\infty}x_n$.

\noindent\textbf{思想方法}\quad 写出 $x_n$ 的通项式，并行递推，变成级数形式求和.

\noindent\textbf{解}\quad
\[
  x_{n+1}=x_n+cv_n^2=x_{n-1}+cv_{n-1}^2+cv_n^2
  =\cdots=x_1+c\sum_{k=1}^n v_k^2.
\]
\[
\begin{aligned}
v_k^2
&=a^2\frac{k-1}{k+1}v_{k-1}^2
=a^2\frac{k-1}{k+1}\cdot a^2\frac{k-2}{k}v_{k-2}^2=\cdots\\
&=a^{2(k-1)}\frac{k-1}{k+1}\cdot\frac{k-2}{k}\cdot\frac{k-3}{k-1}\cdots\frac2{4}\cdot\frac1{3}v_1^2
=a^{2(k-1)}\frac2{k+1}\cdot\frac1k\,v_1^2.
\end{aligned}
\]
因此
\[
\begin{aligned}
x_{n+1}
&=x_1+c\sum_{k=1}^n v_k^2
=x_1+c\sum_{k=1}^n a^{2(k-1)}\frac2{k+1}\cdot\frac1k\,v_1^2\\
&=x_1+2cv_1^2\sum_{k=1}^n a^{2(k-1)}\left(\frac1k-\frac1{k+1}\right).
\end{aligned}
\]
于是
\[
  \lim_{n\to\infty}x_n
  =x_1+2cv_1^2\sum_{k=1}^{\infty}a^{2(k-1)}
  \left(\frac1k-\frac1{k+1}\right), \qquad(1)
\]
其中
\[
  \sum_{k=1}^{\infty}\frac{a^{2(k-1)}}{k+1}
  =\frac1{a^4}\sum_{k=1}^{\infty}\frac{a^{2(k+1)}}{k+1}.
\]
记 $f(a)=\displaystyle\sum_{k=1}^{\infty}\dfrac{a^{2(k+1)}}{k+1}$，则
\[
  f'(a)=\sum_{k=1}^{\infty}\frac{2(k+1)a^{2(k+1)-1}}{k+1}
  =2\sum_{k=1}^{\infty}a^{2k+1}
  =\frac{2a^3}{1-a^2}.
\]
于是
\[
\begin{aligned}
  \sum_{k=1}^{\infty}\frac{a^{2(k+1)}}{k+1}
  =f(a)
  &=\int_0^a\frac{2t^3}{1-t^2}\,dt
  =\int_0^a\frac{2t^3-2t+2t}{1-t^2}\,dt\\
  &=-a^2+\int_0^a\frac{2t}{1-t^2}\,dt
  =-a^2-\ln(1-a^2). \qquad(2)
\end{aligned}
\]
由此得
\[
  \sum_{k=1}^{\infty}\frac{a^{2(k-1)}}{k+1}
  =\frac1{a^4}\sum_{k=1}^{\infty}\frac{a^{2(k+1)}}{k+1}
  =-\frac1{a^2}-\frac1{a^4}\ln(1-a^2). \qquad(3)
\]
式 (1) 中的另一项
\[
\begin{aligned}
  \sum_{k=1}^{\infty}\frac{a^{2(k-1)}}k
  &=1+\sum_{k=2}^{\infty}\frac{a^{2(k-1)}}k
  =1+\frac1{a^2}\sum_{k=2}^{\infty}\frac{a^{2k}}k\\
  &\overset{\text{式 }(2)}{=}1+\frac1{a^2}\left[-a^2-\ln(1-a^2)\right]
  =-\frac1{a^2}\ln(1-a^2). \qquad(4)
\end{aligned}
\]
(3)，(4) 代入 (1) 得
\[
\begin{aligned}
  \lim_{n\to\infty}x_n
  &=x_1+2cv_1^2\left[-\frac1{a^2}\ln(1-a^2)+\frac1{a^2}
  +\frac1{a^4}\ln(1-a^2)\right]\\
  &=x_1+2cv_1^2\frac1{a^4}\left[a^2+(1-a^2)\ln(1-a^2)\right].
\end{aligned}
\]

\noindent\textbf{\(\star\)5.3.8}\quad 设 $a_n\ge0(n=1,2,\cdots)$，$\displaystyle\sum_{n=1}^{\infty}a_nx^n$ 当 $-1<x<1$ 时收敛并且有上界，证明：

1）$\displaystyle\lim_{x\to1^-}\sum_{n=1}^{\infty}a_nx^n$ 存在；\quad
2）$\displaystyle\sum_{n=1}^{\infty}a_n$ 收敛；\quad
3）$\displaystyle\lim_{x\to1^-}\sum_{n=1}^{\infty}a_nx^n=\sum_{n=1}^{\infty}a_n$.（厦门大学）

\noindent\textbf{提示}\quad 1）$x\to1^-$ 时，$\displaystyle\sum_{k=1}^n a_kx^k$ 有一致上界.

2）$\displaystyle\sum a_k$ 的部分和有界.

3）在 $[0,1]$ 上 $\displaystyle\sum a_nx^n$ 一致收敛.

\noindent\textbf{5.3.9}\quad 设 $f(x)=\displaystyle\sum_{n=0}^{\infty}a_nx^n$ 的收敛半径为 $R=+\infty$，令 $f_n(x)=\displaystyle\sum_{k=0}^n a_kx^k$. 求证：当 $n\to\infty$ 时
\[
  f(f_n(x))\rightrightarrows f(f(x))\quad(a\le x\le b).
\]

\noindent\textbf{提示}\quad 收敛半径 $R=+\infty\Rightarrow\forall[a,b]$ 上 $f_n\rightrightarrows f$，并且 $\{f_n\},f$ 一致有界，$\exists A>0$，使 $f(x)\in[-A,A]\ (\forall x\in\mathbb R)$.

\noindent\textbf{再提示}\quad $f$ 在 $[-A,A]$ 上一致连续，$f_n\rightrightarrows f\Longrightarrow f(f_n(x))\rightrightarrows f(f(x))\ (n\to\infty)$ 关于 $x\in[a,b]$ 一致.

\noindent\textbf{\(\star\)5.3.10}\quad 求级数 $\displaystyle\sum_{n=1}^{\infty}n\cdot2^{\frac n2}x^{3n-1}$ 的收敛区间与和函数.（华中师范大学）

\hfill 《收敛区间 $(-1,1)$，端点发散，和 $=\dfrac{2^{1/2}x^2}{(1-x^3)^2}$》

\noindent\textbf{\(\star\)5.3.11}\quad 证明：
\[
  \sum_{n=1}^{\infty}\frac{(-1)^{n-1}}{(2n-1)(2n+1)3^n}
  =\sqrt3\int_0^{1/\sqrt3}x\arctan x\,dx
  =\frac{2\pi\sqrt3-9}{18}.
\]
（西南大学）

\noindent\textbf{提示}\quad 参看例 3.1.19，可写出 $x\arctan x=\displaystyle\sum_{n=1}^{\infty}(-1)^{n-1}\dfrac{x^{2n}}{2n-1},\ |x|\le1$.

\noindent\textbf{再提示}\quad 上式逐项积分可得左边等式，分部积分可得右边等式.

\noindent\textbf{*5.3.12}\quad 验证积分 $\displaystyle\int_0^1\left|\ln\frac{1+x}{1-x}\right|\cdot\frac{dx}{x}$ 存在且等于 $2\displaystyle\sum_{n=1}^{\infty}\dfrac1{(2n-1)^2}$.（湘潭大学）

\noindent\textbf{提示}\quad $x^{-1}\ln\dfrac{1+x}{1-x}=2\displaystyle\sum_{n=1}^{\infty}\dfrac{x^{2n-2}}{2n-1}$ 在 $[0,1-\varepsilon]$ 上逐项积分，然后令 $\varepsilon\to0^+$.

\noindent\textbf{另解}\quad 令 $u=\dfrac{1+x}{1-x}$ 作变换，原积分 $I=2\displaystyle\int_1^{+\infty}\dfrac{\ln u}{u^2-1}\,du$. 再令 $v=\dfrac1u$，则 $I=-2\displaystyle\int_0^1\dfrac{\ln v}{1-v^2}\,dv$，见下面 5.3.14：
\[
  \frac{\pi^2}{4}=2\sum_{n=1}^{\infty}\frac1{(2n-1)^2}\quad(\text{见例 5.4.8}).
\]

\noindent\textbf{*5.3.13}\quad 试证：
\[
  \int_0^x\frac{\arctan t}{t}\ln\frac{x}{t}\,dt
  =\sum_{n=0}^{\infty}\frac{(-1)^nx^{2n+1}}{(2n+1)^3}\quad(|x|\le1).
\]
（四川大学）

\noindent\textbf{提示}\quad 代入 $\arctan t=\displaystyle\sum_{n=1}^{\infty}(-1)^{n-1}\dfrac{t^{2n-1}}{2n-1}\ (|t|\le1)$. 令 $0<\varepsilon<x<1$，在 $[\varepsilon,x]$ 上逐项积分，然后令 $\varepsilon\to0^+$.

\noindent\textbf{**5.3.14}\quad 证明：$A=\displaystyle\int_0^1\dfrac{\ln x}{1-x^2}\,dx$，$B=\displaystyle\int_0^1\dfrac{\ln x}{1-x}\,dx$，$C=\displaystyle\int_0^1\dfrac{\ln x}{1+x}\,dx$ 收敛，并求其值.

\noindent\textbf{提示}\quad $x=0$ 为奇点，当 $0\le x\le\dfrac12$ 时，有
\[
  \left|\frac{\ln x}{1+x}\right|\le|\ln x|,\quad
  \left|\frac{\ln x}{1-x}\right|\le2|\ln x|,\quad
  \left|\frac{\ln x}{1-x^2}\right|\le\frac43|\ln x|.
\]
可证 $A,B,C$ 收敛. 利用分部积分，$\ln(1+x)$ 的展开式在 $[0,1]$ 上可逐项积分.
$C=-\dfrac{\pi^2}{12},\ A=\dfrac32C=-\dfrac{\pi^2}{8},\ B=2C=-\dfrac{\pi^2}{6}$.

\noindent\textbf{*5.3.15}\quad 设幂级数 $\displaystyle\sum_{n=1}^{\infty}a_nx^n$ 的收敛半径大于 0，证明：

1）$\displaystyle\lim_{x\to+0}\sum_{n=1}^{\infty}a_nx^n=0$；

2）如果 $a_1\ne0$，并且在原点的一个邻域里
\[
  \left|\sum_{n=1}^{\infty}a_nx^n\right|\ge |a_1||x|-2x^2
\]
逐点成立，那么 $|a_2|\le2$.（厦门大学）

\noindent\textbf{证}\quad 1）收敛半径 $>0$，则函数 $S(x)=\displaystyle\sum_{n=1}^{\infty}a_nx^n$ 在 $x=0$ 处连续，故
\[
  \lim_{x\to+0}\sum_{n=1}^{\infty}a_nx^n
  =\left.\sum_{n=1}^{\infty}a_nx^n\right|_{x=0}=0.
\]

2）当 $a_1>0$ 时（$a_1<0$ 时类似可证），在 0 点附近，$\displaystyle\sum_{n=1}^{\infty}a_nx^n$ 的符号跟 $a_1$ 相同，因此 $a_1>0$ 时：对充分小的 $\delta>0$，当 $x\in(0,\delta)$ 时，
\[
  \left|\sum_{n=1}^{\infty}a_nx^n\right|
  =\sum_{n=1}^{\infty}a_nx^n
  =a_1x+a_2x^2+o(x^2).
\]
根据已知条件，$\displaystyle\sum_{n=1}^{\infty}a_nx^n\ge a_1x-2x^2,\ x\in(0,\delta)$. 因此
\[
  a_2x^2+o(x^2)\ge-2x^2,\quad a_2+o(1)\ge-2.
\]
令 $x\to0^+$，得
\[
  a_2\ge-2. \qquad(1)
\]
同理，当 $x\in(-\delta,0)$ 时，
\[
  \left|\sum_{n=1}^{\infty}a_nx^n\right|
  =-\left[a_1x+a_2x^2+o(x^2)\right]
  \overset{\text{已知条件}}{\ge}-a_1x-2x^2,\quad x\in(-\delta,0).
\]
故 $-a_2-o(1)\ge-2$，即 $a_2+o(1)\le2$. 令 $x\to0^-$，得 $a_2\le2$.

联系 (1)，得 $|a_2|\le2$.

\section*{§5.4\quad Fourier 级数}

\noindent\textbf{导读}\quad Fourier 级数是函数项级数的一种特殊形式. 它是研究和表示周期函数的有力工具. 本节我们将围绕展开与收敛问题进行讨论.

Fourier 级数相当重要. 可能由于计算过长，考题相对较少.

\subsection*{一、正交系}

\noindent\textbf{要点}\quad 要证明 $\{\varphi_n(x)\}$ 在 $[a,b]$ 上正交，即要证明
\[
  \int_a^b \varphi_n(x)\varphi_m(x)\,dx=0\quad(\text{当 }m\ne n\text{ 时}).
\]

\noindent\textbf{※例 5.4.1}\quad 试证明 Legendre 多项式
\[
  X_n(x)=\frac1{2^nn!}\frac{d^n(x^2-1)^n}{dx^n}
  \equiv\frac1{(2n)!!}\frac{d^n(x^2-1)^n}{dx^n}
\]
在 $[-1,1]$ 上为正交的，并求 $\displaystyle\int_{-1}^1X_n^2(x)\,dx$ 之值.

\noindent\textbf{证}\quad 记
\[
  u=u(x)=(x^2-1)^n=(x-1)^n(x+1)^n,
\]
则 $u,u',u'',\cdots,u^{(n-1)}$ 当 $x=\pm1$ 时为 0，且 $X_n(x)=\dfrac1{(2n)!!}u^{(n)}$ 为 $x$ 的 $n$ 次多项式. 设 $m<n$，则 $X_m(x)$ 为 $m$ 次多项式. 反复应用分部积分法（$n$ 次），有
\[
\begin{aligned}
(2n)!!\int_{-1}^1X_n(x)X_m(x)\,dx
&=\int_{-1}^1u^{(n)}(x)X_m(x)\,dx
=\int_{-1}^1X_m(x)\,du^{(n-1)}(x)\\
&=X_m(x)u^{(n-1)}(x)\bigg|_{-1}^1
-\int_{-1}^1u^{(n-1)}(x)X'_m(x)\,dx\\
&=-\int_{-1}^1u^{(n-1)}(x)X'_m(x)\,dx
=\cdots\\
&=(-1)^n\int_{-1}^1u(x)X_m^{(n)}(x)\,dx=0.
\end{aligned}
\]
（因为 $X_m(x)$ 是 $m$ 次多项式，$n>m$，可知 $X_m^{(n)}(x)\equiv0$.）这就证明 $\{X_n(x)\}$ 有正交性. 用类似的方法可求出
\[
  \int_{-1}^1X_n^2(x)\,dx=\frac2{2n+1}.
\]

\noindent\textbf{※例 5.4.2}\quad 设序列 $\{y_n(x)\}$ 满足方程
\[
  \frac d{dx}\left[p(x)\frac{dy_n}{dx}\right]+\lambda_ny_n=0\quad
  (\forall x\in[a,b])\quad(n=1,2,\cdots) \qquad(1)
\]
（其中 $n\ne m$ 时 $\lambda_n\ne\lambda_m$）及边界条件 $y_n(a)=y_n(b)=0$. 试证 $\{y_n(x)\}$ 在 $[a,b]$ 上为正交系.

\noindent\textbf{证}\quad 由 (1)，$\lambda_ny_n=-\dfrac d{dx}\left[p(x)\dfrac{dy_n}{dx}\right]$，于是
\[
\begin{aligned}
\lambda_n\int_a^b y_ny_m\,dx
&=\int_a^b\lambda_ny_ny_m\,dx
=-\int_a^b\left\{\frac d{dx}\left[p(x)\frac{dy_n}{dx}\right]\right\}y_m\,dx\\
&=-p(x)\frac{dy_n}{dx}y_m\bigg|_a^b
+\int_a^b p(x)\frac{dy_n}{dx}\frac{dy_m}{dx}\,dx
=\int_a^b p(x)\frac{dy_n}{dx}\frac{dy_m}{dx}\,dx. \qquad(2)
\end{aligned}
\]
注意此式积分号下的 $y_n,y_m$ 地位是平等的，将 $m,n$ 互换得
\[
  \lambda_m\int_a^b y_ny_m\,dx
  =\int_a^b p(x)\frac{dy_n}{dx}\frac{dy_m}{dx}\,dx. \qquad(3)
\]
(2) $-$ (3) 得
\[
  (\lambda_n-\lambda_m)\int_a^b y_ny_m\,dx=0,
\]
故
\[
  \int_a^b y_ny_m\,dx=0\quad(\text{当 }n\ne m\text{ 时}).\quad\text{证毕}
\]

\noindent\textbf{※例 5.4.3}\quad 设 $0<\lambda_1<\lambda_2<\cdots<\lambda_n<\cdots$ 满足
\[
  \sigma\sin\sqrt{\lambda_n}l+\sqrt{\lambda_n}\cos\sqrt{\lambda_n}l=0\quad(\sigma>0),
\]
求证：$\{y_n(x)\}=\{\sin\sqrt{\lambda_n}x\}$ 在 $[0,l]$ 上为正交系，并求 $\displaystyle\int_0^l y_n^2(x)\,dx$（用 $\sigma,\lambda_n,l$ 表示）.

\noindent\textbf{提示}\quad 利用三角公式计算积分 $\displaystyle\int_0^l y_m(x)y_n(x)\,dx(m\ne n)$ 及 $\displaystyle\int_0^l y_n^2(x)\,dx$，并对结果应用题设的关系式.

\subsection*{二、Fourier 系数}

\noindent\textbf{要点}\quad 若 $f(x)$ 以 $2\pi$ 为周期，在区间 $[-\pi,\pi]$ 上可积，则
\[
  a_n=\frac1\pi\int_{-\pi}^{\pi}f(x)\cos nx\,dx,\quad n=0,1,2,\cdots,
\]
\[
  b_n=\frac1\pi\int_{-\pi}^{\pi}f(x)\sin nx\,dx,\quad n=1,2,\cdots,
\]
称为 $f(x)$ 的 Fourier 系数. $\dfrac{a_0}2+\displaystyle\sum_{n=1}^{\infty}(a_n\cos nx+b_n\sin nx)$ 称为 $f(x)$ 在 $[-\pi,\pi]$ 上的 Fourier 级数. 由此易知 Fourier 系数有如下性质：

1）若 $f_1(x),f_2(x)$ 的 Fourier 系数为 $a_n^{(1)},b_n^{(1)}$ 与 $a_n^{(2)},b_n^{(2)}$，则函数 $f(x)=\alpha f_1(x)\pm\beta f_2(x)$ 的 Fourier 系数为
\[
  a_n=\alpha a_n^{(1)}\pm\beta a_n^{(2)}\quad(n=0,1,2,\cdots),
\]
\[
  b_n=\alpha b_n^{(1)}\pm\beta b_n^{(2)}\quad(n=1,2,\cdots).
\]

2）若 $f(x)$ 在 $[-\pi,\pi]$ 上连续，分段光滑，$a_n,b_n$ 是 $f(x)$ 的 Fourier 系数，利用分部积分法，易知 $f'(x)$ 的 Fourier 系数（用 $a'_n,b'_n$ 表示）：
\[
  b'_n=-na_n\quad(n=1,2,\cdots),
\]
当 $f(-\pi)=f(\pi)$ 时，$a'_0=0,\ a'_n=nb_n\ (n=1,2,\cdots)$.

3）若 $f(x)$ 在 $[-\pi,\pi]$ 上分段连续，$a_n,b_n$ 是它的 Fourier 系数，则不难验证，
\[
  F(x)\equiv\int_0^x\left(f(t)-\frac{a_0}2\right)dt
\]
的 Fourier 系数（用 $A_n,B_n$ 表示）：
\[
  A_n=-\frac{b_n}{n},\quad B_n=\frac{a_n}{n}\quad(n=1,2,\cdots),\quad
  A_0=2\sum_{n=1}^{\infty}\frac{b_n}{n}.
\]

下面我们再来证明一些性质.

\noindent\textbf{例 5.4.4}\quad 1）设 $f(x)$ 以 $2\pi$ 为周期，在 $(0,2\pi)$ 内有界，试证：若 $f(x)$ 单调减，则系数 $b_n\ge0$；若 $f(x)$ 单调增，则 $b_n\le0\ (n=1,2,\cdots)$；

2）设 $f(x)$ 在 $(0,2\pi)$ 内 $f'(x)$ 有界，试证：若 $f'(x)$ 单调增，则 $a_n\ge0$；若 $f'(x)$ 单调减，则 $a_n\le0\ (n=1,2,\cdots)$；

3）设 $f(x)$ 在 $(0,2\pi)$ 上可积，试证：若 $F(x)\equiv\displaystyle\int_0^x\left(f(t)-\dfrac{a_0}2\right)dt$ 单调减，则 $f$ 的系数 $a_n\ge0$；若 $F(x)$ 单调增，则 $a_n\le0\ (n=1,2,\cdots)$.

\noindent\textbf{证}\quad 1）将 $[0,2\pi]$ $n$ 等分，则
\[
  b_n=\frac1\pi\int_0^{2\pi}f(x)\sin nx\,dx
  =\frac1\pi\sum_{i=1}^n\int_{(i-1)\frac{2\pi}n}^{i\frac{2\pi}n}f(x)\sin nx\,dx
\]
\[
  =\frac1\pi\sum_{i=1}^n\left[
  \int_{(i-1)\frac{2\pi}n}^{(i-\frac12)\frac{2\pi}n}f(x)\sin nx\,dx
  +\int_{(i-\frac12)\frac{2\pi}n}^{i\frac{2\pi}n}f(x)\sin nx\,dx\right]
\]
\[
  =\frac1\pi\sum_{i=1}^n\left[
  \int_{(i-1)\frac{2\pi}n}^{(i-\frac12)\frac{2\pi}n}f(x)\sin nx\,dx
  -\int_{(i-1)\frac{2\pi}n}^{(i-\frac12)\frac{2\pi}n}f\left(t+\frac\pi n\right)\sin nt\,dt\right]\quad\left(\text{其中 }t=x-\frac\pi n\right)
\]
\[
  =\frac1\pi\sum_{i=1}^n
  \int_{(i-1)\frac{2\pi}n}^{(i-\frac12)\frac{2\pi}n}
  \left[f(x)-f\left(x+\frac\pi n\right)\right]\sin nx\,dx\ge0.
\]
（这是因为 $f(x)$ 单调减，$f(x)-f\left(x+\dfrac\pi n\right)\ge0$；又因在
$\left[(i-1)\dfrac{2\pi}n,\left(i-\dfrac12\right)\dfrac{2\pi}n\right]$ 上 $\sin nx\ge0$.）

2）与 3）可利用要点 2）和 3）的关系得到.

\noindent\textbf{例 5.4.5}\quad 设 $f(x)$ 是以 $2\pi$ 为周期的函数，满足 $\alpha$ 阶 H\"older 条件（亦称 $\alpha$ 阶的 Lipschitz 条件）：
\[
  |f(x)-f(y)|\le L|x-y|^{\alpha}\quad(0<\alpha\le1).
\]
证明：$a_n=O\left(\dfrac1{n^{\alpha}}\right)$，$b_n=O\left(\dfrac1{n^{\alpha}}\right)$（当 $n\to\infty$ 时）.

\noindent\textbf{证}\quad
\[
  a_n=\frac1\pi\int_{-\pi}^{\pi}f(x)\cos nx\,dx \qquad(1)
\]
\[
\begin{aligned}
&=\frac1\pi\int_{-\pi-\frac\pi n}^{\pi-\frac\pi n}
  f\left(t+\frac\pi n\right)\cos(nt+\pi)\,dt
  \quad\left(x=t+\frac\pi n\right)\\
&=-\frac1\pi\int_{-\pi}^{\pi}f\left(t+\frac\pi n\right)\cos nt\,dt\\
&=-\frac1\pi\int_{-\pi}^{\pi}f\left(x+\frac\pi n\right)\cos nx\,dx. \qquad(2)
\end{aligned}
\]
(1)，(2) 平均得
\[
  a_n=\frac1{2\pi}\int_{-\pi}^{\pi}
  \left[f(x)-f\left(x+\frac\pi n\right)\right]\cos nx\,dx,
\]
\[
\begin{aligned}
|a_n|
&\le\frac1{2\pi}\int_{-\pi}^{\pi}
\left|f(x)-f\left(x+\frac\pi n\right)\right||\cos nx|\,dx\\
&\le\frac1{2\pi}L\left(\frac\pi n\right)^{\alpha}
\int_{-\pi}^{\pi}|\cos nx|\,dx
\le L\left(\frac\pi n\right)^{\alpha}.
\end{aligned}
\]
因此 $|a_n|=O\left(\dfrac1{n^\alpha}\right)$. 类似可证 $|b_n|=O\left(\dfrac1{n^\alpha}\right)$.

\noindent\textbf{例 5.4.6}\quad 设 $f(x)$ 有界，周期为 $2\pi$，并在 $(-\pi,\pi)$ 内逐段单调，试证
\[
  a_n=O\left(\frac1n\right),\quad b_n=O\left(\frac1n\right)\quad(\text{当 }n\to\infty\text{ 时}).
\]

\noindent\textbf{证 I}\quad
\[
  a_n=\frac1\pi\int_{-\pi}^{\pi}f(x)\cos nx\,dx
  \quad\left(\text{令 }x=t+\frac kn\pi\right)
\]
\[
  =\frac{(-1)^k}{\pi}\int_{-\pi}^{\pi}f\left(x+\frac{k\pi}n\right)\cos nx\,dx
  \quad(k=-n+1,\cdots,0,1,\cdots,n).
\]
同理，
\[
  a_n=\frac{(-1)^{k-1}}{\pi}\int_{-\pi}^{\pi}
  f\left(x+\frac{k-1}{n}\pi\right)\cos nx\,dx.
\]
两式平均得
\[
  a_n=\frac{(-1)^k}{2\pi}\int_{-\pi}^{\pi}
  \left[f\left(x+\frac kn\pi\right)-f\left(x+\frac{k-1}{n}\pi\right)\right]\cos nx\,dx
  \quad(k=-n+1,\cdots,0,1,\cdots,n). \qquad(1)
\]
(1) 中的 $2n$ 个式子相加求平均，得
\[
\begin{aligned}
|a_n|
&=\left|\frac1{2\pi}\int_{-\pi}^{\pi}\frac1{2n}
\sum_{k=-n+1}^{n}(-1)^k
\left[f\left(x+\frac kn\pi\right)-f\left(x+\frac{k-1}{n}\pi\right)\right]\cos nx\,dx\right|\\
&\le\frac1{4n\pi}\int_{-\pi}^{\pi}
\sum_{k=-n+1}^{n}\left|
f\left(x+\frac kn\pi\right)-f\left(x+\frac{k-1}{n}\pi\right)
\right||\cos nx|\,dx\\
&\le\frac1{4n\pi}\bigvee_{-\pi}^{\pi}f\cdot
\int_{-\pi}^{\pi}|\cos nx|\,dx
\le\frac1{2n}\bigvee_{-\pi}^{\pi}f. \qquad(2)
\end{aligned}
\]
（其中 $\displaystyle\bigvee_{-\pi}^{\pi}f=\sup\left\{\sum_{i=1}^{n}|f(x_i)-f(x_{i-1})|\mid -\pi=x_0<x_1<\cdots<x_n=\pi\right\}$ 表示 $f$ 在 $[-\pi,\pi]$ 上的全变差. 因 $f(x)$ 有界，且分段单调，故 $\displaystyle\bigvee_{-\pi}^{\pi}f<+\infty$.）式 (2) 表明 $a_n=O\left(\dfrac1n\right)$. 同理可证 $b_n=O\left(\dfrac1n\right)$.

\noindent\textbf{证 II}\quad （借助于 Stieltjes 积分）.
\[
\begin{aligned}
a_n&=\frac1\pi\int_{-\pi}^{\pi}f(x)\cos nx\,dx
=\frac1{n\pi}\int_{-\pi}^{\pi}f(x)\,d\sin nx\\
&=\frac1{n\pi}f(x)\sin nx\bigg|_{-\pi}^{\pi}
-\frac1{n\pi}\int_{-\pi}^{\pi}\sin nx\,df(x)\footnotemark\\
&=-\frac1{n\pi}\int_{-\pi}^{\pi}\sin nx\,df(x). \qquad(3)
\end{aligned}
\]
因此，
\[
  |a_n|\le\frac1{n\pi}\left|\int_{-\pi}^{\pi}\sin nx\,df(x)\right|
  \le\frac1{n\pi}\max_{-\pi\le x\le\pi}|\sin nx|\cdot\bigvee_{-\pi}^{\pi}f
  \le\frac1{n\pi}\bigvee_{-\pi}^{\pi}f,
\]
所以 $a_n=O\left(\dfrac1n\right)$. 同理可证 $b_n=O\left(\dfrac1n\right)$.

\footnotetext{对于 Stieltjes 积分，分部积分公式亦成立. 设 $-\pi=x_0<x_1<\cdots<x_k=\pi$ 为任一分划，则
\[
\sum_{i=0}^{k-1}f(\xi_i)\Delta\sin nx_i
=\sum_{i=0}^{k-1}f(\xi_i)(\sin nx_{i+1}-\sin nx_i)
=\sum_{i=1}^{k}f(\xi_{i-1})\sin nx_i-\sum_{i=0}^{k-1}f(\xi_i)\sin nx_i
\]
\[
=f(\xi_{k-1})\sin n\pi+\sum_{i=1}^{k-1}[f(\xi_{i-1})-f(\xi_i)]\sin nx_i-f(\xi_0)\sin n(-\pi),
\]
其中 $x_i\le\xi_i\le x_{i+1}$. 令 $\lambda=\max_{0\le i\le k-1}|x_{i+1}-x_i|\to0$，取极限，即得式 (3). 因
\[
  \left|\sum_i\sin n\xi_i\Delta f(x_i)\right|
  \le\max_{-\pi\le x\le\pi}|\sin nx|\cdot\sum_i|f(x_{i+1})-f(x_i)|
  \le\max_{-\pi\le x\le\pi}|\sin nx|\cdot\bigvee_{-\pi}^{\pi}f.
\]
令 $\lambda=\max|x_{i+1}-x_i|\to0$ 取极限，得
\[
  \left|\int_{-\pi}^{\pi}\sin nx\,df(x)\right|
  \le\max_{-\pi\le x\le\pi}|\sin nx|\cdot\bigvee_{-\pi}^{\pi}f.
\]}

\subsection*{三、求 Fourier 展开式}

\noindent\textbf{a.\quad 求 Fourier 展开式的基本方法}

\noindent\textbf{要点}\quad （1）将 $[a,b]$ 上可积函数展开为 Fourier 级数，最基本的方法是

i）按系数公式计算系数：
\[
  a_n=\frac1l\int_a^b f(x)\cos\frac{n\pi x}{l}\,dx,\quad n=0,1,2,\cdots,
\]
\[
  b_n=\frac1l\int_a^b f(x)\sin\frac{n\pi x}{l}\,dx,\quad n=1,2,\cdots,
\]
其中 $l=\dfrac{b-a}{2}$.

ii）将算出的系数代入级数：
\[
  f(x)\sim\frac{a_0}2+\sum_{n=1}^{\infty}a_n\cos\frac{n\pi x}{l}
  +b_n\sin\frac{n\pi x}{l}.
\]
若 $f'(x)$ 在 $[a,b]$ 上存在并连续，$f'(x)$ 的 Fourier 系数记为 $a'_n,b'_n$，那么
\[
  a'_0=0,\quad a'_n=\frac{n\pi}{l}b_n\quad(\text{当 }f(a)=f(b)\text{ 时})\quad(n=1,2,\cdots),
\]
\[
  b'_n=-\frac{n\pi}{l}a_n\quad(n=1,2,\cdots).
\]

iii）根据收敛定理，判定“$\sim$”可改为等号的范围. 若 $f(x)$ 在 $[a,b]$ 上分段光滑\footnote{分段光滑意指：可将区间 $[a,b]$ 分成有限段，在每个小区间内部函数有连续的导数，在小区间端点函数 $f$ 与 $f'$ 有单侧极限.}，则级数的和函数
\[
S(x)=
\begin{cases}
f(x),&\text{当 }x\in(a,b)\text{ 为 }f(x)\text{ 的连续点时},\\[3pt]
\dfrac{f(x-0)+f(x+0)}2,&\text{当 }x\in(a,b)\text{ 为 }f(x)\text{ 的间断点时},\\[6pt]
\dfrac12[f(a+0)+f(b-0)],&\text{当 }x=a,b\text{ 时},\\[6pt]
\text{呈周期},&\text{其他}.
\end{cases}
\]

特别，若 $f(x)$ 以 $2l$ 为周期，或只在 $[-l,l]$ 上有定义，则在上面系数公式里应取区间 $[a,b]=[-l,l]$. 此时，若 $f(x)$ 为奇（偶）函数，则 $a_n=0\ (b_n=0)$；此外还不难验证：若 $f(x)$ 为奇（偶）函数且 $(0,l)$ 上 $f(l-x)=f(x)$，则 $a_n=b_{2n}=0\ (b_n=a_{2n+1}=0)$；若 $f(x)$ 为奇（偶）函数且 $(0,l)$ 上 $f(l-x)=-f(x)$，则 $a_n=b_{2n+1}=0\ (b_n=a_{2n}=0)$.

值得注意的是，可积函数在指定区间上的 Fourier 展开式是唯一的，而三角展开式是随意的\footnote{其系数不要求是此函数在此区间上的 Fourier 系数.}. 将 $f(x)$ 以不同的方式延拓到比 $[a,b]$ 大的区间上，在较大区间上求延拓后的 Fourier 展开式，就可得 $f(x)$ 在 $[a,b]$ 上不同的三角展开式. 例如 $[0,l]$ 上给定的函数 $f(x)$，若将 $f(x)$ 奇延拓到 $[-l,0)$ 上，便可获得 $a_n=0$（即只含正弦）的展开式；若将 $f(x)$ 偶延拓到 $[-l,0)$ 上，便可获得 $b_n=0$（即只含余弦项）的展开式.

（2）由 Fourier 级数的定义及积分的性质易知 Fourier 级数具有可加性：两函数之和的 Fourier 级数等于它们的 Fourier 级数之和（同类项合并）.

（3）由 Fourier 级数的定义及正弦、余弦函数系的正交性易知，三角多项式
\[
  \frac{a_0}2+\sum_{k=1}^{n}(a_k\cos kx+b_k\sin kx)
\]
的 Fourier 级数是它本身.

（4）若 $f(x)$ 在 $[-\pi,\pi]$ 上正常可积，或有奇点但绝对可积，有如下的 Fourier 级数：
\[
  f(x)\sim\frac{a_0}2+\sum_{n=1}^{\infty}(a_n\cos nx+b_n\sin nx),
\]
则不论此级数是否收敛，若收敛，也不管是否收敛于 $f(x)$，恒可逐项积分：
\[
  \int_0^x\left(f(t)-\frac{a_0}2\right)dt
  =\sum_{n=1}^{\infty}\int_0^x(a_n\cos nt+b_n\sin nt)\,dt,\quad x\in[-\pi,\pi],
\]
并且此式必是函数 $\varphi(x)\equiv\displaystyle\int_0^x\left(f(t)-\dfrac{a_0}2\right)dt$ 在 $[-\pi,\pi]$ 上的 Fourier 展开式.

（5）若 $f(x)$ 连续、分段光滑，$f(\pi)=f(-\pi)$，有 Fourier 展开式：
\[
  f(x)=\frac{a_0}2+\sum_{n=1}^{\infty}(a_n\cos nx+b_n\sin nx),\quad x\in[-\pi,\pi].
\]
则逐项求导之后，便得到 $f'(x)$ 的 Fourier 级数
\[
  f'(x)\sim\sum_{n=1}^{\infty}(a_n\cos nx+b_n\sin nx)'.
\]
若附加 $f'(x)$ 分段光滑的条件，则 $f'$ 的 Fourier 级数收敛于
\[
  \frac{f'(x+0)+f'(x-0)}2,\quad x\in(-\pi,\pi).
\]
若再加上 $f''(x)$ 连续的条件，则得 $f'(x)$ 的 Fourier 展开式：
\[
  f'(x)=\sum_{n=1}^{\infty}(a_n\cos nx+b_n\sin nx)',\quad x\in(-\pi,\pi).
\]
若 $f'(-\pi)=f'(\pi)$，则此展开式对于 $\pm\pi$ 也成立.

\noindent\(\star\)\textbf{例 5.4.7}\quad 设 $f(x)=\pi-x,\ x\in(0,\pi)$.

1）将 $f(x)$ 展开为正弦级数；

2）写出和函数的表达式，绘出和函数图形；

3）该级数在 $(0,\pi)$ 内是否一致收敛.（厦门大学）

\noindent\textbf{解}\quad 1）将 $f(x)$ 作奇延拓到 $[-\pi,0]$ 上，求延拓后的函数在 $[-\pi,\pi]$ 上的 Fourier 级数. 这时
\[
  a_n=0\quad(n=0,1,2,\cdots),
\]
\[
  b_n=\frac2\pi\int_0^\pi(\pi-x)\sin nx\,dx
  =-\frac2{n\pi}(\pi-x)\cos nx\bigg|_0^\pi
  -\frac2{n\pi}\int_0^\pi\cos nx\,dx
  =\frac2n.
\]
因延拓后的函数分段光滑，根据收敛定理，
\[
  f(x)=\sum_{n=1}^{\infty}\frac2n\sin nx,\quad\text{当 }x\in(0,\pi)\text{ 时}.
\]

2）级数和函数
\[
\sum_{n=1}^{\infty}\frac2n\sin nx=
\begin{cases}
\pi-x,&\text{在 }(0,\pi)\text{ 内},\\
-\pi+x,&\text{在 }(-\pi,0)\text{ 内},\\
0,&\text{当 }x=0,\ \pm\pi\text{ 时},\\
\text{呈周期},&\text{在 }[-\pi,\pi]\text{ 外}.
\end{cases}
\]
其图形如图 5.4.1 所示.

\begin{center}
  \includegraphics[width=0.58\textwidth]{5-4-1.png}
\end{center}

3）该级数在 $(0,\pi)$ 内非一致收敛. 因为在区间端点 $x=0,\pi$ 上级数收敛，假若级数在 $(0,\pi)$ 内一致收敛，则级数在 $[0,\pi]$ 上一致收敛，和函数应在 $[0,\pi]$ 上连续. 矛盾.

\noindent\textbf{注}\quad 此例给我们很多启示：

（1）所得的展开式，也是函数 $f(x)=\pi-x$ 在区间 $(0,2\pi)$ 内的 Fourier 展开式（想想为什么？）：
\[
  \pi-x=\sum_{n=1}^{\infty}\frac2n\sin nx,\quad x\in(0,2\pi).
\]

（2）级数 $\displaystyle\sum_{n=1}^{\infty}\dfrac2n\sin nx$，尽管通项连续，但和函数在 $x=2k\pi$ 处不连续（$k\in\mathbb Z$）.

（3）（根据 Dirichlet 判别法）在（不含点 $x=2k\pi\ (k\in\mathbb Z)$ 的）任一内闭区间上，此级数总一致收敛，且和函数 $f(x)=\pi-x$ 有导数. 但是，此级数不能逐项求导，逐项求导后级数 $2\displaystyle\sum_{n=1}^{\infty}\cos nx$ 根本不收敛！

\noindent\textbf{*练习}\quad 设 $f(x)=\displaystyle\sum_{n=1}^{\infty}a_n\cos nx$，$\displaystyle\sum_{n=1}^{\infty}|a_n|$ 收敛，若 $\displaystyle\sum_{n=1}^{\infty}B_n\sin nx$ 是 $f(x)$ 在区间 $[0,\pi]$ 上的正弦级数，试求 $B_n(n\in\mathbb N)$.（华中科技大学）

\noindent\textbf{提示}\quad 参看例 5.4.7.

\noindent\textbf{解}\quad
\[
  B_n=\frac2\pi\int_0^\pi f(x)\sin nx\,dx
  =\frac2\pi\int_0^\pi\sum_{k=1}^{\infty}a_k\cos kx\sin nx\,dx,
\]
其中 $|a_k\cos kx\sin nx|\le |a_k|\ (\forall k,n\in\mathbb N)$. 而 $\displaystyle\sum_{k=1}^{\infty}|a_k|$ 收敛，故级数一致收敛，可逐项积分. 于是
\[
\begin{aligned}
B_n
&=\frac2\pi\sum_{k=1}^{\infty}a_k\int_0^\pi\cos kx\sin nx\,dx\\
&=\frac1\pi\sum_{k=1}^{\infty}a_k\int_0^\pi[\sin(n+k)x+\sin(n-k)x]\,dx\\
&=\frac1\pi\sum_{k=1}^{\infty}a_k\left[
-\frac{\cos(n+k)x}{n+k}-\frac{\cos(n-k)x}{n-k}\right]_0^\pi\\
&=\frac1\pi\sum_{k=1}^{\infty}a_k\left[
\frac{(-1)^{n+k+1}+1}{n+k}
+\frac{(-1)^{n-k+1}+1}{n-k}\right].
\end{aligned}
\]
其中右端方括号里的两个分式的分子：
\[
  (-1)^{n+k+1}+1=(-1)^{n-k+1}+1=0\quad(\text{当 }n,k\text{ 奇偶性相同时}),
\]
\[
  (-1)^{n+k+1}+1=(-1)^{n-k+1}+1=2\quad(\text{当 }n,k\text{ 奇偶性相反时}).
\]
故
\[
B_n=
\begin{cases}
\dfrac2\pi\displaystyle\sum_{i=0}^{\infty}a_{2i+1}
\left(\dfrac1{n+2i+1}+\dfrac1{n-2i-1}\right)
=\dfrac4\pi\displaystyle\sum_{i=0}^{\infty}\dfrac{n}{n^2-(2i+1)^2}a_{2i+1},& n\text{ 为偶数},\\[12pt]
\dfrac2\pi\displaystyle\sum_{i=1}^{\infty}a_{2i}
\left(\dfrac1{n+2i}+\dfrac1{n-2i}\right)
=\dfrac4\pi\displaystyle\sum_{i=1}^{\infty}\dfrac{n}{n^2-4i^2}a_{2i},& n\text{ 为奇数}.
\end{cases}
\]

为了计算简便，有时可以更换区间求展开式.

\noindent\(\star\)\textbf{例 5.4.8}\quad 试将
\[
  f(x)=
  \begin{cases}
  1-x,&0\le x\le2,\\
  x-3,&2<x\le4
  \end{cases}
\]
在 $[0,4]$ 上展开为余弦级数，并证明：
\[
  \sum_{n=0}^{\infty}\frac1{(2n+1)^2}=\frac{\pi^2}{8}.
\]

\noindent\textbf{解}\quad 将 $f(x)$ 作偶延拓到 $[-4,0]$ 上，从图 5.4.2 易知，求延拓后函数的 Fourier 级数，最小周期实为 4. 因此我们只要在 $[-2,2]$ 上求展开式即可. 这时

\begin{center}
  \includegraphics[width=0.42\textwidth]{5-4-2.png}
\end{center}

\[
  b_n=0\quad(n=1,2,\cdots),
\]
\[
  a_0=\frac22\int_0^2(1-x)\,dx=0,
\]
\[
\begin{aligned}
a_n&=\frac22\int_0^2(1-x)\cos\frac{n\pi x}{2}\,dx
=-\frac4{n^2\pi^2}[(-1)^n-1]\\
&=
\begin{cases}
0,&n\text{ 为偶数},\\[3pt]
\dfrac8{n^2\pi^2},&n\text{ 为奇数}.
\end{cases}
\end{aligned}
\]
故
\[
  f(x)=\frac8{\pi^2}\sum_{n=0}^{\infty}\frac1{(2n+1)^2}
  \cos\frac{(2n+1)\pi x}{2},\quad x\in[0,4].
\]
在此式里令 $x=0$，即得
\[
  \sum_{n=0}^{\infty}\frac1{(2n+1)^2}=\frac{\pi^2}{8}.
\]

\noindent\(\star\)\textbf{例 5.4.9}\quad 求
\[
  f(x)=
  \begin{cases}
  x,&0\le x<1,\\
  2-x,&1\le x\le2
  \end{cases}
\]
在 $[0,2]$ 上的 Fourier 展开式.

\noindent\textbf{提示}\quad 求 $|x|$ 在 $(-1,1)$ 内的 Fourier 级数.

\noindent\(\star\)\textbf{例 5.4.10}\quad 求
\[
  f(x)=
  \begin{cases}
  0,&-\pi\le x\le0,\\
  \sin x,&0<x\le\pi
  \end{cases}
\]
在 $[-\pi,\pi]$ 上的 Fourier 展开式.

\noindent\textbf{解}\quad 将 $f(x)$ 改写成
\[
  f(x)=\frac{\sin x+|\sin x|}{2},
\]
$\dfrac12|\sin x|$ 为偶函数，且在 $[0,\pi]$ 上满足 $f(\pi-x)=f(x)$，因此
\[
  \frac12|\sin x|=\frac{a_0}2+\sum_{n=1}^{\infty}a_{2n}\cos2nx,
\]
其中
\[
  a_{2n}=\frac4\pi\int_0^{\pi/2}\frac12\sin x\cos2nx\,dx
  =-\frac2{(4n^2-1)\pi}.
\]
所以
\[
  f(x)=\frac{\sin x}{2}+\frac1\pi+\frac2\pi
  \sum_{n=1}^{\infty}\frac{\cos2nx}{1-4n^2}\quad(0<|x|<\pi).
\]

\noindent\(\star\)\textbf{例 5.4.11}\quad 1）将周期为 $2\pi$ 的函数 $f(x)=\dfrac14x(2\pi-x),\ x\in[0,2\pi]$ 展开为 Fourier 级数，并由此求出 $\displaystyle\sum_{n=1}^{\infty}\dfrac1{n^2}$；

2）通过 Fourier 级数的逐项积分求出 $\displaystyle\sum_{n=1}^{\infty}\dfrac1{n^4}$.（复旦大学）

\noindent\textbf{解}\quad 1）
\[
  a_0=\frac1\pi\int_0^{2\pi}\frac14x(2\pi-x)\,dx=\frac13\pi^2,
\]
\[
  a_n=\frac1\pi\int_0^{2\pi}\frac14x(2\pi-x)\cos nx\,dx=-\frac1{n^2},
\]
\[
  b_n=\frac1\pi\int_0^{2\pi}\frac14x(2\pi-x)\sin nx\,dx=0\quad(n=1,2,\cdots),
\]
且 $f(x)=\dfrac14x(2\pi-x)$ 在 $[0,2\pi]$ 上处处可导，$f(0+0)=f(2\pi-0)$，根据收敛定理，有
\[
  \frac14x(2\pi-x)=\frac16\pi^2-\sum_{n=1}^{\infty}\frac1{n^2}\cos nx\quad(x\in[0,2\pi]). \qquad(1)
\]
令 $x=0$，可得
\[
  \sum_{n=1}^{\infty}\frac1{n^2}=\frac{\pi^2}{6}\quad(\text{此结果多次被用到}).
\]

2）将 (1) 中常数项移至左端，根据 Fourier 级数逐项积分的定理，有
\[
  \int_0^x\left[\frac14t(2\pi-t)-\frac16\pi^2\right]dt
  =-\sum_{n=1}^{\infty}\frac1{n^2}\int_0^x\cos nt\,dt,
\]
即
\[
  \frac16\pi^2x-\frac14\pi x^2+\frac{x^3}{12}
  =\sum_{n=1}^{\infty}\frac1{n^3}\sin nx.
\]
此式为左端函数的 Fourier 展开式.

同理继续逐项积分两次，得
\[
  -\frac1{36}\pi^2x^3+\frac1{48}\pi x^4-\frac{x^5}{240}
  =\sum_{n=1}^{\infty}\frac1{n^4}\left(\frac{\sin nx}{n}-x\right),\quad x\in[0,2\pi].
\]
令 $x=2\pi$，可得
\[
  \sum_{n=1}^{\infty}\frac1{n^4}=\frac1{90}\pi^4.
\]

\noindent\textbf{例 5.4.12}\quad 设 $f(x)$ 是以 $2\pi$ 为周期的连续偶函数，它的 Fourier 级数为
\[
  f(x)\sim\frac{a_0}{2}+\sum_{n=1}^{\infty}a_n\cos nx.
\]
求证：

1）函数
\[
  H(x)=\frac1\pi\int_{-\pi}^{\pi}f(x+t)f(t)\,dt
\]
是以 $2\pi$ 为周期的连续偶函数，它的 Fourier 级数是
\[
  H(x)\sim\frac{a_0^2}{2}+\sum_{n=1}^{\infty}a_n^2\cos nx;
\]

2）$\displaystyle\frac{a_0^2}{2}+\sum_{n=1}^{\infty}a_n^2\cos nx$ 一致收敛于 $H(x)$.（同济大学）

\noindent\textbf{证}\quad 1）$\forall x_0\in(-\infty,+\infty)$，$f(x+t)f(t)$ 作为二元函数在
$x_0-1\le x\le x_0+1$，$-\pi\le t\le\pi$ 上连续，所以
\[
  H(x)=\frac1\pi\int_{-\pi}^{\pi}f(x+t)f(t)\,dt
\]
在 $x_0$ 处连续，即 $H(x)$ 在 $(-\infty,+\infty)$ 内处处连续. 又因
\[
\begin{aligned}
H(-x)&=\frac1\pi\int_{-\pi}^{\pi}f(-x+t)f(t)\,dt
\xlongequal{\text{令 }-x+t=u}
\frac1\pi\int_{-x-\pi}^{-x+\pi}f(u)f(u+x)\,du\\
&=\frac1\pi\int_{-\pi}^{\pi}f(u+x)f(u)\,du=H(x)\quad(\forall x\in\mathbb R),
\end{aligned}
\]
所以 $H(x)$ 为偶函数. 用 $a_n',b_n'$ 表示 $H(x)$ 的 Fourier 系数，则
\[
  b_n'=0\ (n=1,2,\cdots),\quad
  a_0'=\frac1\pi\int_{-\pi}^{\pi}H(x)\,dx
  =\frac1\pi\int_{-\pi}^{\pi}\left(\frac1\pi\int_{-\pi}^{\pi}f(x+t)f(t)\,dt\right)dx.
\]
因 $f(x+t)f(t)$ 作为二元函数，在 $-\pi\le x\le\pi$，$-\pi\le t\le\pi$ 上连续，积分可以交换次序. 因此令 $u=x+t$，则
\[
  \frac1\pi\int_{-\pi}^{\pi}f(x+t)\,dx
  =\frac1\pi\int_{-\pi+t}^{\pi+t}f(u)\,du
  =\frac1\pi\int_{-\pi}^{\pi}f(u)\,du=a_0,
\]
于是
\[
  a_0'=\frac1\pi\int_{-\pi}^{\pi}a_0f(t)\,dt
  =a_0\cdot\frac1\pi\int_{-\pi}^{\pi}f(t)\,dt=a_0^2,
\]
\[
\begin{aligned}
a_n'&=\frac1\pi\int_{-\pi}^{\pi}H(x)\cos nx\,dx
=\frac1\pi\int_{-\pi}^{\pi}\left(\frac1\pi\int_{-\pi}^{\pi}f(x+t)f(t)\,dt\right)\cos nx\,dx\\
&=\frac1\pi\int_{-\pi}^{\pi}\left(\frac1\pi\int_{-\pi}^{\pi}f(x+t)\cos nx\,dx\right)f(t)\,dt.
\end{aligned}
\]
令 $x+t=u$，则
\[
\begin{aligned}
\frac1\pi\int_{-\pi}^{\pi}f(x+t)\cos nx\,dx
&=\frac1\pi\int_{-\pi+t}^{\pi+t}f(u)\cos n(u-t)\,du\\
&=\left(\frac1\pi\int_{-\pi}^{\pi}f(u)\cos nu\,du\right)\cos nt
+\left(\frac1\pi\int_{-\pi}^{\pi}f(u)\sin nu\,du\right)\sin nt\\
&=a_n\cos nt.
\end{aligned}
\]
于是
\[
  a_n'=\frac1\pi\int_{-\pi}^{\pi}a_nf(t)\cos nt\,dt=a_n^2,\quad(n=1,2,\cdots).
\]
故
\[
  H(x)\sim\frac{a_0^2}{2}+\sum_{n=1}^{\infty}a_n^2\cos nx.
\]

2）因 $\displaystyle f(x)\sim\frac{a_0}{2}+\sum_{n=1}^{\infty}a_n\cos nx$，由 Bessel 不等式（见例 5.4.20 之注）知
\[
  \frac{a_0^2}{2}+\sum_{n=1}^{\infty}a_n^2
\]
收敛，于是由 $|a_n^2\cos nx|\le a_n^2,\ \forall x\in\mathbb R$ 知级数
$\displaystyle\frac{a_0^2}{2}+\sum_{n=1}^{\infty}a_n^2\cos nx$ 在 $(-\infty,+\infty)$ 内一致收敛.
记和函数为 $S(x)$，如此 $S(x)$ 连续，且
\[
  S(x)=\frac{a_0^2}{2}+\sum_{n=1}^{\infty}a_n^2\cos nx
\]
是 $S(x)$ 的 Fourier 级数. 即 $H(x),S(x)$ 两连续函数有相同的 Fourier 级数，因此
$H(x)\equiv S(x)$，即
\[
  \frac{a_0^2}{2}+\sum_{n=1}^{\infty}a_n^2\cos nx
  \xlongequal{\text{一致收敛}}H(x).
\]

\noindent\textbf{b.\quad 求 Fourier 展开式的一些其他方法}

\noindent\(\star\)\textbf{例 5.4.13}\quad 已知 $f(x)=x$ 在 $(-\pi,\pi)$ 上的 Fourier 展开式为
\[
  x=\sum_{n=1}^{\infty}\frac{2(-1)^{n-1}}{n}\sin nx,\quad |x|<\pi. \qquad(1)
\]
试求函数 $\varphi(x)=x\sin x$ 在 $[-\pi,\pi]$ 上的 Fourier 展开式.

\noindent\textbf{解}\quad 利用三角公式，由 (1) 得
\[
\begin{aligned}
x\sin x
&=\sum_{n=1}^{\infty}\frac{(-1)^{n-1}}{n}2\sin nx\sin x\\
&=1-\frac{\cos x}{2}+2\sum_{n=2}^{\infty}\frac{(-1)^{n-1}}{n^2-1}\cos nx\quad(|x|<\pi).
\end{aligned}
\]
因所得级数当 $|x|<\pi$ 时一致收敛，故为 Fourier 展开式. 又因
$x\sin x\big|_{-\pi}^{\pi}=0$，所以展开式在 $x=\pm\pi$ 时亦成立.

值得注意的是：上例使用的方法，具有普遍性. 我们有

\noindent\textbf{例 5.4.14}\quad 设 $f(x)$ 有如下的 Fourier 级数：
\[
  f(x)\sim\frac{a_0}{2}+\sum_{n=1}^{\infty}a_n\cos nx+b_n\sin nx.
\]
试证：将级数逐项乘以 $\sin x$，可得 $f(x)\sin x$ 的 Fourier 级数：
\[
  f(x)\sin x\sim\frac{b_1}{2}
  +\sum_{n=1}^{\infty}\left(\frac{b_{n+1}-b_{n-1}}2\cos nx
  +\frac{a_{n-1}-a_{n+1}}2\sin nx\right).
\]

\noindent\textbf{证}\quad 逐项乘 $\sin x$ 之后，利用三角公式可得
\[
\begin{aligned}
&\frac{a_0}{2}\sin x+\sum_{n=1}^{\infty}a_n\frac12[\sin(n+1)x-\sin(n-1)x]\\
&\quad+\sum_{n=1}^{\infty}b_n\frac12[\cos(n-1)x-\cos(n+1)x]\\
={}&\frac{a_0}{2}\sin x+\frac12\sum_{n=2}^{\infty}a_{n-1}\sin nx
-\frac12\sum_{n=1}^{\infty}a_{n+1}\sin nx\\
&+\frac12\sum_{n=2}^{\infty}b_{n-1}\cos nx
+\frac12\sum_{n=1}^{\infty}b_{n+1}\cos nx+\frac{b_1}{2}\quad(\text{记 }b_0=0)\\
={}&\frac12b_1+\sum_{n=1}^{\infty}\left(\frac{b_{n+1}-b_{n-1}}2\cos nx
+\frac{a_{n-1}-a_{n+1}}2\sin nx\right).
\end{aligned}
\]
此级数实为 $f(x)\sin x$ 在 $[-\pi,\pi]$ 上的 Fourier 级数.

因为若用 $\alpha_n,\beta_n$ 表示 $f(x)\sin x$ 的 Fourier 系数，则
\[
  \alpha_0=\frac1\pi\int_{-\pi}^{\pi}f(x)\sin x\,dx=b_1,
\]
\[
\begin{aligned}
\alpha_n&=\frac1\pi\int_{-\pi}^{\pi}f(x)\sin x\cos nx\,dx\\
&=\frac1{2\pi}\int_{-\pi}^{\pi}f(x)\sin(n+1)x\,dx
-\frac1{2\pi}\int_{-\pi}^{\pi}f(x)\sin(n-1)x\,dx\\
&=\frac12(b_{n+1}-b_{n-1}),\quad n=1,2,\cdots.
\end{aligned}
\]
同理
\[
  \beta_n=\frac12(a_{n-1}-a_{n+1}),\quad n=1,2,\cdots.
\]

\noindent\textbf{例 5.4.15}\quad 求函数
\[
  f(x)\equiv\sum_{n=1}^{\infty}\beta^n\frac{\sin nx}{\sin x} \qquad(1)
\]
的 Fourier 展开式（$|\beta|<1$）.

\noindent\textbf{解 I}\quad 因 $0<|x|\le\dfrac{\pi}{2}$ 时，
\[
  \left|\beta^n\frac{\sin nx}{\sin x}\right|\le |\beta|^n\frac{2n}{\pi},
\]
且由比式判别法知 $\displaystyle\sum_{n=1}^{\infty}|\beta|^n\dfrac{2n}{\pi}$ 收敛，故级数 (1) 在
$\left[-\dfrac{\pi}{2},\dfrac{\pi}{2}\right]\setminus\{0\}$ 上一致收敛，从而 $f(x)$ 在其上连续，在 $x=0$ 处为可去间断.

类似可知，$f(x)$ 除在 $x=k\pi(k=0,\pm1,\pm2,\cdots)$ 处可去间断之外，其余处处连续，$f(x)$ 在 $[-\pi,\pi]$ 上可积. 由式 (1)，
\[
  f(x)\sin x=\sum_{n=1}^{\infty}\beta^n\sin nx. \qquad(2)
\]
显然，级数 (2) 在 $(-\infty,+\infty)$ 上一致收敛. 以 $a_n,b_n$ 表示 $f(x)$ 的 Fourier 系数，利用上例结果，有
$\dfrac12b_1=0$，$\dfrac12(b_{n+1}-b_{n-1})=0$，从而可知 $b_n=0\ (n=1,2,\cdots)$. 又
\[
  \frac12(a_{n-1}-a_{n+1})=\beta^n \qquad(3)
\]
（这是 $\{a_n\}$ 中间隔一项的递推关系，因此下面只要求出 $a_0$，便可依次求出
$a_2,a_4,\cdots,a_{2n},\cdots$；若求出了 $a_1$，便可求出 $a_3,a_5,\cdots,a_{2n+1},\cdots$），从而
$a_{2n-2}-a_{2n}=2\beta^{2n-1}$，于是
\[
  a_0=\sum_{n=1}^{\infty}(a_{2n-2}-a_{2n})
  =2\sum_{n=1}^{\infty}\beta^{2n-1}=\frac{2\beta}{1-\beta^2}. \qquad(4)
\]
同理，由 $a_{2n-1}-a_{2n+1}=2\beta^{2n}$ 得
\[
  a_1=\frac{2\beta^2}{1-\beta^2}. \qquad(5)
\]
由式 (3)、(4)、(5) 可得
\[
  a_n=\frac{2\beta^{n+1}}{1-\beta^2},\quad n=0,1,2,\cdots.
\]
故
\[
  f(x)\sim\frac{\beta}{1-\beta^2}
  +\frac{2\beta}{1-\beta^2}\sum_{n=1}^{\infty}\beta^n\cos nx\quad(-\infty<x<+\infty).
\]
上面级数一致收敛，故为 $f(x)$ 的 Fourier 展开式，``$\sim$'' 可改写成 ``$=$''.

\noindent\textbf{解 II}\quad （复数解法）若 $z$ 为复数，令 $\operatorname{Im}z$ 表示 $z$ 的虚部. 于是
\[
\begin{aligned}
f(x)&=\frac1{\sin x}\sum_{n=1}^{\infty}\beta^n\sin nx
=\frac1{\sin x}\operatorname{Im}\left[\sum_{n=1}^{\infty}(\beta e^{ix})^n\right]\\
&=\frac1{\sin x}\operatorname{Im}\left(\frac{\beta e^{ix}}{1-\beta e^{ix}}\right)
=\frac{\beta}{\sin x}\cdot
\frac{\operatorname{Im}[e^{ix}(1-\beta e^{-ix})]}
{(1-\beta e^{ix})(1-\beta e^{-ix})}\\
&=\frac{\beta}{(1-\beta e^{ix})(1-\beta e^{-ix})}\\
&=\frac{1-\beta e^{ix}+1-\beta e^{-ix}-(1-\beta e^{ix}-\beta e^{-ix}+\beta^2)}
{(1-\beta e^{ix})(1-\beta e^{-ix})}\cdot\frac{\beta}{1-\beta^2}\\
&=\frac{\beta}{1-\beta^2}\left(\frac1{1-\beta e^{-ix}}+\frac1{1-\beta e^{ix}}-1\right)\\
&=\frac{\beta}{1-\beta^2}\left[\sum_{n=0}^{\infty}\beta^n(e^{-inx}+e^{inx})-1\right]\\
&=\frac{\beta}{1-\beta^2}+\frac{2\beta}{1-\beta^2}\sum_{n=1}^{\infty}\beta^n\cos nx.
\end{aligned}
\]
因上面级数在 $[-\pi,\pi]$ 上一致收敛，故为 Fourier 展开式.

\noindent\textbf{注}\quad 本题如直接用系数公式计算积分，也可以做出，但十分麻烦.

\noindent\textbf{例 5.4.16}\quad 求
\[
  f(x)=\frac{q\sin x}{1-2q\cos x+q^2}\quad(|q|<1)
\]
的 Fourier 展开式.

\noindent\textbf{解}\quad （复数解法）令 $z=e^{ix}$，$\bar z=e^{-ix}$，则
\[
  \sin x=\frac1{2i}(z-\bar z),\quad \cos x=\frac12(z+\bar z).
\]
于是
\[
\begin{aligned}
f(x)&=\frac1{2i}\left(\frac1{1-qz}-\frac1{1-q\bar z}\right)
=\frac1{2i}\left[\sum_{n=0}^{\infty}(qz)^n-\sum_{n=0}^{\infty}(q\bar z)^n\right]\\
&=\sum_{n=0}^{\infty}\frac{z^n-\bar z^n}{2i}\cdot q^n
=\sum_{n=0}^{\infty}q^n\sin nx\quad(-\infty<x<+\infty).
\end{aligned}
\]

\noindent\textbf{例 5.4.17}\quad 求 $f(x)=\ln(1-2q\cos x+q^2)$（$|q|<1$）的 Fourier 展开式，并计算积分
\[
  \int_0^{\pi}\ln(1-2\alpha\cos x+\alpha^2)\,dx\quad(-\infty<\alpha<+\infty).
\]

\noindent\textbf{提示}\quad 可用复数解法，或先求 $f'(x)$ 的 Fourier 级数（利用上题结果），然后逐项积分. 最后，
\[
  \int_0^{\pi}\ln(1-2\alpha\cos x+\alpha^2)\,dx
  =
  \begin{cases}
  0,&|\alpha|\le1,\\
  2\pi\ln|\alpha|,&|\alpha|>1.
  \end{cases}
\]

\noindent\textbf{四、综合性问题}

\noindent\textbf{例 5.4.18}\quad 设 $f(x)$ 在 $[a,b]$ 上可积，试用 Fourier 展开式证明：
\[
  \lim_{n\to\infty}\int_a^b f(x)|\sin nx|\,dx
  =\frac2\pi\int_a^b f(x)\,dx.
\]

\noindent\textbf{方法}\quad 先将 $|\sin nx|$ 展开成（一致收敛的）Fourier 级数，逐项乘 $f(x)$，再逐项积分，逐项取极限. 利用一般书中的 Riemann 引理\footnote{若 $f(x)$ 在区间 $[a,b]$ 上可积，或无界但绝对可积，则
\[
  \lim_{p\to+\infty}\int_a^b f(x)\sin px\,dx=0,\quad
  \lim_{p\to+\infty}\int_a^b f(x)\cos px\,dx=0.
\]
}，便知极限为 $\dfrac2\pi\displaystyle\int_a^b f(x)\,dx$.

\noindent\textbf{证}\quad $1^\circ$\quad $|\sin x|$ 有一致收敛的 Fourier 展开式：
\[
  |\sin x|=\frac2\pi-\frac4\pi\sum_{k=1}^{\infty}\frac{\cos 2kx}{(2k)^2-1},\quad x\in(-\infty,+\infty).
\]
因此
\[
  |\sin nx|=\frac2\pi-\frac4\pi\sum_{k=1}^{\infty}\frac{\cos 2knx}{(2k)^2-1},\quad x\in[-\pi,\pi]. \qquad(1)
\]

$2^\circ$\quad 用有界函数遍乘一致收敛级数的各项，所得级数仍一致收敛. 故
\[
  f(x)|\sin nx|
  =\frac2\pi f(x)-\frac4\pi\sum_{k=1}^{\infty}f(x)\frac{\cos 2knx}{4k^2-1},\quad x\in[-\pi,\pi]
\]
一致收敛，可逐项积分：
\[
  \int_a^b f(x)|\sin nx|\,dx
  =\frac2\pi\int_a^b f(x)\,dx
  -\frac4\pi\sum_{k=1}^{\infty}\int_a^b\frac{f(x)\cos 2knx}{4k^2-1}\,dx. \qquad(2)
\]

$3^\circ$\quad 设 $|f(x)|\le M\ (\forall x\in[a,b])$，级数 (2) 的通项
\[
  \left|\int_a^b\frac{f(x)\cos 2knx}{4k^2-1}\,dx\right|
  \le\frac{M\cdot(b-a)}{4k^2-1},
\]
且 $\displaystyle\sum_{k=1}^{\infty}\frac{M(b-a)}{4k^2-1}$ 收敛，因而级数 (2) 一致收敛（关于 $n$）.

$4^\circ$\quad 在级数 (2) 中，令 $n\to\infty$，逐项取极限，并利用 Riemann 引理，得
\[
\begin{aligned}
\lim_{n\to\infty}\int_a^b f(x)|\sin nx|\,dx
&=\frac2\pi\int_a^b f(x)\,dx
-\frac4\pi\sum_{k=1}^{\infty}\lim_{n\to\infty}
\frac{\displaystyle\int_a^b f(x)\cos 2knx\,dx}{4k^2-1}\\
&=\frac2\pi\int_a^b f(x)\,dx.
\end{aligned}
\]

\noindent\textbf{注}\quad 请跟例 4.1.9 及例 4.1.10 进行比较.

由 Fourier 展开式可得到别的展开式. 如

\noindent\textbf{例 5.4.19}\quad 已知 $f(x)=\cos\alpha x(\alpha\ne\text{整数})$ 在 $(-\pi,\pi)$ 内的 Fourier 展开式为
\[
  \cos\alpha x=\frac2\pi\sin\alpha\pi
  \left[\frac1{2\alpha}+\sum_{n=1}^{\infty}(-1)^n\frac{\alpha\cos nx}{\alpha^2-n^2}\right], \qquad(1)
\]
试证下列展开式成立：
\[
  \frac1{\sin x}=\frac1x+\sum_{n=1}^{\infty}(-1)^n\frac{2x}{x^2-n^2\pi^2}, \qquad(2)
\]
\[
  \cot x=\frac1x+\sum_{n=1}^{\infty}\left(\frac1{x-n\pi}+\frac1{x+n\pi}\right), \qquad(3)
\]
\[
  \tan x=-\sum_{n=1}^{\infty}\left(\frac1{x-\dfrac{2n+1}{2}\pi}
  +\frac1{x+\dfrac{2n-1}{2}\pi}\right)-\frac2{2x-\pi}. \qquad(4)
\]

\noindent\textbf{解}\quad 将式 (1) 改写成
\[
  \frac{\pi}{2}\frac{\cos\alpha x}{\sin\alpha\pi}
  =\frac1{2\alpha}+\sum_{n=1}^{\infty}(-1)^n\frac{\alpha\cos nx}{\alpha^2-n^2}. \qquad(5)
\]
令 $x=0$，再把 $\alpha\pi$ 改写成 $x$，即可得式 (2).

在式 (5) 中令 $x=\pi$，再将 $\alpha\pi$ 改写成 $x$，便可得式 (3). 最后，利用
$\tan x=-\cot\left(x-\dfrac{\pi}{2}\right)$ 可得式 (4). 解毕.

下例表明，在均方意义下，采用 Fourier 系数，可使三角多项式逼近达到最佳.

\noindent\(\star\)\textbf{例 5.4.20}\quad 假设函数 $f(x)$ 在闭区间 $[-\pi,\pi]$ 上可积，$T_n(x)$ 为三角多项式：
\[
  T_n(x)=\frac{\alpha_0}{2}+\sum_{k=1}^{n}(\alpha_k\cos kx+\beta_k\sin kx), \qquad(1)
\]
试求系数 $\alpha_k,\beta_k$ 使均方误差
\[
  \delta_n=\frac1{2\pi}\int_{-\pi}^{\pi}|f(x)-T_n(x)|^2\,dx \qquad(2)
\]
最小.（中山大学，中国人民大学）

\noindent\textbf{解}\quad 将式 (1) 代入 (2)，乘开并利用三角函数的正交性，可得
\[
\begin{aligned}
0\le\delta_n
&=\frac1{2\pi}\int_{-\pi}^{\pi}|f(x)-T_n(x)|^2\,dx\\
&=\frac1{2\pi}\int_{-\pi}^{\pi}
\left[f(x)-\frac{\alpha_0}{2}-\sum_{k=1}^{n}(\alpha_k\cos kx+\beta_k\sin kx)\right]^2dx\\
&=\frac1{2\pi}\int_{-\pi}^{\pi}f^2(x)\,dx
-\frac12\left[\frac{a_0^2}{2}+\sum_{k=1}^{n}(a_k^2+b_k^2)\right]\\
&\quad+\frac12\left[\frac12(\alpha_0-a_0)^2
+\sum_{k=1}^{n}(\alpha_k-a_k)^2+(\beta_k-b_k)^2\right],
\end{aligned}
\]
其中 $a_k,b_k$ 是 $f$ 的 Fourier 系数. 可见当
$\alpha_k=a_k(k=0,1,2,\cdots)$，$\beta_k=b_k(k=1,2,\cdots)$ 时
$\delta_n$ 最小，此时
\[
  \delta_n=\hat\delta_n=\frac1{2\pi}\int_{-\pi}^{\pi}f^2(x)\,dx
  -\frac12\left[\frac{a_0^2}{2}+\sum_{k=1}^{n}(a_k^2+b_k^2)\right]. \qquad(3)
\]

\noindent\textbf{注}\quad 根据 $\delta_n$ 的定义式 (2)，恒有 $\delta_n\ge0$. 因此利用本例的结果，我们实际也证明了：若 $f(x)$ 在 $[-\pi,\pi]$ 上可积，则 Bessel 不等式
\[
  \frac1\pi\int_{-\pi}^{\pi}f^2(x)\,dx\ge
  \frac{a_0^2}{2}+\sum_{n=1}^{\infty}(a_n^2+b_n^2)
\]
成立，其中 $a_n,b_n$ 是 $f(x)$ 在 $[-\pi,\pi]$ 上的 Fourier 系数. 另外式 (3) 表明当采用 Fourier 系数时，
$\delta_n=\hat\delta_n\downarrow$（当 $n$ 增大时），利用这一点，容易证明新的结论.

\noindent\textbf{例 5.4.21}\quad 设 $f(x)$ 在 $[-\pi,\pi]$ 上连续，则有 Parseval 等式成立（用 $a_n,b_n$ 表示 $f(x)$ 在 $[-\pi,\pi]$ 上的 Fourier 系数）：
\[
  \frac1\pi\int_{-\pi}^{\pi}f^2(x)\,dx
  =\frac{a_0^2}{2}+\sum_{n=1}^{\infty}(a_n^2+b_n^2). \qquad(1)
\]

\noindent\textbf{证}\quad 根据 Weierstrass 逼近定理，$\forall\varepsilon>0$，存在三角多项式
$T_N(x)$，使得
\[
  |f(x)-T_N(x)|<\varepsilon^{1/2}\quad(\forall x\in[-\pi,\pi]).
\]
从而
\[
  \delta_N=\frac1{2\pi}\int_{-\pi}^{\pi}|f(x)-T_N(x)|^2\,dx<\varepsilon.
\]
根据上例，$T_N(x)$ 的系数采用 Fourier 系数时，$\delta_N$ 达到最小，且这时的
$\delta_N=\hat\delta_N$. 取上例式 (3) 的形式，当 $n$ 增大时，$\hat\delta_n\downarrow$. 故 $n\ge N$ 时，有
\[
  0\le\frac1{2\pi}\int_{-\pi}^{\pi}f^2(x)\,dx
  -\frac12\left[\frac{a_0^2}{2}+\sum_{k=1}^{n}(a_k^2+b_k^2)\right]<\varepsilon.
\]
式 (1) 获证.

\noindent\textbf{注}\quad 可以证明，只要 $f(x)$ 在 $[-\pi,\pi]$ 上正常可积，或者 $[-\pi,\pi]$ 上有奇点，但
$f(x)$ 平方可积，则 Parseval 等式仍保持成立.

下面讨论 Parseval 等式的若干应用.

\noindent\textbf{例 5.4.22}\quad 试利用 Parseval 等式证明：

1）若 $f$ 与 $g$ 在 $[-\pi,\pi]$ 上符合 Parseval 等式的条件，它们的 Fourier 系数分别为
$a_n,b_n$ 与 $\alpha_n,\beta_n$，试证：
\[
  \frac1\pi\int_{-\pi}^{\pi}f(x)g(x)\,dx
  =\frac{a_0\alpha_0}{2}+\sum_{n=1}^{\infty}(a_n\alpha_n+b_n\beta_n).
\]

2）就 $[-\pi,\pi]$ 上全体连续函数而言，三角函数系
\[
  \{1,\cos x,\sin x,\cdots,\cos nx,\sin nx,\cdots\}
\]
为完全的. 即：$[-\pi,\pi]$ 上任意连续函数 $f(x)$，若与三角函数系的每个元正交，则
$f(x)\equiv0$ 于 $[-\pi,\pi]$ 上.

\noindent\textbf{证}\quad 1）写出 $f(x)+g(x)$ 与 $f(x)-g(x)$ 的 Parseval 等式，相减即得.

2）由已知条件可知 $a_n=b_n=0$，从而利用 Parseval 等式，知
$\displaystyle\int_{-\pi}^{\pi}f^2(x)\,dx=0$. 因 $f(x)$ 连续，故 $f(x)\equiv0$.

\noindent\textbf{例 5.4.23}\quad 设函数 $f(x)$ 在区间 $[0,2\pi]$ 上可积，$\varphi(x)$ 在 $[0,2\pi]$ 上连续，且在
$(0,2\pi)$ 内可展开为它的 Fourier 级数：
\[
  \varphi(x)=\frac{a_0}{2}+\sum_{k=1}^{\infty}(a_k\cos kx+b_k\sin kx),\quad x\in(0,2\pi).
\]
试证：
\[
  f(x)\varphi(x)=\frac{a_0}{2}f(x)+
  \sum_{k=1}^{\infty}(a_kf(x)\cos kx+b_kf(x)\sin kx) \qquad(1)
\]
可在 $[0,2\pi]$ 上逐项积分.

\noindent\textbf{分析}\quad 要证明
\[
  \int_0^{2\pi}f(x)\varphi(x)\,dx
  =\int_0^{2\pi}\frac{a_0}{2}f(x)\,dx+
  \sum_{k=1}^{\infty}\int_0^{2\pi}(a_kf(x)\cos kx+b_kf(x)\sin kx)\,dx,
\]
即要证明：
\[
  \int_0^{2\pi}f(x)\varphi(x)\,dx-\int_0^{2\pi}\frac{a_0}{2}f(x)\,dx
  =\lim_{n\to\infty}\sum_{k=1}^{n}\int_0^{2\pi}(a_kf(x)\cos kx+b_kf(x)\sin kx)\,dx.
\]
记
\[
  S_n(x)=\frac{a_0}{2}+\sum_{k=1}^{n}(a_k\cos kx+b_k\sin kx),
\]
即要证明
\[
  \int_0^{2\pi}f(x)(\varphi(x)-S_n(x))\,dx\to0\quad(n\to\infty).
\]
根据 Cauchy 不等式，
\[
  \left|\int_0^{2\pi}f(x)(\varphi(x)-S_n(x))\,dx\right|
  \le\left[\int_0^{2\pi}f^2(x)\,dx\cdot
  \int_0^{2\pi}(\varphi(x)-S_n(x))^2\,dx\right]^{1/2}.
\]
由此式可知，问题进一步归结为证明：
\[
  \int_0^{2\pi}(\varphi(x)-S_n(x))^2\,dx\to0\quad(n\to\infty).
\]
用 Fourier 系数公式，不难验证有
\[
  \frac1{2\pi}\int_0^{2\pi}|\varphi(x)-S_n(x)|^2\,dx
  =\frac1{2\pi}\int_0^{2\pi}\varphi^2(x)\,dx
  -\frac12\left[\frac{a_0^2}{2}+\sum_{k=1}^{n}(a_k^2+b_k^2)\right].
\]
由于 $\varphi(x)$ 连续，根据 Parseval 等式，上式右端当 $n\to\infty$ 时趋于零，问题获证.

\noindent\textbf{注}\quad 这里我们看到，若 $\varphi(x)$ 连续，则
\[
  \frac1{2\pi}\int_0^{2\pi}(\varphi(x)-S_n(x))^2\,dx\to0\quad(n\to\infty),
\]
即 $S_n(x)$ 在平方平均意义下收敛于 $\varphi(x)$.

\noindent\(\star\)\textbf{例 5.4.24}\quad 设函数 $f(x)$ 在区间 $[0,2\pi]$ 上可积，证明：
\[
  \frac1{2\pi}\int_0^{2\pi}f(x)(\pi-x)\,dx=\sum_{n=1}^{\infty}\frac{b_n}{n}, \qquad(1)
\]
其中
\[
  b_n=\frac1\pi\int_0^{2\pi}f(x)\sin nx\,dx\quad(n=1,2,\cdots).
\]
（南京航空学院）

\noindent\textbf{证 I}\quad $\varphi(x)=\pi-x$ 在 $[0,2\pi]$ 上连续可微，按收敛定理可得 Fourier 展开式：
\[
  \pi-x=\sum_{n=1}^{\infty}\frac2n\sin nx,\quad x\in(0,2\pi).
\]
根据上题，可在 $[0,2\pi]$ 上对
\[
  f(x)(\pi-x)=\sum_{n=1}^{\infty}\frac2n f(x)\sin nx
\]
逐项积分，于是式 (1) 获证.

\noindent\textbf{证 II}\quad 用 $a_n,b_n;\alpha_n,\beta_n$ 分别表示两个可积函数 $f(x),\varphi(x)$ 的 Fourier 系数，则已知有关系（例 5.4.22）：
\[
  \frac1\pi\int_0^{2\pi}f(x)\varphi(x)\,dx
  =\frac{a_0\alpha_0}{2}+\sum_{n=1}^{\infty}(a_n\alpha_n+b_n\beta_n).
\]
令 $\varphi(x)=\pi-x$，可得 $\alpha_n=0,\ \beta_n=\dfrac2n$，这便得到式 (1).

\noindent\textbf{例 5.4.25}\quad 设 $f(x)$ 是以 $2\pi$ 为周期的连续函数，
\[
  V_n(x)=\frac{(2n)!!}{2\pi(2n-1)!!}\int_{-\pi}^{\pi}f(t)\cos^{2n}\frac{t-x}{2}\,dt,
\]
证明：$V_n(x)\rightrightarrows f(x)$（当 $n\to\infty$ 时）于 $[-\pi,\pi]$ 上.

\noindent\textbf{方法}\quad 与收敛定理的证法类似.

\noindent\textbf{证}\quad 令 $t-x=u$，则
\[
\begin{aligned}
V_n(x)&=\frac{(2n)!!}{2\pi(2n-1)!!}
\int_{-\pi-x}^{\pi-x}f(x+u)\cos^{2n}\frac u2\,du\\
&=\frac{(2n)!!}{2\pi(2n-1)!!}
\int_{-\pi}^{\pi}f(x+u)\cos^{2n}\frac u2\,du\\
&=\frac{(2n)!!}{2\pi(2n-1)!!}
\int_0^{\pi}[f(x+u)+f(x-u)]\cos^{2n}\frac u2\,du.
\end{aligned}
\]
利用 Wallis 公式或分部积分可得
\[
  \int_0^{\pi}\cos^{2n}\frac x2\,dx=\frac{(2n-1)!!}{(2n)!!}\pi.
\]
因此有单位分解：
\[
  1=\frac{(2n)!!}{\pi(2n-1)!!}\int_0^{\pi}\cos^{2n}\frac u2\,du,
\]
故
\[
  f(x)=\frac{(2n)!!}{2\pi(2n-1)!!}
  \int_0^{\pi}2f(x)\cos^{2n}\frac u2\,du,
\]
从而
\[
  |V_n(x)-f(x)|
  =\frac{(2n)!!}{2\pi(2n-1)!!}
  \left|\int_0^{\pi}[f(x+u)+f(x-u)-2f(x)]\cos^{2n}\frac u2\,du\right|. \qquad(1)
\]
记 $\varphi(x,u)=f(x+u)+f(x-u)-2f(x)$. 因 $f(x)$ 在 $[-\pi,2\pi]$ 上连续，所以一致连续. 故
$\forall\varepsilon>0$，$\exists\delta>0$（$\delta<\pi$），使得当 $|x'-x''|<\delta$ 时有
$|f(x')-f(x'')|<\dfrac{\varepsilon}{4}$. 于是当 $x\in[0,\pi]$，$0<|u|<\delta$ 时，有
\[
  |\varphi(x,u)|\le |f(x+u)-f(x)|+|f(x-u)-f(x)|<\frac{\varepsilon}{2}.
\]
由此
\[
\begin{aligned}
&\frac{(2n)!!}{2\pi(2n-1)!!}
\left|\int_0^{\delta}\varphi(x,u)\cos^{2n}\frac u2\,du\right|\\
&\le\frac{(2n)!!}{2\pi(2n-1)!!}
\int_0^{\delta}|\varphi(x,u)|\cos^{2n}\frac u2\,du
<\frac{\varepsilon}{2}\frac{(2n)!!}{2\pi(2n-1)!!}
\int_0^{\pi}\cos^{2n}\frac u2\,du<\frac{\varepsilon}{2}. \qquad(2)
\end{aligned}
\]
因为 $f(x)$ 是连续周期函数，所以存在常数 $M$ 使得处处有 $|\varphi(x,u)|\le M$. 因而
\[
\begin{aligned}
&\frac{(2n)!!}{2\pi(2n-1)!!}
\left|\int_{\delta}^{\pi}\varphi(x,u)\cos^{2n}\frac u2\,du\right|\\
&\le M\frac{(2n)!!}{2\pi(2n-1)!!}
\int_{\delta}^{\pi}\cos^{2n}\frac u2\,du
\quad\left(\text{令 }v=\frac u2\right)\\
&=M\frac{(2n)!!}{2\pi(2n-1)!!}
\cdot2\int_{\delta/2}^{\pi/2}\cos^{2n}v\,dv
\le M\frac{(2n)!!}{2\pi(2n-1)!!}\cos^{2n}\frac{\delta}{2}. \qquad(3)
\end{aligned}
\]
记 $q=\cos\dfrac{\delta}{2}$，因 $0<\delta<\pi$，所以 $0<q<1$，
$\displaystyle\sum_{n=1}^{\infty}\dfrac{(2n)!!}{(2n-1)!!}q^{2n}$ 收敛，
\[
  \frac{(2n)!!}{(2n-1)!!}q^{2n}\to0\quad(n\to\infty).
\]
故对 $\varepsilon>0$，$\exists N>0$，当 $n>N$ 时，
\[
  M\frac{(2n)!!}{2\pi(2n-1)!!}\cos^{2n}\frac{\delta}{2}<\frac{\varepsilon}{2}. \qquad(4)
\]
最后由 (1)、(2)、(3)、(4) 得
\[
\begin{aligned}
|V_n(x)-f(x)|
&<\frac{(2n)!!}{2\pi(2n-1)!!}
\left|\int_0^{\delta}\varphi(x,u)\cos^{2n}\frac u2\,du\right|\\
&\quad+\frac{(2n)!!}{2\pi(2n-1)!!}
\left|\int_{\delta}^{\pi}\varphi(x,u)\cos^{2n}\frac u2\,du\right|\\
&<\frac{\varepsilon}{2}+\frac{\varepsilon}{2}=\varepsilon.
\end{aligned}
\]
即 $V_n(x)\rightrightarrows f(x)$（当 $n\to\infty$ 时）于 $[-\pi,\pi]$ 上.

\noindent\(\star\)\textbf{例 5.4.26}\quad 设 $f(x)$ 在 $[0,\pi]$ 上有连续导数，且 $f'(x)$ 在 $[0,\pi]$ 上分段光滑，
$\displaystyle\int_0^{\pi}f(x)\,dx=0$，试证：
\[
  \int_0^{\pi}f'^2(x)\,dx\ge\int_0^{\pi}f^2(x)\,dx.
\]

\noindent\textbf{证}\quad 将 $f(x)$ 偶延拓到 $[-\pi,0]$ 上，由已知条件，延拓后的函数能在 $[-\pi,\pi]$ 展开为 Fourier 级数，
$a_0=0,b_n=0$，且可以逐项微分：
\[
  f(x)=\sum_{n=1}^{\infty}a_n\cos nx,\ x\in[0,\pi],\quad
  f'(x)=-\sum_{n=1}^{\infty}na_n\sin nx,\ x\in[0,\pi].
\]
两者均为 Fourier 展开式. 利用 Parseval 等式知
\[
  \frac2\pi\int_0^{\pi}f^2(x)\,dx=\sum_{n=1}^{\infty}a_n^2,
\]
\[
  \frac2\pi\int_0^{\pi}f'^2(x)\,dx
  =\sum_{n=1}^{\infty}n^2a_n^2\ge\sum_{n=1}^{\infty}a_n^2
  =\frac2\pi\int_0^{\pi}f^2(x)\,dx,
\]
故
\[
  \int_0^{\pi}f'^2(x)\,dx\ge\int_0^{\pi}f^2(x)\,dx.
\]

\noindent\(\star\)\textbf{练习}\quad 设 $f(x)$ 在 $\mathbb R$ 上有连续导数，有周期 $\ell(>0)$，且
$\displaystyle\int_0^{\ell}f(x)\,dx=0$. 试利用 Fourier 级数展开式证明：
\[
  \int_0^{\ell}|f'(x)|^2\,dx\ge\frac{4\pi^2}{\ell^2}\int_0^{\ell}|f(x)|^2\,dx, \qquad(1)
\]
其中等号成立当且仅当 $f(x)=a_1\cos\dfrac{2\pi x}{\ell}+b_1\sin\dfrac{2\pi x}{\ell}$.（浙江大学）

\noindent\textbf{提示}\quad
\[
  f(x)\sim\frac{a_0}{2}+\sum_{n=1}^{\infty}
  \left(a_n\cos\frac{2n\pi x}{\ell}+b_n\sin\frac{2n\pi x}{\ell}\right),
\]
\[
  f'(x)\sim\sum_{n=1}^{\infty}
  \left(a_n'\cos\frac{2n\pi x}{\ell}+b_n'\sin\frac{2n\pi x}{\ell}\right),
\]
其中
\[
  a_0'=0,\quad a_n'=\frac{2\pi}{\ell}nb_n,\quad b_n'=-\frac{2\pi}{\ell}na_n
  \quad(n=1,2,\cdots). \qquad(2)
\]
再对 $f(x),f'(x)$ 分别利用 Parseval 等式，可推出（欲证的）不等式 (1).

\noindent\textbf{证}\quad （证明等式 (2).）
\[
\begin{aligned}
a_n&=\frac2{\ell}\int_0^{\ell}f(x)\cos\frac{2n\pi x}{\ell}\,dx\quad(\text{利用分部积分})\\
&=-\frac{\ell}{2n\pi}\cdot\frac2{\ell}\int_0^{\ell}f'(x)\sin\frac{2n\pi x}{\ell}\,dx
=-\frac{\ell}{2n\pi}b_n'\quad(n=1,2,\cdots).
\end{aligned}
\]
同理，有 $b_n=\dfrac{\ell}{2n\pi}a_n'\ (n=1,2,\cdots)$，
$a_0'=\dfrac1{\ell}\int_0^{\ell}f'(x)\,dx=0$（因周期为 $\ell$：$f(\ell)-f(0)=0$）.

检查前面的 Parseval 等式的推导过程，将 $(2\pi)$ 改为 $\ell$，保持有效（半周期 $\pi$ 对应 $\dfrac{\ell}{2}$）. 应用 Parseval 等式：
\[
  \frac2{\ell}\int_0^{\ell}f^2(x)\,dx
  =\frac{a_0}{2}+\sum_{n=1}^{\infty}(a_n^2+b_n^2)
  \xlongequal{a_0=0}\sum_{n=1}^{\infty}(a_n^2+b_n^2), \qquad(3)
\]
\[
\begin{aligned}
\frac2{\ell}\int_0^{\ell}|f'(x)|^2\,dx
&=\sum_{n=1}^{\infty}[(a_n')^2+(b_n')^2]
\xlongequal{\text{式 }(2)}
\sum_{n=1}^{\infty}\left[\left(\frac{2\pi}{\ell}a_n\right)^2+
\left(-\frac{2\pi}{\ell}b_n\right)^2\right]n^2\\
&=\frac{4\pi^2}{\ell^2}\sum_{n=1}^{\infty}(n^2a_n^2+n^2b_n^2)
\ge\frac{4\pi^2}{\ell^2}\sum_{n=1}^{\infty}(a_n^2+b_n^2)\\
&\xlongequal{\text{式 }(3)}
\frac{4\pi^2}{\ell^2}\left(\frac2{\ell}\int_0^{\ell}f^2(x)\,dx\right).
\end{aligned}
\]
不等式 (1) 获证. 此式明显看出，``$\ge$'' 变成 ``$=$'' 的充分必要条件是：
（当 $n\ge2$ 时）$a_n=b_n=0$（否则为 ``$>$''）. 亦即：等号成立的充分必要条件是：
\[
  f(x)=a_1\cos\frac{2\pi x}{\ell}+b_1\sin\frac{2\pi x}{\ell}.
\]

\noindent\textbf{关于 Fourier 级数的一致收敛问题}

因为 Fourier 级数如果一致收敛，其和函数必在 $\mathbb R$ 上连续. 故以 $2\pi$ 为周期的函数要想在 $\mathbb R$ 上展开成一致收敛的 Fourier 级数，必要条件是它在全数轴上连续. 若
$f(x)$ 仅在 $[-\pi,\pi]$ 上给出，则 $f(x)$ 必须在 $[-\pi,\pi]$ 上连续，而且
$f(-\pi)=f(\pi)$. 下例给出该问题的一个充分条件.

\noindent\(\star\)\textbf{例 5.4.27}\quad 证明：若 $f(x)$ 是周期为 $2\pi$ 的连续函数，在 $[-\pi,\pi]$ 上分段光滑，则 $f(x)$ 的 Fourier 级数一致收敛于 $f(x)$.

\noindent\textbf{证}\quad
\[
  f(x)=\frac{a_0}{2}+\sum_{n=1}^{\infty}(a_n\cos nx+b_n\sin nx) \quad(\text{收敛定理}). \qquad(1)
\]
这时 $f'(x)$ 分段连续，其 Fourier 系数 $a_n',b_n'$ 有关系
$a_n'=nb_n,\ b_n'=-na_n$.
（所谓 Fourier 级数逐项求导性，见例 5.4.4 前的要点中性质 2）和例 5.4.7 前的要点 5. 在 $f'$ 分段连续，$f$ 连续，$f(-\pi)=f(\pi)$ 的条件下，该性质可直接通过分部积分得到.）因此
\[
  |a_n|+|b_n|\le\left|\frac{a_n'}{n}\right|+\left|\frac{b_n'}{n}\right|
  \le\frac12(a_n'^2+b_n'^2)+\frac1{n^2}.
\]
（此处应用了平均值不等式，如：
$\left|\dfrac{a_n'}{n}\right|=\sqrt{\dfrac{a_n'^2}{n^2}}\le
\dfrac12\left(a_n'^2+\dfrac1{n^2}\right)$.）根据 Bessel 不等式，可知
\[
  \sum_{n=1}^{\infty}(|a_n|+|b_n|)
  \le\frac12\sum_{n=1}^{\infty}\left(a_n'^2+b_n'^2+\frac2{n^2}\right)<+\infty,
\]
故原 Fourier 级数 (1) 不仅一致收敛，而且绝对一致收敛（于 $\mathbb R$ 上）.

\noindent\textbf{注}\quad （任意区间的情况）若 $f(x)$ 在 $[a,b]$ 上连续，分段光滑，$f(a+0)=f(b-0)$，则
$f(x)$ 可在 $[a,b]$ 上展开成一致收敛的 Fourier 级数.

下例给出一个应用.

\noindent\(\star\)\textbf{例 5.4.28}\quad 设 $f(x)$ 为以 $2\pi$ 为周期，且在 $[-\pi,\pi]$ 上可积的函数，$a_n,b_n$ 是
$f(x)$ 的 Fourier 系数.

1）试求（延迟函数）$f(x+t)$ 的 Fourier 系数；

2）若 $f$ 连续，在 $[-\pi,\pi]$ 上分段光滑，试求卷积函数
\[
  F(x)=\frac1\pi\int_{-\pi}^{\pi}f(t)f(x+t)\,dt
\]
的 Fourier 展开式，并由此推出 Parseval 等式.（哈尔滨工业大学）

\noindent\textbf{解}\quad 1）将 $f(x+t)$ 的 Fourier 系数记作 $A_n,B_n$，则
\[
\begin{aligned}
A_n&=\frac1\pi\int_{-\pi}^{\pi}f(x+t)\cos nx\,dx
\xlongequal{\text{令 }x+t=u}
\frac1\pi\int_{-\pi+t}^{\pi+t}f(u)\cos n(u-t)\,du\\
&=\frac1\pi\int_{-\pi+t}^{\pi+t}f(u)(\cos nu\cos nt+\sin nu\sin nt)\,du\\
&=\cos nt\cdot\frac1\pi\int_{-\pi}^{\pi}f(u)\cos nu\,du
+\sin nt\cdot\frac1\pi\int_{-\pi}^{\pi}f(u)\sin nu\,du\\
&=a_n\cos nt+b_n\sin nt\quad(n=1,2,\cdots).
\end{aligned}
\]
类似地，
\[
  A_0=a_0,\quad B_n=b_n\cos nt-a_n\sin nt\quad(n=1,2,\cdots).
\]

2）根据上例，这时
\[
  f(x)=\frac{a_0}{2}+\sum_{n=1}^{\infty}a_n\cos nx+b_n\sin nx
\]
在 $\mathbb R$ 上一致收敛. 连续周期函数 $f(x)$ 必有界. 上式两边同乘有界函数 $f(x+t)$，仍一致收敛，可以逐项积分. 故
\[
\begin{aligned}
F(x)&=\frac1\pi\int_{-\pi}^{\pi}f(t)f(x+t)\,dt\\
&=\frac1\pi\int_{-\pi}^{\pi}\frac{a_0}{2}f(x+t)\,dt
+\sum_{n=1}^{\infty}\frac1\pi\int_{-\pi}^{\pi}(a_n\cos nt+b_n\sin nt)f(x+t)\,dt\\
&=\frac{a_0}{2}A_0+\sum_{n=1}^{\infty}(a_nA_n+b_nB_n)\quad(\text{利用 1）的结果})\\
&=\frac{a_0^2}{2}+\sum_{n=1}^{\infty}
[a_n(a_n\cos nx+b_n\sin nx)+b_n(b_n\cos nx-a_n\sin nx)]\\
&=\frac{a_0^2}{2}+\sum_{n=1}^{\infty}(a_n^2+b_n^2)\cos nx\quad(-\infty<t<+\infty). \qquad(1)
\end{aligned}
\]
在此式中令 $x=0$，即得 Parseval 等式：
\[
  \frac1\pi\int_{-\pi}^{\pi}f^2(x)\,dx
  =\frac{a_0^2}{2}+\sum_{n=1}^{\infty}(a_n^2+b_n^2).
\]
由此即知式 (1) 一致收敛，进而 (1) 为 $F(x)$ 的 Fourier 级数（留作练习，见练习 5.4.1）.

\noindent\textbf{注}\quad $1^\circ$\quad 求卷积的 Fourier 级数，本例与例 5.4.12 方法、途径完全不同.

$2^\circ$\quad 从式 (1) 可知 $F(x)$ 应是偶函数. 通过变量替换也容易验证（见例 5.4.12）.

\noindent\(\star\)\textbf{例 5.4.29}\quad 构造两个以 $2\pi$ 为周期的连续函数，使其 Fourier 级数在 $[0,\pi]$ 上一致收敛于 0.（北京大学）

\noindent\textbf{提示}\quad 题目只要求以 $2\pi$ 为周期，在 $[0,\pi]$ 上一致收敛于 0，因此 $(-\pi,0)$ 上的值可适当选取.

\noindent\textbf{再提示}\quad 例如，在 $[0,\pi]$ 上令 $f(x)\equiv0$，在 $(-\pi,0)$ 上连续，分段光滑，且
$f(-\pi)=0$.

如在 $[0,\pi]$ 上，设 $f(x)\equiv0$；在 $[-\pi,0)$ 上，设 $f(x)=x(x+\pi)$（或 $f(x)=\sin x$）. 然后以 $2\pi$ 周期，延拓到全数轴，写出 $f(x)$ 的 Fourier 级数，利用例 5.4.27 即知：此 Fourier 级数一致收敛，且在 $[0,\pi]$ 上一致收敛于 0.

\noindent\textbf{注}\quad 任意给定的级数
$\displaystyle\frac{a_0}{2}+\sum_{n=1}^{\infty}a_n\cos nx+b_n\sin nx$
未必一定是某函数的 Fourier 级数. 请看

\noindent\(\star\)\textbf{例 5.4.30}\quad 试证：级数
\[
  \sum_{n=2}^{\infty}\frac1{\ln n}\sin nx
\]
不可能是某函数的 Fourier 级数.

（Fourier 级数比较，这里 $b_n=\dfrac1{\ln n}(n=2,3,\cdots)$，其余系数都为零.）

\noindent\textbf{提示}\quad 作为 Fourier 级数，必要条件是：
$\displaystyle\sum_{n=1}^{\infty}\dfrac{b_n}{n}$ 收敛（见例 5.4.24），但这里，利用 Cauchy 积分判别法，知
$\displaystyle\sum_{n=1}^{\infty}\dfrac{b_n}{n}
=\sum_{n=1}^{\infty}\dfrac1{n\ln n}$ 发散.

\noindent\(\star\)\textbf{例 5.4.31}\quad 设 $f(x),g(x)$ 在 $[a,b]$ 上 Riemann 可积，证明：$f(x),g(x)$ 具有相同 Fourier 系数的充分必要条件是
\[
  \int_a^b|f(x)-g(x)|\,dx=0. \qquad(1)
\]
（北京大学）

\noindent\textbf{证}\quad 分别用 $a_n,b_n;\alpha_n,\beta_n;A_n,B_n$ 表示 $f(x),g(x)$ 及
$F(x)=f(x)-g(x)$ 的 Fourier 系数.

（1）（充分性）若式 (1) 成立，那么
\[
  0\le |a_n-\alpha_n|=
  \left|\frac1{\ell}\int_a^b\left(f(x)\cos\frac{n\pi x}{\ell}
  -g(x)\cos\frac{n\pi x}{\ell}\right)dx\right|
  \le\frac1{\ell}\int_a^b|f(x)-g(x)|\,dx=0
  \quad\left(\ell=\frac{b-a}{2}\right).
\]
推知：$a_n=\alpha_n\ (n=0,1,2,\cdots)$. 同理可证：$b_n=\beta_n\ (n=1,2,\cdots)$.

（2）（必要性）若 Fourier 系数相等，则 $a_n=\alpha_n,b_n=\beta_n$，故
\[
  A_n=a_n-\alpha_n=0,\quad B_n=b_n-\beta_n=0. \qquad(2)
\]
又因为 $f(x),g(x)$ 可积，故 $f(x)-g(x)$ 可积，可用 Parseval 等式：
\[
\begin{aligned}
0\le\left(\int_a^b|f(x)-g(x)|\,dx\right)^2
&\le\left(\int_a^b1^2\,dx\cdot\int_a^b|f(x)-g(x)|^2\,dx\right)\\
&\xlongequal{\text{Parseval 等式}}
(b-a)\left[\frac{A_0^2}{2}+\sum_1^{\infty}(A_n^2+B_n^2)\right]
\xlongequal{\text{式 }(2)}0.
\end{aligned}
\]
因此，式 (1) 成立.

\noindent\(\star\)\textbf{例 5.4.32（Fourier 展开式的唯一性）}\quad 在区间 $[-\pi,\pi]$ 上，函数 $f(x)$ 和
$g(x)$ 分别有 Fourier 级数如下：
\[
  f(x)\sim\frac{a_0}{2}+\sum_{n=1}^{\infty}a_n\cos nx+b_n\sin nx, \qquad(1)
\]
\[
  g(x)\sim\frac{\alpha_0}{2}+\sum_{n=1}^{\infty}\alpha_n\cos nx+\beta_n\sin nx. \qquad(2)
\]
若两级数都在 $[-\pi,\pi]$ 上收敛，两和函数都在 $[-\pi,\pi]$ 上连续并且相等，问：对应的系数是否保持相等？即，是否
$a_n=\alpha_n\ (n=0,1,2,\cdots),\ b_n=\beta_n\ (n=1,2,\cdots)$？如果成立，请给出证明；如果不成立，加上什么条件就能成立？请说明理由.（北京大学）

（该题有较大灵活性，能用不同方式回答，下面 5 种回答，仅供参考.）

\noindent\textbf{回答方式 1}\quad （证明：假定两级数都在 $[-\pi,\pi]$ 上一致收敛，则命题成立.）

用 $S(x)$ 表示公共的和函数，则
\[
  S(x)=\frac{a_0}{2}+\sum_{n=1}^{\infty}a_n\cos nx+b_n\sin nx.
\]
两端同乘（有界函数）$\cos nx$，仍一致收敛，且可逐项求积分，根据三角系的正交性，得
\[
  a_n=\frac1\pi\int_{-\pi}^{\pi}S(x)\cos nx\,dx\quad(n=0,1,2,\cdots).
\]
同理，式 (2) 中的 $\alpha_n$ 也如此. 于是
\[
  a_n=\frac1\pi\int_{-\pi}^{\pi}S(x)\cos nx\,dx=\alpha_n\quad(n=0,1,2,\cdots).
\]
类似可证
\[
  b_n=\beta_n,\quad(n=1,2,\cdots).
\]

\noindent\textbf{回答方式 2}\quad （证明：假设 $f(x),g(x)$ 在 $[-\pi,\pi]$ 上分段光滑，则命题成立.）

若 $f(x),g(x)$ 是以 $2\pi$ 为周期的连续函数，在 $[-\pi,\pi]$ 上分段光滑，则它们的 Fourier 级数一致收敛（如例 5.4.27 所证），接下来用回答方式 1 的方法，得
$a_n=\alpha_n,\ b_n=\beta_n$.

\noindent\textbf{※回答方式 3}\quad （应用 Heine-Cantor 定理.）该定理指出：若两个三角级数：
\[
  \frac{a_0}{2}+\sum_{n=1}^{\infty}a_n\cos nx+b_n\sin nx, \qquad(1)
\]
\[
  \frac{\alpha_0}{2}+\sum_{n=1}^{\infty}\alpha_n\cos nx+\beta_n\sin nx \qquad(2)
\]
在区间 $[-\pi,\pi]$ 上（最多除有限个点外）处处收敛于同一函数 $f(x)$，则两级数恒等，即
$a_n=\alpha_n\ (n=0,1,2,\cdots),\ b_n=\beta_n\ (n=1,2,\cdots)$. 该定理对问题作了肯定的回答.

该定理的证明十分复杂，这里将证明的思路做几点概括：

$1^\circ$\quad 问题等价于：``形如 Fourier 级数的三角级数 (1)，若它在 $[-\pi,\pi]$ 上收敛于 0（可以存在有限个点例外），则式 (1) 的系数必都为零.'' 接下来应用三个定理：

设三角级数为
\[
  \frac{a_0}{2}+\sum_{n=1}^{\infty}a_n\cos nx+b_n\sin nx.
\]
i）（Cantor 引理）若此级数在某有限闭区间上收敛，那么其系数
\[
  a_n\to0,\quad b_n\to0\quad(\text{当 }n\to\infty). \qquad(3)
\]
ii）（Riemann 定理）若此级数的系数 $a_n\to0,b_n\to0$（当 $n\to\infty$），则此三角级数逐项积分两次所得的级数收敛. 记其和为
\[
  F(x)=\frac{a_0x^2}{4}-\sum_{n=1}^{\infty}\frac{a_n\cos nx+b_n\sin nx}{n^2}, \qquad(4)
\]
则 $F(x)$ 的二阶对称导数必为零：$F^{(2)}(x)\equiv0$.

（注\quad 1）关于二阶对称导数的基本知识，参见本书第三章习题 3.2.34 和 3.2.35. 2）上面 $1^\circ$ 中已述：若级数 (1) 收敛于 0（除有限个点外），则 (1) 的系数必为 0. 其实，后来的研究发现（见回答方式 4）：此条件下级数 (1) 必为 Fourier 级数，而 Fourier 级数每点之值等于左、右极限的平均值. 故在这有限个点上级数值只能为 0（尽管我们原先不知道）.）

iii）（Schwarz 定理）二阶对称导数 $F^{(2)}(x)\equiv0$，则 $F$ 必是线性函数：$F(x)=cx+d$.

$2^\circ$\quad （回到本题）式 (4) 可改写为
\[
  \frac{a_0x^2}{4}-\sum_{n=1}^{\infty}\frac{a_n\cos nx+b_n\sin nx}{n^2}=cx+d
\]
或
\[
  \frac{a_0x^2}{4}-cx=d+\sum_{n=1}^{\infty}\frac{a_n\cos nx+b_n\sin nx}{n^2}.
\]
因此只能 $a_0=c=0$. 即得
\[
  0=d+\sum_{n=1}^{\infty}\frac{a_n\cos nx+b_n\sin nx}{n^2}.
\]
可见此级数一致收敛（和为 0）. 根据本题回答方式 1，此式的系数应为 0. 故系数的分子
$a_n=\alpha_n\ (n=0,1,2,\cdots),\ b_n=\beta_n\ (n=1,2,\cdots)$. 证毕.

\noindent\textbf{※回答方式 4}\quad （P. du Bois Reymond 定理）三角级数
\[
  \frac{a_0}{2}+\sum_{n=1}^{\infty}a_n\cos nx+b_n\sin nx
\]
若收敛于有界可积函数，那么此级数必是 $s(x)$ 的 Fourier 级数，下式成了 Fourier 展开式
\[
  s(x)=\frac{a_0}{2}+\sum_{n=1}^{\infty}a_n\cos nx+b_n\sin nx.
\]
（注\quad 此定理被 C. J. de la Vallee Poussin 作了推广：收敛到无界、绝对可积函数也行.）
因此
\[
  a_n=\frac1\pi\int_{-\pi}^{\pi}S(x)\cos nx\,dx=\alpha_n,\quad(n=0,1,2,\cdots).
\]
同理，有
\[
  b_n=\beta_n\quad(n=1,2,\cdots).
\]

\noindent\textbf{注}\quad 第 3、4 两种回答的详细证明，可参看菲赫金哥尔茨的《微积分学教程》第三卷第三分册 732 小节.

\noindent\textbf{※回答方式 5}\quad （应用 Fejer 定理.）已知式 (1) 的和函数为 $S(x)$，即 $\forall x$，级数的部分和
$S_n(x)\to S(x)$（当 $n\to\infty$），因此有
\[
  \sigma_n(x)=\frac{S_1(x)+S_2(x)+\cdots+S_n(x)}{n}\to S(x)\quad(n\to\infty).
\]
题称：$f(x),g(x)$ 有 Fourier 级数，暗示它们为可积函数（或无界但绝对可积的函数），应用 Fejer 定理，在
$f(x)$ 的连续点上 $S(x)=f(x)$，不连续点最多是一个零测集，对积分值没有影响，因此
\[
  a_n=\frac1\pi\int_{-\pi}^{\pi}f(x)\cos nx\,dx
  =\frac1\pi\int_{-\pi}^{\pi}S(x)\cos nx\,dx.
\]
同理有
\[
  \alpha_n=\frac1\pi\int_{-\pi}^{\pi}f(x)\cos nx\,dx
  =\frac1\pi\int_{-\pi}^{\pi}S(x)\cos nx\,dx.
\]
于是有
\[
  a_n=\alpha_n\ (n=0,1,2,\cdots).
\]
类似可证
\[
  b_n=\beta_n\quad(n=1,2,\cdots).
\]

\noindent\textbf{注}\quad Fejer 定理的证明可参看何琛、史济怀、徐森林的《数学分析》第三册；或常庚哲、史济怀的《数学分析教程》下册.

第 5 个回答，虽然简洁明了，但“不连续点最多是一个零测集，对积分值没有影响”是实变函数的知识，超出平常数学分析课程范围.

\subsection*{单元练习 5.4}

\noindent\(\star\)\textbf{5.4.1}\quad 设
$\displaystyle\frac{a_0}{2}+\sum_{k=1}^{\infty}(a_k\cos kx+b_k\sin kx)$
在 $[-\pi,\pi]$ 上一致收敛，试证它必是 $[-\pi,\pi]$ 上其和函数的 Fourier 级数.（西北师范大学）

\noindent\textbf{提示}\quad 用 $\tilde a_n,\tilde b_n$ 表示和函数的 Fourier 系数，利用逐项积分，容易验证
$\tilde a_n=a_n,\tilde b_n=b_n$.

\noindent\(\star\)\textbf{5.4.2}\quad 设
\[
  f(x)=
  \begin{cases}
  0,&-\pi\le x<0,\\
  1,&0\le x\le\pi,
  \end{cases}
\]
1）求 $f(x)$ 的 Fourier 级数；

2）这级数收敛吗？收敛于 $f(x)$ 吗？为什么？

3）这级数在区间 $(-\pi,\pi)$ 里一致收敛吗？为什么？（厦门大学）

\noindent\textbf{提示}\quad 直接用公式和收敛定理.
\[
  f(x)\sim\frac12+\frac2\pi\sum_{n=1}^{\infty}\frac1{2n-1}\sin(2n-1)x
  =
  \begin{cases}
  f(x),&0<|x|<\pi,\\
  \dfrac12,&x=-\pi,0,\pi,\\
  \text{周期 }2\pi,&\text{其他}.
  \end{cases}
\]
在 $(-\pi,\pi)$ 上非一致收敛，因和函数已不连续.

\noindent\textbf{5.4.3}\quad 已知 $f$ 是以 $2\pi$ 为周期的可积函数，它的 Fourier 系数为 $a_n,b_n(n\ge0)$，求函数
\[
  f_1(x)=\frac1{2h}\int_{x-h}^{x+h}f(\xi)\,d\xi\quad(h\ne0)
\]
的 Fourier 系数 $A_n,B_n(n\ge0)$.（西北师范大学，合肥工业大学）
\[
  \left\langle A_0=a_0,\quad A_n=\frac{a_n}{nh}\sin nh,\quad
  B_n=\frac{b_n}{nh}\sin nh\ (n=1,2,\cdots)\right\rangle
\]

\noindent\textbf{5.4.4}\quad 试将 $f(x)=-\pi-x$ 在 $(-\pi,0)$ 内展开成正弦级数，并判断此级数在 $(-\pi,0)$ 是否一致收敛.（河北师范大学）

\noindent\textbf{提示}\quad 作奇延拓即变成例 5.4.7.

\noindent\(\star\)\textbf{5.4.5}\quad 试将周期函数 $f(x)=\arcsin(\sin x)$ 展为 Fourier 级数.（哈尔滨工业大学）
\[
  \left\langle
  f(x)=\frac4\pi\sum_{n=0}^{\infty}\frac{(-1)^n}{(2n+1)^2}\sin(2n+1)x,
  \text{ 在 }\mathbb R\text{ 上处处成立}
  \right\rangle
\]

\noindent\textbf{提示}\quad $f$ 为连续奇函数，以 $2\pi$ 为周期，在 $[0,\pi]$ 上
\[
  f(x)=
  \begin{cases}
  x,&x\in\left[0,\dfrac{\pi}{2}\right],\\
  \pi-x,&x\in\left[\dfrac{\pi}{2},\pi\right].
  \end{cases}
\]

\noindent\textbf{5.4.6}\quad 已知
\[
  f(x)=\frac{\pi}{2}\frac{e^x+e^{-x}}{e^{\pi}-e^{-\pi}},
\]
1）在 $[-\pi,\pi]$ 上将 $f(x)$ 展为 Fourier 级数；
\[
  \left\langle \frac12+\sum_{n=1}^{\infty}\frac{(-1)^n}{1+n^2}\cos nx\right\rangle
\]
2）求级数 $\displaystyle\sum_{n=1}^{\infty}\frac{(-1)^n}{1+(2n)^2}$ 之和.（天津大学）
\[
  \left\langle \frac{\pi}{2}\left(e^{\pi/2}-e^{-\pi/2}\right)^{-1}-\frac12\right\rangle
\]

\noindent\textbf{提示}\quad $f$ 为偶函数，$b_n=0$. 注意
\[
  \int e^{ax}\cos bx\,dx=\frac{e^{ax}}{a^2+b^2}(a\cos bx+b\sin bx),
\]
\[
  \int e^{ax}\sin bx\,dx=\frac{e^{ax}}{a^2+b^2}(a\sin bx-b\cos bx)
\]
不时地被用到.

\noindent\(\star\)\textbf{5.4.7}\quad 设 $f(x)$ 是以 $2\pi$ 为周期的周期函数，且 $f(x)=x,\ -\pi<x<\pi$，求 $f(x)$ 与
$|f(x)|$ 的 Fourier 级数，它们的 Fourier 级数是否一致收敛（给出证明）？（北京大学）

\noindent\textbf{提示}\quad $1^\circ$\quad
\[
  f(x)\sim2\sum_{n=1}^{\infty}\frac{(-1)^{n-1}}{n}\sin nx
  =
  \begin{cases}
  x,&-\pi<x<\pi,\\
  0,&x=-\pi,\pi,\\
  \text{呈周期},&\text{其他}.
  \end{cases}
\]
若级数在 $(-\pi,\pi)$ 内一致收敛，加上 $-\pi,\pi$ 处收敛，可知在 $[-\pi,\pi]$ 上一致收敛. 和函数应当在
$[-\pi,\pi]$ 上连续. 与结果矛盾.

$2^\circ$\quad
\[
  |f(x)|=\frac{\pi}{2}-\frac4\pi\sum_{n=1}^{\infty}
  \frac1{(2n-1)^2}\cos(2n-1)x,\quad x\in[-\pi,\pi]
\]
一致收敛，因 $\displaystyle\sum_{n=1}^{\infty}\frac1{(2n-1)^2}$ 收敛，或利用例 5.4.27 的结论.

\noindent\textbf{5.4.8}\quad 在 $[0,\pi]$ 上将 $f(x)=x+\cos x$ 展开为余弦级数.（华中科技大学）
\[
  \left\langle f(x)=\frac{\pi}{2}+\cos x-\sum_{n=1}^{\infty}
  \frac4{\pi(2n-1)^2}\cos(2n-1)x\right\rangle
\]

\noindent\textbf{提示}\quad $\cos x$ 的 Fourier 级数是它自己，只要求出 $g(x)=x$ 的展开式，两者相加即得.

\noindent\(\star\)\textbf{5.4.9}\quad 试利用 5.4.2 题的结果，求出 $g(x)=\operatorname{sgn}x$，$h(x)=|x|$ 在
$(-\pi,\pi)$ 内的 Fourier 展开式.

\noindent\textbf{提示}\quad $g(x)=2\left(f(x)-\dfrac12\right)$. 于是由 $f$ 可写出 $g$ 的展开式，逐项积分可得
$h(x)$ 的展开式（注意何处可写等号）.

\noindent\textbf{5.4.10}\quad 设 $f(x)=x,\ x\in\left[0,\dfrac{\pi}{2}\right]$，试将 $f(x)$ 展开成
$\displaystyle\sum_{n=1}^{\infty}b_{2n-1}\sin(2n-1)x$ 型的三角级数.

\noindent\textbf{提示}\quad 参考对比 5.4.5 题的结果，寻查做法，下题给出一般结果.

\noindent\(\star\)\textbf{5.4.11}\quad 设 $f(x)$ 以 $2\pi$ 为周期，$[-\pi,\pi]$ 上可积，$a_n,b_n$ 是它的 Fourier 系数. 试证：
\[
\begin{array}{ll}
1)\ f(-x)=f(x),\ f(\pi-x)=-f(x)
&\Rightarrow
\begin{cases}
b_n=0,&n=1,2,\cdots,\\
a_{2n}=0,&n=0,1,2,\cdots;
\end{cases}\\[10pt]
2)\ f(-x)=f(x),\ f(\pi-x)=f(x)
&\Rightarrow
\begin{cases}
b_n=0,&n=1,2,\cdots,\\
a_{2n-1}=0,&n=1,2,\cdots;
\end{cases}\\[10pt]
3)\ f(-x)=-f(x),\ f(\pi-x)=-f(x)
&\Rightarrow
\begin{cases}
a_n=0,&n=0,1,2,\cdots,\\
b_{2n-1}=0,&n=1,2,\cdots;
\end{cases}\\[10pt]
4)\ f(-x)=-f(x),\ f(\pi-x)=f(x)
&\Rightarrow
\begin{cases}
a_n=0,&n=0,1,\cdots,\\
b_{2n}=0,&n=1,2,\cdots.
\end{cases}
\end{array}
\]

\noindent\textbf{提示}\quad 可用系数公式直接验证. 注意：$f(-x)=f(x)$（或 $-f(x)$）是偶（或奇）性条件，导致
$b_n=0$（或 $a_n=0$）. 这是共知，好记.

$f(\pi-x)=-f(x)$ 表明图形关于点 $\left(\dfrac{\pi}{2},0\right)$ 中心对称.

$f(\pi-x)=f(x)$ 表明图形关于直线 $x=\dfrac{\pi}{2}$ 轴对称.

该例的结果可帮助我们预料和校验计算结果. 例如 5.4.5 题.
感兴趣的读者，不妨用“偶心奇轴皆无偶，奇心偶轴皆无奇”两句口诀来记忆. 意即：偶函数在对点
$\left(\dfrac{\pi}{2},0\right)$ 作中心延拓，或奇函数对直线 $x=\dfrac{\pi}{2}$ 作轴对称延拓时，系数下标就不会有偶数出现. 第二句类似.

回头再做上题就容易了.

\noindent\textbf{5.4.12}\quad 求下列函数在指定区间上的 Fourier 级数：

1）
\[
  f(x)=
  \begin{cases}
  x,&x\in[0,\pi],\\
  2,&x\in[-\pi,0),
  \end{cases}
\]
于 $[-\pi,\pi]$ 上；（中山大学）

2）$f(x)=x+x^2$，于 $[-\pi,\pi]$ 上，并求
$\displaystyle\sum_{n=1}^{\infty}\frac1{n^2}$；（中南大学）

3）$f(x)=\left(\dfrac{\pi-x}{2}\right)^2$，于 $(0,2\pi]$ 上，并求
$\displaystyle\sum_{n=1}^{\infty}\frac1{n^2}$；（复旦大学）

4）
\[
  f(x)=
  \begin{cases}
  e^x,&\left[0,\dfrac{\pi}{2}\right],\\
  0,&\left[-\dfrac{\pi}{2},0\right),
  \end{cases}
\]
于 $\left[-\dfrac{\pi}{2},\dfrac{\pi}{2}\right]$，并求和函数；（湘潭大学）

5）$f(x)=x$，于 $(0,2)$ 上，按余弦展开；（国防科技大学）

6）$f(x)=1$，于 $(0,\pi]$ 上，按正弦展开，并求和函数.（南京大学）

\noindent\textbf{**5.4.13}\quad 求函数
$f(x)=\ln\left(2\cos\dfrac{x}{2}\right)$ 在 $(-\pi,\pi)$ 内的 Fourier 级数展开式.
\[
  \left\langle \sum_{n=1}^{\infty}(-1)^{n-1}\frac{\cos nx}{n}\quad(-\pi<x<\pi)\right\rangle
\]

\noindent\textbf{*5.4.14}\quad 证明级数
\[
  \sum_{n=1}^{\infty}\frac{\sin nx}{\ln(n+1)}
\]
不可能是某个可积函数 $f(x)$ 的 Fourier 级数.

\noindent\textbf{提示}\quad 可用反证法及 Fourier 级数逐项积分定理.

\noindent\textbf{*5.4.15}\quad 写出
\[
  f(x)=
  \begin{cases}
  1,&\text{当 }|x|\le\alpha,\\
  0,&\text{当 }\alpha<|x|\le\pi
  \end{cases}
\]
的 Fourier 级数，并根据 Parseval 等式求和：
\[
  \left\langle f(x)\sim\frac{\alpha}{\pi}+
  \sum_{n=1}^{\infty}\frac{2\sin n\alpha}{n\pi}\cos nx\right\rangle
\]
1）$\displaystyle\sum_{n=1}^{\infty}\frac{\sin^2n\alpha}{n^2}$；
\[
  \left\langle\frac{\alpha(\pi-\alpha)}{2}\right\rangle
\]
2）$\displaystyle\sum_{n=1}^{\infty}\frac{\cos^2n\alpha}{n^2}$（已知 $\displaystyle\sum_{n=1}^{\infty}\frac1{n^2}=\frac{\pi^2}{6}$）.
\[
  \left\langle\frac16(\pi^2-3\pi\alpha+3\alpha^2)\right\rangle
\]

\noindent\textbf{※5.4.16}\quad 设 $f(x)$ 是以 $2\pi$ 为周期的函数，在 $[-\pi,\pi]$ 上可积，则已知它的 Fourier 级数的部分和
$S_n(x)$ 可表示为 Dirichlet 积分：
\[
  S_n(x)=\frac1\pi\int_{-\pi}^{\pi}f(x+t)
  \frac{\sin\left(n+\dfrac12\right)t}{2\sin\dfrac t2}\,dt,
\]
其中
\[
  \frac{\sin\left(n+\dfrac12\right)t}{2\sin\dfrac t2}
  =\frac12+\cos t+\cos2t+\cdots+\cos nt=D_n(t)
\]
称为 Dirichlet 核. $S_n(x)$ 的平均值
$\displaystyle\sigma_n(x)=\frac1n\sum_{k=0}^{n-1}S_k(x)$ 称为 Cesaro 和. 试证：

1）
\[
  D_0(x)+\cdots+D_{n-1}(x)=\frac12\left(\frac{\sin\dfrac n2x}{\sin\dfrac x2}\right)^2;
\]
2）
\[
  \frac1{2n\pi}\int_{-\pi}^{\pi}
  \left(\frac{\sin\dfrac n2x}{\sin\dfrac x2}\right)^2\,dx=1;
\]
3）$\forall\delta>0$，
\[
  \frac1{n\pi}\int_{\delta}^{\pi}
  \left(\frac{\sin\dfrac n2x}{\sin\dfrac x2}\right)^2\,dx\to0
  \quad(\text{当 }n\to\infty);
\]
4）若 $f(x)$ 是以 $2\pi$ 为周期的连续函数，则当 $n\to\infty$ 时
$\sigma_n(x)\rightrightarrows f(x)$ 于 $[-\pi,\pi]$ 上.

\noindent\textbf{证}\quad 1）
\[
  2\sin\frac x2\sum_{k=0}^{n-1}\sin\left(k+\frac12\right)x
  =\sum_{k=0}^{n-1}[\cos kx-\cos(k+1)x]
  =1-\cos nx=2\sin^2\frac{nx}{2}.
\]
所以
\[
  \sum_{k=0}^{n-1}D_k(x)
  =\frac{\displaystyle\sum_{k=0}^{n-1}\sin\left(k+\frac12\right)x}{2\sin\dfrac x2}
  =\frac{\sin^2\dfrac{nx}{2}}{2\sin^2\dfrac x2}. \qquad(1)
\]
2）
\[
\begin{aligned}
\frac1{2n\pi}\int_{-\pi}^{\pi}
\left(\frac{\sin\dfrac{nx}{2}}{\sin\dfrac x2}\right)^2dx
&\xlongequal{\text{式 }(1)}
\frac1{2n\pi}\int_{-\pi}^{\pi}2\sum_{k=0}^{n-1}D_k(x)\,dx\\
&=\frac1{2n\pi}\sum_{k=0}^{n-1}\int_{-\pi}^{\pi}2\left(\frac12+\sum_{i=1}^{k}\cos ix\right)dx=1. \qquad(2)
\end{aligned}
\]
3）$\forall\delta:0<\delta<\pi$，
\[
  \frac1{n\pi}\int_{\delta}^{\pi}
  \left(\frac{\sin\dfrac{nx}{2}}{\sin\dfrac x2}\right)^2dx
  \le\frac1{n\pi}\cdot\frac{\pi-\delta}{\sin^2\dfrac{\delta}{2}}\to0
  \quad(n\to\infty). \qquad(3)
\]
4）
\[
  S_n(x)=\frac1\pi\int_{-\pi}^{\pi}f(x+t)
  \frac{\sin\left(n+\dfrac12\right)t}{2\sin\dfrac t2}\,dt
  =\frac1\pi\int_{-\pi}^{\pi}f(x+t)D_n(t)\,dt,
\]
\[
\begin{aligned}
\sigma_n(x)&=\frac1n\sum_{k=0}^{n-1}S_k(x)
=\frac1{n\pi}\sum_{k=0}^{n-1}\int_{-\pi}^{\pi}f(x+t)D_k(t)\,dt\\
&=\frac1{n\pi}\int_{-\pi}^{\pi}f(x+t)\sum_{k=0}^{n-1}D_k(t)\,dt
\xlongequal{\text{式 }(1)}
\frac1{2n\pi}\int_{-\pi}^{\pi}f(x+t)
\left(\frac{\sin\dfrac{nt}{2}}{\sin\dfrac t2}\right)^2dt. \qquad(4)
\end{aligned}
\]
另一方面，（利用拟合法的思想）由式 (2) 两端同乘 $f(x)$，可得
\[
  f(x)=\frac1{2n\pi}\int_{-\pi}^{\pi}f(x)
  \left(\frac{\sin\dfrac{nt}{2}}{\sin\dfrac t2}\right)^2dt. \qquad(5)
\]
（4）、（5）两式相减得
\[
  \sigma_n(x)-f(x)=\frac1{2n\pi}\int_{-\pi}^{\pi}
  [f(x+t)-f(x)]
  \left(\frac{\sin\dfrac{nt}{2}}{\sin\dfrac t2}\right)^2dt. \qquad(6)
\]
因为 $f(x)$ 在 $[-\pi,\pi]$ 连续，知（有界）：$\exists M>0$，当 $x\in[-\pi,\pi]$ 时，恒有
\[
  |f(x)|\le M. \qquad(7)
\]
根据 Cantor 定理，$f(x)$ 在 $[-\pi,\pi]$ 一致连续，故 $\forall\varepsilon>0$，$\exists\delta:\pi>\delta>0$，当
$|t|<\delta$ 时，只要 $x,x+t\in[-\pi,\pi]$，则恒有
\[
  |f(x+t)-f(x)|\le\frac{\varepsilon}{3}. \qquad(8)
\]
为了证明：（当 $n\to\infty$ 时）$\sigma_n(x)\rightrightarrows f(x)$，只需证明：
$\forall\varepsilon>0$，$\exists N>0$，当 $n>N$ 时，

\[
  \left|\frac1{2n\pi}\int_{-\pi}^{\pi}[f(x+t)-f(x)]
  \left(\frac{\sin\dfrac{nt}{2}}{\sin\dfrac t2}\right)^2dt\right|
  <\varepsilon\quad(\forall x\in[-\pi,\pi]). \qquad(9)
\]
为此，将式 (6) 右端的积分析为三段：
\[
  \sigma_n(x)-f(x)=
  \left(\frac1{2n\pi}\int_{-\pi}^{-\delta}
  +\frac1{2n\pi}\int_{-\delta}^{\delta}
  +\frac1{2n\pi}\int_{\delta}^{\pi}\right)
  [f(x+t)-f(x)]
  \left(\frac{\sin\dfrac{nt}{2}}{\sin\dfrac t2}\right)^2dt
  =I_1+I_2+I_3.
\]
那么
\[
\begin{aligned}
|I_2|
&=\left|\frac1{2n\pi}\int_{-\delta}^{\delta}[f(x+t)-f(x)]
\left(\frac{\sin\dfrac{nt}{2}}{\sin\dfrac t2}\right)^2dt\right|\\
&\le\frac1{2n\pi}\int_{-\delta}^{\delta}|f(x+t)-f(x)|
\left(\frac{\sin\dfrac{nt}{2}}{\sin\dfrac t2}\right)^2dt\\
&\xlongequal{\text{式 }(8)}
\frac{\varepsilon}{3}\cdot\frac1{2n\pi}\int_{-\delta}^{\delta}
\left(\frac{\sin\dfrac{nt}{2}}{\sin\dfrac t2}\right)^2dt
\le\frac{\varepsilon}{3}\cdot\frac1{2n\pi}\int_{-\pi}^{\pi}
\left(\frac{\sin\dfrac{nt}{2}}{\sin\dfrac t2}\right)^2dt
=\frac{\varepsilon}{3}
\end{aligned}
\]
（利用 2）的结果）.
\[
\begin{aligned}
|I_3|
&=\left|\frac1{2n\pi}\int_{\delta}^{\pi}[f(x+t)-f(x)]
\left(\frac{\sin\dfrac{nt}{2}}{\sin\dfrac t2}\right)^2dt\right|\\
&\xlongequal{\text{式 }(7)}
\le\frac{M}{n\pi}\int_{\delta}^{\pi}
\left(\frac{\sin\dfrac{nt}{2}}{\sin\dfrac t2}\right)^2dt\to0\quad(n\to\infty)
\end{aligned}
\]
（利用 3）的结果）.

故对上面的 $\varepsilon>0$，$\exists N>0$，当 $n>N$ 时，能使得
$|I_3|\le\dfrac{\varepsilon}{3}$.

而
\[
\begin{aligned}
|I_1|
&=\left|\frac1{2n\pi}\int_{-\pi}^{-\delta}[f(x+t)-f(x)]
\left(\frac{\sin\dfrac{nt}{2}}{\sin\dfrac t2}\right)^2dt\right|\\
&\le\frac{M}{n\pi}\int_{-\pi}^{-\delta}
\left(\frac{\sin\dfrac{nt}{2}}{\sin\dfrac t2}\right)^2dt
\xlongequal{\text{偶性}}
\frac{M}{n\pi}\int_{\delta}^{\pi}
\left(\frac{\sin\dfrac{nt}{2}}{\sin\dfrac t2}\right)^2dt
\le\frac{\varepsilon}{3}.
\end{aligned}
\]
因此式 (9) 成立：
\[
  |\sigma_n(x)-f(x)|\le |I_1|+|I_2|+|I_3|
  \le\frac{\varepsilon}{3}+\frac{\varepsilon}{3}+\frac{\varepsilon}{3}=\varepsilon.
\]
一致收敛获证.

\noindent\textbf{※5.4.17}\quad 设 $f(x)$ 是以 $2\pi$ 为周期的连续函数，$S_n(x)$ 是 $f(x)$ 的 Fourier 级数的部分和，
\[
  g_n(x)=\int_{-\pi}^{\pi}
  \frac{\cos(x-u)}{\sqrt{1+\sin^2(x+u)}}S_n(u)\,du.
\]
试证：

1）存在与 $x,n$ 无关的数 $K$，使得 $|g_n(x)|\le K\ (x\in[-\pi,\pi])$；

2）当 $n\to\infty$ 时，
\[
  g_n(x)\to\int_{-\pi}^{\pi}
  \frac{\cos(x-u)}{\sqrt{1+\sin^2(x+u)}}f(u)\,du
\]
于 $[-\pi,\pi]$ 上.

\noindent\textbf{提示}\quad 利用 Schwarz 不等式及 Parseval 等式.

\noindent\textbf{※5.4.18}\quad 设 $T_n(x)$ 为 $n$ 阶三角多项式：
\[
  T_n(x)=\frac{\alpha_0}{2}+\sum_{k=1}^{n}(\alpha_k\cos kx+\beta_k\sin kx).
\]
试证：

1）
\[
  \max_{-\pi\le x\le\pi}|T_n'(x)|
  \le n^2\max_{-\pi\le x\le\pi}|T_n(x)|\quad(n>1),
\]
当 $\alpha_0\ge0$ 时，对 $n=1$ 也成立；

2）若 $\alpha_{n-1}=1$，则
\[
  \max_{-\pi\le x\le\pi}|T_n(x)|\ge\frac{\pi}{4}\quad(n>1).
\]

\noindent\textbf{提示}\quad 1）可应用 Cauchy 不等式. 2）可考虑积分
$\displaystyle\int_{-\pi}^{\pi}T_n(x)\cos(n-1)x\,dx$.

\noindent\textbf{证}\quad 1）
\[
  |T_n'(x)|=\left|\sum_{k=1}^{n}k(-\alpha_k\sin kx+\beta_k\cos kx)\right|
  \le n\sum_{k=1}^{n}(|\alpha_k||\sin kx|+|\beta_k||\cos kx|).
\]
可用多种方法证明下面的不等式，例如应用 Cauchy 不等式：
\[
  |\alpha_k||\sin kx|+|\beta_k||\cos kx|
  \le(\alpha_k^2+\beta_k^2)^{1/2}(\sin^2kx+\cos^2kx)^{1/2}
  =(\alpha_k^2+\beta_k^2)^{1/2}.
\]
\[
\begin{aligned}
|T_n'(x)|
&\le n\sum_{k=1}^{n}1\cdot(\alpha_k^2+\beta_k^2)^{1/2}
\le n\sqrt n\left[\sum_{k=1}^{n}(\alpha_k^2+\beta_k^2)\right]^{1/2}\\
&\le n\sqrt n\left[\frac{\alpha_0^2}{2}
+\sum_{k=1}^{n}(\alpha_k^2+\beta_k^2)\right]^{1/2}
\le n\sqrt n\left(\frac1\pi\int_{-\pi}^{\pi}T_n^2(x)\,dx\right)^{1/2}\\
&\le n\sqrt n\max_{-\pi\le x\le\pi}|T_n(x)|\sqrt2.
\end{aligned}
\]
当 $n\ge2$ 时，
\[
  |T_n'(x)|\le n^2\max_{-\pi\le x\le\pi}|T_n(x)|\quad(\forall x\in[-\pi,\pi]).
\]
因此，当 $n\ge2$ 时，有
\[
  \max_{-\pi\le x\le\pi}|T_n'(x)|
  \le n^2\max_{-\pi\le x\le\pi}|T_n(x)|.
\]
当 $n=1$ 时，可直接验证：当 $\alpha_0\ge0$ 时，不等式成立. 因为
\[
\begin{aligned}
T_1'(x)&=-\alpha_1\sin x+\beta_1\cos x\\
&=\sqrt{\alpha_1^2+\beta_1^2}
\left(-\frac{\alpha_1}{\sqrt{\alpha_1^2+\beta_1^2}}\sin x
+\frac{\beta_1}{\sqrt{\alpha_1^2+\beta_1^2}}\cos x\right)\\
&=\sqrt{\alpha_1^2+\beta_1^2}\cos(x+\theta_1),
\quad \text{其中 }\sin\theta_1=\frac{\alpha_1}{\sqrt{\alpha_1^2+\beta_1^2}}.
\end{aligned}
\]
\[
  T_1(x)=\frac{\alpha_0}{2}+\alpha_1\cos x+\beta_1\sin x
  =\frac{\alpha_0}{2}+\sqrt{\alpha_1^2+\beta_1^2}\sin(x+\theta_1).
\]
所以当 $\alpha_0\ge0$ 时，
\[
  \max_{-\pi\le x\le\pi}|T_1'(x)|
  =\sqrt{\alpha_1^2+\beta_1^2}
  \le\left|\frac{\alpha_0}{2}\right|+\sqrt{\alpha_1^2+\beta_1^2}
  =\max_{-\pi\le x\le\pi}|T_1(x)|.
\]

2）设 $n>1$，因已知 $\alpha_{n-1}=1$，那么
\[
\begin{aligned}
1=\alpha_{n-1}
&=\frac1\pi\int_{-\pi}^{\pi}T_{n-1}(x)\cos(n-1)x\,dx\\
&\le\frac1\pi\int_{-\pi}^{\pi}|T_{n-1}(x)||\cos(n-1)x|\,dx\\
&\le\max_{-\pi\le x\le\pi}|T_{n-1}(x)|\cdot
\frac1\pi\int_{-\pi}^{\pi}|\cos(n-1)x|\,dx
\xlongequal{\text{见下}}
\max_{-\pi\le x\le\pi}|T_{n-1}(x)|\cdot\frac4\pi. \qquad(1)
\end{aligned}
\]
（因为其中：
\[
\begin{aligned}
\int_{-\pi}^{\pi}|\cos(n-1)x|\,dx
&\xlongequal{\text{令 }(n-1)x=t}
\frac1{n-1}\int_{-(n-1)\pi}^{(n-1)\pi}|\cos t|\,dt\\
&=\frac2{n-1}\int_0^{(n-1)\pi}|\cos t|\,dt
\xlongequal{|\cos t|\text{ 周期为 }\pi}
2\int_0^{\pi}|\cos t|\,dt=4.
\end{aligned}
\]
）
因此，式 (1) 最后的等号成立. 结论 2）获证.

\end{document}
